\documentclass{article}
\usepackage{z-eves}
\begin{document}

\section*{Unary primitive recursive functions in Z/EVES}

Formalization of Section 6 of D. E. Severin, \emph{Unary primitive recursive functions},
J. Symbolic Logic 73(4), 2008, at the level of closures (as in the Mizar version), and of the
two open questions answered in the README.

Conventions: one-argument functions are the elements of $\nat \fun \nat$.  A function or
operator $X$ of the paper is introduced by an axiomatic box (\texttt{funX} or \texttt{opX})
whose defining equations are labelled rewrite rules.  Every axiomatic box is followed by a
consistency theorem \texttt{XEx}, which exhibits a witness (\texttt{XWit}, defined by an
explicit term) satisfying the box; so no definition adds an axiom that could be inconsistent.

\section{Building functions}

Z/EVES has no rules showing that a $\lambda$-term is a total function; these four lemmas
provide them.  A relation that is single-valued and total on $X$ is a total function.

\begin[disabled]{theorem}{rule RelPow}[X, Y]
  X \rel Y = \power (X \cross Y)
\end{theorem}

\begin{zproof}
use relDefinition[X, Y];
prove;
\end{zproof}

\begin[disabled]{theorem}{rule PfunIntro}[X, Y]
  f \in X \rel Y \land (\forall x: X; y, z: Y | (x, y) \in f \land (x, z) \in f @ y = z)
  \implies (f \in X \pfun Y \iff true)
\end{theorem}

\begin{zproof}
prove;
use pfunDef[X, Y][R := f];
prove;
with enabled (inPower, compDef) prove;
with enabled (invDef, inId) prove by reduce;
instantiate x\_\_0 == y, y\_\_0 == x, z\_\_0 == z;
prove;
\end{zproof}

\begin[disabled]{theorem}{rule DomTotal}[X, Y]
  f \in X \rel Y \land (\forall x: X @ \exists y: Y @ (x, y) \in f) \implies \dom f = X
\end{theorem}

\begin{zproof}
apply extensionality;
prove;
with enabled (inDom) prove;
instantiate x == y;
prove;
instantiate y\_\_2 == y\_\_0;
prove;
\end{zproof}

\begin[disabled]{theorem}{rule FunIntro}[X, Y]
  f \in X \rel Y \land (\forall x: X @ \exists y: Y @ (x, y) \in f)
     \land (\forall x: X; y, z: Y | (x, y) \in f \land (x, z) \in f @ y = z)
  \implies (f \in X \fun Y \iff true)
\end{theorem}

\begin{zproof}
use pfunIsFun[X, Y][A := X, B := Y, f := f];
with enabled (PfunIntro, DomTotal) prove;
\end{zproof}

\begin{theorem}{rule FunIsPfun}[X, Y]
  f \in X \fun Y \implies (f \in X \pfun Y \iff true)
\end{theorem}

\begin{zproof}
use fun\_type[X, Y];
with enabled (inPower) prove;
\end{zproof}

Relational composition and iteration of one-argument functions (the witness for iteration
$F^{\square(a)}$ uses the Toolkit's iteration $F^n$).

\begin{theorem}{rule FunIsRel}[X, Y]
  f \in X \fun Y \implies (f \in X \rel Y \iff true)
\end{theorem}

\begin{zproof}
use inPowerTransitive[X := X \fun Y, Y := X \pfun Y, Z := X \rel Y];
use fun\_type[X, Y];
use pfun\_type[X, Y];
prove;
\end{zproof}

\begin{theorem}{DomFunGen}[X, Y]
  \forall f: X \fun Y @ \dom f = X
\end{theorem}

\begin{zproof}
prove;
\end{zproof}

\begin{theorem}{BigcapSub}[X]
  \forall T: \power (\power X); Y: \power X | Y \in T @ \bigcap T \subseteq Y
\end{theorem}

\begin{zproof}
prove;
\end{zproof}

\begin{theorem}{SubAntisym}[X]
  \forall A, B: \power X | A \subseteq B \land B \subseteq A @ A = B
\end{theorem}

\begin{zproof}
prove;
apply extensionality to predicate A = B;
with enabled (subsetDef, inPower) prove;
\end{zproof}

\begin{theorem}{InPair}[X]
  \forall u, p, q: X | u \in \{ p, q \} @ u = p \lor u = q
\end{theorem}

\begin{zproof}
prove;
\end{zproof}

\begin{theorem}{rule CompRelFun}
  \forall f, g: \nat \fun \nat @ f \comp[\nat, \nat, \nat] g \in \nat \fun \nat
\end{theorem}
\begin{zproof}
prove;
with enabled (easyDomComp) prove;
\end{zproof}

\begin{theorem}{rule CompRelApp}
  \forall f, g: \nat \fun \nat; x: \nat @ (f \comp[\nat, \nat, \nat] g)~x = g~(f~x)
\end{theorem}

\begin{zproof}
prove;
\end{zproof}

\begin{zed}
  IterSet == \{ n: \nat | \forall F: \nat \fun \nat @ iter[\nat]~n~F \in \nat \fun \nat \}
\end{zed}

\begin{zproof}[IterSet\$domainCheck]
prove;
use applyInRanFun[\num, \power (\nat \cross \nat) \fun \power (\nat \cross \nat)][f := iter[\nat], a := n];
use DomFunGen[\power (\nat \cross \nat), \power (\nat \cross \nat)][f := iter[\nat]~n];
prove;
equality substitute;
prove;
\end{zproof}

\begin{theorem}{rule IterBase}
  0 \in IterSet
\end{theorem}

\begin{zproof}
invoke IterSet;
prove;
\end{zproof}

\begin{theorem}{rule IterStep}
  \forall n: IterSet @ 1 + n \in IterSet
\end{theorem}

\begin{zproof}
invoke IterSet;
prove;
use iterPositive[\nat][R\_\_1 := F, n\_\_0 := n];
instantiate F\_\_0 == F;
with enabled (easyDomComp) prove;
\end{zproof}

\begin{theorem}{IterSub}
  \nat \subseteq IterSet
\end{theorem}

\begin{zproof}
apply natInduction to predicate \nat \subseteq IterSet;
prove;
invoke IterSet;
prove;
\end{zproof}

\begin{theorem}{rule IterFun}
  \forall F: \nat \fun \nat; n: \nat @ iter[\nat]~n~F \in \nat \fun \nat
\end{theorem}

\begin{zproof}
prove;
use IterSub;
with enabled (inPower) prove;
instantiate e == n;
invoke IterSet;
prove;
instantiate F\_\_1 == F;
prove;
\end{zproof}

\begin{theorem}{IterSucR}
  \forall F: \nat \fun \nat; k, a: \nat @ iter[\nat]~(1 + k)~F~a = iter[\nat]~k~F~(F~a)
\end{theorem}

\begin{zproof}
prove;
use composePositiveIterates[\nat][n := 1, k := k, R := F, Y := \nat];
use CompRelApp[f := F, g := iter[\nat]~k~F, x := a];
use IterFun[F := F, n := k];
prove;
\end{zproof}

\begin{theorem}{rule IterInDom}
  \forall F: \nat \fun \nat; n: \nat @ F \in \dom (iter[\nat]~n)
\end{theorem}

\begin{zproof}
prove;
use applyInRanFun[\num, \power (\nat \cross \nat) \fun \power (\nat \cross \nat)][f := iter[\nat], a := n];
use DomFunGen[\power (\nat \cross \nat), \power (\nat \cross \nat)][f := iter[\nat]~n];
prove;
\end{zproof}

\begin{theorem}{rule IterZero}
  \forall F: \nat \fun \nat; a: \nat @ iter[\nat]~0~F~a = a
\end{theorem}

\begin{zproof}
prove;
\end{zproof}

\begin{theorem}{rule IterSuc}
  \forall F: \nat \fun \nat; n, a: \nat @ iter[\nat]~(1 + n)~F~a = F~(iter[\nat]~n~F~a)
\end{theorem}

\begin{zproof}
prove;
use composePositiveIterates[\nat][n := n, k := 1, R := F, Y := \nat];
use CompRelApp[f := iter[\nat]~n~F, g := F, x := a];
use IterFun[F := F, n := n];
prove;
\end{zproof}

The domain of a one-argument function, and extensionality.

\begin{theorem}{rule DomUF}
  \forall f: \nat \fun \nat @ \dom f = \nat
\end{theorem}

\begin{zproof}
prove;
\end{zproof}

\begin{theorem}{rule UFAppNat}
  \forall f: \nat \fun \nat; x: \nat @ f~x \in \nat
\end{theorem}

\begin{zproof}
prove;
use applyInRanFun[\nat, \nat][f := f, a := x];
prove;
\end{zproof}

\begin{theorem}{LeOrGt}
  \forall a, b: \num @ a \leq b \lor b < a
\end{theorem}

\begin{zproof}
prove;
\end{zproof}

\begin{theorem}{UFNonneg}
  \forall f: \nat \fun \nat; x: \nat @ f~x \geq 0
\end{theorem}

\begin{zproof}
prove;
use applyInRanFun[\nat, \nat][f := f, a := x];
with disabled (UFAppNat) prove;
\end{zproof}

\begin{theorem}{rule NatAdd}
  \forall a, b: \nat @ a + b \in \nat
\end{theorem}

\begin{zproof}
prove;
\end{zproof}

\begin{theorem}{rule NatMul}
  \forall a, b: \nat @ a * b \in \nat
\end{theorem}

\begin{zproof}
prove;
\end{zproof}

\begin{theorem}{rule DomUOP}
  \forall u: (\nat \fun \nat) \fun (\nat \fun \nat) @ \dom u = \nat \fun \nat
\end{theorem}

\begin{zproof}
prove;
\end{zproof}

\begin{theorem}{rule DomBOP}
  \forall b: (\nat \fun \nat) \cross (\nat \fun \nat) \fun (\nat \fun \nat) @
    \dom b = (\nat \fun \nat) \cross (\nat \fun \nat)
\end{theorem}

\begin{zproof}
prove;
\end{zproof}

\begin{theorem}{rule AppliesUOP}
  \forall u: (\nat \fun \nat) \fun (\nat \fun \nat); F: \nat \fun \nat @ u applies\$to F
\end{theorem}

\begin{zproof}
prove;
\end{zproof}

\begin{theorem}{rule AppliesBOP}
  \forall b: (\nat \fun \nat) \cross (\nat \fun \nat) \fun (\nat \fun \nat); F, G: \nat \fun \nat @
    b applies\$to (F, G)
\end{theorem}

\begin{zproof}
prove;
\end{zproof}

\begin{theorem}{rule AppliesUF}
  \forall f: \nat \fun \nat; x: \nat @ f applies\$to x
\end{theorem}

\begin{zproof}
prove;
\end{zproof}

\begin{theorem}{rule UOPApp}
  \forall u: (\nat \fun \nat) \fun (\nat \fun \nat); F: \nat \fun \nat @ u~F \in \nat \fun \nat
\end{theorem}

\begin{zproof}
prove;
\end{zproof}

\begin{theorem}{rule BOPApp}
  \forall b: (\nat \fun \nat) \cross (\nat \fun \nat) \fun (\nat \fun \nat); F, G: \nat \fun \nat @
    b(F, G) \in \nat \fun \nat
\end{theorem}

\begin{zproof}
prove;
\end{zproof}

\begin{theorem}{rule NatUOPApp}
  \forall b: \nat \fun ((\nat \fun \nat) \fun (\nat \fun \nat)); a: \nat @
    b~a \in (\nat \fun \nat) \fun (\nat \fun \nat)
\end{theorem}

\begin{zproof}
prove;
\end{zproof}

\begin[disabled]{theorem}{rule UFExt}
  \forall F, G: \nat \fun \nat @ F = G \iff (\forall x: \nat @ F~x = G~x)
\end{theorem}

\begin{zproof}
use funExtensionality[\nat, \nat][f := F, g := G];
prove;
\end{zproof}

\section{Section 2: functions and operators}

\subsection{The identity and the constants}

\begin{axdef}
  funI: \nat \fun \nat
\where
  \Label{rule IVal} \forall x: \nat @ funI~x = x
\end{axdef}

\begin{zproof}[funI\$domainCheck]
prove;
\end{zproof}

\begin{zed}
  IWit == (\lambda x: \nat @ x)
\end{zed}

\begin{theorem}{IWitFun}
  IWit \in \nat \fun \nat
\end{theorem}

\begin{zproof}
invoke IWit;
apply FunIntro to predicate (\lambda x: \nat @ x) \in \nat \fun \nat;
prove by reduce;
with enabled (RelPow) prove by reduce;
\end{zproof}

\begin{theorem}{rule IWitVal}
  \forall x: \nat @ IWit~x = x
\end{theorem}

\begin{zproof}
prove;
use IWitFun;
use pairInFunction[\nat, \nat][f := IWit, y := x];
invoke IWit;
prove by reduce;
\end{zproof}

\begin{theorem}{IEx}
  \exists funI: \nat \fun \nat @ \forall x: \nat @ funI~x = x
\end{theorem}

\begin{zproof}
instantiate funI == IWit;
use IWitFun;
prove;
\end{zproof}

\begin{axdef}
  cst: \nat \fun (\nat \fun \nat)
\where
  \Label{rule CstVal} \forall n, x: \nat @ cst~n~x = n
\end{axdef}

\begin{zproof}[cst\$domainCheck]
prove;
\end{zproof}

\begin{zed}
  CstWit == (\lambda n: \nat @ (\lambda x: \nat @ n))
\end{zed}

\begin{theorem}{CstWitFun}
  CstWit \in \nat \fun (\nat \fun \nat)
\end{theorem}

\begin{zproof}
invoke CstWit;
apply FunIntro to predicate (\lambda n: \nat @ (\lambda x: \nat @ n)) \in \nat \fun (\nat \fun \nat);
prove by reduce;
with enabled (RelPow) prove by reduce;
\end{zproof}

\begin{theorem}{CstWitApp}
  \forall n: \nat @ CstWit~n = (\lambda x: \nat @ n)
\end{theorem}

\begin{zproof}
prove;
use CstWitFun;
use pairInFunction[\nat, \nat \fun \nat][f := CstWit, x := n, y := (\lambda x: \nat @ n)];
invoke CstWit;
prove by reduce;
\end{zproof}

\begin{theorem}{rule CstWitVal}
  \forall n, x: \nat @ CstWit~n~x = n
\end{theorem}

\begin{zproof}
prove;
use CstWitApp;
prove;
equality substitute CstWit~n;
prove;
\end{zproof}

\begin{theorem}{CstEx}
  \exists cst: \nat \fun (\nat \fun \nat) @ \forall n, x: \nat @ cst~n~x = n
\end{theorem}

\begin{zproof}
instantiate cst == CstWit;
use CstWitFun;
prove;
\end{zproof}

\subsection{Arithmetic functions}

\begin{axdef}
  funS: \nat \fun \nat
\where
  \Label{rule SVal} \forall x: \nat @ funS~x = x + 1
\end{axdef}

\begin{zproof}[funS\$domainCheck]
prove;
\end{zproof}

\begin{zed}
  SWit == (\lambda x: \nat @ x + 1)
\end{zed}

\begin{theorem}{SWitFun}
  SWit \in \nat \fun \nat
\end{theorem}

\begin{zproof}
invoke SWit;
apply FunIntro to predicate (\lambda x: \nat @ x + 1) \in \nat \fun \nat;
prove by reduce;
with enabled (RelPow) prove by reduce;
\end{zproof}

\begin{theorem}{rule SWitVal}
  \forall x: \nat @ SWit~x = x + 1
\end{theorem}

\begin{zproof}
prove;
use SWitFun;
use pairInFunction[\nat, \nat][f := SWit, x := x, y := x + 1];
invoke SWit;
prove by reduce;
\end{zproof}

\begin{theorem}{SEx}
  \exists funS: \nat \fun \nat @ \forall x: \nat @ funS~x = x + 1
\end{theorem}

\begin{zproof}
instantiate funS == SWit;
use SWitFun;
prove;
\end{zproof}

\begin{axdef}
  funD: \nat \fun \nat
\where
  \Label{rule DVal} \forall x: \nat @ funD~x = 2 * x
\end{axdef}

\begin{zproof}[funD\$domainCheck]
prove;
\end{zproof}

\begin{zed}
  DWit == (\lambda x: \nat @ 2 * x)
\end{zed}

\begin{theorem}{DWitFun}
  DWit \in \nat \fun \nat
\end{theorem}

\begin{zproof}
invoke DWit;
apply FunIntro to predicate (\lambda x: \nat @ 2 * x) \in \nat \fun \nat;
prove by reduce;
with enabled (RelPow) prove by reduce;
\end{zproof}

\begin{theorem}{rule DWitVal}
  \forall x: \nat @ DWit~x = 2 * x
\end{theorem}

\begin{zproof}
prove;
use DWitFun;
use pairInFunction[\nat, \nat][f := DWit, x := x, y := 2 * x];
invoke DWit;
prove by reduce;
\end{zproof}

\begin{theorem}{DEx}
  \exists funD: \nat \fun \nat @ \forall x: \nat @ funD~x = 2 * x
\end{theorem}

\begin{zproof}
instantiate funD == DWit;
use DWitFun;
prove;
\end{zproof}

\begin{axdef}
  funSq: \nat \fun \nat
\where
  \Label{rule SqVal} \forall x: \nat @ funSq~x = x * x
\end{axdef}

\begin{zproof}[funSq\$domainCheck]
prove;
\end{zproof}

\begin{zed}
  SqWit == (\lambda x: \nat @ x * x)
\end{zed}

\begin{theorem}{SqWitFun}
  SqWit \in \nat \fun \nat
\end{theorem}

\begin{zproof}
invoke SqWit;
apply FunIntro to predicate (\lambda x: \nat @ x * x) \in \nat \fun \nat;
prove by reduce;
with enabled (RelPow) prove by reduce;
\end{zproof}

\begin{theorem}{rule SqWitVal}
  \forall x: \nat @ SqWit~x = x * x
\end{theorem}

\begin{zproof}
prove;
use SqWitFun;
use pairInFunction[\nat, \nat][f := SqWit, x := x, y := x * x];
invoke SqWit;
prove by reduce;
\end{zproof}

\begin{theorem}{SqEx}
  \exists funSq: \nat \fun \nat @ \forall x: \nat @ funSq~x = x * x
\end{theorem}

\begin{zproof}
instantiate funSq == SqWit;
use SqWitFun;
prove;
\end{zproof}

\begin{axdef}
  funHf: \nat \fun \nat
\where
  \Label{rule HfVal} \forall x: \nat @ funHf~x = x \div 2
\end{axdef}

\begin{zproof}[funHf\$domainCheck]
prove;
\end{zproof}

\begin{zed}
  HfWit == (\lambda x: \nat @ x \div 2)
\end{zed}

\begin{zproof}[HfWit\$domainCheck]
prove;
\end{zproof}

\begin{theorem}{HfWitFun}
  HfWit \in \nat \fun \nat
\end{theorem}

\begin{zproof}
invoke HfWit;
apply FunIntro to predicate (\lambda x: \nat @ x \div 2) \in \nat \fun \nat;
prove by reduce;
with enabled (RelPow) prove by reduce;
\end{zproof}

\begin{theorem}{rule HfWitVal}
  \forall x: \nat @ HfWit~x = x \div 2
\end{theorem}

\begin{zproof}
prove;
use HfWitFun;
use pairInFunction[\nat, \nat][f := HfWit, x := x, y := x \div 2];
invoke HfWit;
prove by reduce;
\end{zproof}

\begin{theorem}{HfEx}
  \exists funHf: \nat \fun \nat @ \forall x: \nat @ funHf~x = x \div 2
\end{theorem}

\begin{zproof}
instantiate funHf == HfWit;
use HfWitFun;
prove;
\end{zproof}

\begin{axdef}
  funO: \nat \fun \nat
\where
  \Label{rule OVal} \forall x: \nat @ funO~x = \IF x = 0 \THEN 1 \ELSE 0
\end{axdef}

\begin{zproof}[funO\$domainCheck]
prove;
\end{zproof}

\begin{zed}
  OWit == (\lambda x: \nat @ 0) \oplus \{ 0 \mapsto 1 \}
\end{zed}

\begin{theorem}{OWitFun}
  OWit \in \nat \fun \nat
\end{theorem}

\begin{zproof}
invoke OWit;
prove by reduce;
\end{zproof}

\begin{theorem}{rule OWitVal}
  \forall x: \nat @ OWit~x = \IF x = 0 \THEN 1 \ELSE 0
\end{theorem}

\begin{zproof}
invoke OWit;
prove by reduce;
\end{zproof}

\begin{theorem}{OEx}
  \exists funO: \nat \fun \nat @ \forall x: \nat @ funO~x = \IF x = 0 \THEN 1 \ELSE 0
\end{theorem}

\begin{zproof}
instantiate funO == OWit;
use OWitFun;
prove;
\end{zproof}

\begin{axdef}
  funSgn: \nat \fun \nat
\where
  \Label{rule SgnVal} \forall x: \nat @ funSgn~x = \IF x = 0 \THEN 0 \ELSE 1
\end{axdef}

\begin{zproof}[funSgn\$domainCheck]
prove;
\end{zproof}

\begin{zed}
  SgnWit == (\lambda x: \nat @ 1) \oplus \{ 0 \mapsto 0 \}
\end{zed}

\begin{theorem}{SgnWitFun}
  SgnWit \in \nat \fun \nat
\end{theorem}

\begin{zproof}
invoke SgnWit;
prove by reduce;
\end{zproof}

\begin{theorem}{rule SgnWitVal}
  \forall x: \nat @ SgnWit~x = \IF x = 0 \THEN 0 \ELSE 1
\end{theorem}

\begin{zproof}
invoke SgnWit;
prove by reduce;
\end{zproof}

\begin{theorem}{SgnEx}
  \exists funSgn: \nat \fun \nat @ \forall x: \nat @ funSgn~x = \IF x = 0 \THEN 0 \ELSE 1
\end{theorem}

\begin{zproof}
instantiate funSgn == SgnWit;
use SgnWitFun;
prove;
\end{zproof}

\subsection{Integer square root and the excess over a square}

\begin{axdef}
  funRt: \nat \fun \nat
\where
  \Label{RtSpec} \forall x: \nat @ funRt~x * funRt~x \leq x \land x < (funRt~x + 1) * (funRt~x + 1)
\end{axdef}

\begin{zproof}[funRt\$domainCheck]
prove;
\end{zproof}

\begin{theorem}{RtUnique}
  \forall r, s, x: \nat | r * r \leq x \land x < (r + 1) * (r + 1)
     \land s * s \leq x \land x < (s + 1) * (s + 1) @ r = s
\end{theorem}

\begin{zproof}
prove;
split r < s;
cases;
use timesMonotonic2[a := r + 1, b := r + 1, c := s, d := s];
prove;
next;
split s < r;
cases;
use timesMonotonic2[a := s + 1, b := s + 1, c := r, d := r];
prove;
next;
prove;
next;
\end{zproof}

\begin{zed}
  RtExSet == \{ x: \nat | \exists r: \nat @ r * r \leq x \land x < (r + 1) * (r + 1) \}
\end{zed}

\begin{theorem}{rule RtExBase}
  0 \in RtExSet
\end{theorem}

\begin{zproof}
invoke RtExSet;
prove;
instantiate r == 0;
prove;
\end{zproof}

\begin{theorem}{rule RtExStep}
  \forall x: RtExSet @ 1 + x \in RtExSet
\end{theorem}

\begin{zproof}
invoke RtExSet;
prove;
split x + 1 < (r + 1) * (r + 1);
cases;
instantiate r\_\_0 == r;
prove;
next;
instantiate r\_\_1 == r + 1;
prove;
next;
\end{zproof}

\begin{theorem}{RtExSub}
  \nat \subseteq RtExSet
\end{theorem}

\begin{zproof}
apply natInduction to predicate \nat \subseteq RtExSet;
prove;
invoke RtExSet;
prove;
\end{zproof}

\begin{theorem}{RtExists}
  \forall x: \nat @ \exists r: \nat @ r * r \leq x \land x < (r + 1) * (r + 1)
\end{theorem}

\begin{zproof}
prove;
use RtExSub;
with enabled (inPower) prove;
instantiate e == x;
invoke RtExSet;
prove;
\end{zproof}

\begin{zed}
  RtWit == \{ x, r: \nat | r * r \leq x \land x < (r + 1) * (r + 1) @ (x, r) \}
\end{zed}

\begin{theorem}{RtWitFun}
  RtWit \in \nat \fun \nat
\end{theorem}

\begin{zproof}
invoke RtWit;
apply FunIntro to predicate \{ x, r: \nat | r * r \leq x \land x < (r + 1) * (r + 1) @ (x, r) \}
  \in \nat \fun \nat;
with enabled (RelPow) prove by reduce;
use RtExists[x := x];
use RtUnique[r := y, s := z, x := x\_\_0];
prove;
instantiate y\_\_0 == r;
prove;
\end{zproof}

\begin{theorem}{rule RtWitLe}
  \forall x: \nat @ RtWit~x * RtWit~x \leq x
\end{theorem}

\begin{zproof}
prove;
use RtExists[x := x];
use RtWitFun;
prove;
use pairInFunction[\nat, \nat][f := RtWit, x := x, y := r];
invoke RtWit;
prove by reduce;
\end{zproof}

\begin{theorem}{rule RtWitLt}
  \forall x: \nat @ x < (1 + RtWit~x) * (1 + RtWit~x)
\end{theorem}

\begin{zproof}
prove;
use RtExists[x := x];
use RtWitFun;
prove;
use pairInFunction[\nat, \nat][f := RtWit, x := x, y := r];
invoke RtWit;
prove by reduce;
\end{zproof}

\begin{theorem}{RtEx}
  \exists funRt: \nat \fun \nat @
    \forall x: \nat @ funRt~x * funRt~x \leq x \land x < (funRt~x + 1) * (funRt~x + 1)
\end{theorem}

\begin{zproof}
instantiate funRt == RtWit;
use RtWitFun;
prove;
\end{zproof}

$\mathrm{Rt}(x) = r$ whenever $r^2 \le x < (r+1)^2$.

\begin{theorem}{RtI}
  \forall r, x: \nat | r * r \leq x \land x < (r + 1) * (r + 1) @ funRt~x = r
\end{theorem}

\begin{zproof}
prove;
use RtSpec[x := x];
use RtUnique[r := funRt~x, s := r, x := x];
use UFNonneg[f := funRt, x := x];
prove;
\end{zproof}

\begin{axdef}
  funE: \nat \fun \nat
\where
  \Label{disabled rule EVal} \forall x: \nat @ funE~x = x - funRt~x * funRt~x
\end{axdef}

\begin{zproof}[funE\$domainCheck]
prove;
\end{zproof}

\begin{zed}
  EWit == (\lambda x: \nat @ x - funRt~x * funRt~x)
\end{zed}

\begin{zproof}[EWit\$domainCheck]
prove;
\end{zproof}

\begin{theorem}{EWitRel}
  (\lambda x: \nat @ x - funRt~x * funRt~x) \in \power (\nat \cross \nat)
\end{theorem}

\begin{zproof}
with enabled (inPower) prove by reduce;
use RtSpec[x := x];
prove;
\end{zproof}

\begin{theorem}{EWitFun}
  EWit \in \nat \fun \nat
\end{theorem}

\begin{zproof}
invoke EWit;
use EWitRel;
apply FunIntro to predicate (\lambda x: \nat @ x - funRt~x * funRt~x) \in \nat \fun \nat;
with enabled (RelPow) prove by reduce;
use RtSpec[x := x];
prove;
\end{zproof}

\begin{theorem}{rule EWitVal}
  \forall x: \nat @ EWit~x = x - funRt~x * funRt~x
\end{theorem}

\begin{zproof}
prove;
use EWitFun;
use pairInFunction[\nat, \nat][f := EWit, x := x, y := x - funRt~x * funRt~x];
invoke EWit;
prove by reduce;
\end{zproof}

\begin{theorem}{EEx}
  \exists funE: \nat \fun \nat @ \forall x: \nat @ funE~x = x - funRt~x * funRt~x
\end{theorem}

\begin{zproof}
instantiate funE == EWit;
use EWitFun;
prove;
\end{zproof}

\subsection{Operators}

$F^+(x) = F(x) + x$.

\begin{axdef}
  opPlusI: (\nat \fun \nat) \fun (\nat \fun \nat)
\where
  \Label{rule PlusIVal} \forall F: \nat \fun \nat; x: \nat @ opPlusI~F~x = F~x + x
\end{axdef}

\begin{zproof}[opPlusI\$domainCheck]
prove;
\end{zproof}

\begin{theorem}{rule PlusIInner}
  \forall F: \nat \fun \nat @ (\lambda x: \nat @ F~x + x) \in \nat \fun \nat
\end{theorem}

\begin{zproof}
apply FunIntro to predicate (\lambda x: \nat @ F~x + x) \in \nat \fun \nat;
with enabled (RelPow, inPower) prove by reduce;
\end{zproof}

\begin{zed}
  PlusIWit == (\lambda F: \nat \fun \nat @ (\lambda x: \nat @ F~x + x))
\end{zed}

\begin{zproof}[PlusIWit\$domainCheck]
prove;
\end{zproof}

\begin{theorem}{PlusIWitFun}
  PlusIWit \in (\nat \fun \nat) \fun (\nat \fun \nat)
\end{theorem}

\begin{zproof}
invoke PlusIWit;
apply FunIntro to predicate (\lambda F: \nat \fun \nat @ (\lambda x: \nat @ F~x + x)) \in (\nat \fun \nat) \fun (\nat \fun \nat);
with enabled (RelPow, inPower) prove by reduce;
\end{zproof}

\begin{theorem}{PlusIWitApp}
  \forall F: \nat \fun \nat @ PlusIWit~F = (\lambda x: \nat @ F~x + x)
\end{theorem}

\begin{zproof}
use PlusIWitFun;
use pairInFunction[\nat \fun \nat, \nat \fun \nat][f := PlusIWit, x := F, y := (\lambda x: \nat @ F~x + x)];
invoke PlusIWit;
prove by reduce;
\end{zproof}

\begin{theorem}{rule PlusIWitVal}
  \forall F: \nat \fun \nat; x: \nat @ PlusIWit~F~x = F~x + x
\end{theorem}

\begin{zproof}
prove;
use PlusIWitApp[F := F];
use pairInFunction[\nat, \nat][f := (\lambda x: \nat @ F~x + x), x := x, y := F~x + x];
prove;
\end{zproof}

\begin{theorem}{PlusIEx}
  \exists opPlusI: (\nat \fun \nat) \fun (\nat \fun \nat) @ \forall F: \nat \fun \nat; x: \nat @ opPlusI~F~x = F~x + x
\end{theorem}

\begin{zproof}
instantiate opPlusI == PlusIWit;
use PlusIWitFun;
prove;
\end{zproof}

Binary operators are functions on pairs of one-argument functions: $F + G$ and the composition $FG$.

\begin{axdef}
  opPlus: (\nat \fun \nat) \cross (\nat \fun \nat) \fun (\nat \fun \nat)
\where
  \Label{rule PlusVal} \forall F, G: \nat \fun \nat; x: \nat @ opPlus(F, G)~x = F~x + G~x
\end{axdef}

\begin{zproof}[opPlus\$domainCheck]
prove;
\end{zproof}

\begin{theorem}{rule PlusInner}
  \forall F, G: \nat \fun \nat @ (\lambda x: \nat @ F~x + G~x) \in \nat \fun \nat
\end{theorem}

\begin{zproof}
apply FunIntro to predicate (\lambda x: \nat @ F~x + G~x) \in \nat \fun \nat;
with enabled (RelPow, inPower) prove by reduce;
\end{zproof}

\begin{zed}
  PlusWit == (\lambda F, G: \nat \fun \nat @ (\lambda x: \nat @ F~x + G~x))
\end{zed}

\begin{zproof}[PlusWit\$domainCheck]
prove;
\end{zproof}

\begin{theorem}{PlusWitFun}
  PlusWit \in (\nat \fun \nat) \cross (\nat \fun \nat) \fun (\nat \fun \nat)
\end{theorem}

\begin{zproof}
invoke PlusWit;
apply FunIntro to predicate (\lambda F, G: \nat \fun \nat @ (\lambda x: \nat @ F~x + G~x)) \in (\nat \fun \nat) \cross (\nat \fun \nat) \fun (\nat \fun \nat);
with enabled (RelPow, inPower) prove by reduce;
instantiate F == x.1, G == x.2;
prove;
\end{zproof}

\begin{theorem}{PlusWitApp}
  \forall F, G: \nat \fun \nat @ PlusWit(F, G) = (\lambda x: \nat @ F~x + G~x)
\end{theorem}

\begin{zproof}
use PlusWitFun;
use pairInFunction[(\nat \fun \nat) \cross (\nat \fun \nat), \nat \fun \nat][f := PlusWit, x := (F, G), y := (\lambda x: \nat @ F~x + G~x)];
invoke PlusWit;
prove by reduce;
\end{zproof}

\begin{theorem}{rule PlusWitVal}
  \forall F, G: \nat \fun \nat; x: \nat @ PlusWit(F, G)~x = F~x + G~x
\end{theorem}

\begin{zproof}
prove;
use PlusWitApp[F := F, G := G];
use pairInFunction[\nat, \nat][f := (\lambda x: \nat @ F~x + G~x), x := x, y := F~x + G~x];
prove;
\end{zproof}

\begin{theorem}{PlusEx}
  \exists opPlus: (\nat \fun \nat) \cross (\nat \fun \nat) \fun (\nat \fun \nat) @ \forall F, G: \nat \fun \nat; x: \nat @ opPlus(F, G)~x = F~x + G~x
\end{theorem}

\begin{zproof}
instantiate opPlus == PlusWit;
use PlusWitFun;
prove;
\end{zproof}

\begin{axdef}
  comp: (\nat \fun \nat) \cross (\nat \fun \nat) \fun (\nat \fun \nat)
\where
  \Label{rule CompVal} \forall F, G: \nat \fun \nat; x: \nat @ comp(F, G)~x = F~(G~x)
\end{axdef}

\begin{zproof}[comp\$domainCheck]
prove;
\end{zproof}

\begin{theorem}{rule CompInner}
  \forall F, G: \nat \fun \nat @ (\lambda x: \nat @ F~(G~x)) \in \nat \fun \nat
\end{theorem}

\begin{zproof}
apply FunIntro to predicate (\lambda x: \nat @ F~(G~x)) \in \nat \fun \nat;
with enabled (RelPow, inPower) prove by reduce;
\end{zproof}

\begin{zed}
  CompWit == (\lambda F, G: \nat \fun \nat @ (\lambda x: \nat @ F~(G~x)))
\end{zed}

\begin{zproof}[CompWit\$domainCheck]
prove;
\end{zproof}

\begin{theorem}{CompWitFun}
  CompWit \in (\nat \fun \nat) \cross (\nat \fun \nat) \fun (\nat \fun \nat)
\end{theorem}

\begin{zproof}
invoke CompWit;
apply FunIntro to predicate (\lambda F, G: \nat \fun \nat @ (\lambda x: \nat @ F~(G~x))) \in (\nat \fun \nat) \cross (\nat \fun \nat) \fun (\nat \fun \nat);
with enabled (RelPow, inPower) prove by reduce;
instantiate F == x.1, G == x.2;
prove;
\end{zproof}

\begin{theorem}{CompWitApp}
  \forall F, G: \nat \fun \nat @ CompWit(F, G) = (\lambda x: \nat @ F~(G~x))
\end{theorem}

\begin{zproof}
use CompWitFun;
use pairInFunction[(\nat \fun \nat) \cross (\nat \fun \nat), \nat \fun \nat][f := CompWit, x := (F, G), y := (\lambda x: \nat @ F~(G~x))];
invoke CompWit;
prove by reduce;
\end{zproof}

\begin{theorem}{rule CompWitVal}
  \forall F, G: \nat \fun \nat; x: \nat @ CompWit(F, G)~x = F~(G~x)
\end{theorem}

\begin{zproof}
prove;
use CompWitApp[F := F, G := G];
use pairInFunction[\nat, \nat][f := (\lambda x: \nat @ F~(G~x)), x := x, y := F~(G~x)];
prove;
\end{zproof}

\begin{theorem}{CompEx}
  \exists comp: (\nat \fun \nat) \cross (\nat \fun \nat) \fun (\nat \fun \nat) @ \forall F, G: \nat \fun \nat; x: \nat @ comp(F, G)~x = F~(G~x)
\end{theorem}

\begin{zproof}
instantiate comp == CompWit;
use CompWitFun;
prove;
\end{zproof}

Iteration from $a$: $F^{\square(a)}(0) = a$ and $F^{\square(a)}(x+1) = F(F^{\square(a)}(x))$
(in the defining equations Z/EVES's normal form $1 + x$ is used for $x + 1$).

\begin{axdef}
  opBoxA: \nat \fun ((\nat \fun \nat) \fun (\nat \fun \nat))
\where
  \Label{rule BoxA0} \forall a: \nat; F: \nat \fun \nat @ opBoxA~a~F~0 = a \\
  \Label{rule BoxASuc} \forall a: \nat; F: \nat \fun \nat; x: \nat @ opBoxA~a~F~(1 + x) = F~(opBoxA~a~F~x)
\end{axdef}

\begin{zproof}[opBoxA\$domainCheck]
prove;
\end{zproof}

\begin{theorem}{rule BoxAInner}
  \forall a: \nat; F: \nat \fun \nat @ (\lambda n: \nat @ iter[\nat]~n~F~a) \in \nat \fun \nat
\end{theorem}

\begin{zproof}
apply FunIntro to predicate (\lambda n: \nat @ iter[\nat]~n~F~a) \in \nat \fun \nat;
with enabled (RelPow, inPower) prove by reduce;
\end{zproof}

\begin{theorem}{rule BoxAMid}
  \forall a: \nat @ (\lambda F: \nat \fun \nat @ (\lambda n: \nat @ iter[\nat]~n~F~a)) \in (\nat \fun \nat) \fun (\nat \fun \nat)
\end{theorem}

\begin{zproof}
apply FunIntro to predicate (\lambda F: \nat \fun \nat @ (\lambda n: \nat @ iter[\nat]~n~F~a)) \in (\nat \fun \nat) \fun (\nat \fun \nat);
with enabled (RelPow, inPower) prove by reduce;
\end{zproof}

\begin{zed}
  BoxAWit == (\lambda a: \nat @ (\lambda F: \nat \fun \nat @ (\lambda n: \nat @ iter[\nat]~n~F~a)))
\end{zed}

\begin{zproof}[BoxAWit\$domainCheck]
prove;
\end{zproof}

\begin{theorem}{BoxAWitFun}
  BoxAWit \in \nat \fun ((\nat \fun \nat) \fun (\nat \fun \nat))
\end{theorem}

\begin{zproof}
invoke BoxAWit;
apply FunIntro to predicate (\lambda a: \nat @ (\lambda F: \nat \fun \nat @ (\lambda n: \nat @ iter[\nat]~n~F~a))) \in \nat \fun ((\nat \fun \nat) \fun (\nat \fun \nat));
with enabled (RelPow, inPower) prove by reduce;
\end{zproof}

\begin{theorem}{BoxAWitApp}
  \forall a: \nat @ BoxAWit~a = (\lambda F: \nat \fun \nat @ (\lambda n: \nat @ iter[\nat]~n~F~a))
\end{theorem}

\begin{zproof}
use BoxAWitFun;
use pairInFunction[\nat, (\nat \fun \nat) \fun (\nat \fun \nat)][f := BoxAWit, x := a, y := (\lambda F: \nat \fun \nat @ (\lambda n: \nat @ iter[\nat]~n~F~a))];
invoke BoxAWit;
prove by reduce;
\end{zproof}

\begin{theorem}{BoxAWitApp2}
  \forall a: \nat; F: \nat \fun \nat @ BoxAWit~a~F = (\lambda n: \nat @ iter[\nat]~n~F~a)
\end{theorem}

\begin{zproof}
prove;
use BoxAWitApp[a := a];
use pairInFunction[\nat \fun \nat, \nat \fun \nat][f := (\lambda F: \nat \fun \nat @ (\lambda n: \nat @ iter[\nat]~n~F~a)), x := F, y := (\lambda n: \nat @ iter[\nat]~n~F~a)];
prove;
\end{zproof}

\begin{theorem}{rule BoxAWitVal}
  \forall a: \nat; F: \nat \fun \nat; n: \nat @ BoxAWit~a~F~n = iter[\nat]~n~F~a
\end{theorem}

\begin{zproof}
prove;
use BoxAWitApp2[a := a, F := F];
use pairInFunction[\nat, \nat][f := (\lambda n: \nat @ iter[\nat]~n~F~a), x := n, y := iter[\nat]~n~F~a];
prove;
\end{zproof}

\begin{theorem}{BoxAEx}
  \exists opBoxA: \nat \fun ((\nat \fun \nat) \fun (\nat \fun \nat)) @
    (\forall a: \nat; F: \nat \fun \nat @ opBoxA~a~F~0 = a)
    \land (\forall a: \nat; F: \nat \fun \nat; x: \nat @ opBoxA~a~F~(1 + x) = F~(opBoxA~a~F~x))
\end{theorem}

\begin{zproof}
instantiate opBoxA == BoxAWit;
use BoxAWitFun;
prove;
\end{zproof}

Pure iteration $F^\square$, i.e.\ iteration from $0$ (its witness is $F^{\square(0)}$).

\begin{axdef}
  opBox: (\nat \fun \nat) \fun (\nat \fun \nat)
\where
  \Label{rule Box0} \forall F: \nat \fun \nat @ opBox~F~0 = 0 \\
  \Label{rule BoxSuc} \forall F: \nat \fun \nat; x: \nat @ opBox~F~(1 + x) = F~(opBox~F~x)
\end{axdef}

\begin{zproof}[opBox\$domainCheck]
prove;
\end{zproof}

\begin{theorem}{BoxEx}
  \exists opBox: (\nat \fun \nat) \fun (\nat \fun \nat) @
    (\forall F: \nat \fun \nat @ opBox~F~0 = 0)
    \land (\forall F: \nat \fun \nat; x: \nat @ opBox~F~(1 + x) = F~(opBox~F~x))
\end{theorem}

\begin{zproof}
instantiate opBox == opBoxA~0;
prove;
\end{zproof}

\subsection{Bases and closures}

A basis is given by a set $I$ of initial functions, a set $U$ of unary operators and a set
$B$ of binary operators; its closure is the least set containing $I$ and closed under $U$
and $B$, i.e.\ the intersection of all such sets.  First, the sets closed under $U$ and
those closed under $B$.

\begin{axdef}
  UClosed: \power ((\nat \fun \nat) \fun (\nat \fun \nat)) \fun \power (\power (\nat \fun \nat))
\where
  \Label{disabled rule UClosedDef} \forall U: \power ((\nat \fun \nat) \fun (\nat \fun \nat)) @ UClosed~U = \{ X: \power (\nat \fun \nat) | \forall u: U; F: X @ u~F \in X \}
\end{axdef}

\begin{zproof}[UClosed\$domainCheck]
prove;
\end{zproof}

\begin{zed}
  UClosedWit == (\lambda U: \power ((\nat \fun \nat) \fun (\nat \fun \nat)) @ \{ X: \power (\nat \fun \nat) | \forall u: U; F: X @ u~F \in X \})
\end{zed}

\begin{zproof}[UClosedWit\$domainCheck]
prove;
\end{zproof}

\begin{theorem}{UClosedWitFun}
  UClosedWit \in \power ((\nat \fun \nat) \fun (\nat \fun \nat)) \fun \power (\power (\nat \fun \nat))
\end{theorem}

\begin{zproof}
invoke UClosedWit;
apply FunIntro to predicate (\lambda U: \power ((\nat \fun \nat) \fun (\nat \fun \nat)) @ \{ X: \power (\nat \fun \nat) | \forall u: U; F: X @ u~F \in X \}) \in \power ((\nat \fun \nat) \fun (\nat \fun \nat)) \fun \power (\power (\nat \fun \nat));
with enabled (RelPow, inPower) prove by reduce;
\end{zproof}

\begin[disabled]{theorem}{rule UClosedWitVal}
  \forall U: \power ((\nat \fun \nat) \fun (\nat \fun \nat)) @ UClosedWit~U = \{ X: \power (\nat \fun \nat) | \forall u: U; F: X @ u~F \in X \}
\end{theorem}

\begin{zproof}
use UClosedWitFun;
use pairInFunction[\power ((\nat \fun \nat) \fun (\nat \fun \nat)), \power (\power (\nat \fun \nat))][f := UClosedWit, x := U, y := \{ X: \power (\nat \fun \nat) | \forall u: U; F: X @ u~F \in X \}];
invoke UClosedWit;
prove by reduce;
\end{zproof}

\begin{theorem}{UClosedEx}
  \exists UClosed: \power ((\nat \fun \nat) \fun (\nat \fun \nat)) \fun \power (\power (\nat \fun \nat)) @ \forall U: \power ((\nat \fun \nat) \fun (\nat \fun \nat)) @ UClosed~U = \{ X: \power (\nat \fun \nat) | \forall u: U; F: X @ u~F \in X \}
\end{theorem}

\begin{zproof}
instantiate UClosed == UClosedWit;
use UClosedWitFun;
with enabled (UClosedWitVal) prove;
\end{zproof}



\begin{axdef}
  BClosed: \power ((\nat \fun \nat) \cross (\nat \fun \nat) \fun (\nat \fun \nat)) \fun \power (\power (\nat \fun \nat))
\where
  \Label{disabled rule BClosedDef} \forall B: \power ((\nat \fun \nat) \cross (\nat \fun \nat) \fun (\nat \fun \nat)) @ BClosed~B = \{ X: \power (\nat \fun \nat) | \forall b: B; F, G: X @ b(F, G) \in X \}
\end{axdef}

\begin{zproof}[BClosed\$domainCheck]
prove;
\end{zproof}

\begin{zed}
  BClosedWit == (\lambda B: \power ((\nat \fun \nat) \cross (\nat \fun \nat) \fun (\nat \fun \nat)) @ \{ X: \power (\nat \fun \nat) | \forall b: B; F, G: X @ b(F, G) \in X \})
\end{zed}

\begin{zproof}[BClosedWit\$domainCheck]
prove;
\end{zproof}

\begin{theorem}{BClosedWitFun}
  BClosedWit \in \power ((\nat \fun \nat) \cross (\nat \fun \nat) \fun (\nat \fun \nat)) \fun \power (\power (\nat \fun \nat))
\end{theorem}

\begin{zproof}
invoke BClosedWit;
apply FunIntro to predicate (\lambda B: \power ((\nat \fun \nat) \cross (\nat \fun \nat) \fun (\nat \fun \nat)) @ \{ X: \power (\nat \fun \nat) | \forall b: B; F, G: X @ b(F, G) \in X \}) \in \power ((\nat \fun \nat) \cross (\nat \fun \nat) \fun (\nat \fun \nat)) \fun \power (\power (\nat \fun \nat));
with enabled (RelPow, inPower) prove by reduce;
\end{zproof}

\begin[disabled]{theorem}{rule BClosedWitVal}
  \forall B: \power ((\nat \fun \nat) \cross (\nat \fun \nat) \fun (\nat \fun \nat)) @ BClosedWit~B = \{ X: \power (\nat \fun \nat) | \forall b: B; F, G: X @ b(F, G) \in X \}
\end{theorem}

\begin{zproof}
use BClosedWitFun;
use pairInFunction[\power ((\nat \fun \nat) \cross (\nat \fun \nat) \fun (\nat \fun \nat)), \power (\power (\nat \fun \nat))][f := BClosedWit, x := B, y := \{ X: \power (\nat \fun \nat) | \forall b: B; F, G: X @ b(F, G) \in X \}];
invoke BClosedWit;
prove by reduce;
\end{zproof}

\begin{theorem}{BClosedEx}
  \exists BClosed: \power ((\nat \fun \nat) \cross (\nat \fun \nat) \fun (\nat \fun \nat)) \fun \power (\power (\nat \fun \nat)) @ \forall B: \power ((\nat \fun \nat) \cross (\nat \fun \nat) \fun (\nat \fun \nat)) @ BClosed~B = \{ X: \power (\nat \fun \nat) | \forall b: B; F, G: X @ b(F, G) \in X \}
\end{theorem}

\begin{zproof}
instantiate BClosed == BClosedWit;
use BClosedWitFun;
with enabled (BClosedWitVal) prove;
\end{zproof}

Domains and applications of the curried closure function.

\begin{theorem}{rule ClosApp1}
  \forall c: \power (\nat \fun \nat) \fun (\power ((\nat \fun \nat) \fun (\nat \fun \nat)) \fun (\power ((\nat \fun \nat) \cross (\nat \fun \nat) \fun (\nat \fun \nat)) \fun \power (\nat \fun \nat))); I: \power (\nat \fun \nat) @ c~I \in \power ((\nat \fun \nat) \fun (\nat \fun \nat)) \fun (\power ((\nat \fun \nat) \cross (\nat \fun \nat) \fun (\nat \fun \nat)) \fun \power (\nat \fun \nat))
\end{theorem}

\begin{zproof}
prove;
\end{zproof}

\begin{theorem}{rule ClosApp2}
  \forall g: \power ((\nat \fun \nat) \fun (\nat \fun \nat)) \fun (\power ((\nat \fun \nat) \cross (\nat \fun \nat) \fun (\nat \fun \nat)) \fun \power (\nat \fun \nat)); U: \power ((\nat \fun \nat) \fun (\nat \fun \nat)) @ g~U \in \power ((\nat \fun \nat) \cross (\nat \fun \nat) \fun (\nat \fun \nat)) \fun \power (\nat \fun \nat)
\end{theorem}

\begin{zproof}
prove;
\end{zproof}

\begin{theorem}{rule ClosDom2}
  \forall g: \power ((\nat \fun \nat) \fun (\nat \fun \nat)) \fun (\power ((\nat \fun \nat) \cross (\nat \fun \nat) \fun (\nat \fun \nat)) \fun \power (\nat \fun \nat)) @ \dom g = \power ((\nat \fun \nat) \fun (\nat \fun \nat))
\end{theorem}

\begin{zproof}
prove;
\end{zproof}

\begin{theorem}{rule ClosApplies3}
  \forall h: \power ((\nat \fun \nat) \cross (\nat \fun \nat) \fun (\nat \fun \nat)) \fun \power (\nat \fun \nat); B: \power ((\nat \fun \nat) \cross (\nat \fun \nat) \fun (\nat \fun \nat)) @ h applies\$to B
\end{theorem}

\begin{zproof}
prove;
\end{zproof}

The closure, as a curried function of $I$, $U$ and $B$.

\begin{axdef}
  clos: \power (\nat \fun \nat) \fun (\power ((\nat \fun \nat) \fun (\nat \fun \nat)) \fun (\power ((\nat \fun \nat) \cross (\nat \fun \nat) \fun (\nat \fun \nat)) \fun \power (\nat \fun \nat)))
\where
  \Label{ClosDef} \forall I: \power (\nat \fun \nat); U: \power ((\nat \fun \nat) \fun (\nat \fun \nat)); B: \power ((\nat \fun \nat) \cross (\nat \fun \nat) \fun (\nat \fun \nat)) @ clos~I~U~B = \bigcap \{ X: UClosed~U \cap BClosed~B | I \subseteq X \}
\end{axdef}

\begin{zproof}[clos\$domainCheck]
prove;
\end{zproof}

Consistency: the intersection is a set of one-argument functions, because the set of all
one-argument functions belongs to the family.

\begin{theorem}{rule ClosFamUF}
  \forall I: \power (\nat \fun \nat); U: \power ((\nat \fun \nat) \fun (\nat \fun \nat)); B: \power ((\nat \fun \nat) \cross (\nat \fun \nat) \fun (\nat \fun \nat)) @
    \nat \fun \nat \in \{ X: UClosed~U \cap BClosed~B | I \subseteq X \}
\end{theorem}

\begin{zproof}
with enabled (UClosedDef, BClosedDef) prove by reduce;
\end{zproof}

\begin{theorem}{rule ClosInPower}
  \forall I: \power (\nat \fun \nat); U: \power ((\nat \fun \nat) \fun (\nat \fun \nat)); B: \power ((\nat \fun \nat) \cross (\nat \fun \nat) \fun (\nat \fun \nat)) @ \bigcap \{ X: UClosed~U \cap BClosed~B | I \subseteq X \} \in \power (\nat \fun \nat)
\end{theorem}

\begin{zproof}
prove;
with enabled (inPower, inBigcap) prove;
\end{zproof}

\begin{theorem}{rule ClosL3}
  \forall I: \power (\nat \fun \nat); U: \power ((\nat \fun \nat) \fun (\nat \fun \nat)) @ (\lambda B: \power ((\nat \fun \nat) \cross (\nat \fun \nat) \fun (\nat \fun \nat)) @ \bigcap \{ X: UClosed~U \cap BClosed~B | I \subseteq X \}) \in \power ((\nat \fun \nat) \cross (\nat \fun \nat) \fun (\nat \fun \nat)) \fun \power (\nat \fun \nat)
\end{theorem}

\begin{zproof}
apply FunIntro to predicate (\lambda B: \power ((\nat \fun \nat) \cross (\nat \fun \nat) \fun (\nat \fun \nat)) @ \bigcap \{ X: UClosed~U \cap BClosed~B | I \subseteq X \}) \in \power ((\nat \fun \nat) \cross (\nat \fun \nat) \fun (\nat \fun \nat)) \fun \power (\nat \fun \nat);
with enabled (RelPow, inPower) prove;
\end{zproof}

\begin{theorem}{rule ClosL2}
  \forall I: \power (\nat \fun \nat) @ (\lambda U: \power ((\nat \fun \nat) \fun (\nat \fun \nat)) @ (\lambda B: \power ((\nat \fun \nat) \cross (\nat \fun \nat) \fun (\nat \fun \nat)) @ \bigcap \{ X: UClosed~U \cap BClosed~B | I \subseteq X \})) \in \power ((\nat \fun \nat) \fun (\nat \fun \nat)) \fun (\power ((\nat \fun \nat) \cross (\nat \fun \nat) \fun (\nat \fun \nat)) \fun \power (\nat \fun \nat))
\end{theorem}

\begin{zproof}
apply FunIntro to predicate (\lambda U: \power ((\nat \fun \nat) \fun (\nat \fun \nat)) @ (\lambda B: \power ((\nat \fun \nat) \cross (\nat \fun \nat) \fun (\nat \fun \nat)) @ \bigcap \{ X: UClosed~U \cap BClosed~B | I \subseteq X \})) \in \power ((\nat \fun \nat) \fun (\nat \fun \nat)) \fun (\power ((\nat \fun \nat) \cross (\nat \fun \nat) \fun (\nat \fun \nat)) \fun \power (\nat \fun \nat));
with enabled (RelPow, inPower) prove;
\end{zproof}

\begin{zed}
  ClosWit == (\lambda I: \power (\nat \fun \nat) @ (\lambda U: \power ((\nat \fun \nat) \fun (\nat \fun \nat)) @ (\lambda B: \power ((\nat \fun \nat) \cross (\nat \fun \nat) \fun (\nat \fun \nat)) @ \bigcap \{ X: UClosed~U \cap BClosed~B | I \subseteq X \})))
\end{zed}

\begin{zproof}[ClosWit\$domainCheck]
prove;
\end{zproof}

\begin{theorem}{ClosWitFun}
  ClosWit \in \power (\nat \fun \nat) \fun (\power ((\nat \fun \nat) \fun (\nat \fun \nat)) \fun (\power ((\nat \fun \nat) \cross (\nat \fun \nat) \fun (\nat \fun \nat)) \fun \power (\nat \fun \nat)))
\end{theorem}

\begin{zproof}
invoke ClosWit;
apply FunIntro to predicate (\lambda I: \power (\nat \fun \nat) @ (\lambda U: \power ((\nat \fun \nat) \fun (\nat \fun \nat)) @ (\lambda B: \power ((\nat \fun \nat) \cross (\nat \fun \nat) \fun (\nat \fun \nat)) @ \bigcap \{ X: UClosed~U \cap BClosed~B | I \subseteq X \}))) \in \power (\nat \fun \nat) \fun (\power ((\nat \fun \nat) \fun (\nat \fun \nat)) \fun (\power ((\nat \fun \nat) \cross (\nat \fun \nat) \fun (\nat \fun \nat)) \fun \power (\nat \fun \nat)));
with enabled (RelPow, inPower) prove;
\end{zproof}

\begin{theorem}{ClosWitApp1}
  \forall I: \power (\nat \fun \nat) @ ClosWit~I = (\lambda U: \power ((\nat \fun \nat) \fun (\nat \fun \nat)) @ (\lambda B: \power ((\nat \fun \nat) \cross (\nat \fun \nat) \fun (\nat \fun \nat)) @ \bigcap \{ X: UClosed~U \cap BClosed~B | I \subseteq X \}))
\end{theorem}

\begin{zproof}
use ClosWitFun;
use pairInFunction[\power (\nat \fun \nat), \power ((\nat \fun \nat) \fun (\nat \fun \nat)) \fun (\power ((\nat \fun \nat) \cross (\nat \fun \nat) \fun (\nat \fun \nat)) \fun \power (\nat \fun \nat))][f := ClosWit, x := I, y := (\lambda U: \power ((\nat \fun \nat) \fun (\nat \fun \nat)) @ (\lambda B: \power ((\nat \fun \nat) \cross (\nat \fun \nat) \fun (\nat \fun \nat)) @ \bigcap \{ X: UClosed~U \cap BClosed~B | I \subseteq X \}))];
invoke ClosWit;
prove;
\end{zproof}

\begin{theorem}{ClosWitApp2}
  \forall I: \power (\nat \fun \nat); U: \power ((\nat \fun \nat) \fun (\nat \fun \nat)) @ ClosWit~I~U = (\lambda B: \power ((\nat \fun \nat) \cross (\nat \fun \nat) \fun (\nat \fun \nat)) @ \bigcap \{ X: UClosed~U \cap BClosed~B | I \subseteq X \})
\end{theorem}

\begin{zproof}
use ClosWitApp1[I := I];
use pairInFunction[\power ((\nat \fun \nat) \fun (\nat \fun \nat)), \power ((\nat \fun \nat) \cross (\nat \fun \nat) \fun (\nat \fun \nat)) \fun \power (\nat \fun \nat)][f := (\lambda U: \power ((\nat \fun \nat) \fun (\nat \fun \nat)) @ (\lambda B: \power ((\nat \fun \nat) \cross (\nat \fun \nat) \fun (\nat \fun \nat)) @ \bigcap \{ X: UClosed~U \cap BClosed~B | I \subseteq X \})), x := U, y := (\lambda B: \power ((\nat \fun \nat) \cross (\nat \fun \nat) \fun (\nat \fun \nat)) @ \bigcap \{ X: UClosed~U \cap BClosed~B | I \subseteq X \})];
prove;
\end{zproof}

\begin{theorem}{rule ClosWitVal}
  \forall I: \power (\nat \fun \nat); U: \power ((\nat \fun \nat) \fun (\nat \fun \nat)); B: \power ((\nat \fun \nat) \cross (\nat \fun \nat) \fun (\nat \fun \nat)) @ ClosWit~I~U~B = \bigcap \{ X: UClosed~U \cap BClosed~B | I \subseteq X \}
\end{theorem}

\begin{zproof}
use ClosWitApp2[I := I, U := U];
use pairInFunction[\power ((\nat \fun \nat) \cross (\nat \fun \nat) \fun (\nat \fun \nat)), \power (\nat \fun \nat)][f := (\lambda B: \power ((\nat \fun \nat) \cross (\nat \fun \nat) \fun (\nat \fun \nat)) @ \bigcap \{ X: UClosed~U \cap BClosed~B | I \subseteq X \}), x := B, y := \bigcap \{ X: UClosed~U \cap BClosed~B | I \subseteq X \}];
prove;
\end{zproof}

\begin{theorem}{ClosEx}
  \exists clos: \power (\nat \fun \nat) \fun (\power ((\nat \fun \nat) \fun (\nat \fun \nat)) \fun (\power ((\nat \fun \nat) \cross (\nat \fun \nat) \fun (\nat \fun \nat)) \fun \power (\nat \fun \nat))) @ \forall I: \power (\nat \fun \nat); U: \power ((\nat \fun \nat) \fun (\nat \fun \nat)); B: \power ((\nat \fun \nat) \cross (\nat \fun \nat) \fun (\nat \fun \nat)) @ clos~I~U~B = \bigcap \{ X: UClosed~U \cap BClosed~B | I \subseteq X \}
\end{theorem}

\begin{zproof}
instantiate clos == ClosWit;
use ClosWitFun;
prove;
\end{zproof}

\subsection{The closure rules}

\begin{theorem}{rule ClosInit}
  \forall I: \power (\nat \fun \nat); U: \power ((\nat \fun \nat) \fun (\nat \fun \nat)); B: \power ((\nat \fun \nat) \cross (\nat \fun \nat) \fun (\nat \fun \nat)); F: \nat \fun \nat | F \in I @ F \in clos~I~U~B
\end{theorem}

\begin{zproof}
prove;
use ClosDef[I := I, U := U, B := B];
prove;
with enabled (inBigcap, subsetDef, inPower) prove;
instantiate e\_\_2 == F;
prove;
\end{zproof}

\begin[disabled]{theorem}{rule ClosMem}
  \forall I: \power (\nat \fun \nat); U: \power ((\nat \fun \nat) \fun (\nat \fun \nat)); B: \power ((\nat \fun \nat) \cross (\nat \fun \nat) \fun (\nat \fun \nat)); F: \power (\num \cross \num) @
    F \in clos~I~U~B \iff (\forall Y: \{ X: UClosed~U \cap BClosed~B | I \subseteq X \} @ F \in Y)
\end{theorem}

\begin{zproof}
use ClosDef[I := I, U := U, B := B];
prove;
\end{zproof}

\begin{theorem}{rule ClosUnop}
  \forall I: \power (\nat \fun \nat); U: \power ((\nat \fun \nat) \fun (\nat \fun \nat)); B: \power ((\nat \fun \nat) \cross (\nat \fun \nat) \fun (\nat \fun \nat)); u: (\nat \fun \nat) \fun (\nat \fun \nat); F: \nat \fun \nat | u \in U \land F \in clos~I~U~B @
    u~F \in clos~I~U~B
\end{theorem}

\begin{zproof}
with enabled (ClosMem) prove;
instantiate Y\_\_0 == Y;
with enabled (UClosedDef) prove;
instantiate u\_\_0 == u, F\_\_0 == F;
prove;
\end{zproof}

\begin{theorem}{rule ClosBinop}
  \forall I: \power (\nat \fun \nat); U: \power ((\nat \fun \nat) \fun (\nat \fun \nat)); B: \power ((\nat \fun \nat) \cross (\nat \fun \nat) \fun (\nat \fun \nat)); b: (\nat \fun \nat) \cross (\nat \fun \nat) \fun (\nat \fun \nat); F, G: \nat \fun \nat |
    b \in B \land F \in clos~I~U~B \land G \in clos~I~U~B @ b(F, G) \in clos~I~U~B
\end{theorem}

\begin{zproof}
with enabled (ClosMem) prove;
instantiate Y\_\_0 == Y, Y\_\_1 == Y;
with enabled (BClosedDef) prove;
instantiate b\_\_0 == b, F\_\_0 == F, G\_\_0 == G;
prove;
\end{zproof}

The closure is contained in every closed set containing $I$.

\begin{theorem}{ClosInPow}
  \forall I: \power (\nat \fun \nat); U: \power ((\nat \fun \nat) \fun (\nat \fun \nat)); B: \power ((\nat \fun \nat) \cross (\nat \fun \nat) \fun (\nat \fun \nat)) @ clos~I~U~B \in \power (\nat \fun \nat)
\end{theorem}

\begin{zproof}
use applyInRanFun[\power ((\nat \fun \nat) \cross (\nat \fun \nat) \fun (\nat \fun \nat)), \power (\nat \fun \nat)][f := clos~I~U, a := B];
prove;
\end{zproof}

\begin{theorem}{ClosElemUF}
  \forall I: \power (\nat \fun \nat); U: \power ((\nat \fun \nat) \fun (\nat \fun \nat)); B: \power ((\nat \fun \nat) \cross (\nat \fun \nat) \fun (\nat \fun \nat)); F: \power (\num \cross \num) | F \in clos~I~U~B @ F \in \nat \fun \nat
\end{theorem}

\begin{zproof}
use applyInRanFun[\power ((\nat \fun \nat) \cross (\nat \fun \nat) \fun (\nat \fun \nat)), \power (\nat \fun \nat)][f := clos~I~U, a := B];
prove;
with enabled (inPower) prove;
\end{zproof}

\begin{theorem}{ClosLeast}
  \forall I: \power (\nat \fun \nat); U: \power ((\nat \fun \nat) \fun (\nat \fun \nat)); B: \power ((\nat \fun \nat) \cross (\nat \fun \nat) \fun (\nat \fun \nat)); Z: \power (\nat \fun \nat) |
    Z \in UClosed~U \land Z \in BClosed~B \land I \subseteq Z @ clos~I~U~B \subseteq Z
\end{theorem}

\begin{zproof}
use ClosDef[I := I, U := U, B := B];
use BigcapSub[\nat \fun \nat][T := \{ X: UClosed~U \cap BClosed~B | I \subseteq X \}, Y := Z];
prove;
\end{zproof}

\subsection{Arithmetic facts about Rt and E used silently in the paper}

\begin{theorem}{RtSqAdd}
  \forall m, k: \nat | k \leq 2 * m @ funRt~(m * m + k) = m
\end{theorem}

\begin{zproof}
prove;
use RtI[r := m, x := m * m + k];
prove;
\end{zproof}

\begin{theorem}{ESqAdd}
  \forall m, k: \nat | k \leq 2 * m @ funE~(m * m + k) = k
\end{theorem}

\begin{zproof}
prove;
use RtSqAdd[m := m, k := k];
with enabled (EVal) prove;
\end{zproof}

\begin{theorem}{rule RtSq}
  \forall m: \nat @ funRt~(m * m) = m
\end{theorem}

\begin{zproof}
prove;
use RtSqAdd[m := m, k := 0];
prove;
\end{zproof}

\begin{theorem}{ELe}
  \forall x: \nat @ funE~x \leq x
\end{theorem}

\begin{zproof}
prove;
use RtSpec[x := x];
with enabled (EVal) prove;
\end{zproof}

\begin{theorem}{rule E0}
  funE~0 = 0
\end{theorem}

\begin{zproof}
use ESqAdd[m := 0, k := 0];
prove;
\end{zproof}

\begin{theorem}{rule E1}
  funE~1 = 0
\end{theorem}

\begin{zproof}
use ESqAdd[m := 1, k := 0];
prove;
\end{zproof}

\begin{theorem}{rule E2}
  funE~2 = 1
\end{theorem}

\begin{zproof}
use ESqAdd[m := 1, k := 1];
prove;
\end{zproof}

\begin{theorem}{rule E3}
  funE~3 = 2
\end{theorem}

\begin{zproof}
use ESqAdd[m := 1, k := 2];
prove;
\end{zproof}

\begin{theorem}{rule E4}
  funE~4 = 0
\end{theorem}

\begin{zproof}
use ESqAdd[m := 2, k := 0];
prove;
\end{zproof}


\section{Section 6: iteration and the unary operator $F^+$}

The closure of the basis $\mathcal{F} = \langle \bar 1, E, F^+, FG, F^\square\rangle$.

\begin{zed}
  C6 == clos~\{ cst~1, funE \}~\{ opPlusI, opBox \}~\{ comp \}
\end{zed}

\begin{zproof}[C6\$domainCheck]
prove;
\end{zproof}

\begin[disabled]{theorem}{rule C6Def}
  C6 = clos~\{ cst~1, funE \}~\{ opPlusI, opBox \}~\{ comp \}
\end{theorem}

\begin{zproof}
invoke C6;
prove;
\end{zproof}

\begin[disabled]{theorem}{rule C6UF}
  \forall F: \power (\num \cross \num) | F \in C6 @ F \in \nat \fun \nat
\end{theorem}

\begin{zproof}
prove;
use ClosElemUF[I := \{ cst~1, funE \}, U := \{ opPlusI, opBox \}, B := \{ comp \}, F := F];
with enabled (C6Def) prove;
\end{zproof}

\begin{theorem}{rule C6One}
  cst~1 \in C6
\end{theorem}

\begin{zproof}
with enabled (C6Def) prove;
\end{zproof}

\begin{theorem}{rule C6E}
  funE \in C6
\end{theorem}

\begin{zproof}
with enabled (C6Def) prove;
\end{zproof}

\begin{theorem}{rule C6PlusI}
  \forall F: \nat \fun \nat | F \in C6 @ opPlusI~F \in C6
\end{theorem}

\begin{zproof}
with enabled (C6Def) prove;
\end{zproof}

\begin{theorem}{rule C6Box}
  \forall F: \nat \fun \nat | F \in C6 @ opBox~F \in C6
\end{theorem}

\begin{zproof}
with enabled (C6Def) prove;
\end{zproof}

\begin{theorem}{rule C6Comp}
  \forall F, G: \nat \fun \nat | F \in C6 \land G \in C6 @ comp(F, G) \in C6
\end{theorem}

\begin{zproof}
with enabled (C6Def) prove;
\end{zproof}



\subsection{Auxiliary functions of Section 6}

\begin{axdef}
  funM3: \nat \fun \nat
\where
  \Label{rule M3Val} \forall x: \nat @ funM3~x = 3 * x
\end{axdef}

\begin{zproof}[funM3\$domainCheck]
prove;
\end{zproof}

\begin{zed}
  M3Wit == (\lambda x: \nat @ 3 * x)
\end{zed}

\begin{theorem}{M3WitFun}
  M3Wit \in \nat \fun \nat
\end{theorem}

\begin{zproof}
invoke M3Wit;
apply FunIntro to predicate (\lambda x: \nat @ 3 * x) \in \nat \fun \nat;
prove by reduce;
with enabled (RelPow) prove by reduce;
\end{zproof}

\begin{theorem}{rule M3WitVal}
  \forall x: \nat @ M3Wit~x = 3 * x
\end{theorem}

\begin{zproof}
prove;
use M3WitFun;
use pairInFunction[\nat, \nat][f := M3Wit, x := x, y := 3 * x];
invoke M3Wit;
prove by reduce;
\end{zproof}

\begin{theorem}{M3Ex}
  \exists funM3: \nat \fun \nat @ \forall x: \nat @ funM3~x = 3 * x
\end{theorem}

\begin{zproof}
instantiate funM3 == M3Wit;
use M3WitFun;
prove;
\end{zproof}

$R(x) = x + 2\,\mathrm{Rt}(x)$, as in Lemma 5.7.

\begin{axdef}
  funR: \nat \fun \nat
\where
  \Label{rule RVal} \forall x: \nat @ funR~x = x + 2 * funRt~x
\end{axdef}

\begin{zproof}[funR\$domainCheck]
prove;
\end{zproof}

\begin{zed}
  RWit == (\lambda x: \nat @ x + 2 * funRt~x)
\end{zed}

\begin{zproof}[RWit\$domainCheck]
prove;
\end{zproof}

\begin{theorem}{RWitFun}
  RWit \in \nat \fun \nat
\end{theorem}

\begin{zproof}
invoke RWit;
apply FunIntro to predicate (\lambda x: \nat @ x + 2 * funRt~x) \in \nat \fun \nat;
prove by reduce;
with enabled (RelPow) prove by reduce;
\end{zproof}

\begin{theorem}{rule RWitVal}
  \forall x: \nat @ RWit~x = x + 2 * funRt~x
\end{theorem}

\begin{zproof}
prove;
use RWitFun;
use pairInFunction[\nat, \nat][f := RWit, x := x, y := x + 2 * funRt~x];
invoke RWit;
prove by reduce;
\end{zproof}

\begin{theorem}{REx}
  \exists funR: \nat \fun \nat @ \forall x: \nat @ funR~x = x + 2 * funRt~x
\end{theorem}

\begin{zproof}
instantiate funR == RWit;
use RWitFun;
prove;
\end{zproof}

The excess $E(x)$ vanishes exactly on the squares.

\begin{theorem}{ESquareIf}
  \forall k: \nat @ funE~(k * k) = 0
\end{theorem}

\begin{zproof}
prove;
use ESqAdd[m := k, k := 0];
prove;
\end{zproof}

\begin{theorem}{ESquareOnly}
  \forall x: \nat | funE~x = 0 @ x = funRt~x * funRt~x
\end{theorem}

\begin{zproof}
with enabled (EVal) prove;
\end{zproof}

\begin{theorem}{ESquare}
  \forall x: \nat @ funE~x = 0 \iff (\exists k: \nat @ x = k * k)
\end{theorem}

\begin{zproof}
prove;
split funE~x = 0;
cases;
use ESquareOnly[x := x];
prove;
instantiate k == funRt~x;
prove;
next;
prove;
use ESquareIf[k := k];
prove;
next;
\end{zproof}

\begin{theorem}{rule OEVal}
  \forall x: \nat @ funO~(funE~x) = \IF (\exists k: \nat @ x = k * k) \THEN 1 \ELSE 0
\end{theorem}

\begin{zproof}
prove;
use ESquare[x := x];
prove;
instantiate k\_\_1 == k;
prove;
\end{zproof}

$Q$ is the characteristic function of the squares.

\begin{axdef}
  funQ: \nat \fun \nat
\where
  \Label{rule QVal} \forall x: \nat @ funQ~x = \IF (\exists k: \nat @ x = k * k) \THEN 1 \ELSE 0
\end{axdef}

\begin{zproof}[funQ\$domainCheck]
prove;
\end{zproof}

\begin{theorem}{QEx}
  \exists funQ: \nat \fun \nat @
    \forall x: \nat @ funQ~x = \IF (\exists k: \nat @ x = k * k) \THEN 1 \ELSE 0
\end{theorem}

\begin{zproof}
instantiate funQ == comp(funO, funE);
prove;
\end{zproof}

\subsection{First functions in $\mathrm{clos}\,\mathcal{F}$}

\begin{theorem}{SExpr}
  opPlusI~(cst~1) = funS
\end{theorem}

\begin{zproof}
apply UFExt to predicate opPlusI~(cst~1) = funS;
prove;
\end{zproof}

\begin{theorem}{ZeroExpr}
  comp(funE, cst~1) = cst~0
\end{theorem}

\begin{zproof}
apply UFExt to predicate comp(funE, cst~1) = cst~0;
prove;
\end{zproof}

\begin{theorem}{DExpr}
  opPlusI~(opPlusI~(cst~0)) = funD
\end{theorem}

\begin{zproof}
apply UFExt to predicate opPlusI~(opPlusI~(cst~0)) = funD;
prove;
\end{zproof}

\begin{theorem}{M3Expr}
  opPlusI~funD = funM3
\end{theorem}

\begin{zproof}
apply UFExt to predicate opPlusI~funD = funM3;
prove;
\end{zproof}

\begin{theorem}{OExpr}
  comp(comp(comp(funE, funD), funS), funSgn) = funO
\end{theorem}

\begin{zproof}
apply UFExt to predicate comp(comp(comp(funE, funD), funS), funSgn) = funO;
prove;
\end{zproof}

\begin{theorem}{QExpr}
  comp(funO, funE) = funQ
\end{theorem}

\begin{zproof}
apply UFExt to predicate comp(funO, funE) = funQ;
prove;
\end{zproof}

Robinson's difference formula: if $b \le a$ then $a - b = E((a+b)^2 + 3a + b + 1)$.

\begin{theorem}{RobDiffE}
  \forall a, b: \nat | b \leq a @ funE~((a + b) * (a + b) + 3 * a + b + 1) = a - b
\end{theorem}

\begin{zproof}
prove;
use ESqAdd[m := a + b + 1, k := a - b];
prove;
\end{zproof}

\subsection{Division with remainder}

\begin{theorem}{DivLe}
  \forall x, d, q: \nat | 0 < d \land x < d * q + d @ x \div d \leq q
\end{theorem}

\begin{zproof}
prove;
\end{zproof}

\begin{theorem}{DivGe}
  \forall x, d, q: \nat | 0 < d \land d * q \leq x @ q \leq x \div d
\end{theorem}

\begin{zproof}
prove;
\end{zproof}

\begin{theorem}{DivMod}
  \forall x, d, q, r: \nat | 0 < d \land r < d \land x = d * q + r @ x \div d = q \land x \mod d = r
\end{theorem}

\begin{zproof}
prove;
use DivLe[x := x, d := d, q := q];
use DivGe[x := x, d := d, q := q];
use divModRelation[x := x, y := d];
with disabled (divLowerBound, divUpperBound) prove;
\end{zproof}

$x \bmod 3$.

\begin{axdef}
  funMod3: \nat \fun \nat
\where
  \Label{rule Mod3Val} \forall x: \nat @ funMod3~x = x \mod 3
\end{axdef}

\begin{zproof}[funMod3\$domainCheck]
prove;
\end{zproof}

\begin{zed}
  Mod3Wit == (\lambda x: \nat @ x \mod 3)
\end{zed}

\begin{zproof}[Mod3Wit\$domainCheck]
prove;
\end{zproof}

\begin{theorem}{Mod3WitFun}
  Mod3Wit \in \nat \fun \nat
\end{theorem}

\begin{zproof}
invoke Mod3Wit;
apply FunIntro to predicate (\lambda x: \nat @ x \mod 3) \in \nat \fun \nat;
prove by reduce;
with enabled (RelPow) prove by reduce;
\end{zproof}

\begin{theorem}{rule Mod3WitVal}
  \forall x: \nat @ Mod3Wit~x = x \mod 3
\end{theorem}

\begin{zproof}
prove;
use Mod3WitFun;
use pairInFunction[\nat, \nat][f := Mod3Wit, x := x, y := x \mod 3];
invoke Mod3Wit;
prove by reduce;
\end{zproof}

\begin{theorem}{Mod3Ex}
  \exists funMod3: \nat \fun \nat @ \forall x: \nat @ funMod3~x = x \mod 3
\end{theorem}

\begin{zproof}
instantiate funMod3 == Mod3Wit;
use Mod3WitFun;
prove;
\end{zproof}

\begin{theorem}{DivModOf}
  \forall x, d: \nat | 0 < d @ x = d * (x \div d) + x \mod d \land 0 \leq x \mod d \land x \mod d < d
\end{theorem}

\begin{zproof}
prove;
\end{zproof}

\begin{theorem}{Mod3Cases}
  \forall n: \nat @ (n \mod 3 = 0 \land (1 + n) \mod 3 = 1) \lor (n \mod 3 = 1 \land (1 + n) \mod 3 = 2)
    \lor (n \mod 3 = 2 \land (1 + n) \mod 3 = 0)
\end{theorem}

\begin{zproof}
prove;
use DivModOf[x := n, d := 3];
split n \mod 3 = 2;
cases;
use DivMod[x := 1 + n, d := 3, q := n \div 3 + 1, r := 0];
prove;
next;
use DivMod[x := 1 + n, d := 3, q := n \div 3, r := n \mod 3 + 1];
prove;
next;
\end{zproof}

\subsection{First functions in $\mathrm{clos}\,\mathcal{F}$, continued}

\begin{theorem}{SgnExpr}
  opBox~(cst~1) = funSgn
\end{theorem}

\begin{zproof}
apply UFExt to predicate opBox~(cst~1) = funSgn;
prove;
split x = 0;
cases;
prove;
next;
use BoxSuc[F\_\_0 := cst~1, x := x - 1];
prove;
next;
\end{zproof}

\begin{zed}
  Mod3Set == \{ n: \nat | opBox~(comp(comp(funE, funS), funS))~n = n \mod 3 \}
\end{zed}

\begin{zproof}[Mod3Set\$domainCheck]
prove;
\end{zproof}

\begin{theorem}{rule Mod3Base}
  0 \in Mod3Set
\end{theorem}

\begin{zproof}
invoke Mod3Set;
prove;
\end{zproof}

\begin{theorem}{rule Mod3Step}
  \forall n: Mod3Set @ 1 + n \in Mod3Set
\end{theorem}

\begin{zproof}
invoke Mod3Set;
prove;
use Mod3Cases[n := n];
split n \mod 3 = 0;
cases;
prove;
next;
split n \mod 3 = 1;
cases;
prove;
next;
prove;
next;
\end{zproof}

\begin{theorem}{Mod3Sub}
  \nat \subseteq Mod3Set
\end{theorem}

\begin{zproof}
apply natInduction to predicate \nat \subseteq Mod3Set;
prove;
invoke Mod3Set;
prove;
\end{zproof}

\begin{theorem}{Mod3All}
  \forall n: \nat @ opBox~(comp(comp(funE, funS), funS))~n = n \mod 3
\end{theorem}

\begin{zproof}
prove;
use Mod3Sub;
invoke Mod3Set;
with enabled (inPower) prove;
instantiate e == n;
prove;
\end{zproof}

\begin{theorem}{Mod3Expr}
  opBox~(comp(comp(funE, funS), funS)) = funMod3
\end{theorem}

\begin{zproof}
apply UFExt to predicate opBox~(comp(comp(funE, funS), funS)) = funMod3;
prove;
use Mod3All[n := x];
prove;
\end{zproof}

\subsection{Membership in $\mathrm{clos}\,\mathcal{F}$}

\begin{theorem}{SInExpr}
  opPlusI~(cst~1) \in C6
\end{theorem}

\begin{zproof}
prove;
\end{zproof}

\begin{theorem}{rule SIn}
  funS \in C6
\end{theorem}

\begin{zproof}
use SInExpr;
use SExpr;
prove;
\end{zproof}

\begin{theorem}{SgnInExpr}
  opBox~(cst~1) \in C6
\end{theorem}

\begin{zproof}
prove;
\end{zproof}

\begin{theorem}{rule SgnIn}
  funSgn \in C6
\end{theorem}

\begin{zproof}
use SgnInExpr;
use SgnExpr;
prove;
\end{zproof}

\begin{theorem}{ZeroInExpr}
  comp(funE, cst~1) \in C6
\end{theorem}

\begin{zproof}
prove;
\end{zproof}

\begin{theorem}{rule ZeroIn}
  cst~0 \in C6
\end{theorem}

\begin{zproof}
use ZeroInExpr;
use ZeroExpr;
prove;
\end{zproof}

\begin{theorem}{DInExpr}
  opPlusI~(opPlusI~(cst~0)) \in C6
\end{theorem}

\begin{zproof}
prove;
\end{zproof}

\begin{theorem}{rule DIn}
  funD \in C6
\end{theorem}

\begin{zproof}
use DInExpr;
use DExpr;
prove;
\end{zproof}

\begin{theorem}{M3InExpr}
  opPlusI~funD \in C6
\end{theorem}

\begin{zproof}
prove;
\end{zproof}

\begin{theorem}{rule M3In}
  funM3 \in C6
\end{theorem}

\begin{zproof}
use M3InExpr;
use M3Expr;
prove;
\end{zproof}

\begin{theorem}{OInExpr}
  comp(comp(comp(funE, funD), funS), funSgn) \in C6
\end{theorem}

\begin{zproof}
prove;
\end{zproof}

\begin{theorem}{rule OIn}
  funO \in C6
\end{theorem}

\begin{zproof}
use OInExpr;
use OExpr;
prove;
\end{zproof}

\begin{theorem}{QInExpr}
  comp(funO, funE) \in C6
\end{theorem}

\begin{zproof}
prove;
\end{zproof}

\begin{theorem}{rule QIn}
  funQ \in C6
\end{theorem}

\begin{zproof}
use QInExpr;
use QExpr;
prove;
\end{zproof}

\begin{theorem}{Mod3InExpr}
  opBox~(comp(comp(funE, funS), funS)) \in C6
\end{theorem}

\begin{zproof}
prove;
\end{zproof}

\begin{theorem}{rule Mod3In}
  funMod3 \in C6
\end{theorem}

\begin{zproof}
use Mod3InExpr;
use Mod3Expr;
prove;
\end{zproof}

\subsection{$R(x) = x + 2\,\mathrm{Rt}(x)$ and $\mathrm{Sq}$ (Lemma 5.7 technique)}

\begin{theorem}{QSq}
  \forall x, k: \nat | x = k * k @ funQ~x = 1
\end{theorem}

\begin{zproof}
prove;
instantiate k\_\_0 == k;
prove;
\end{zproof}

\begin{theorem}{NotSquare}
  \forall m, x: \nat | m * m < x \land x < (m + 1) * (m + 1) @ \lnot (\exists k: \nat @ x = k * k)
\end{theorem}

\begin{zproof}
prove;
split k \leq m;
cases;
use timesMonotonic2[a := k, b := k, c := m, d := m];
prove;
next;
use timesMonotonic2[a := m + 1, b := m + 1, c := k, d := k];
prove;
next;
\end{zproof}

\begin{theorem}{QNonSq}
  \forall m, x: \nat | m * m < x \land x < (m + 1) * (m + 1) @ funQ~x = 0
\end{theorem}

\begin{zproof}
prove;
use NotSquare[m := m, x := x];
prove;
instantiate k\_\_0 == k;
prove;
\end{zproof}

\begin{theorem}{RStep}
  \forall n: \nat @ opPlusI~(comp(comp(comp(comp(funD, funQ), funS), funS), funS))~(funS~(funR~n)) = funR~(1 + n)
\end{theorem}

\begin{zproof}
prove;
use RtSpec[x := n];
split (funRt~n + 1) * (funRt~n + 1) = n + 1;
cases;
use RtI[r := funRt~n + 1, x := n + 1];
prove;
instantiate k == 2 + funRt~n;
prove;
next;
use RtI[r := funRt~n, x := n + 1];
use QNonSq[m := funRt~n + 1, x := 4 + n + 2 * funRt~n];
prove;
next;
\end{zproof}


\begin{zed}
  RBoxSet == \{ n: \nat | opBox~(comp(opPlusI~(comp(comp(comp(comp(funD, funQ), funS), funS), funS)), funS))~n = funR~n \}
\end{zed}

\begin{zproof}[RBoxSet\$domainCheck]
prove;
\end{zproof}

\begin{theorem}{rule RBoxBase}
  0 \in RBoxSet
\end{theorem}

\begin{zproof}
invoke RBoxSet;
prove;
use RtSq[m := 0];
prove;
\end{zproof}

\begin{theorem}{rule RBoxStep}
  \forall n: RBoxSet @ 1 + n \in RBoxSet
\end{theorem}

\begin{zproof}
invoke RBoxSet;
prove;
use RStep[n := n];
prove;
\end{zproof}

\begin{theorem}{RBoxSub}
  \nat \subseteq RBoxSet
\end{theorem}

\begin{zproof}
apply natInduction to predicate \nat \subseteq RBoxSet;
prove;
invoke RBoxSet;
prove;
\end{zproof}

\begin{theorem}{RBoxAll}
  \forall n: \nat @ opBox~(comp(opPlusI~(comp(comp(comp(comp(funD, funQ), funS), funS), funS)), funS))~n = funR~n
\end{theorem}

\begin{zproof}
prove;
use RBoxSub;
invoke RBoxSet;
with enabled (inPower) prove;
instantiate e == n;
prove;
\end{zproof}

\begin{theorem}{RExpr}
  opBox~(comp(opPlusI~(comp(comp(comp(comp(funD, funQ), funS), funS), funS)), funS)) = funR
\end{theorem}

\begin{zproof}
apply UFExt to predicate opBox~(comp(opPlusI~(comp(comp(comp(comp(funD, funQ), funS), funS), funS)), funS)) = funR;
prove;
use RBoxAll[n := x];
prove;
\end{zproof}

\begin{zed}
  SqBoxSet == \{ n: \nat | opBox~(comp(funS, funR))~n = n * n \}
\end{zed}

\begin{zproof}[SqBoxSet\$domainCheck]
prove;
\end{zproof}

\begin{theorem}{rule SqBoxBase}
  0 \in SqBoxSet
\end{theorem}

\begin{zproof}
invoke SqBoxSet;
prove;
\end{zproof}

\begin{theorem}{rule SqBoxStep}
  \forall n: SqBoxSet @ 1 + n \in SqBoxSet
\end{theorem}

\begin{zproof}
invoke SqBoxSet;
prove;
equality substitute opBox~(comp(funS, funR))~n;
rewrite;
\end{zproof}

\begin{theorem}{SqBoxSub}
  \nat \subseteq SqBoxSet
\end{theorem}

\begin{zproof}
apply natInduction to predicate \nat \subseteq SqBoxSet;
prove;
invoke SqBoxSet;
prove;
\end{zproof}

\begin{theorem}{SqBoxAll}
  \forall n: \nat @ opBox~(comp(funS, funR))~n = n * n
\end{theorem}

\begin{zproof}
prove;
use SqBoxSub;
invoke SqBoxSet;
with enabled (inPower) prove;
instantiate e == n;
prove;
\end{zproof}

\begin{theorem}{SqExpr}
  opBox~(comp(funS, funR)) = funSq
\end{theorem}

\begin{zproof}
apply UFExt to predicate opBox~(comp(funS, funR)) = funSq;
prove;
use SqBoxAll[n := x];
prove;
\end{zproof}

\begin{theorem}{RInExpr}
  opBox~(comp(opPlusI~(comp(comp(comp(comp(funD, funQ), funS), funS), funS)), funS)) \in C6
\end{theorem}

\begin{zproof}
prove;
\end{zproof}

\begin{theorem}{rule RIn}
  funR \in C6
\end{theorem}

\begin{zproof}
use RInExpr;
use RExpr;
prove;
\end{zproof}

\begin{theorem}{SqInExpr}
  opBox~(comp(funS, funR)) \in C6
\end{theorem}

\begin{zproof}
prove;
\end{zproof}

\begin{theorem}{rule SqIn}
  funSq \in C6
\end{theorem}

\begin{zproof}
use SqInExpr;
use SqExpr;
prove;
\end{zproof}

Inflationary functions ($F(x) \ge x$ for all $x$), and pairs with $G(x)^2 \le F(x)$.

\begin{zed}
  Infl == \{ F: \nat \fun \nat | \forall x: \nat @ x \leq F~x \}
\end{zed}

\begin{zproof}[Infl\$domainCheck]
prove;
\end{zproof}

\begin{theorem}{rule InflLe}
  \forall F: \nat \fun \nat; x: \nat | F \in Infl @ x \leq F~x
\end{theorem}

\begin{zproof}
prove;
invoke Infl;
prove;
instantiate x\_\_0 == x;
prove;
\end{zproof}

\begin{zed}
  SqBelow == \{ F, G: \nat \fun \nat | \forall x: \nat @ G~x * G~x \leq F~x \}
\end{zed}

\begin{zproof}[SqBelow\$domainCheck]
prove;
\end{zproof}

\begin{theorem}{rule SqBelowLe}
  \forall F, G: \nat \fun \nat; x: \nat | (F, G) \in SqBelow @ G~x * G~x \leq F~x
\end{theorem}

\begin{zproof}
prove;
invoke SqBelow;
prove;
instantiate x\_\_0 == x;
prove;
\end{zproof}

\subsection{Lemma 6.2: $F^- = E\,S\,(M_3 F E)^+ (\mathrm{Sq}\,F^+)^+$}

Definition 6.1 defines $F^-(x) = F(x) - x$, meaningful when $F(x) \ge x$; Lemmas 6.2 and 6.4
are stated accordingly: the closure contains a function with those values.

\begin{zed}
  exprMinusI == (\lambda F: \nat \fun \nat @ comp(comp(comp(funE, funS), opPlusI~(comp(comp(funM3, F), funE))), opPlusI~(comp(funSq, opPlusI~F))))
\end{zed}

\begin{zproof}[exprMinusI\$domainCheck]
prove;
\end{zproof}

\begin{theorem}{rule exprMinusIFun}
  exprMinusI \in (\nat \fun \nat) \fun (\nat \fun \nat)
\end{theorem}

\begin{zproof}
invoke exprMinusI;
apply FunIntro to predicate (\lambda F: \nat \fun \nat @ comp(comp(comp(funE, funS), opPlusI~(comp(comp(funM3, F), funE))), opPlusI~(comp(funSq, opPlusI~F)))) \in (\nat \fun \nat) \fun (\nat \fun \nat);
with enabled (RelPow, inPower) prove by reduce;
\end{zproof}

\begin[disabled]{theorem}{rule exprMinusIApp}
  \forall F: \nat \fun \nat @ exprMinusI~F = comp(comp(comp(funE, funS), opPlusI~(comp(comp(funM3, F), funE))), opPlusI~(comp(funSq, opPlusI~F)))
\end{theorem}

\begin{zproof}
use exprMinusIFun;
use pairInFunction[\nat \fun \nat, \nat \fun \nat][f := exprMinusI, x := F, y := comp(comp(comp(funE, funS), opPlusI~(comp(comp(funM3, F), funE))), opPlusI~(comp(funSq, opPlusI~F)))];
invoke exprMinusI;
prove by reduce;
\end{zproof}

\begin{theorem}{rule MinusIIn}
  \forall F: \nat \fun \nat | F \in C6 @ exprMinusI~F \in C6
\end{theorem}

\begin{zproof}
with enabled (exprMinusIApp) prove;
\end{zproof}

\begin{theorem}{rule MinusIVal}
  \forall F: \nat \fun \nat; x: \nat | x \leq F~x @ exprMinusI~F~x = F~x - x
\end{theorem}

\begin{zproof}
with enabled (exprMinusIApp) prove;
use ESqAdd[m := F~x + x, k := x];
use RobDiffE[a := F~x, b := x];
prove;
\end{zproof}

\begin{theorem}{Lemma62}
  \forall F: \nat \fun \nat | F \in C6 \land F \in Infl @
    \exists G: \nat \fun \nat | G \in C6 @ \forall x: \nat @ G~x = F~x - x
\end{theorem}

\begin{zproof}
prove;
instantiate G == exprMinusI~F;
prove;
\end{zproof}

\subsection{Lemma 6.3: $\mathrm{Hf} = ((S\,\mathrm{Mod}_3^+)^\square)^-$ and $\mathrm{Rt} = \mathrm{Hf}\,R^-$}

\begin{zed}
  HfAuxSet == \{ q: \nat | opBox~(comp(funS, opPlusI~funMod3))~(2 * q) = 3 * q \land opBox~(comp(funS, opPlusI~funMod3))~(1 + 2 * q) = 1 + 3 * q \}
\end{zed}

\begin{zproof}[HfAuxSet\$domainCheck]
prove;
\end{zproof}

\begin{theorem}{rule HfAuxBase}
  0 \in HfAuxSet
\end{theorem}

\begin{zproof}
invoke HfAuxSet;
prove;
use DivMod[x := 0, d := 3, q := 0, r := 0];
prove;
\end{zproof}

\begin{theorem}{rule HfAuxStep}
  \forall q: HfAuxSet @ 1 + q \in HfAuxSet
\end{theorem}

\begin{zproof}
invoke HfAuxSet;
prove;
use BoxSuc[F\_\_0 := comp(funS, opPlusI~funMod3), x := 1 + 2 * q];
use BoxSuc[F\_\_0 := comp(funS, opPlusI~funMod3), x := 2 + 2 * q];
use DivMod[x := 1 + 3 * q, d := 3, q := q, r := 1];
use DivMod[x := 3 + 3 * q, d := 3, q := q + 1, r := 0];
prove;
\end{zproof}

\begin{theorem}{HfAuxSub}
  \nat \subseteq HfAuxSet
\end{theorem}

\begin{zproof}
apply natInduction to predicate \nat \subseteq HfAuxSet;
prove;
invoke HfAuxSet;
prove;
\end{zproof}

\begin{theorem}{HfAuxAll}
  \forall q: \nat @ opBox~(comp(funS, opPlusI~funMod3))~(2 * q) = 3 * q \land opBox~(comp(funS, opPlusI~funMod3))~(1 + 2 * q) = 1 + 3 * q
\end{theorem}

\begin{zproof}
prove;
use HfAuxSub;
invoke HfAuxSet;
with enabled (inPower) prove;
instantiate e == q;
prove;
\end{zproof}

\begin{theorem}{HfBoxVal}
  \forall x: \nat @ opBox~(comp(funS, opPlusI~funMod3))~x = x + x \div 2
\end{theorem}

\begin{zproof}
prove;
use DivModOf[x := x, d := 2];
use HfAuxAll[q := x \div 2];
split x \mod 2 = 0;
cases;
prove;
next;
prove;
next;
\end{zproof}

\begin{theorem}{HfExpr}
  exprMinusI~(opBox~(comp(funS, opPlusI~funMod3))) = funHf
\end{theorem}

\begin{zproof}
apply UFExt to predicate exprMinusI~(opBox~(comp(funS, opPlusI~funMod3))) = funHf;
prove;
use HfBoxVal[x := x];
use MinusIVal[F := opBox~(comp(funS, opPlusI~funMod3)), x := x];
prove;
\end{zproof}

\begin{theorem}{RtExpr}
  comp(funHf, exprMinusI~funR) = funRt
\end{theorem}

\begin{zproof}
apply UFExt to predicate comp(funHf, exprMinusI~funR) = funRt;
prove;
use MinusIVal[F := funR, x := x];
use DivMod[x := 2 * funRt~x, d := 2, q := funRt~x, r := 0];
use UFNonneg[f := funRt, x := x];
prove;
\end{zproof}

\begin{theorem}{HfInExpr}
  exprMinusI~(opBox~(comp(funS, opPlusI~funMod3))) \in C6
\end{theorem}

\begin{zproof}
prove;
\end{zproof}

\begin{theorem}{rule HfIn}
  funHf \in C6
\end{theorem}

\begin{zproof}
use HfInExpr;
use HfExpr;
prove;
\end{zproof}
\begin{theorem}{RtInExpr}
  comp(funHf, exprMinusI~funR) \in C6
\end{theorem}

\begin{zproof}
prove;
\end{zproof}

\begin{theorem}{rule RtIn}
  funRt \in C6
\end{theorem}

\begin{zproof}
use RtInExpr;
use RtExpr;
prove;
\end{zproof}

\subsection{Lemma 6.4: $F \oplus G = (F^- E)^+ (\mathrm{Sq}\,G)^+$ and $F \ominus G = (F^+ E)^- (\mathrm{Sq}\,G)^+$}

$(F \oplus G)(x) = F(x) + G(x)^2$ (Definition 6.1); $F \ominus G$, meaningful when $F(x) \ge G(x)^2$, is treated like $F^-$.

\begin{axdef}
  opOplus: (\nat \fun \nat) \cross (\nat \fun \nat) \fun (\nat \fun \nat)
\where
  \Label{rule OplusVal} \forall F, G: \nat \fun \nat; x: \nat @ opOplus(F, G)~x = F~x + G~x * G~x
\end{axdef}

\begin{zproof}[opOplus\$domainCheck]
prove;
\end{zproof}

\begin{theorem}{rule OplusInner}
  \forall F, G: \nat \fun \nat @ (\lambda x: \nat @ F~x + G~x * G~x) \in \nat \fun \nat
\end{theorem}

\begin{zproof}
apply FunIntro to predicate (\lambda x: \nat @ F~x + G~x * G~x) \in \nat \fun \nat;
with enabled (RelPow, inPower) prove by reduce;
\end{zproof}

\begin{zed}
  OplusWit == (\lambda F, G: \nat \fun \nat @ (\lambda x: \nat @ F~x + G~x * G~x))
\end{zed}

\begin{zproof}[OplusWit\$domainCheck]
prove;
\end{zproof}

\begin{theorem}{OplusWitFun}
  OplusWit \in (\nat \fun \nat) \cross (\nat \fun \nat) \fun (\nat \fun \nat)
\end{theorem}

\begin{zproof}
invoke OplusWit;
apply FunIntro to predicate (\lambda F, G: \nat \fun \nat @ (\lambda x: \nat @ F~x + G~x * G~x)) \in (\nat \fun \nat) \cross (\nat \fun \nat) \fun (\nat \fun \nat);
with enabled (RelPow, inPower) prove by reduce;
instantiate F == x.1, G == x.2;
prove;
\end{zproof}

\begin{theorem}{OplusWitApp}
  \forall F, G: \nat \fun \nat @ OplusWit(F, G) = (\lambda x: \nat @ F~x + G~x * G~x)
\end{theorem}

\begin{zproof}
use OplusWitFun;
use pairInFunction[(\nat \fun \nat) \cross (\nat \fun \nat), \nat \fun \nat][f := OplusWit, x := (F, G), y := (\lambda x: \nat @ F~x + G~x * G~x)];
invoke OplusWit;
prove by reduce;
\end{zproof}

\begin{theorem}{rule OplusWitVal}
  \forall F, G: \nat \fun \nat; x: \nat @ OplusWit(F, G)~x = F~x + G~x * G~x
\end{theorem}

\begin{zproof}
prove;
use OplusWitApp[F := F, G := G];
use pairInFunction[\nat, \nat][f := (\lambda x: \nat @ F~x + G~x * G~x), x := x, y := F~x + G~x * G~x];
prove;
\end{zproof}

\begin{theorem}{OplusEx}
  \exists opOplus: (\nat \fun \nat) \cross (\nat \fun \nat) \fun (\nat \fun \nat) @ \forall F, G: \nat \fun \nat; x: \nat @ opOplus(F, G)~x = F~x + G~x * G~x
\end{theorem}

\begin{zproof}
instantiate opOplus == OplusWit;
use OplusWitFun;
prove;
\end{zproof}

\begin{zed}
  exprOplus == (\lambda F, G: \nat \fun \nat @ comp(opPlusI~(comp(exprMinusI~F, funE)), opPlusI~(comp(funSq, G))))
\end{zed}

\begin{zproof}[exprOplus\$domainCheck]
prove;
\end{zproof}

\begin{theorem}{rule exprOplusFun}
  exprOplus \in (\nat \fun \nat) \cross (\nat \fun \nat) \fun (\nat \fun \nat)
\end{theorem}

\begin{zproof}
invoke exprOplus;
apply FunIntro to predicate (\lambda F, G: \nat \fun \nat @ comp(opPlusI~(comp(exprMinusI~F, funE)), opPlusI~(comp(funSq, G)))) \in (\nat \fun \nat) \cross (\nat \fun \nat) \fun (\nat \fun \nat);
with enabled (RelPow, inPower) prove;
instantiate F == x.1, G == x.2;
prove;
\end{zproof}

\begin[disabled]{theorem}{rule exprOplusApp}
  \forall F, G: \nat \fun \nat @ exprOplus(F, G) = comp(opPlusI~(comp(exprMinusI~F, funE)), opPlusI~(comp(funSq, G)))
\end{theorem}

\begin{zproof}
use exprOplusFun;
use pairInFunction[(\nat \fun \nat) \cross (\nat \fun \nat), \nat \fun \nat][f := exprOplus, x := (F, G), y := comp(opPlusI~(comp(exprMinusI~F, funE)), opPlusI~(comp(funSq, G)))];
invoke exprOplus;
prove by reduce;
\end{zproof}

\begin{zed}
  exprOminus == (\lambda F, G: \nat \fun \nat @ comp(exprMinusI~(comp(opPlusI~F, funE)), opPlusI~(comp(funSq, G))))
\end{zed}

\begin{zproof}[exprOminus\$domainCheck]
prove;
\end{zproof}

\begin{theorem}{rule exprOminusFun}
  exprOminus \in (\nat \fun \nat) \cross (\nat \fun \nat) \fun (\nat \fun \nat)
\end{theorem}

\begin{zproof}
invoke exprOminus;
apply FunIntro to predicate (\lambda F, G: \nat \fun \nat @ comp(exprMinusI~(comp(opPlusI~F, funE)), opPlusI~(comp(funSq, G)))) \in (\nat \fun \nat) \cross (\nat \fun \nat) \fun (\nat \fun \nat);
with enabled (RelPow, inPower) prove;
instantiate F == x.1, G == x.2;
prove;
\end{zproof}

\begin[disabled]{theorem}{rule exprOminusApp}
  \forall F, G: \nat \fun \nat @ exprOminus(F, G) = comp(exprMinusI~(comp(opPlusI~F, funE)), opPlusI~(comp(funSq, G)))
\end{theorem}

\begin{zproof}
use exprOminusFun;
use pairInFunction[(\nat \fun \nat) \cross (\nat \fun \nat), \nat \fun \nat][f := exprOminus, x := (F, G), y := comp(exprMinusI~(comp(opPlusI~F, funE)), opPlusI~(comp(funSq, G)))];
invoke exprOminus;
prove by reduce;
\end{zproof}

\begin{theorem}{rule OplusIn}
  \forall F, G: \nat \fun \nat | F \in C6 \land G \in C6 @ exprOplus(F, G) \in C6
\end{theorem}

\begin{zproof}
with enabled (exprOplusApp) prove;
\end{zproof}

\begin{theorem}{rule OminusIn}
  \forall F, G: \nat \fun \nat | F \in C6 \land G \in C6 @ exprOminus(F, G) \in C6
\end{theorem}

\begin{zproof}
with enabled (exprOminusApp) prove;
\end{zproof}

\begin{theorem}{rule OplusExprVal}
  \forall F, G: \nat \fun \nat; x: \nat | x \leq F~x \land x \leq G~x @
    exprOplus(F, G)~x = F~x + G~x * G~x
\end{theorem}

\begin{zproof}
with enabled (exprOplusApp) prove;
use ESqAdd[m := G~x, k := x];
use MinusIVal[F := F, x := x];
prove;
\end{zproof}

\begin{theorem}{rule OminusExprVal}
  \forall F, G: \nat \fun \nat; x: \nat | x \leq G~x \land G~x * G~x \leq F~x @
    exprOminus(F, G)~x = F~x - G~x * G~x
\end{theorem}

\begin{zproof}
with enabled (exprOminusApp) prove;
use ESqAdd[m := G~x, k := x];
use MinusIVal[F := comp(opPlusI~F, funE), x := G~x * G~x + x];
prove;
\end{zproof}

\begin{theorem}{OplusEq}
  \forall F, G: \nat \fun \nat | F \in Infl \land G \in Infl @ exprOplus(F, G) = opOplus(F, G)
\end{theorem}

\begin{zproof}
apply UFExt to predicate exprOplus(F, G) = opOplus(F, G);
prove;
\end{zproof}

\begin{theorem}{Lemma64Plus}
  \forall F, G: \nat \fun \nat | F \in C6 \land G \in C6 \land F \in Infl \land G \in Infl @
    opOplus(F, G) \in C6
\end{theorem}

\begin{zproof}
use OplusEq[F := F, G := G];
prove;
equality substitute opOplus(F, G);
rewrite;
\end{zproof}

\begin{theorem}{Lemma64Minus}
  \forall F, G: \nat \fun \nat | F \in C6 \land G \in C6 \land G \in Infl \land (F, G) \in SqBelow @
    \exists H: \nat \fun \nat | H \in C6 @ \forall x: \nat @ H~x = F~x - G~x * G~x
\end{theorem}

\begin{zproof}
prove;
instantiate H == exprOminus(F, G);
prove;
\end{zproof}

\begin{theorem}{Lemma64}
  \forall F, G: \nat \fun \nat | F \in C6 \land G \in C6 \land F \in Infl \land G \in Infl @
    opOplus(F, G) \in C6 \land ((F, G) \in SqBelow \implies
      (\exists H: \nat \fun \nat | H \in C6 @ \forall x: \nat @ H~x = F~x - G~x * G~x))
\end{theorem}

\begin{zproof}
use Lemma64Plus[F := F, G := G];
use Lemma64Minus[F := F, G := G];
prove;
instantiate H\_\_0 == H;
prove;
instantiate x\_\_0 == x;
prove;
\end{zproof}

\subsection{Lemma 6.5: the predecessor $\hat P$ and the product $F \otimes G$}

\begin{theorem}{MulNonneg}
  \forall a, b: \nat @ a * b \geq 0
\end{theorem}

\begin{zproof}
prove;
\end{zproof}

\begin{theorem}{SqProd}
  \forall a, b: \num @ (a * b) * (a * b) = a * (a * (b * b))
\end{theorem}

\begin{zproof}
prove;
\end{zproof}

\begin{theorem}{RtSqProd}
  \forall a, b: \nat @ funRt~(a * (a * (b * b))) = a * b
\end{theorem}

\begin{zproof}
prove;
use RtI[r := a * b, x := a * (a * (b * b))];
use MulNonneg[a := a, b := b];
prove;
\end{zproof}

\begin{theorem}{LeSq}
  \forall n: \nat @ n \leq n * n
\end{theorem}

\begin{zproof}
prove;
split n = 0;
cases;
prove;
next;
use timesMonotonic1[a := n, b := 1, c := n];
prove;
next;
\end{zproof}

$\hat P = ((\mathrm{Sq}^+)^+)^+ \ominus S$ is a predecessor: $\hat P(x+1) = x$.

\begin{zed}
  exprPhat == exprOminus(opPlusI~(opPlusI~(opPlusI~funSq)), funS)
\end{zed}

\begin{zproof}[exprPhat\$domainCheck]
prove;
\end{zproof}

\begin{theorem}{rule PhatFun}
  exprPhat \in \nat \fun \nat
\end{theorem}

\begin{zproof}
invoke exprPhat;
prove;
\end{zproof}

\begin{theorem}{rule PhatIn}
  exprPhat \in C6
\end{theorem}

\begin{zproof}
invoke exprPhat;
prove;
\end{zproof}

\begin{theorem}{rule PhatVal}
  \forall x: \nat @ exprPhat~(1 + x) = x
\end{theorem}

\begin{zproof}
prove;
invoke exprPhat;
use OminusExprVal[F := opPlusI~(opPlusI~(opPlusI~funSq)), G := funS, x := 1 + x];
prove;
\end{zproof}

$(F \otimes G)(x) = F(x)\,G(x)$ (Definition 6.1), and its expression in Lemma 6.5.

\begin{axdef}
  opOtimes: (\nat \fun \nat) \cross (\nat \fun \nat) \fun (\nat \fun \nat)
\where
  \Label{rule OtimesVal} \forall F, G: \nat \fun \nat; x: \nat @ opOtimes(F, G)~x = F~x * G~x
\end{axdef}

\begin{zproof}[opOtimes\$domainCheck]
prove;
\end{zproof}

\begin{theorem}{rule OtimesInner}
  \forall F, G: \nat \fun \nat @ (\lambda x: \nat @ F~x * G~x) \in \nat \fun \nat
\end{theorem}

\begin{zproof}
apply FunIntro to predicate (\lambda x: \nat @ F~x * G~x) \in \nat \fun \nat;
with enabled (RelPow, inPower) prove by reduce;
\end{zproof}

\begin{zed}
  OtimesWit == (\lambda F, G: \nat \fun \nat @ (\lambda x: \nat @ F~x * G~x))
\end{zed}

\begin{zproof}[OtimesWit\$domainCheck]
prove;
\end{zproof}

\begin{theorem}{OtimesWitFun}
  OtimesWit \in (\nat \fun \nat) \cross (\nat \fun \nat) \fun (\nat \fun \nat)
\end{theorem}

\begin{zproof}
invoke OtimesWit;
apply FunIntro to predicate (\lambda F, G: \nat \fun \nat @ (\lambda x: \nat @ F~x * G~x)) \in (\nat \fun \nat) \cross (\nat \fun \nat) \fun (\nat \fun \nat);
with enabled (RelPow, inPower) prove by reduce;
instantiate F == x.1, G == x.2;
prove;
\end{zproof}

\begin{theorem}{OtimesWitApp}
  \forall F, G: \nat \fun \nat @ OtimesWit(F, G) = (\lambda x: \nat @ F~x * G~x)
\end{theorem}

\begin{zproof}
use OtimesWitFun;
use pairInFunction[(\nat \fun \nat) \cross (\nat \fun \nat), \nat \fun \nat][f := OtimesWit, x := (F, G), y := (\lambda x: \nat @ F~x * G~x)];
invoke OtimesWit;
prove by reduce;
\end{zproof}

\begin{theorem}{rule OtimesWitVal}
  \forall F, G: \nat \fun \nat; x: \nat @ OtimesWit(F, G)~x = F~x * G~x
\end{theorem}

\begin{zproof}
prove;
use OtimesWitApp[F := F, G := G];
use pairInFunction[\nat, \nat][f := (\lambda x: \nat @ F~x * G~x), x := x, y := F~x * G~x];
prove;
\end{zproof}

\begin{theorem}{OtimesEx}
  \exists opOtimes: (\nat \fun \nat) \cross (\nat \fun \nat) \fun (\nat \fun \nat) @ \forall F, G: \nat \fun \nat; x: \nat @ opOtimes(F, G)~x = F~x * G~x
\end{theorem}

\begin{zproof}
instantiate opOtimes == OtimesWit;
use OtimesWitFun;
prove;
\end{zproof}

\begin{zed}
  exprOtimes == (\lambda F, G: \nat \fun \nat @ comp(funRt, exprOminus(comp(comp(funHf, funHf), exprOminus(comp(comp(funSq, exprPhat), exprOminus(exprOplus(comp(comp(comp(funSq, funS), funSq), F), G), comp(funSq, F))), comp(funSq, G))), comp(funSq, F))))
\end{zed}

\begin{zproof}[exprOtimes\$domainCheck]
prove;
\end{zproof}

\begin{theorem}{rule exprOtimesFun}
  exprOtimes \in (\nat \fun \nat) \cross (\nat \fun \nat) \fun (\nat \fun \nat)
\end{theorem}

\begin{zproof}
invoke exprOtimes;
apply FunIntro to predicate (\lambda F, G: \nat \fun \nat @ comp(funRt, exprOminus(comp(comp(funHf, funHf), exprOminus(comp(comp(funSq, exprPhat), exprOminus(exprOplus(comp(comp(comp(funSq, funS), funSq), F), G), comp(funSq, F))), comp(funSq, G))), comp(funSq, F)))) \in (\nat \fun \nat) \cross (\nat \fun \nat) \fun (\nat \fun \nat);
with enabled (RelPow, inPower) prove;
instantiate F == x.1, G == x.2;
prove;
\end{zproof}

\begin[disabled]{theorem}{rule exprOtimesApp}
  \forall F, G: \nat \fun \nat @ exprOtimes(F, G) = comp(funRt, exprOminus(comp(comp(funHf, funHf), exprOminus(comp(comp(funSq, exprPhat), exprOminus(exprOplus(comp(comp(comp(funSq, funS), funSq), F), G), comp(funSq, F))), comp(funSq, G))), comp(funSq, F)))
\end{theorem}

\begin{zproof}
use exprOtimesFun;
use pairInFunction[(\nat \fun \nat) \cross (\nat \fun \nat), \nat \fun \nat][f := exprOtimes, x := (F, G), y := comp(funRt, exprOminus(comp(comp(funHf, funHf), exprOminus(comp(comp(funSq, exprPhat), exprOminus(exprOplus(comp(comp(comp(funSq, funS), funSq), F), G), comp(funSq, F))), comp(funSq, G))), comp(funSq, F)))];
invoke exprOtimes;
prove by reduce;
\end{zproof}

\begin{theorem}{rule OtimesIn}
  \forall F, G: \nat \fun \nat | F \in C6 \land G \in C6 @ exprOtimes(F, G) \in C6
\end{theorem}

\begin{zproof}
with enabled (exprOtimesApp) prove;
\end{zproof}

\begin{theorem}{rule OtimesExprVal}
  \forall F, G: \nat \fun \nat; x: \nat | x \leq F~x \land x \leq G~x @
    exprOtimes(F, G)~x = F~x * G~x
\end{theorem}

\begin{zproof}
use LeSq[n := F~x];
use LeSq[n := G~x];
use LeSq[n := 1 + F~x * F~x];
use MulNonneg[a := F~x, b := F~x];
use MulNonneg[a := G~x, b := G~x];
use MulNonneg[a := F~x * F~x, b := G~x * G~x];
use MulNonneg[a := F~x * F~x, b := F~x * F~x];
with enabled (exprOtimesApp) prove;
use DivMod[x := 4 * (F~x * F~x * F~x * F~x + F~x * F~x * G~x * G~x), d := 2,
  q := 2 * (F~x * F~x * F~x * F~x + F~x * F~x * G~x * G~x), r := 0];
prove;
use DivMod[x := 2 * (F~x * F~x * F~x * F~x + F~x * F~x * G~x * G~x), d := 2,
  q := F~x * F~x * F~x * F~x + F~x * F~x * G~x * G~x, r := 0];
prove;
use RtSqProd[a := F~x, b := G~x];
prove;
\end{zproof}

\begin{theorem}{OtimesEq}
  \forall F, G: \nat \fun \nat | F \in Infl \land G \in Infl @ exprOtimes(F, G) = opOtimes(F, G)
\end{theorem}

\begin{zproof}
apply UFExt to predicate exprOtimes(F, G) = opOtimes(F, G);
prove;
\end{zproof}

\begin{theorem}{Lemma65}
  \forall F, G: \nat \fun \nat | F \in C6 \land G \in C6 \land F \in Infl \land G \in Infl @
    opOtimes(F, G) \in C6
\end{theorem}

\begin{zproof}
use OtimesEq[F := F, G := G];
prove;
equality substitute opOtimes(F, G);
rewrite;
\end{zproof}

\subsection{Lemma 6.6: $F + G$ and the difference $F - G$}

\begin{theorem}{LeProd}
  \forall x, a, b: \nat | x \leq a \land x \leq b @ x \leq a * b
\end{theorem}

\begin{zproof}
prove;
split x = 0;
cases;
use MulNonneg[a := a, b := b];
prove;
next;
use timesMonotonic1[a := a, b := 1, c := b];
prove;
next;
\end{zproof}

\begin{theorem}{RtSqSum}
  \forall a, b: \nat @ funRt~(2 * (a * b) + a * a + b * b) = a + b
\end{theorem}

\begin{zproof}
prove;
use RtI[r := a + b, x := 2 * (a * b) + a * a + b * b];
prove;
\end{zproof}

$F + G = ((\mathrm{Rt}((D(F^+ \otimes G^+) \oplus F^+) \oplus G^+))^-)^-$.

\begin{zed}
  exprPlus == (\lambda F, G: \nat \fun \nat @ exprMinusI~(exprMinusI~(comp(funRt, exprOplus(exprOplus(comp(funD, exprOtimes(opPlusI~F, opPlusI~G)), opPlusI~F), opPlusI~G)))))
\end{zed}

\begin{zproof}[exprPlus\$domainCheck]
prove;
\end{zproof}

\begin{theorem}{rule exprPlusFun}
  exprPlus \in (\nat \fun \nat) \cross (\nat \fun \nat) \fun (\nat \fun \nat)
\end{theorem}

\begin{zproof}
invoke exprPlus;
apply FunIntro to predicate (\lambda F, G: \nat \fun \nat @ exprMinusI~(exprMinusI~(comp(funRt, exprOplus(exprOplus(comp(funD, exprOtimes(opPlusI~F, opPlusI~G)), opPlusI~F), opPlusI~G))))) \in (\nat \fun \nat) \cross (\nat \fun \nat) \fun (\nat \fun \nat);
with enabled (RelPow, inPower) prove;
instantiate F == x.1, G == x.2;
prove;
\end{zproof}

\begin[disabled]{theorem}{rule exprPlusApp}
  \forall F, G: \nat \fun \nat @ exprPlus(F, G) = exprMinusI~(exprMinusI~(comp(funRt, exprOplus(exprOplus(comp(funD, exprOtimes(opPlusI~F, opPlusI~G)), opPlusI~F), opPlusI~G))))
\end{theorem}

\begin{zproof}
use exprPlusFun;
use pairInFunction[(\nat \fun \nat) \cross (\nat \fun \nat), \nat \fun \nat][f := exprPlus, x := (F, G), y := exprMinusI~(exprMinusI~(comp(funRt, exprOplus(exprOplus(comp(funD, exprOtimes(opPlusI~F, opPlusI~G)), opPlusI~F), opPlusI~G))))];
invoke exprPlus;
prove by reduce;
\end{zproof}

\begin{theorem}{rule PlusExprIn}
  \forall F, G: \nat \fun \nat | F \in C6 \land G \in C6 @ exprPlus(F, G) \in C6
\end{theorem}

\begin{zproof}
with enabled (exprPlusApp) prove;
\end{zproof}

\begin{theorem}{rule PlusExprVal}
  \forall F, G: \nat \fun \nat; x: \nat @ exprPlus(F, G)~x = F~x + G~x
\end{theorem}

\begin{zproof}
use UFNonneg[f := F, x := x];
use UFNonneg[f := G, x := x];
use LeProd[x := x, a := F~x + x, b := G~x + x];
use MulNonneg[a := F~x + x, b := G~x + x];
use MulNonneg[a := F~x + x, b := F~x + x];
use RtSqSum[a := F~x + x, b := G~x + x];
with enabled (exprPlusApp) prove;
\end{zproof}

\begin{theorem}{PlusExprEq}
  \forall F, G: \nat \fun \nat @ exprPlus(F, G) = opPlus(F, G)
\end{theorem}

\begin{zproof}
apply UFExt to predicate exprPlus(F, G) = opPlus(F, G);
prove;
\end{zproof}

The difference function of footnote 8: it equals $F - G$ where $F \ge G$.

\begin{axdef}
  opDiffE: (\nat \fun \nat) \cross (\nat \fun \nat) \fun (\nat \fun \nat)
\where
  \Label{rule DiffEVal} \forall F, G: \nat \fun \nat; x: \nat @ opDiffE(F, G)~x = funE~((F~x + G~x) * (F~x + G~x) + 3 * F~x + G~x + 1)
\end{axdef}

\begin{zproof}[opDiffE\$domainCheck]
prove;
\end{zproof}

\begin{theorem}{rule DiffEInner}
  \forall F, G: \nat \fun \nat @ (\lambda x: \nat @ funE~((F~x + G~x) * (F~x + G~x) + 3 * F~x + G~x + 1)) \in \nat \fun \nat
\end{theorem}

\begin{zproof}
apply FunIntro to predicate (\lambda x: \nat @ funE~((F~x + G~x) * (F~x + G~x) + 3 * F~x + G~x + 1)) \in \nat \fun \nat;
with enabled (RelPow, inPower) prove by reduce;
\end{zproof}

\begin{zed}
  DiffEWit == (\lambda F, G: \nat \fun \nat @ (\lambda x: \nat @ funE~((F~x + G~x) * (F~x + G~x) + 3 * F~x + G~x + 1)))
\end{zed}

\begin{zproof}[DiffEWit\$domainCheck]
prove;
\end{zproof}

\begin{theorem}{DiffEWitFun}
  DiffEWit \in (\nat \fun \nat) \cross (\nat \fun \nat) \fun (\nat \fun \nat)
\end{theorem}

\begin{zproof}
invoke DiffEWit;
apply FunIntro to predicate (\lambda F, G: \nat \fun \nat @ (\lambda x: \nat @ funE~((F~x + G~x) * (F~x + G~x) + 3 * F~x + G~x + 1))) \in (\nat \fun \nat) \cross (\nat \fun \nat) \fun (\nat \fun \nat);
with enabled (RelPow, inPower) prove by reduce;
instantiate F == x.1, G == x.2;
prove;
\end{zproof}

\begin{theorem}{DiffEWitApp}
  \forall F, G: \nat \fun \nat @ DiffEWit(F, G) = (\lambda x: \nat @ funE~((F~x + G~x) * (F~x + G~x) + 3 * F~x + G~x + 1))
\end{theorem}

\begin{zproof}
use DiffEWitFun;
use pairInFunction[(\nat \fun \nat) \cross (\nat \fun \nat), \nat \fun \nat][f := DiffEWit, x := (F, G), y := (\lambda x: \nat @ funE~((F~x + G~x) * (F~x + G~x) + 3 * F~x + G~x + 1))];
invoke DiffEWit;
prove by reduce;
\end{zproof}

\begin{theorem}{rule DiffEWitVal}
  \forall F, G: \nat \fun \nat; x: \nat @ DiffEWit(F, G)~x = funE~((F~x + G~x) * (F~x + G~x) + 3 * F~x + G~x + 1)
\end{theorem}

\begin{zproof}
prove;
use DiffEWitApp[F := F, G := G];
use pairInFunction[\nat, \nat][f := (\lambda x: \nat @ funE~((F~x + G~x) * (F~x + G~x) + 3 * F~x + G~x + 1)), x := x, y := funE~((F~x + G~x) * (F~x + G~x) + 3 * F~x + G~x + 1)];
prove;
\end{zproof}

\begin{theorem}{DiffEEx}
  \exists opDiffE: (\nat \fun \nat) \cross (\nat \fun \nat) \fun (\nat \fun \nat) @ \forall F, G: \nat \fun \nat; x: \nat @ opDiffE(F, G)~x = funE~((F~x + G~x) * (F~x + G~x) + 3 * F~x + G~x + 1)
\end{theorem}

\begin{zproof}
instantiate opDiffE == DiffEWit;
use DiffEWitFun;
prove;
\end{zproof}

$F - G = E\,S\,(\mathrm{Sq}(F + G) + M_3 F + G)$.

\begin{zed}
  exprMinus == (\lambda F, G: \nat \fun \nat @ comp(comp(funE, funS), exprPlus(exprPlus(comp(funSq, exprPlus(F, G)), comp(funM3, F)), G)))
\end{zed}

\begin{zproof}[exprMinus\$domainCheck]
prove;
\end{zproof}

\begin{theorem}{rule exprMinusFun}
  exprMinus \in (\nat \fun \nat) \cross (\nat \fun \nat) \fun (\nat \fun \nat)
\end{theorem}

\begin{zproof}
invoke exprMinus;
apply FunIntro to predicate (\lambda F, G: \nat \fun \nat @ comp(comp(funE, funS), exprPlus(exprPlus(comp(funSq, exprPlus(F, G)), comp(funM3, F)), G))) \in (\nat \fun \nat) \cross (\nat \fun \nat) \fun (\nat \fun \nat);
with enabled (RelPow, inPower) prove;
instantiate F == x.1, G == x.2;
prove;
\end{zproof}

\begin[disabled]{theorem}{rule exprMinusApp}
  \forall F, G: \nat \fun \nat @ exprMinus(F, G) = comp(comp(funE, funS), exprPlus(exprPlus(comp(funSq, exprPlus(F, G)), comp(funM3, F)), G))
\end{theorem}

\begin{zproof}
use exprMinusFun;
use pairInFunction[(\nat \fun \nat) \cross (\nat \fun \nat), \nat \fun \nat][f := exprMinus, x := (F, G), y := comp(comp(funE, funS), exprPlus(exprPlus(comp(funSq, exprPlus(F, G)), comp(funM3, F)), G))];
invoke exprMinus;
prove by reduce;
\end{zproof}

\begin{theorem}{rule MinusExprIn}
  \forall F, G: \nat \fun \nat | F \in C6 \land G \in C6 @ exprMinus(F, G) \in C6
\end{theorem}

\begin{zproof}
with enabled (exprMinusApp) prove;
\end{zproof}

\begin{theorem}{MinusExprEq}
  \forall F, G: \nat \fun \nat @ exprMinus(F, G) = opDiffE(F, G)
\end{theorem}

\begin{zproof}
apply UFExt to predicate exprMinus(F, G) = opDiffE(F, G);
with enabled (exprMinusApp) prove;
\end{zproof}

\begin{theorem}{DiffEIsDifference}
  \forall F, G: \nat \fun \nat; x: \nat | G~x \leq F~x @ opDiffE(F, G)~x = F~x - G~x
\end{theorem}

\begin{zproof}
prove;
use RobDiffE[a := F~x, b := G~x];
prove;
\end{zproof}

\begin{theorem}{Lemma66}
  \forall F, G: \nat \fun \nat | F \in C6 \land G \in C6 @ opPlus(F, G) \in C6 \land opDiffE(F, G) \in C6
\end{theorem}

\begin{zproof}
use PlusExprEq[F := F, G := G];
use MinusExprEq[F := F, G := G];
prove;
equality substitute opPlus(F, G);
equality substitute opDiffE(F, G);
rewrite;
\end{zproof}

\begin{theorem}{rule C6Plus}
  \forall F, G: \nat \fun \nat | F \in C6 \land G \in C6 @ opPlus(F, G) \in C6
\end{theorem}

\begin{zproof}
use Lemma66[F := F, G := G];
prove;
\end{zproof}

\begin{theorem}{rule C6DiffE}
  \forall F, G: \nat \fun \nat | F \in C6 \land G \in C6 @ opDiffE(F, G) \in C6
\end{theorem}

\begin{zproof}
use Lemma66[F := F, G := G];
prove;
\end{zproof}

\section{Theorem 6.7 (closure part)}

Robinson [15, Sec.~7, Thm.~3] proved that $\langle S, E, F + G, FG, F^\square\rangle$ is
suitable; here we show that it generates exactly $\mathrm{clos}\,\mathcal{F}$, so
$\mathrm{clos}\,\mathcal{F} = \mathrm{Prim}(\mathbb{N},\mathbb{N})$ follows from Robinson's theorem.

\begin{zed}
  RobinsonClos == clos~\{ funS, funE \}~\{ opBox \}~\{ opPlus, comp \}
\end{zed}

\begin{zproof}[RobinsonClos\$domainCheck]
prove;
\end{zproof}

\begin[disabled]{theorem}{rule RobDef}
  RobinsonClos = clos~\{ funS, funE \}~\{ opBox \}~\{ opPlus, comp \}
\end{theorem}

\begin{zproof}
invoke RobinsonClos;
prove;
\end{zproof}

\begin[disabled]{theorem}{rule RobUF}
  \forall F: \power (\num \cross \num) | F \in RobinsonClos @ F \in \nat \fun \nat
\end{theorem}

\begin{zproof}
prove;
use ClosElemUF[I := \{ funS, funE \}, U := \{ opBox \}, B := \{ opPlus, comp \}, F := F];
with enabled (RobDef) prove;
\end{zproof}

\begin{theorem}{rule RobS}
  funS \in RobinsonClos
\end{theorem}

\begin{zproof}
with enabled (RobDef) prove;
\end{zproof}

\begin{theorem}{rule RobE}
  funE \in RobinsonClos
\end{theorem}

\begin{zproof}
with enabled (RobDef) prove;
\end{zproof}

\begin{theorem}{rule RobBox}
  \forall F: \nat \fun \nat | F \in RobinsonClos @ opBox~F \in RobinsonClos
\end{theorem}

\begin{zproof}
with enabled (RobDef) prove;
\end{zproof}

\begin{theorem}{rule RobPlus}
  \forall F, G: \nat \fun \nat | F \in RobinsonClos \land G \in RobinsonClos @ opPlus(F, G) \in RobinsonClos
\end{theorem}

\begin{zproof}
with enabled (RobDef) prove;
\end{zproof}

\begin{theorem}{rule RobComp}
  \forall F, G: \nat \fun \nat | F \in RobinsonClos \land G \in RobinsonClos @ comp(F, G) \in RobinsonClos
\end{theorem}

\begin{zproof}
with enabled (RobDef) prove;
\end{zproof}

\begin{zed}
  BoxESet == \{ n: \nat | opBox~funE~n = 0 \}
\end{zed}

\begin{zproof}[BoxESet\$domainCheck]
prove;
\end{zproof}

\begin{theorem}{rule BoxEBase}
  0 \in BoxESet
\end{theorem}

\begin{zproof}
invoke BoxESet;
prove;
\end{zproof}

\begin{theorem}{rule BoxEStep}
  \forall n: BoxESet @ 1 + n \in BoxESet
\end{theorem}

\begin{zproof}
invoke BoxESet;
prove;
\end{zproof}

\begin{theorem}{BoxESub}
  \nat \subseteq BoxESet
\end{theorem}

\begin{zproof}
apply natInduction to predicate \nat \subseteq BoxESet;
prove;
invoke BoxESet;
prove;
\end{zproof}

\begin{theorem}{BoxEAll}
  \forall n: \nat @ opBox~funE~n = 0
\end{theorem}

\begin{zproof}
prove;
use BoxESub;
invoke BoxESet;
with enabled (inPower) prove;
instantiate e == n;
prove;
\end{zproof}

\begin{theorem}{BoxEExpr}
  opBox~funE = cst~0
\end{theorem}

\begin{zproof}
apply UFExt to predicate opBox~funE = cst~0;
prove;
use BoxEAll[n := x];
prove;
\end{zproof}

\begin{zed}
  BoxSSet == \{ n: \nat | opBox~funS~n = n \}
\end{zed}

\begin{zproof}[BoxSSet\$domainCheck]
prove;
\end{zproof}

\begin{theorem}{rule BoxSBase}
  0 \in BoxSSet
\end{theorem}

\begin{zproof}
invoke BoxSSet;
prove;
\end{zproof}

\begin{theorem}{rule BoxSStep}
  \forall n: BoxSSet @ 1 + n \in BoxSSet
\end{theorem}

\begin{zproof}
invoke BoxSSet;
prove;
\end{zproof}

\begin{theorem}{BoxSSub}
  \nat \subseteq BoxSSet
\end{theorem}

\begin{zproof}
apply natInduction to predicate \nat \subseteq BoxSSet;
prove;
invoke BoxSSet;
prove;
\end{zproof}

\begin{theorem}{BoxSAll}
  \forall n: \nat @ opBox~funS~n = n
\end{theorem}

\begin{zproof}
prove;
use BoxSSub;
invoke BoxSSet;
with enabled (inPower) prove;
instantiate e == n;
prove;
\end{zproof}

\begin{theorem}{BoxSExpr}
  opBox~funS = funI
\end{theorem}

\begin{zproof}
apply UFExt to predicate opBox~funS = funI;
prove;
use BoxSAll[n := x];
prove;
\end{zproof}

\begin{theorem}{OneExpr}
  comp(funS, cst~0) = cst~1
\end{theorem}

\begin{zproof}
apply UFExt to predicate comp(funS, cst~0) = cst~1;
prove;
\end{zproof}

\begin{theorem}{PlusIEq}
  \forall F: \nat \fun \nat @ opPlusI~F = opPlus(F, funI)
\end{theorem}

\begin{zproof}
apply UFExt to predicate opPlusI~F = opPlus(F, funI);
prove;
\end{zproof}

\begin{theorem}{RobIExpr}
  opBox~funS \in RobinsonClos
\end{theorem}

\begin{zproof}
prove;
\end{zproof}

\begin{theorem}{rule RobI}
  funI \in RobinsonClos
\end{theorem}

\begin{zproof}
use RobIExpr;
use BoxSExpr;
prove;
\end{zproof}
\begin{theorem}{RobZeroExpr}
  opBox~funE \in RobinsonClos
\end{theorem}

\begin{zproof}
prove;
\end{zproof}

\begin{theorem}{rule RobZero}
  cst~0 \in RobinsonClos
\end{theorem}

\begin{zproof}
use BoxEExpr;
equality substitute cst~0;
rewrite;
\end{zproof}
\begin{theorem}{RobOneExpr}
  comp(funS, cst~0) \in RobinsonClos
\end{theorem}

\begin{zproof}
prove;
\end{zproof}

\begin{theorem}{rule RobOne}
  cst~1 \in RobinsonClos
\end{theorem}

\begin{zproof}
use RobOneExpr;
use OneExpr;
prove;
\end{zproof}

\begin{theorem}{rule RobPlusI}
  \forall F: \nat \fun \nat | F \in RobinsonClos @ opPlusI~F \in RobinsonClos
\end{theorem}

\begin{zproof}
use PlusIEq[F := F];
prove;
\end{zproof}



\begin{theorem}{C6InPower}
  C6 \in \power (\nat \fun \nat)
\end{theorem}

\begin{zproof}
use ClosInPow[I := \{ cst~1, funE \}, U := \{ opPlusI, opBox \}, B := \{ comp \}];
with enabled (C6Def) prove;
\end{zproof}

\begin{theorem}{RobInPower}
  RobinsonClos \in \power (\nat \fun \nat)
\end{theorem}

\begin{zproof}
use ClosInPow[I := \{ funS, funE \}, U := \{ opBox \}, B := \{ opPlus, comp \}];
with enabled (RobDef) prove;
\end{zproof}

\begin{theorem}{RobUClosed}
  RobinsonClos \in UClosed~\{ opPlusI, opBox \}
\end{theorem}

\begin{zproof}
use RobInPower;
with enabled (UClosedDef, C6UF, RobUF) prove;
split u = opPlusI;
cases;
prove;
next;
prove;
next;
\end{zproof}

\begin{theorem}{RobBClosed}
  RobinsonClos \in BClosed~\{ comp \}
\end{theorem}

\begin{zproof}
use RobInPower;
with enabled (BClosedDef, C6UF, RobUF) prove;
\end{zproof}

\begin{theorem}{C6UClosed}
  C6 \in UClosed~\{ opBox \}
\end{theorem}

\begin{zproof}
use C6InPower;
with enabled (UClosedDef, C6UF, RobUF) prove;
\end{zproof}

\begin{theorem}{C6BClosed}
  C6 \in BClosed~\{ opPlus, comp \}
\end{theorem}

\begin{zproof}
use C6InPower;
with enabled (BClosedDef, C6UF, RobUF) prove;
split b = opPlus;
cases;
prove;
next;
prove;
next;
\end{zproof}

\begin{theorem}{C6SubRob}
  C6 \subseteq RobinsonClos
\end{theorem}

\begin{zproof}
use ClosLeast[I := \{ cst~1, funE \}, U := \{ opPlusI, opBox \}, B := \{ comp \}, Z := RobinsonClos];
use RobInPower;
use RobUClosed;
use RobBClosed;
with enabled (C6Def) prove;
\end{zproof}

\begin{theorem}{RobSubC6}
  RobinsonClos \subseteq C6
\end{theorem}

\begin{zproof}
use ClosLeast[I := \{ funS, funE \}, U := \{ opBox \}, B := \{ opPlus, comp \}, Z := C6];
use C6InPower;
use C6UClosed;
use C6BClosed;
with enabled (RobDef) prove;
\end{zproof}

\begin{theorem}{Theorem67}
  C6 = RobinsonClos
\end{theorem}

\begin{zproof}
use C6SubRob;
use RobSubC6;
use C6InPower;
use RobInPower;
use SubAntisym[\nat \fun \nat][A := C6, B := RobinsonClos];
prove;
\end{zproof}

\section{Cantor pairing functions}

$A(x) = x(x+1)/2$, $V$ its inverse ($A(V(x)) \le x < A(V(x)+1)$),
$K(x) = x - A(V(x))$ and $L(x) = A(V(x)+1) - x - 1$.

\begin{axdef}
  funA: \nat \fun \nat
\where
  \Label{disabled rule AVal} \forall x: \nat @ funA~x = (x * x + x) \div 2
\end{axdef}

\begin{zproof}[funA\$domainCheck]
prove;
\end{zproof}

\begin{zed}
  AWit == (\lambda x: \nat @ (x * x + x) \div 2)
\end{zed}

\begin{theorem}{AWitFun}
  AWit \in \nat \fun \nat
\end{theorem}

\begin{zproof}
invoke AWit;
apply FunIntro to predicate (\lambda x: \nat @ (x * x + x) \div 2) \in \nat \fun \nat;
prove by reduce;
with enabled (RelPow) prove by reduce;
\end{zproof}

\begin{theorem}{rule AWitVal}
  \forall x: \nat @ AWit~x = (x * x + x) \div 2
\end{theorem}

\begin{zproof}
prove;
use AWitFun;
use pairInFunction[\nat, \nat][f := AWit, x := x, y := (x * x + x) \div 2];
invoke AWit;
prove by reduce;
\end{zproof}

\begin{theorem}{AEx}
  \exists funA: \nat \fun \nat @ \forall x: \nat @ funA~x = (x * x + x) \div 2
\end{theorem}

\begin{zproof}
instantiate funA == AWit;
use AWitFun;
prove;
\end{zproof}

\begin{zproof}[AWit\$domainCheck]
prove;
\end{zproof}

\begin{zed}
  TwoASet == \{ n: \nat | 2 * funA~n = n * n + n \}
\end{zed}

\begin{zproof}[TwoASet\$domainCheck]
prove;
\end{zproof}

\begin{theorem}{rule TwoABase}
  0 \in TwoASet
\end{theorem}

\begin{zproof}
invoke TwoASet;
prove;
use DivMod[x := 0, d := 2, q := 0, r := 0];
with enabled (AVal) prove;
\end{zproof}

\begin{theorem}{rule TwoAStep}
  \forall n: TwoASet @ 1 + n \in TwoASet
\end{theorem}

\begin{zproof}
invoke TwoASet;
prove;
use DivMod[x := (1 + n) * (1 + n) + (1 + n), d := 2, q := funA~n + n + 1, r := 0];
with enabled (AVal) prove;
\end{zproof}

\begin{theorem}{TwoASub}
  \nat \subseteq TwoASet
\end{theorem}

\begin{zproof}
apply natInduction to predicate \nat \subseteq TwoASet;
prove;
invoke TwoASet;
prove;
\end{zproof}

\begin{theorem}{TwoAAll}
  \forall n: \nat @ 2 * funA~n = n * n + n
\end{theorem}

\begin{zproof}
prove;
use TwoASub;
invoke TwoASet;
with enabled (inPower) prove;
instantiate e == n;
prove;
\end{zproof}

\begin{theorem}{TwoA}
  \forall x: \nat @ 2 * funA~x = x * x + x
\end{theorem}

\begin{zproof}
prove;
use TwoAAll[n := x];
prove;
\end{zproof}

\begin{theorem}{ASuc}
  \forall x: \nat @ funA~(1 + x) = funA~x + x + 1
\end{theorem}

\begin{zproof}
prove;
use TwoA[x := x];
use TwoA[x := 1 + x];
use UFNonneg[f := funA, x := x];
use UFNonneg[f := funA, x := 1 + x];
prove;
\end{zproof}

\begin{theorem}{AZero}
  funA~0 = 0 \land funA~1 = 1
\end{theorem}

\begin{zproof}
use TwoA[x := 0];
use TwoA[x := 1];
use UFNonneg[f := funA, x := 0];
use UFNonneg[f := funA, x := 1];
prove;
\end{zproof}

\begin{theorem}{AMono}
  \forall m, n: \nat | m \leq n @ funA~m \leq funA~n
\end{theorem}

\begin{zproof}
prove;
use TwoA[x := m];
use TwoA[x := n];
use timesMonotonic2[a := m, b := m, c := n, d := n];
use UFNonneg[f := funA, x := m];
use UFNonneg[f := funA, x := n];
prove;
\end{zproof}

\begin{theorem}{AIndexUq}
  \forall r, s, x: \nat | funA~r \leq x \land x < funA~(1 + r) \land funA~s \leq x \land x < funA~(1 + s)
    @ r = s
\end{theorem}

\begin{zproof}
prove;
split r < s;
cases;
use AMono[m := 1 + r, n := s];
prove;
next;
split s < r;
cases;
use AMono[m := 1 + s, n := r];
prove;
next;
prove;
next;
\end{zproof}

\begin{zed}
  AIdxSet == \{ x: \nat | \exists t: \nat @ funA~t \leq x \land x < funA~(1 + t) \}
\end{zed}

\begin{zproof}[AIdxSet\$domainCheck]
prove;
\end{zproof}

\begin{theorem}{rule AIdxBase}
  0 \in AIdxSet
\end{theorem}

\begin{zproof}
invoke AIdxSet;
prove;
use AZero;
instantiate t == 0;
prove;
\end{zproof}

\begin{theorem}{rule AIdxStep}
  \forall x: AIdxSet @ 1 + x \in AIdxSet
\end{theorem}

\begin{zproof}
invoke AIdxSet;
prove;
use ASuc[x := 1 + t];
use UFNonneg[f := funA, x := t];
use UFNonneg[f := funA, x := 1 + t];
use UFNonneg[f := funA, x := 2 + t];
use LeOrGt[a := funA~(1 + t), b := 1 + x];
split funA~(1 + t) \leq 1 + x;
cases;
instantiate t\_\_0 == 1 + t;
prove;
next;
instantiate t\_\_1 == t;
prove;
next;
\end{zproof}

\begin{theorem}{AIdxSub}
  \nat \subseteq AIdxSet
\end{theorem}

\begin{zproof}
apply natInduction to predicate \nat \subseteq AIdxSet;
prove;
invoke AIdxSet;
prove;
\end{zproof}

\begin{theorem}{AIndexEx}
  \forall x: \nat @ \exists t: \nat @ funA~t \leq x \land x < funA~(1 + t)
\end{theorem}

\begin{zproof}
prove;
use AIdxSub;
with enabled (inPower) prove;
instantiate e == x;
invoke AIdxSet;
prove;
\end{zproof}

\begin{axdef}
  funV: \nat \fun \nat
\where
  \Label{VSpec} \forall x: \nat @ funA~(funV~x) \leq x \land x < funA~(1 + funV~x)
\end{axdef}

\begin{zproof}[funV\$domainCheck]
prove;
\end{zproof}

\begin{zed}
  VWit == \{ x, t: \nat | funA~t \leq x \land x < funA~(1 + t) @ (x, t) \}
\end{zed}

\begin{zproof}[VWit\$domainCheck]
prove;
\end{zproof}

\begin{theorem}{VWitFun}
  VWit \in \nat \fun \nat
\end{theorem}

\begin{zproof}
invoke VWit;
apply FunIntro to predicate \{ x, t: \nat | funA~t \leq x \land x < funA~(1 + t) @ (x, t) \}
  \in \nat \fun \nat;
with enabled (RelPow) prove by reduce;
use AIndexEx[x := x];
use AIndexUq[r := y, s := z, x := x\_\_0];
prove;
instantiate y\_\_0 == t;
prove;
\end{zproof}

\begin{theorem}{rule VWitLe}
  \forall x: \nat @ funA~(VWit~x) \leq x
\end{theorem}

\begin{zproof}
prove;
use AIndexEx[x := x];
use VWitFun;
prove;
use pairInFunction[\nat, \nat][f := VWit, x := x, y := t];
invoke VWit;
prove by reduce;
\end{zproof}

\begin{theorem}{rule VWitLt}
  \forall x: \nat @ x < funA~(1 + VWit~x)
\end{theorem}

\begin{zproof}
prove;
use AIndexEx[x := x];
use VWitFun;
prove;
use pairInFunction[\nat, \nat][f := VWit, x := x, y := t];
invoke VWit;
prove by reduce;
\end{zproof}

\begin{theorem}{VEx}
  \exists funV: \nat \fun \nat @ \forall x: \nat @ funA~(funV~x) \leq x \land x < funA~(1 + funV~x)
\end{theorem}

\begin{zproof}
instantiate funV == VWit;
use VWitFun;
prove;
\end{zproof}

\begin{theorem}{VChar}
  \forall t, x: \nat | funA~t \leq x \land x < funA~(1 + t) @ funV~x = t
\end{theorem}

\begin{zproof}
prove;
use VSpec[x := x];
use AIndexUq[r := funV~x, s := t, x := x];
use UFNonneg[f := funV, x := x];
prove;
\end{zproof}

\begin{axdef}
  funK: \nat \fun \nat
\where
  \Label{disabled rule KVal} \forall x: \nat @ funK~x = x - funA~(funV~x)
\end{axdef}

\begin{zproof}[funK\$domainCheck]
prove;
\end{zproof}

\begin{zed}
  KWit == (\lambda x: \nat @ x - funA~(funV~x))
\end{zed}

\begin{zproof}[KWit\$domainCheck]
prove;
\end{zproof}

\begin{theorem}{KWitRel}
  (\lambda x: \nat @ x - funA~(funV~x)) \in \power (\nat \cross \nat)
\end{theorem}

\begin{zproof}
with enabled (inPower) prove by reduce;
use VSpec[x := x];
prove;
\end{zproof}

\begin{theorem}{KWitFun}
  KWit \in \nat \fun \nat
\end{theorem}

\begin{zproof}
invoke KWit;
use KWitRel;
apply FunIntro to predicate (\lambda x: \nat @ x - funA~(funV~x)) \in \nat \fun \nat;
with enabled (RelPow) prove by reduce;
use VSpec[x := x];
prove;
\end{zproof}

\begin{theorem}{rule KWitVal}
  \forall x: \nat @ KWit~x = x - funA~(funV~x)
\end{theorem}

\begin{zproof}
prove;
use KWitFun;
use pairInFunction[\nat, \nat][f := KWit, x := x, y := x - funA~(funV~x)];
invoke KWit;
prove by reduce;
\end{zproof}

\begin{theorem}{KEx}
  \exists funK: \nat \fun \nat @ \forall x: \nat @ funK~x = x - funA~(funV~x)
\end{theorem}

\begin{zproof}
instantiate funK == KWit;
use KWitFun;
prove;
\end{zproof}

\begin{axdef}
  funL: \nat \fun \nat
\where
  \Label{disabled rule LVal} \forall x: \nat @ funL~x = funA~(1 + funV~x) - x - 1
\end{axdef}

\begin{zproof}[funL\$domainCheck]
prove;
\end{zproof}

\begin{zed}
  LWit == (\lambda x: \nat @ funA~(1 + funV~x) - x - 1)
\end{zed}

\begin{zproof}[LWit\$domainCheck]
prove;
\end{zproof}

\begin{theorem}{LWitRel}
  (\lambda x: \nat @ funA~(1 + funV~x) - x - 1) \in \power (\nat \cross \nat)
\end{theorem}

\begin{zproof}
with enabled (inPower) prove by reduce;
use VSpec[x := x];
prove;
\end{zproof}

\begin{theorem}{LWitFun}
  LWit \in \nat \fun \nat
\end{theorem}

\begin{zproof}
invoke LWit;
use LWitRel;
apply FunIntro to predicate (\lambda x: \nat @ funA~(1 + funV~x) - x - 1) \in \nat \fun \nat;
with enabled (RelPow) prove by reduce;
use VSpec[x := x];
prove;
\end{zproof}

\begin{theorem}{rule LWitVal}
  \forall x: \nat @ LWit~x = funA~(1 + funV~x) - x - 1
\end{theorem}

\begin{zproof}
prove;
use LWitFun;
use pairInFunction[\nat, \nat][f := LWit, x := x, y := funA~(1 + funV~x) - x - 1];
invoke LWit;
prove by reduce;
\end{zproof}

\begin{theorem}{LEx}
  \exists funL: \nat \fun \nat @ \forall x: \nat @ funL~x = funA~(1 + funV~x) - x - 1
\end{theorem}

\begin{zproof}
instantiate funL == LWit;
use LWitFun;
prove;
\end{zproof}

Values on a diagonal of the Cantor table: for $k \le t$, $V(A(t)+k) = t$, $K(A(t)+k) = k$ and
$L(A(t)+k) = t - k$.

\begin{theorem}{Diag}
  \forall k, t: \nat | k \leq t @
    funV~(funA~t + k) = t \land funK~(funA~t + k) = k \land funL~(funA~t + k) = t - k
\end{theorem}

\begin{zproof}
prove;
use ASuc[x := t];
use UFNonneg[f := funA, x := t];
use VChar[t := t, x := funA~t + k];
with enabled (KVal, LVal) prove;
\end{zproof}

\begin{theorem}{Decomp}
  \forall x: \nat @ x = funA~(funV~x) + funK~x \land funK~x \leq funV~x
\end{theorem}

\begin{zproof}
prove;
use VSpec[x := x];
use ASuc[x := funV~x];
use UFNonneg[f := funA, x := funV~x];
use UFNonneg[f := funA, x := 1 + funV~x];
use UFNonneg[f := funV, x := x];
with enabled (KVal) prove;
\end{zproof}

\subsection{$A$, $V$, $K$ and $L$ are in $\mathrm{clos}\,\mathcal{F}$}

\begin{theorem}{AExpr}
  comp(funHf, opPlusI~funSq) = funA
\end{theorem}

\begin{zproof}
apply UFExt to predicate comp(funHf, opPlusI~funSq) = funA;
prove;
use TwoA[x := x];
use UFNonneg[f := funA, x := x];
use DivMod[x := x * x + x, d := 2, q := funA~x, r := 0];
prove;
\end{zproof}

\begin{theorem}{RtPos}
  \forall y: \nat | 1 \leq y @ 1 \leq funRt~y
\end{theorem}

\begin{zproof}
prove;
use RtSpec[x := y];
use UFNonneg[f := funRt, x := y];
split funRt~y = 0;
cases;
prove;
next;
prove;
next;
\end{zproof}

\begin{theorem}{PhatPos}
  \forall y: \nat | 1 \leq y @ exprPhat~y = y - 1
\end{theorem}

\begin{zproof}
prove;
use PhatVal[x := y - 1];
prove;
\end{zproof}

\begin{theorem}{VFormula}
  \forall x: \nat @ funV~x = (funRt~(1 + 8 * x) - 1) \div 2
\end{theorem}

\begin{zproof}
prove;
use VSpec[x := x];
use TwoA[x := funV~x];
use TwoA[x := 1 + funV~x];
use UFNonneg[f := funA, x := funV~x];
use UFNonneg[f := funA, x := 1 + funV~x];
use UFNonneg[f := funV, x := x];
use LeOrGt[a := (2 * funV~x + 2) * (2 * funV~x + 2), b := 1 + 8 * x];
split (2 * funV~x + 2) * (2 * funV~x + 2) \leq 1 + 8 * x;
cases;
use RtI[r := 2 * funV~x + 2, x := 1 + 8 * x];
prove;
use DivMod[x := 2 * funV~x + 1, d := 2, q := funV~x, r := 1];
prove;
next;
use RtI[r := 2 * funV~x + 1, x := 1 + 8 * x];
prove;
use DivMod[x := 2 * funV~x, d := 2, q := funV~x, r := 0];
prove;
next;
\end{zproof}

\begin{theorem}{VExpr}
  comp(comp(comp(comp(comp(comp(funHf, exprPhat), funRt), funS), funD), funD), funD) = funV
\end{theorem}

\begin{zproof}
apply UFExt to predicate comp(comp(comp(comp(comp(comp(funHf, exprPhat), funRt), funS), funD), funD), funD) = funV;
prove;
use VFormula[x := x];
use RtPos[y := 1 + 8 * x];
use PhatPos[y := funRt~(1 + 8 * x)];
prove;
\end{zproof}

\begin{theorem}{IInExpr}
  opBox~funS \in C6
\end{theorem}

\begin{zproof}
prove;
\end{zproof}

\begin{theorem}{rule IIn}
  funI \in C6
\end{theorem}

\begin{zproof}
use IInExpr;
use BoxSExpr;
prove;
\end{zproof}
\begin{theorem}{AInExpr}
  comp(funHf, opPlusI~funSq) \in C6
\end{theorem}

\begin{zproof}
prove;
\end{zproof}

\begin{theorem}{rule AIn}
  funA \in C6
\end{theorem}

\begin{zproof}
use AInExpr;
use AExpr;
prove;
\end{zproof}
\begin{theorem}{VInExpr}
  comp(comp(comp(comp(comp(comp(funHf, exprPhat), funRt), funS), funD), funD), funD) \in C6
\end{theorem}

\begin{zproof}
prove;
\end{zproof}

\begin{theorem}{rule VIn}
  funV \in C6
\end{theorem}

\begin{zproof}
use VInExpr;
use VExpr;
prove;
\end{zproof}

\begin{theorem}{KExpr}
  opDiffE(funI, comp(funA, funV)) = funK
\end{theorem}

\begin{zproof}
apply UFExt to predicate opDiffE(funI, comp(funA, funV)) = funK;
prove;
use VSpec[x := x];
use DiffEIsDifference[F := funI, G := comp(funA, funV), x := x];
with enabled (KVal) prove;
\end{zproof}

\begin{theorem}{LExpr}
  opDiffE(comp(funA, comp(funS, funV)), funS) = funL
\end{theorem}

\begin{zproof}
apply UFExt to predicate opDiffE(comp(funA, comp(funS, funV)), funS) = funL;
prove;
use VSpec[x := x];
use DiffEIsDifference[F := comp(funA, comp(funS, funV)), G := funS, x := x];
with enabled (LVal) prove;
\end{zproof}

\begin{theorem}{KInExpr}
  opDiffE(funI, comp(funA, funV)) \in C6
\end{theorem}

\begin{zproof}
prove;
\end{zproof}

\begin{theorem}{rule KIn}
  funK \in C6
\end{theorem}

\begin{zproof}
use KInExpr;
use KExpr;
prove;
\end{zproof}
\begin{theorem}{LInExpr}
  opDiffE(comp(funA, comp(funS, funV)), funS) \in C6
\end{theorem}

\begin{zproof}
prove;
\end{zproof}

\begin{theorem}{rule LIn}
  funL \in C6
\end{theorem}

\begin{zproof}
use LInExpr;
use LExpr;
prove;
\end{zproof}

\section{Theorem 6.8: the bases with $K$ and $L$}

\subsection{Auxiliary functions (Robinson [16, p.~664])}

\begin{theorem}{KAt}
  \forall m, t, k: \nat | m = funA~t + k \land k \leq t @ funK~m = k
\end{theorem}

\begin{zproof}
prove;
use Diag[k := k, t := t];
prove;
\end{zproof}

\begin{theorem}{RtLe}
  \forall x: \nat @ funRt~x \leq x
\end{theorem}

\begin{zproof}
prove;
use RtSpec[x := x];
use LeSq[n := funRt~x];
prove;
\end{zproof}

$Z(x) = x(x+3)/2 = A(x) + x$ and $Y(x) = 2x - \mathrm{Rt}(x)$.

\begin{axdef}
  funZ: \nat \fun \nat
\where
  \Label{rule ZVal} \forall x: \nat @ funZ~x = funA~x + x
\end{axdef}

\begin{zproof}[funZ\$domainCheck]
prove;
\end{zproof}

\begin{zed}
  ZWit == (\lambda x: \nat @ funA~x + x)
\end{zed}

\begin{zproof}[ZWit\$domainCheck]
prove;
\end{zproof}

\begin{theorem}{ZWitFun}
  ZWit \in \nat \fun \nat
\end{theorem}

\begin{zproof}
invoke ZWit;
apply FunIntro to predicate (\lambda x: \nat @ funA~x + x) \in \nat \fun \nat;
prove by reduce;
with enabled (RelPow) prove by reduce;
\end{zproof}

\begin{theorem}{rule ZWitVal}
  \forall x: \nat @ ZWit~x = funA~x + x
\end{theorem}

\begin{zproof}
prove;
use ZWitFun;
use pairInFunction[\nat, \nat][f := ZWit, x := x, y := funA~x + x];
invoke ZWit;
prove by reduce;
\end{zproof}

\begin{theorem}{ZEx}
  \exists funZ: \nat \fun \nat @ \forall x: \nat @ funZ~x = funA~x + x
\end{theorem}

\begin{zproof}
instantiate funZ == ZWit;
use ZWitFun;
prove;
\end{zproof}

\begin{axdef}
  funY: \nat \fun \nat
\where
  \Label{disabled rule YVal} \forall x: \nat @ funY~x = 2 * x - funRt~x
\end{axdef}

\begin{zproof}[funY\$domainCheck]
prove;
\end{zproof}

\begin{zed}
  YWit == (\lambda x: \nat @ 2 * x - funRt~x)
\end{zed}

\begin{zproof}[YWit\$domainCheck]
prove;
\end{zproof}

\begin{theorem}{YWitRel}
  (\lambda x: \nat @ 2 * x - funRt~x) \in \power (\nat \cross \nat)
\end{theorem}

\begin{zproof}
with enabled (inPower) prove by reduce;
use RtLe[x := x];
prove;
\end{zproof}

\begin{theorem}{YWitFun}
  YWit \in \nat \fun \nat
\end{theorem}

\begin{zproof}
invoke YWit;
use YWitRel;
apply FunIntro to predicate (\lambda x: \nat @ 2 * x - funRt~x) \in \nat \fun \nat;
with enabled (RelPow) prove by reduce;
use RtLe[x := x];
prove;
\end{zproof}

\begin{theorem}{rule YWitVal}
  \forall x: \nat @ YWit~x = 2 * x - funRt~x
\end{theorem}

\begin{zproof}
prove;
use YWitFun;
use pairInFunction[\nat, \nat][f := YWit, x := x, y := 2 * x - funRt~x];
invoke YWit;
prove by reduce;
\end{zproof}

\begin{theorem}{YEx}
  \exists funY: \nat \fun \nat @ \forall x: \nat @ funY~x = 2 * x - funRt~x
\end{theorem}

\begin{zproof}
instantiate funY == YWit;
use YWitFun;
prove;
\end{zproof}

\begin{zed}
  ZBoxSet == \{ n: \nat | opBox~(comp(comp(funS, funS), opPlusI~funK))~n = funZ~n \}
\end{zed}

\begin{zproof}[ZBoxSet\$domainCheck]
prove;
\end{zproof}

\begin{theorem}{rule ZBoxBase}
  0 \in ZBoxSet
\end{theorem}

\begin{zproof}
invoke ZBoxSet;
prove;
use AZero;
prove;
\end{zproof}

\begin{theorem}{rule ZBoxStep}
  \forall n: ZBoxSet @ 1 + n \in ZBoxSet
\end{theorem}

\begin{zproof}
invoke ZBoxSet;
prove;
use Diag[k := n, t := n];
use ASuc[x := n];
prove;
\end{zproof}

\begin{theorem}{ZBoxSub}
  \nat \subseteq ZBoxSet
\end{theorem}

\begin{zproof}
apply natInduction to predicate \nat \subseteq ZBoxSet;
prove;
invoke ZBoxSet;
prove;
\end{zproof}

\begin{theorem}{ZBoxAll}
  \forall n: \nat @ opBox~(comp(comp(funS, funS), opPlusI~funK))~n = funZ~n
\end{theorem}

\begin{zproof}
prove;
use ZBoxSub;
invoke ZBoxSet;
with enabled (inPower) prove;
instantiate e == n;
prove;
\end{zproof}

\begin{theorem}{ZExpr}
  opBox~(comp(comp(funS, funS), opPlusI~funK)) = funZ
\end{theorem}

\begin{zproof}
apply UFExt to predicate opBox~(comp(comp(funS, funS), opPlusI~funK)) = funZ;
prove;
use ZBoxAll[n := x];
prove;
\end{zproof}

Robinson's formula: $K(Z(a+b) + 2a + 3) = a - b$ if $a \ge b$.

\begin{theorem}{RobDiffK}
  \forall a, b: \nat | b \leq a @ funK~(funZ~(a + b) + 2 * a + 3) = a - b
\end{theorem}

\begin{zproof}
prove;
use ASuc[x := a + b];
use ASuc[x := 1 + a + b];
use KAt[m := funA~(a + b) + a + b + 2 * a + 3, t := a + b + 2, k := a - b];
prove;
\end{zproof}

$I - F = K\,S\,S\,S\,((Z F^+)^+)^+$ where $F(x) \le x$, and
$D - F = K\,S\,S\,S\,((((Z (F^+)^+)^+)^+)^+)^+$ where $F(x) \le 2x$.

\begin{zed}
  exprIminusF == (\lambda F: \nat \fun \nat @ comp(comp(comp(comp(funK, funS), funS), funS), opPlusI~(opPlusI~(comp(funZ, opPlusI~F)))))
\end{zed}

\begin{zproof}[exprIminusF\$domainCheck]
prove;
\end{zproof}

\begin{theorem}{rule exprIminusFFun}
  exprIminusF \in (\nat \fun \nat) \fun (\nat \fun \nat)
\end{theorem}

\begin{zproof}
invoke exprIminusF;
apply FunIntro to predicate (\lambda F: \nat \fun \nat @ comp(comp(comp(comp(funK, funS), funS), funS), opPlusI~(opPlusI~(comp(funZ, opPlusI~F))))) \in (\nat \fun \nat) \fun (\nat \fun \nat);
with enabled (RelPow, inPower) prove by reduce;
\end{zproof}

\begin[disabled]{theorem}{rule exprIminusFApp}
  \forall F: \nat \fun \nat @ exprIminusF~F = comp(comp(comp(comp(funK, funS), funS), funS), opPlusI~(opPlusI~(comp(funZ, opPlusI~F))))
\end{theorem}

\begin{zproof}
use exprIminusFFun;
use pairInFunction[\nat \fun \nat, \nat \fun \nat][f := exprIminusF, x := F, y := comp(comp(comp(comp(funK, funS), funS), funS), opPlusI~(opPlusI~(comp(funZ, opPlusI~F))))];
invoke exprIminusF;
prove by reduce;
\end{zproof}

\begin{zed}
  exprDminusF == (\lambda F: \nat \fun \nat @ comp(comp(comp(comp(funK, funS), funS), funS), opPlusI~(opPlusI~(opPlusI~(opPlusI~(comp(funZ, opPlusI~(opPlusI~F))))))))
\end{zed}

\begin{zproof}[exprDminusF\$domainCheck]
prove;
\end{zproof}

\begin{theorem}{rule exprDminusFFun}
  exprDminusF \in (\nat \fun \nat) \fun (\nat \fun \nat)
\end{theorem}

\begin{zproof}
invoke exprDminusF;
apply FunIntro to predicate (\lambda F: \nat \fun \nat @ comp(comp(comp(comp(funK, funS), funS), funS), opPlusI~(opPlusI~(opPlusI~(opPlusI~(comp(funZ, opPlusI~(opPlusI~F)))))))) \in (\nat \fun \nat) \fun (\nat \fun \nat);
with enabled (RelPow, inPower) prove by reduce;
\end{zproof}

\begin[disabled]{theorem}{rule exprDminusFApp}
  \forall F: \nat \fun \nat @ exprDminusF~F = comp(comp(comp(comp(funK, funS), funS), funS), opPlusI~(opPlusI~(opPlusI~(opPlusI~(comp(funZ, opPlusI~(opPlusI~F)))))))
\end{theorem}

\begin{zproof}
use exprDminusFFun;
use pairInFunction[\nat \fun \nat, \nat \fun \nat][f := exprDminusF, x := F, y := comp(comp(comp(comp(funK, funS), funS), funS), opPlusI~(opPlusI~(opPlusI~(opPlusI~(comp(funZ, opPlusI~(opPlusI~F)))))))];
invoke exprDminusF;
prove by reduce;
\end{zproof}

\begin{theorem}{IminusFVal}
  \forall F: \nat \fun \nat; x: \nat | F~x \leq x @ exprIminusF~F~x = x - F~x
\end{theorem}

\begin{zproof}
with enabled (exprIminusFApp) prove;
use RobDiffK[a := x, b := F~x];
prove;
\end{zproof}

\begin{theorem}{DminusFVal}
  \forall F: \nat \fun \nat; x: \nat | F~x \leq 2 * x @ exprDminusF~F~x = 2 * x - F~x
\end{theorem}

\begin{zproof}
with enabled (exprDminusFApp) prove;
use RobDiffK[a := 2 * x, b := F~x];
prove;
\end{zproof}

$Y = ((\mathrm{Sgn}\,K)^+ S)^\square$: $Y(n+1) = Y(n) + 1 + \mathrm{Sgn}\,K(Y(n)+1)$, and
$K(Y(n)+1) = 0$ exactly when $n+1$ is a square ($Y(n)+1 = 2n - r + 1$ with $r = \mathrm{Rt}(n)$
lies on the diagonal of $A(2r+1)$, $A(2r)$ or $A(2r-1)$).

\begin{theorem}{YStep}
  \forall n: \nat @ opPlusI~(comp(funSgn, funK))~(funS~(funY~n)) = funY~(1 + n)
\end{theorem}

\begin{zproof}
with enabled (YVal) prove;
use RtSpec[x := n];
use RtLe[x := n];
use UFNonneg[f := funRt, x := n];
use LeOrGt[a := (1 + funRt~n) * (1 + funRt~n), b := 1 + n];
split (1 + funRt~n) * (1 + funRt~n) \leq 1 + n;
cases;
use RtI[r := 1 + funRt~n, x := 1 + n];
use TwoA[x := 1 + 2 * funRt~n];
use UFNonneg[f := funA, x := 1 + 2 * funRt~n];
use KAt[m := 1 + 2 * n - funRt~n, t := 1 + 2 * funRt~n, k := 0];
prove;
next;
use RtI[r := funRt~n, x := 1 + n];
use LeOrGt[a := funRt~n * funRt~n + funRt~n, b := n];
split funRt~n * funRt~n + funRt~n \leq n;
cases;
use TwoA[x := 2 * funRt~n];
use UFNonneg[f := funA, x := 2 * funRt~n];
use KAt[m := 1 + 2 * n - funRt~n, t := 2 * funRt~n, k := 1 + 2 * n - 2 * funRt~n * funRt~n - 2 * funRt~n];
prove;
next;
use TwoA[x := 2 * funRt~n - 1];
use UFNonneg[f := funA, x := 2 * funRt~n - 1];
use KAt[m := 1 + 2 * n - funRt~n, t := 2 * funRt~n - 1, k := 1 + 2 * n - 2 * funRt~n * funRt~n];
prove;
next;
\end{zproof}

\begin{zed}
  YBoxSet == \{ n: \nat | opBox~(comp(opPlusI~(comp(funSgn, funK)), funS))~n = funY~n \}
\end{zed}

\begin{zproof}[YBoxSet\$domainCheck]
prove;
\end{zproof}

\begin{theorem}{rule YBoxBase}
  0 \in YBoxSet
\end{theorem}

\begin{zproof}
invoke YBoxSet;
prove;
use RtSq[m := 0];
with enabled (YVal) prove;
\end{zproof}

\begin{theorem}{rule YBoxStep}
  \forall n: YBoxSet @ 1 + n \in YBoxSet
\end{theorem}

\begin{zproof}
invoke YBoxSet;
prove;
use YStep[n := n];
prove;
\end{zproof}

\begin{theorem}{YBoxSub}
  \nat \subseteq YBoxSet
\end{theorem}

\begin{zproof}
apply natInduction to predicate \nat \subseteq YBoxSet;
prove;
invoke YBoxSet;
prove;
\end{zproof}

\begin{theorem}{YBoxAll}
  \forall n: \nat @ opBox~(comp(opPlusI~(comp(funSgn, funK)), funS))~n = funY~n
\end{theorem}

\begin{zproof}
prove;
use YBoxSub;
invoke YBoxSet;
with enabled (inPower) prove;
instantiate e == n;
prove;
\end{zproof}

\begin{theorem}{YExpr}
  opBox~(comp(opPlusI~(comp(funSgn, funK)), funS)) = funY
\end{theorem}

\begin{zproof}
apply UFExt to predicate opBox~(comp(opPlusI~(comp(funSgn, funK)), funS)) = funY;
prove;
use YBoxAll[n := x];
prove;
\end{zproof}

$\mathrm{Rt} = D - Y$, $\mathrm{Sq} = (((\bar 1^+)^+ \mathrm{Rt})^+)^\square$ and $E = I - \mathrm{Sq}\,\mathrm{Rt}$.

\begin{theorem}{RtExprK}
  exprDminusF~funY = funRt
\end{theorem}

\begin{zproof}
apply UFExt to predicate exprDminusF~funY = funRt;
prove;
use RtLe[x := x];
use UFNonneg[f := funRt, x := x];
use DminusFVal[F := funY, x := x];
with enabled (YVal) prove;
\end{zproof}

\begin{zed}
  SqKSet == \{ n: \nat | opBox~(opPlusI~(comp(opPlusI~(opPlusI~(cst~1)), funRt)))~n = n * n \}
\end{zed}

\begin{zproof}[SqKSet\$domainCheck]
prove;
\end{zproof}

\begin{theorem}{rule SqKBase}
  0 \in SqKSet
\end{theorem}

\begin{zproof}
invoke SqKSet;
prove;
\end{zproof}

\begin{theorem}{rule SqKStep}
  \forall n: SqKSet @ 1 + n \in SqKSet
\end{theorem}

\begin{zproof}
invoke SqKSet;
prove;
equality substitute opBox~(opPlusI~(comp(opPlusI~(opPlusI~(cst~1)), funRt)))~n;
rewrite;
\end{zproof}

\begin{theorem}{SqKSub}
  \nat \subseteq SqKSet
\end{theorem}

\begin{zproof}
apply natInduction to predicate \nat \subseteq SqKSet;
prove;
invoke SqKSet;
prove;
\end{zproof}

\begin{theorem}{SqKAll}
  \forall n: \nat @ opBox~(opPlusI~(comp(opPlusI~(opPlusI~(cst~1)), funRt)))~n = n * n
\end{theorem}

\begin{zproof}
prove;
use SqKSub;
invoke SqKSet;
with enabled (inPower) prove;
instantiate e == n;
prove;
\end{zproof}

\begin{theorem}{SqExprK}
  opBox~(opPlusI~(comp(opPlusI~(opPlusI~(cst~1)), funRt))) = funSq
\end{theorem}

\begin{zproof}
apply UFExt to predicate opBox~(opPlusI~(comp(opPlusI~(opPlusI~(cst~1)), funRt))) = funSq;
prove;
use SqKAll[n := x];
prove;
\end{zproof}

\begin{theorem}{EExprK}
  exprIminusF~(comp(funSq, funRt)) = funE
\end{theorem}

\begin{zproof}
apply UFExt to predicate exprIminusF~(comp(funSq, funRt)) = funE;
prove;
use RtSpec[x := x];
use IminusFVal[F := comp(funSq, funRt), x := x];
with enabled (EVal) prove;
\end{zproof}

\subsection{The basis $\mathcal{F}' = \langle \bar 1, K, F^+, FG, F^\square\rangle$}

\begin{zed}
  C6K == clos~\{ cst~1, funK \}~\{ opPlusI, opBox \}~\{ comp \}
\end{zed}

\begin{zproof}[C6K\$domainCheck]
prove;
\end{zproof}

\begin[disabled]{theorem}{rule C6KDef}
  C6K = clos~\{ cst~1, funK \}~\{ opPlusI, opBox \}~\{ comp \}
\end{theorem}

\begin{zproof}
invoke C6K;
prove;
\end{zproof}

\begin{theorem}{C6KInPower}
  C6K \in \power (\nat \fun \nat)
\end{theorem}

\begin{zproof}
use ClosInPow[I := \{ cst~1, funK \}, U := \{ opPlusI, opBox \}, B := \{ comp \}];
with enabled (C6KDef) prove;
\end{zproof}

\begin[disabled]{theorem}{rule C6KUF}
  \forall F: \power (\num \cross \num) | F \in C6K @ F \in \nat \fun \nat
\end{theorem}

\begin{zproof}
prove;
use ClosElemUF[I := \{ cst~1, funK \}, U := \{ opPlusI, opBox \}, B := \{ comp \}, F := F];
with enabled (C6KDef) prove;
\end{zproof}

\begin{theorem}{rule C6KInit0}
  cst~1 \in C6K
\end{theorem}

\begin{zproof}
with enabled (C6KDef) prove;
\end{zproof}

\begin{theorem}{rule C6KInit1}
  funK \in C6K
\end{theorem}

\begin{zproof}
with enabled (C6KDef) prove;
\end{zproof}

\begin{theorem}{rule C6KUnop0}
  \forall F: \nat \fun \nat | F \in C6K @ opPlusI~F \in C6K
\end{theorem}

\begin{zproof}
with enabled (C6KDef) prove;
\end{zproof}

\begin{theorem}{rule C6KUnop1}
  \forall F: \nat \fun \nat | F \in C6K @ opBox~F \in C6K
\end{theorem}

\begin{zproof}
with enabled (C6KDef) prove;
\end{zproof}

\begin{theorem}{rule C6KBinop0}
  \forall F, G: \nat \fun \nat | F \in C6K \land G \in C6K @ comp(F, G) \in C6K
\end{theorem}

\begin{zproof}
with enabled (C6KDef) prove;
\end{zproof}

\begin{theorem}{SInKExpr}
  opPlusI~(cst~1) \in C6K
\end{theorem}

\begin{zproof}
prove;
\end{zproof}

\begin{theorem}{rule SInK}
  funS \in C6K
\end{theorem}

\begin{zproof}
use SInKExpr;
use SExpr;
prove;
\end{zproof}
\begin{theorem}{SgnInKExpr}
  opBox~(cst~1) \in C6K
\end{theorem}

\begin{zproof}
prove;
\end{zproof}

\begin{theorem}{rule SgnInK}
  funSgn \in C6K
\end{theorem}

\begin{zproof}
use SgnInKExpr;
use SgnExpr;
prove;
\end{zproof}
\begin{theorem}{ZInKExpr}
  opBox~(comp(comp(funS, funS), opPlusI~funK)) \in C6K
\end{theorem}

\begin{zproof}
prove;
\end{zproof}

\begin{theorem}{rule ZInK}
  funZ \in C6K
\end{theorem}

\begin{zproof}
use ZInKExpr;
use ZExpr;
prove;
\end{zproof}
\begin{theorem}{YInKExpr}
  opBox~(comp(opPlusI~(comp(funSgn, funK)), funS)) \in C6K
\end{theorem}

\begin{zproof}
prove;
\end{zproof}

\begin{theorem}{rule YInK}
  funY \in C6K
\end{theorem}

\begin{zproof}
use YInKExpr;
use YExpr;
prove;
\end{zproof}
\begin{theorem}{RtInKExpr}
  exprDminusF~funY \in C6K
\end{theorem}

\begin{zproof}
with enabled (exprDminusFApp) prove;
\end{zproof}

\begin{theorem}{rule RtInK}
  funRt \in C6K
\end{theorem}

\begin{zproof}
use RtInKExpr;
use RtExprK;
prove;
\end{zproof}
\begin{theorem}{SqInKExpr}
  opBox~(opPlusI~(comp(opPlusI~(opPlusI~(cst~1)), funRt))) \in C6K
\end{theorem}

\begin{zproof}
prove;
\end{zproof}

\begin{theorem}{rule SqInK}
  funSq \in C6K
\end{theorem}

\begin{zproof}
use SqInKExpr;
use SqExprK;
prove;
\end{zproof}
\begin{theorem}{EInKExpr}
  exprIminusF~(comp(funSq, funRt)) \in C6K
\end{theorem}

\begin{zproof}
with enabled (exprIminusFApp) prove;
\end{zproof}

\begin{theorem}{rule EInK}
  funE \in C6K
\end{theorem}

\begin{zproof}
use EInKExpr;
use EExprK;
prove;
\end{zproof}

\begin{theorem}{Theorem68KSub1U}
  C6 \in UClosed~\{ opPlusI, opBox \}
\end{theorem}

\begin{zproof}
use C6InPower;
with enabled (UClosedDef, C6UF) prove;
split u = opPlusI;
cases;
prove;
next;
prove;
next;
\end{zproof}

\begin{theorem}{Theorem68KSub1B}
  C6 \in BClosed~\{ comp \}
\end{theorem}

\begin{zproof}
use C6InPower;
with enabled (BClosedDef, C6UF) prove;
prove;
\end{zproof}

\begin{theorem}{Theorem68KSub1}
  C6K \subseteq C6
\end{theorem}

\begin{zproof}
use ClosLeast[I := \{ cst~1, funK \}, U := \{ opPlusI, opBox \}, B := \{ comp \}, Z := C6];
use C6InPower;
use Theorem68KSub1U;
use Theorem68KSub1B;
with enabled (C6KDef) prove;
\end{zproof}

\begin{theorem}{Theorem68KSub2U}
  C6K \in UClosed~\{ opPlusI, opBox \}
\end{theorem}

\begin{zproof}
use C6KInPower;
with enabled (UClosedDef, C6KUF) prove;
split u = opPlusI;
cases;
prove;
next;
prove;
next;
\end{zproof}

\begin{theorem}{Theorem68KSub2B}
  C6K \in BClosed~\{ comp \}
\end{theorem}

\begin{zproof}
use C6KInPower;
with enabled (BClosedDef, C6KUF) prove;
prove;
\end{zproof}

\begin{theorem}{Theorem68KSub2}
  C6 \subseteq C6K
\end{theorem}

\begin{zproof}
use ClosLeast[I := \{ cst~1, funE \}, U := \{ opPlusI, opBox \}, B := \{ comp \}, Z := C6K];
use C6KInPower;
use Theorem68KSub2U;
use Theorem68KSub2B;
with enabled (C6Def) prove;
\end{zproof}

\begin{theorem}{Theorem68K}
  C6K = C6
\end{theorem}

\begin{zproof}
use Theorem68KSub1;
use Theorem68KSub2;
use C6KInPower;
use C6InPower;
use SubAntisym[\nat \fun \nat][A := C6K, B := C6];
prove;
\end{zproof}

\subsection{The basis $\mathcal{F}'' = \langle \bar 1, L, F^+, FG, F^\square\rangle$}

\begin{theorem}{LAt}
  \forall m, t, k: \nat | m = funA~t + k \land k \leq t @ funL~m = t - k
\end{theorem}

\begin{zproof}
prove;
use Diag[k := k, t := t];
prove;
\end{zproof}

$K = L\,\bar 1^+ ((\bar 1^+)^+ L)^+$, i.e.\ $K(x) = L(x + 2L(x) + 2)$.

\begin{theorem}{KExprL}
  comp(comp(funL, opPlusI~(cst~1)), opPlusI~(comp(opPlusI~(opPlusI~(cst~1)), funL))) = funK
\end{theorem}

\begin{zproof}
apply UFExt to predicate comp(comp(funL, opPlusI~(cst~1)), opPlusI~(comp(opPlusI~(opPlusI~(cst~1)), funL))) = funK;
prove;
use Decomp[x := x];
use UFNonneg[f := funV, x := x];
use UFNonneg[f := funA, x := funV~x];
use ASuc[x := funV~x];
use LAt[m := x, t := funV~x, k := funK~x];
prove;
use UFNonneg[f := funK, x := x];
use LAt[m := 2 + x - 2 * funK~x + 2 * funV~x, t := 1 + funV~x, k := funV~x - funK~x + 1];
prove;
\end{zproof}

\begin{zed}
  C6L == clos~\{ cst~1, funL \}~\{ opPlusI, opBox \}~\{ comp \}
\end{zed}

\begin{zproof}[C6L\$domainCheck]
prove;
\end{zproof}

\begin[disabled]{theorem}{rule C6LDef}
  C6L = clos~\{ cst~1, funL \}~\{ opPlusI, opBox \}~\{ comp \}
\end{theorem}

\begin{zproof}
invoke C6L;
prove;
\end{zproof}

\begin{theorem}{C6LInPower}
  C6L \in \power (\nat \fun \nat)
\end{theorem}

\begin{zproof}
use ClosInPow[I := \{ cst~1, funL \}, U := \{ opPlusI, opBox \}, B := \{ comp \}];
with enabled (C6LDef) prove;
\end{zproof}

\begin[disabled]{theorem}{rule C6LUF}
  \forall F: \power (\num \cross \num) | F \in C6L @ F \in \nat \fun \nat
\end{theorem}

\begin{zproof}
prove;
use ClosElemUF[I := \{ cst~1, funL \}, U := \{ opPlusI, opBox \}, B := \{ comp \}, F := F];
with enabled (C6LDef) prove;
\end{zproof}

\begin{theorem}{rule C6LInit0}
  cst~1 \in C6L
\end{theorem}

\begin{zproof}
with enabled (C6LDef) prove;
\end{zproof}

\begin{theorem}{rule C6LInit1}
  funL \in C6L
\end{theorem}

\begin{zproof}
with enabled (C6LDef) prove;
\end{zproof}

\begin{theorem}{rule C6LUnop0}
  \forall F: \nat \fun \nat | F \in C6L @ opPlusI~F \in C6L
\end{theorem}

\begin{zproof}
with enabled (C6LDef) prove;
\end{zproof}

\begin{theorem}{rule C6LUnop1}
  \forall F: \nat \fun \nat | F \in C6L @ opBox~F \in C6L
\end{theorem}

\begin{zproof}
with enabled (C6LDef) prove;
\end{zproof}

\begin{theorem}{rule C6LBinop0}
  \forall F, G: \nat \fun \nat | F \in C6L \land G \in C6L @ comp(F, G) \in C6L
\end{theorem}

\begin{zproof}
with enabled (C6LDef) prove;
\end{zproof}

\begin{theorem}{KInLExpr}
  comp(comp(funL, opPlusI~(cst~1)), opPlusI~(comp(opPlusI~(opPlusI~(cst~1)), funL))) \in C6L
\end{theorem}

\begin{zproof}
prove;
\end{zproof}

\begin{theorem}{rule KInL}
  funK \in C6L
\end{theorem}

\begin{zproof}
use KInLExpr;
use KExprL;
prove;
\end{zproof}

\begin{theorem}{Theorem68LSub1U}
  C6 \in UClosed~\{ opPlusI, opBox \}
\end{theorem}

\begin{zproof}
use C6InPower;
with enabled (UClosedDef, C6UF) prove;
split u = opPlusI;
cases;
prove;
next;
prove;
next;
\end{zproof}

\begin{theorem}{Theorem68LSub1B}
  C6 \in BClosed~\{ comp \}
\end{theorem}

\begin{zproof}
use C6InPower;
with enabled (BClosedDef, C6UF) prove;
prove;
\end{zproof}

\begin{theorem}{Theorem68LSub1}
  C6L \subseteq C6
\end{theorem}

\begin{zproof}
use ClosLeast[I := \{ cst~1, funL \}, U := \{ opPlusI, opBox \}, B := \{ comp \}, Z := C6];
use C6InPower;
use Theorem68LSub1U;
use Theorem68LSub1B;
with enabled (C6LDef) prove;
\end{zproof}

\begin{theorem}{Theorem68LSubKU}
  C6L \in UClosed~\{ opPlusI, opBox \}
\end{theorem}

\begin{zproof}
use C6LInPower;
with enabled (UClosedDef, C6LUF) prove;
split u = opPlusI;
cases;
prove;
next;
prove;
next;
\end{zproof}

\begin{theorem}{Theorem68LSubKB}
  C6L \in BClosed~\{ comp \}
\end{theorem}

\begin{zproof}
use C6LInPower;
with enabled (BClosedDef, C6LUF) prove;
prove;
\end{zproof}

\begin{theorem}{Theorem68LSubK}
  C6K \subseteq C6L
\end{theorem}

\begin{zproof}
use ClosLeast[I := \{ cst~1, funK \}, U := \{ opPlusI, opBox \}, B := \{ comp \}, Z := C6L];
use C6LInPower;
use Theorem68LSubKU;
use Theorem68LSubKB;
with enabled (C6KDef) prove;
\end{zproof}

\begin{theorem}{Theorem68LSub2}
  C6 \subseteq C6L
\end{theorem}

\begin{zproof}
use Theorem68LSubK;
use Theorem68K;
prove;
\end{zproof}

\begin{theorem}{Theorem68L}
  C6L = C6
\end{theorem}

\begin{zproof}
use Theorem68LSub1;
use Theorem68LSub2;
use C6LInPower;
use C6InPower;
use SubAntisym[\nat \fun \nat][A := C6L, B := C6];
prove;
\end{zproof}

\begin{theorem}{Theorem68}
  C6K = C6 \land C6L = C6
\end{theorem}

\begin{zproof}
use Theorem68K;
use Theorem68L;
prove;
\end{zproof}

\section{Remark 6.9 and iteration from any $a \ge 1$}

\subsection{The product, multiples and iteration from $a$ in $\mathrm{clos}\,\mathcal{F}$}

$FG = \mathrm{Hf}((F+G)^2 - (F^2 + G^2))$ for arbitrary $F, G$.

\begin{theorem}{TimesEq}
  \forall F, G: \nat \fun \nat @ comp(funHf, opDiffE(comp(funSq, opPlus(F, G)), opPlus(comp(funSq, F), comp(funSq, G)))) = opOtimes(F, G)
\end{theorem}

\begin{zproof}
apply UFExt to predicate comp(funHf, opDiffE(comp(funSq, opPlus(F, G)), opPlus(comp(funSq, F), comp(funSq, G)))) = opOtimes(F, G);
prove;
use MulNonneg[a := F~x, b := G~x];
use DiffEIsDifference[F := comp(funSq, opPlus(F, G)), G := opPlus(comp(funSq, F), comp(funSq, G)), x := x];
prove;
use DivMod[x := 2 * (F~x * G~x), d := 2, q := F~x * G~x, r := 0];
prove;
\end{zproof}

\begin{theorem}{rule C6Times}
  \forall F, G: \nat \fun \nat | F \in C6 \land G \in C6 @ opOtimes(F, G) \in C6
\end{theorem}

\begin{zproof}
use TimesEq[F := F, G := G];
prove;
\end{zproof}

Multiples $M_r(x) = r x$, and $M_r = (\dots(\bar 0^+)^+\dots)^+$ ($r$ times) in every closure
containing $\bar 0$ and closed under $F^+$.

\begin{axdef}
  funM: \nat \fun (\nat \fun \nat)
\where
  \Label{rule MVal} \forall r, x: \nat @ funM~r~x = r * x
\end{axdef}

\begin{zproof}[funM\$domainCheck]
prove;
\end{zproof}

\begin{theorem}{rule MInner}
  \forall r: \nat @ (\lambda x: \nat @ r * x) \in \nat \fun \nat
\end{theorem}

\begin{zproof}
apply FunIntro to predicate (\lambda x: \nat @ r * x) \in \nat \fun \nat;
with enabled (RelPow, inPower) prove by reduce;
\end{zproof}

\begin{zed}
  MWit == (\lambda r: \nat @ (\lambda x: \nat @ r * x))
\end{zed}

\begin{theorem}{MWitFun}
  MWit \in \nat \fun (\nat \fun \nat)
\end{theorem}

\begin{zproof}
invoke MWit;
apply FunIntro to predicate (\lambda r: \nat @ (\lambda x: \nat @ r * x)) \in \nat \fun (\nat \fun \nat);
with enabled (RelPow, inPower) prove by reduce;
\end{zproof}

\begin{theorem}{MWitApp}
  \forall r: \nat @ MWit~r = (\lambda x: \nat @ r * x)
\end{theorem}

\begin{zproof}
use MWitFun;
use pairInFunction[\nat, \nat \fun \nat][f := MWit, x := r, y := (\lambda x: \nat @ r * x)];
invoke MWit;
prove by reduce;
\end{zproof}

\begin{theorem}{rule MWitVal}
  \forall r, x: \nat @ MWit~r~x = r * x
\end{theorem}

\begin{zproof}
prove;
use MWitApp[r := r];
use pairInFunction[\nat, \nat][f := (\lambda x: \nat @ r * x), x := x, y := r * x];
prove;
\end{zproof}

\begin{theorem}{MEx}
  \exists funM: \nat \fun (\nat \fun \nat) @ \forall r, x: \nat @ funM~r~x = r * x
\end{theorem}

\begin{zproof}
instantiate funM == MWit;
use MWitFun;
prove;
\end{zproof}

\begin{theorem}{rule NatUFApp}
  \forall g: \nat \fun (\nat \fun \nat); r: \nat @ g~r \in \nat \fun \nat
\end{theorem}

\begin{zproof}
prove;
\end{zproof}

\begin{theorem}{MZero}
  funM~0 = cst~0
\end{theorem}

\begin{zproof}
apply UFExt to predicate funM~0 = cst~0;
prove;
\end{zproof}

\begin{theorem}{MSuc}
  \forall r: \nat @ opPlusI~(funM~r) = funM~(1 + r)
\end{theorem}

\begin{zproof}
apply UFExt to predicate opPlusI~(funM~r) = funM~(1 + r);
prove;
\end{zproof}

\begin{zed}
  MSet == \{ r: \nat | \forall I: \power (\nat \fun \nat); U: \power ((\nat \fun \nat) \fun (\nat \fun \nat)); B: \power ((\nat \fun \nat) \cross (\nat \fun \nat) \fun (\nat \fun \nat)) | cst~0 \in clos~I~U~B \land opPlusI \in U @
    funM~r \in clos~I~U~B \}
\end{zed}

\begin{zproof}[MSet\$domainCheck]
prove;
\end{zproof}

\begin{theorem}{rule MBase}
  0 \in MSet
\end{theorem}

\begin{zproof}
invoke MSet;
use MZero;
prove;
\end{zproof}

\begin{theorem}{rule MStep}
  \forall r: MSet @ 1 + r \in MSet
\end{theorem}

\begin{zproof}
invoke MSet;
prove;
instantiate I\_\_0 == I, U\_\_0 == U, B\_\_0 == B;
use MSuc[r := r];
prove;
equality substitute funM~(1 + r);
rewrite;
\end{zproof}

\begin{theorem}{MSub}
  \nat \subseteq MSet
\end{theorem}

\begin{zproof}
apply natInduction to predicate \nat \subseteq MSet;
prove;
invoke MSet;
prove;
\end{zproof}

\begin{theorem}{MInClos}
  \forall I: \power (\nat \fun \nat); U: \power ((\nat \fun \nat) \fun (\nat \fun \nat)); B: \power ((\nat \fun \nat) \cross (\nat \fun \nat) \fun (\nat \fun \nat)); r: \nat |
    cst~0 \in clos~I~U~B \land opPlusI \in U @ funM~r \in clos~I~U~B
\end{theorem}

\begin{zproof}
prove;
use MSub;
apply subsetDef to predicate \nat \subseteq MSet;
apply inPower to predicate \nat \in \power MSet;
prove;
instantiate e == r;
invoke MSet;
prove;
instantiate I\_\_1 == I, U\_\_1 == U, B\_\_1 == B;
prove;
\end{zproof}

\begin{theorem}{rule MInC6}
  \forall r: \nat @ funM~r \in C6
\end{theorem}

\begin{zproof}
prove;
use MInClos[I := \{ cst~1, funE \}, U := \{ opPlusI, opBox \}, B := \{ comp \}, r := r];
use ZeroIn;
with enabled (C6Def) prove;
\end{zproof}

Soundness for iteration from $a$: $F^{\square(a)} = \hat P\,T^\square S$ with $T(0) = a + 1$ and
$T(y+1) = F(y) + 1$, namely $T = \mathrm{Sgn}\cdot(S F \hat P) + M_{a+1} O$.

\begin{theorem}{BoxATVal}
  \forall a: \nat; F: \nat \fun \nat; y: \nat @ opPlus(opOtimes(funSgn, comp(comp(funS, F), exprPhat)), comp(funM~(1 + a), funO))~0 = 1 + a \land opPlus(opOtimes(funSgn, comp(comp(funS, F), exprPhat)), comp(funM~(1 + a), funO))~(1 + y) = 1 + F~y
\end{theorem}

\begin{zproof}
prove;
\end{zproof}

\begin{zed}
  BoxASet == \{ n: \nat | \forall a: \nat; F: \nat \fun \nat @
    opBox~(opPlus(opOtimes(funSgn, comp(comp(funS, F), exprPhat)), comp(funM~(1 + a), funO)))~(1 + n) = 1 + opBoxA~a~F~n \}
\end{zed}

\begin{zproof}[BoxASet\$domainCheck]
prove;
\end{zproof}

\begin{theorem}{rule BoxABase}
  0 \in BoxASet
\end{theorem}

\begin{zproof}
invoke BoxASet;
prove;
\end{zproof}

\begin{theorem}{rule BoxAStep}
  \forall n: BoxASet @ 1 + n \in BoxASet
\end{theorem}

\begin{zproof}
invoke BoxASet;
with disabled (BoxSuc) prove;
instantiate a\_\_0 == a, F\_\_0 == F;
use BoxSuc[F\_\_0 := opPlus(opOtimes(funSgn, comp(comp(funS, F), exprPhat)), comp(funM~(1 + a), funO)), x := 1 + n];
use BoxATVal[a := a, F := F, y := opBoxA~a~F~n];
with disabled (BoxSuc, PlusVal) prove;
\end{zproof}

\begin{theorem}{BoxASub}
  \nat \subseteq BoxASet
\end{theorem}

\begin{zproof}
apply natInduction to predicate \nat \subseteq BoxASet;
prove;
invoke BoxASet;
prove;
\end{zproof}

\begin{theorem}{BoxAExprEq}
  \forall a: \nat; F: \nat \fun \nat @ comp(comp(exprPhat, opBox~(opPlus(opOtimes(funSgn, comp(comp(funS, F), exprPhat)), comp(funM~(1 + a), funO)))), funS) = opBoxA~a~F
\end{theorem}

\begin{zproof}
apply UFExt to predicate comp(comp(exprPhat, opBox~(opPlus(opOtimes(funSgn, comp(comp(funS, F), exprPhat)), comp(funM~(1 + a), funO)))), funS) = opBoxA~a~F;
with disabled (BoxSuc) prove;
use BoxASub;
apply subsetDef to predicate \nat \subseteq BoxASet;
apply inPower to predicate \nat \in \power BoxASet;
with disabled (BoxSuc) prove;
instantiate e == x;
invoke BoxASet;
with disabled (BoxSuc) prove;
instantiate a\_\_1 == a, F\_\_1 == F;
with disabled (BoxSuc) prove;
equality substitute opBox~(opPlus(opOtimes(funSgn, comp(comp(funS, F), exprPhat)), comp(funM~(1 + a), funO)))~(1 + x);
with disabled (BoxSuc) rewrite;
\end{zproof}

\begin{theorem}{rule BoxAIn}
  \forall a: \nat; F: \nat \fun \nat | F \in C6 @ opBoxA~a~F \in C6
\end{theorem}

\begin{zproof}
use BoxAExprEq[a := a, F := F];
prove;
\end{zproof}

\subsection{Powers under composition}

$F^k$ is the Toolkit's iterate $F^k = $ \texttt{iter k F}.

\begin{zed}
  IterSSet == \{ k: \nat | \forall x: \nat @ iter[\nat]~k~funS~x = x + k \}
\end{zed}

\begin{zproof}[IterSSet\$domainCheck]
prove;
\end{zproof}

\begin{theorem}{rule IterSBase}
  0 \in IterSSet
\end{theorem}

\begin{zproof}
invoke IterSSet;
prove;
\end{zproof}

\begin{theorem}{rule IterSStep}
  \forall k: IterSSet @ 1 + k \in IterSSet
\end{theorem}

\begin{zproof}
invoke IterSSet;
prove;
instantiate x\_\_0 == x;
prove;
\end{zproof}

\begin{theorem}{IterSSub}
  \nat \subseteq IterSSet
\end{theorem}

\begin{zproof}
apply natInduction to predicate \nat \subseteq IterSSet;
prove;
invoke IterSSet;
prove;
\end{zproof}

\begin{theorem}{IterSAll}
  \forall k: \nat @ \forall x: \nat @ iter[\nat]~k~funS~x = x + k
\end{theorem}

\begin{zproof}
prove;
use IterSSub;
invoke IterSSet;
with enabled (inPower) prove;
instantiate e == k;
prove;
instantiate x\_\_1 == x;
prove;
\end{zproof}

Predecessors ($P(x + 1) = x$) undo successors: $P^k(x + k) = x$.

\begin{zed}
  Preds == \{ P: \nat \fun \nat | \forall x: \nat @ P~(1 + x) = x \}
\end{zed}

\begin{zproof}[Preds\$domainCheck]
prove;
\end{zproof}

\begin{theorem}{rule PredsVal}
  \forall P: \nat \fun \nat; x: \nat | P \in Preds @ P~(1 + x) = x
\end{theorem}

\begin{zproof}
prove;
invoke Preds;
prove;
instantiate x\_\_0 == x;
prove;
\end{zproof}

\begin{zed}
  IterPSet == \{ k: \nat | \forall P: \nat \fun \nat; x: \nat | P \in Preds @ iter[\nat]~k~P~(x + k) = x \}
\end{zed}

\begin{zproof}[IterPSet\$domainCheck]
prove;
\end{zproof}

\begin{theorem}{rule IterPBase}
  0 \in IterPSet
\end{theorem}

\begin{zproof}
invoke IterPSet;
prove;
\end{zproof}

\begin{theorem}{rule IterPStep}
  \forall k: IterPSet @ 1 + k \in IterPSet
\end{theorem}

\begin{zproof}
invoke IterPSet;
prove;
instantiate P\_\_0 == P, x\_\_0 == 1 + x;
use IterSuc[F := P, n := k, a := 1 + x + k];
prove;
equality substitute iter[\nat]~k~P~(1 + (k + x));
rewrite;
\end{zproof}

\begin{theorem}{IterPSub}
  \nat \subseteq IterPSet
\end{theorem}

\begin{zproof}
apply natInduction to predicate \nat \subseteq IterPSet;
prove;
invoke IterPSet;
prove;
\end{zproof}

\begin{theorem}{IterPAll}
  \forall k: \nat @ \forall P: \nat \fun \nat; x: \nat | P \in Preds @ iter[\nat]~k~P~(x + k) = x
\end{theorem}

\begin{zproof}
prove;
use IterPSub;
invoke IterPSet;
with enabled (inPower) prove;
instantiate e == k;
prove;
instantiate P\_\_1 == P, x\_\_1 == x;
prove;
\end{zproof}

$F^k G$ stays in every closure containing composition.

\begin{zed}
  PowSet == \{ k: \nat | \forall I: \power (\nat \fun \nat); U: \power ((\nat \fun \nat) \fun (\nat \fun \nat)); B: \power ((\nat \fun \nat) \cross (\nat \fun \nat) \fun (\nat \fun \nat)); F, G: \nat \fun \nat |
    comp \in B \land F \in clos~I~U~B \land G \in clos~I~U~B @ comp(iter[\nat]~k~F, G) \in clos~I~U~B \}
\end{zed}

\begin{zproof}[PowSet\$domainCheck]
prove;
\end{zproof}

\begin{theorem}{IterCompZero}
  \forall F, G: \nat \fun \nat @ comp(iter[\nat]~0~F, G) = G
\end{theorem}

\begin{zproof}
apply UFExt to predicate comp(iter[\nat]~0~F, G) = G;
prove;
\end{zproof}

\begin{theorem}{IterCompSuc}
  \forall F, G: \nat \fun \nat; k: \nat @ comp(iter[\nat]~(1 + k)~F, G) = comp(F, comp(iter[\nat]~k~F, G))
\end{theorem}

\begin{zproof}
apply UFExt to predicate comp(iter[\nat]~(1 + k)~F, G) = comp(F, comp(iter[\nat]~k~F, G));
prove;
\end{zproof}

\begin{theorem}{rule PowBase}
  0 \in PowSet
\end{theorem}

\begin{zproof}
invoke PowSet;
prove;
use IterCompZero[F := F, G := G];
prove;
\end{zproof}

\begin{theorem}{rule PowStep}
  \forall k: PowSet @ 1 + k \in PowSet
\end{theorem}

\begin{zproof}
invoke PowSet;
prove;
instantiate I\_\_0 == I, U\_\_0 == U, B\_\_0 == B, F\_\_0 == F, G\_\_0 == G;
use IterCompSuc[F := F, G := G, k := k];
prove;
\end{zproof}

\begin{theorem}{PowSub}
  \nat \subseteq PowSet
\end{theorem}

\begin{zproof}
apply natInduction to predicate \nat \subseteq PowSet;
prove;
invoke PowSet;
prove;
\end{zproof}

\begin{theorem}{PowIn}
  \forall I: \power (\nat \fun \nat); U: \power ((\nat \fun \nat) \fun (\nat \fun \nat)); B: \power ((\nat \fun \nat) \cross (\nat \fun \nat) \fun (\nat \fun \nat)); F, G: \nat \fun \nat; k: \nat |
    comp \in B \land F \in clos~I~U~B \land G \in clos~I~U~B @ comp(iter[\nat]~k~F, G) \in clos~I~U~B
\end{theorem}

\begin{zproof}
prove;
use PowSub;
apply subsetDef to predicate \nat \subseteq PowSet;
apply inPower to predicate \nat \in \power PowSet;
prove;
instantiate e == k;
invoke PowSet;
prove;
instantiate I\_\_1 == I, U\_\_1 == U, B\_\_1 == B, F\_\_1 == F, G\_\_1 == G;
prove;
\end{zproof}

Recovering iteration from $0$ out of iteration from $a$ (as in Remark 5.10 of the paper):
$F^\square = P^a (S^a F P^a)^{\square(a)}$ for a predecessor $P$.

\begin{zed}
  BoxFromASet == \{ n: \nat | \forall a: \nat; F, P: \nat \fun \nat | P \in Preds @ opBoxA~a~(comp(comp(iter[\nat]~a~funS, F), iter[\nat]~a~P))~n = opBox~F~n + a \}
\end{zed}

\begin{zproof}[BoxFromASet\$domainCheck]
prove;
\end{zproof}

\begin{theorem}{rule BoxFromABase}
  0 \in BoxFromASet
\end{theorem}

\begin{zproof}
invoke BoxFromASet;
prove;
\end{zproof}

\begin{theorem}{rule BoxFromAStep}
  \forall n: BoxFromASet @ 1 + n \in BoxFromASet
\end{theorem}

\begin{zproof}
invoke BoxFromASet;
prove;
instantiate a\_\_0 == a, F\_\_0 == F, P\_\_0 == P;
use IterPAll[k := a, P := P, x := opBox~F~n];
use IterSAll[k := a, x := F~(opBox~F~n)];
prove;
\end{zproof}

\begin{theorem}{BoxFromASub}
  \nat \subseteq BoxFromASet
\end{theorem}

\begin{zproof}
apply natInduction to predicate \nat \subseteq BoxFromASet;
prove;
invoke BoxFromASet;
prove;
\end{zproof}

\begin{theorem}{BoxFromAAll}
  \forall n: \nat @ \forall a: \nat; F, P: \nat \fun \nat | P \in Preds @ opBoxA~a~(comp(comp(iter[\nat]~a~funS, F), iter[\nat]~a~P))~n = opBox~F~n + a
\end{theorem}

\begin{zproof}
prove;
use BoxFromASub;
invoke BoxFromASet;
with enabled (inPower) prove;
instantiate e == n;
prove;
instantiate a\_\_1 == a, F\_\_1 == F, P\_\_1 == P;
prove;
\end{zproof}

\begin{theorem}{BoxFromBoxA}
  \forall a: \nat; F, P: \nat \fun \nat | P \in Preds @ comp(iter[\nat]~a~P, opBoxA~a~(comp(comp(iter[\nat]~a~funS, F), iter[\nat]~a~P))) = opBox~F
\end{theorem}

\begin{zproof}
apply UFExt to predicate comp(iter[\nat]~a~P, opBoxA~a~(comp(comp(iter[\nat]~a~funS, F), iter[\nat]~a~P))) = opBox~F;
prove;
use BoxFromAAll[n := x, a := a, F := F, P := P];
use IterPAll[k := a, P := P, x := opBox~F~x];
prove;
\end{zproof}

\subsection{The bases $\langle X, F^+, FG, F^{\square(a)}\rangle$: the common part}

$X$ admits multipliers if $X(0) = 0$ and for every $i \ge 1$ some $r$ has $X(r(i+1)) = i$;
then $T_j = X M_{r(j)}$ maps $j$ to $j - 1$ and $0$ to $0$.

\begin{zed}
  Mult == \{ X: \nat \fun \nat | X~0 = 0 \land (\forall i: \nat | 1 \leq i @ \exists r: \nat @ X~(r * (1 + i)) = i) \}
\end{zed}

\begin{zproof}[Mult\$domainCheck]
prove;
\end{zproof}

\begin{theorem}{MultZero}
  \forall X: \nat \fun \nat | X \in Mult @ X~0 = 0
\end{theorem}

\begin{zproof}
invoke Mult;
prove;
\end{zproof}

\begin{theorem}{MultEx}
  \forall X: \nat \fun \nat; i: \nat | X \in Mult \land 1 \leq i @ \exists r: \nat @ X~(r * (1 + i)) = i
\end{theorem}

\begin{zproof}
prove;
invoke Mult;
prove;
instantiate i\_\_0 == i;
prove;
instantiate r\_\_1 == r;
prove;
\end{zproof}

The chain $N = T_2 \cdots T_{k+1}$ maps $k+1$ to $1$ and $0$ to $0$.

\begin{zed}
  ChainSet == \{ k: \nat | \forall X: \nat \fun \nat; a: \nat | X \in Mult \land cst~0 \in clos~\{ X \}~\{ opPlusI, opBoxA~a \}~\{ comp \} @
    \exists W: \nat \fun \nat | W \in clos~\{ X \}~\{ opPlusI, opBoxA~a \}~\{ comp \} @ W~(1 + k) = 1 \land W~0 = 0 \}
\end{zed}

\begin{zproof}[ChainSet\$domainCheck]
prove;
\end{zproof}

\begin{theorem}{rule ChainBase}
  0 \in ChainSet
\end{theorem}

\begin{zproof}
invoke ChainSet;
prove;
use MInClos[I := \{ X \}, U := \{ opPlusI, opBoxA~a \}, B := \{ comp \}, r := 1];
prove;
instantiate W == funM~1;
prove;
\end{zproof}

\begin{theorem}{rule ChainStep}
  \forall k: ChainSet @ 1 + k \in ChainSet
\end{theorem}

\begin{zproof}
invoke ChainSet;
prove;
instantiate X\_\_0 == X, a\_\_0 == a;
prove;
use MultEx[X := X, i := 1 + k];
use MultZero[X := X];
prove;
use MInClos[I := \{ X \}, U := \{ opPlusI, opBoxA~a \}, B := \{ comp \}, r := r];
prove;
instantiate W\_\_1 == comp(W, comp(X, funM~r));
prove;
\end{zproof}

\begin{theorem}{ChainSub}
  \nat \subseteq ChainSet
\end{theorem}

\begin{zproof}
apply natInduction to predicate \nat \subseteq ChainSet;
prove;
invoke ChainSet;
prove;
\end{zproof}

\begin{theorem}{ChainAll}
  \forall X: \nat \fun \nat; a, k: \nat | X \in Mult \land cst~0 \in clos~\{ X \}~\{ opPlusI, opBoxA~a \}~\{ comp \} @
    \exists W: \nat \fun \nat | W \in clos~\{ X \}~\{ opPlusI, opBoxA~a \}~\{ comp \} @ W~(1 + k) = 1 \land W~0 = 0
\end{theorem}

\begin{zproof}
prove;
use ChainSub;
apply subsetDef to predicate \nat \subseteq ChainSet;
apply inPower to predicate \nat \in \power ChainSet;
prove;
instantiate e == k;
invoke ChainSet;
prove;
instantiate X\_\_1 == X, a\_\_1 == a;
prove;
instantiate W\_\_2 == W;
prove;
\end{zproof}

\begin{theorem}{MOne}
  funM~1 = funI \land funM~2 = funD
\end{theorem}

\begin{zproof}
use UFExt[F := funM~1, G := funI];
use UFExt[F := funM~2, G := funD];
prove;
\end{zproof}

\begin{zed}
  BoxAISet == \{ n: \nat | \forall a: \nat @ opBoxA~a~funI~n = a \}
\end{zed}

\begin{zproof}[BoxAISet\$domainCheck]
prove;
\end{zproof}

\begin{theorem}{rule BoxAIBase}
  0 \in BoxAISet
\end{theorem}

\begin{zproof}
invoke BoxAISet;
prove;
\end{zproof}

\begin{theorem}{rule BoxAIStep}
  \forall n: BoxAISet @ 1 + n \in BoxAISet
\end{theorem}

\begin{zproof}
invoke BoxAISet;
prove;
\end{zproof}

\begin{theorem}{BoxAISub}
  \nat \subseteq BoxAISet
\end{theorem}

\begin{zproof}
apply natInduction to predicate \nat \subseteq BoxAISet;
prove;
invoke BoxAISet;
prove;
\end{zproof}

\begin{theorem}{BoxAIAll}
  \forall n: \nat @ \forall a: \nat @ opBoxA~a~funI~n = a
\end{theorem}

\begin{zproof}
prove;
use BoxAISub;
invoke BoxAISet;
with enabled (inPower) prove;
instantiate e == n;
prove;
instantiate a\_\_1 == a;
prove;
\end{zproof}

\begin{theorem}{OneFromW}
  \forall W: \nat \fun \nat; a: \nat | W~a = 1 @ comp(W, opBoxA~a~funI) = cst~1
\end{theorem}

\begin{zproof}
apply UFExt to predicate comp(W, opBoxA~a~funI) = cst~1;
prove;
use BoxAIAll[n := x];
prove;
\end{zproof}

\begin{theorem}{OFromW}
  \forall W: \nat \fun \nat; a: \nat | W~a = 1 \land W~0 = 0 @ comp(W, opBoxA~a~(cst~0)) = funO
\end{theorem}

\begin{zproof}
apply UFExt to predicate comp(W, opBoxA~a~(cst~0)) = funO;
prove;
split x = 0;
cases;
prove;
next;
use BoxASuc[a\_\_0 := a, F\_\_0 := cst~0, x := x - 1];
prove;
next;
\end{zproof}

\begin{theorem}{ExtAI}
  \forall X: \nat \fun \nat; a: \nat | X \in Mult \land cst~0 \in clos~\{ X \}~\{ opPlusI, opBoxA~a \}~\{ comp \} @ funI \in clos~\{ X \}~\{ opPlusI, opBoxA~a \}~\{ comp \} \land funD \in clos~\{ X \}~\{ opPlusI, opBoxA~a \}~\{ comp \}
\end{theorem}

\begin{zproof}
prove;
use MInClos[I := \{ X \}, U := \{ opPlusI, opBoxA~a \}, B := \{ comp \}, r := 1];
use MInClos[I := \{ X \}, U := \{ opPlusI, opBoxA~a \}, B := \{ comp \}, r := 2];
use MOne;
prove;
\end{zproof}

\begin{theorem}{OneInCB}
  \forall X, W: \nat \fun \nat; a: \nat | W \in clos~\{ X \}~\{ opPlusI, opBoxA~a \}~\{ comp \} \land funI \in clos~\{ X \}~\{ opPlusI, opBoxA~a \}~\{ comp \} \land W~a = 1 @
    cst~1 \in clos~\{ X \}~\{ opPlusI, opBoxA~a \}~\{ comp \}
\end{theorem}

\begin{zproof}
use OneFromW[W := W, a := a];
prove;
\end{zproof}

\begin{theorem}{OExprInCB}
  \forall X, W: \nat \fun \nat; a: \nat | W \in clos~\{ X \}~\{ opPlusI, opBoxA~a \}~\{ comp \} \land cst~0 \in clos~\{ X \}~\{ opPlusI, opBoxA~a \}~\{ comp \} @
    comp(W, opBoxA~a~(cst~0)) \in clos~\{ X \}~\{ opPlusI, opBoxA~a \}~\{ comp \}
\end{theorem}

\begin{zproof}
prove;
\end{zproof}

\begin{theorem}{OInCB}
  \forall X, W: \nat \fun \nat; a: \nat | W \in clos~\{ X \}~\{ opPlusI, opBoxA~a \}~\{ comp \} \land cst~0 \in clos~\{ X \}~\{ opPlusI, opBoxA~a \}~\{ comp \} \land W~a = 1 \land W~0 = 0 @
    funO \in clos~\{ X \}~\{ opPlusI, opBoxA~a \}~\{ comp \}
\end{theorem}

\begin{zproof}
use OExprInCB[X := X, W := W, a := a];
use OFromW[W := W, a := a];
prove;
\end{zproof}

\begin{theorem}{ExtAOne}
  \forall X: \nat \fun \nat; a: \nat | 1 \leq a \land X \in Mult \land cst~0 \in clos~\{ X \}~\{ opPlusI, opBoxA~a \}~\{ comp \} @
    cst~1 \in clos~\{ X \}~\{ opPlusI, opBoxA~a \}~\{ comp \} \land funO \in clos~\{ X \}~\{ opPlusI, opBoxA~a \}~\{ comp \}
\end{theorem}

\begin{zproof}
prove;
use ExtAI[X := X, a := a];
use ChainAll[X := X, a := a, k := a - 1];
prove;
use OneInCB[X := X, W := W, a := a];
use OInCB[X := X, W := W, a := a];
prove;
\end{zproof}

\begin{theorem}{SExprInCB}
  \forall X: \nat \fun \nat; a: \nat | cst~1 \in clos~\{ X \}~\{ opPlusI, opBoxA~a \}~\{ comp \} @ opPlusI~(cst~1) \in clos~\{ X \}~\{ opPlusI, opBoxA~a \}~\{ comp \}
\end{theorem}

\begin{zproof}
prove;
\end{zproof}

\begin{theorem}{ExtAS}
  \forall X: \nat \fun \nat; a: \nat | 1 \leq a \land X \in Mult \land cst~0 \in clos~\{ X \}~\{ opPlusI, opBoxA~a \}~\{ comp \} @ funS \in clos~\{ X \}~\{ opPlusI, opBoxA~a \}~\{ comp \}
\end{theorem}

\begin{zproof}
prove;
use ExtAOne[X := X, a := a];
prove;
use SExprInCB[X := X, a := a];
use SExpr;
prove;
\end{zproof}

\begin{theorem}{IterCompI}
  \forall F: \nat \fun \nat; k: \nat @ comp(iter[\nat]~k~F, funI) = iter[\nat]~k~F
\end{theorem}

\begin{zproof}
apply UFExt to predicate comp(iter[\nat]~k~F, funI) = iter[\nat]~k~F;
prove;
\end{zproof}

\begin{theorem}{IterIn}
  \forall I: \power (\nat \fun \nat); U: \power ((\nat \fun \nat) \fun (\nat \fun \nat)); B: \power ((\nat \fun \nat) \cross (\nat \fun \nat) \fun (\nat \fun \nat)); F: \nat \fun \nat; k: \nat |
    comp \in B \land F \in clos~I~U~B \land funI \in clos~I~U~B @ iter[\nat]~k~F \in clos~I~U~B
\end{theorem}

\begin{zproof}
prove;
use PowIn[I := I, U := U, B := B, F := F, G := funI, k := k];
use IterCompI[F := F, k := k];
prove;
\end{zproof}

With a predecessor in the closure, iteration from $0$ is available.

\begin{theorem}{ExtABox}
  \forall X, P, F: \nat \fun \nat; a: \nat | 1 \leq a \land X \in Mult \land cst~0 \in clos~\{ X \}~\{ opPlusI, opBoxA~a \}~\{ comp \} \land P \in clos~\{ X \}~\{ opPlusI, opBoxA~a \}~\{ comp \}
      \land P \in Preds \land F \in clos~\{ X \}~\{ opPlusI, opBoxA~a \}~\{ comp \} @ opBox~F \in clos~\{ X \}~\{ opPlusI, opBoxA~a \}~\{ comp \}
\end{theorem}

\begin{zproof}
prove;
use ExtAI[X := X, a := a];
use ExtAS[X := X, a := a];
use IterIn[I := \{ X \}, U := \{ opPlusI, opBoxA~a \}, B := \{ comp \}, F := P, k := a];
use IterIn[I := \{ X \}, U := \{ opPlusI, opBoxA~a \}, B := \{ comp \}, F := funS, k := a];
use BoxFromBoxA[a := a, F := F, P := P];
prove;
\end{zproof}

\subsection{Simulation}

\begin{theorem}{CBUClosed}
  \forall X, P: \nat \fun \nat; a: \nat | 1 \leq a \land X \in Mult \land cst~0 \in clos~\{ X \}~\{ opPlusI, opBoxA~a \}~\{ comp \} \land P \in clos~\{ X \}~\{ opPlusI, opBoxA~a \}~\{ comp \} \land P \in Preds @
    clos~\{ X \}~\{ opPlusI, opBoxA~a \}~\{ comp \} \in UClosed~\{ opPlusI, opBox \}
\end{theorem}

\begin{zproof}
prove;
use ClosInPow[I := \{ X \}, U := \{ opPlusI, opBoxA~a \}, B := \{ comp \}];
with enabled (UClosedDef) prove;
use ClosElemUF[I := \{ X \}, U := \{ opPlusI, opBoxA~a \}, B := \{ comp \}, F := F];
split u = opPlusI;
cases;
prove;
next;
use ExtABox[X := X, P := P, F := F, a := a];
prove;
next;
\end{zproof}

\begin{theorem}{CBBClosed}
  \forall X: \nat \fun \nat; a: \nat @ clos~\{ X \}~\{ opPlusI, opBoxA~a \}~\{ comp \} \in BClosed~\{ comp \}
\end{theorem}

\begin{zproof}
prove;
use ClosInPow[I := \{ X \}, U := \{ opPlusI, opBoxA~a \}, B := \{ comp \}];
with enabled (BClosedDef) prove;
\end{zproof}

\begin{theorem}{ExtASup}
  \forall X, P: \nat \fun \nat; a: \nat | 1 \leq a \land X \in Mult \land cst~0 \in clos~\{ X \}~\{ opPlusI, opBoxA~a \}~\{ comp \} \land P \in clos~\{ X \}~\{ opPlusI, opBoxA~a \}~\{ comp \} \land P \in Preds @
    clos~\{ cst~1, X \}~\{ opPlusI, opBox \}~\{ comp \} \subseteq clos~\{ X \}~\{ opPlusI, opBoxA~a \}~\{ comp \}
\end{theorem}

\begin{zproof}
prove;
use ClosLeast[I := \{ cst~1, X \}, U := \{ opPlusI, opBox \}, B := \{ comp \}, Z := clos~\{ X \}~\{ opPlusI, opBoxA~a \}~\{ comp \}];
use ClosInPow[I := \{ X \}, U := \{ opPlusI, opBoxA~a \}, B := \{ comp \}];
use CBUClosed[X := X, P := P, a := a];
use CBBClosed[X := X, a := a];
use ExtAOne[X := X, a := a];
prove;
\end{zproof}

\begin{theorem}{C6UClosedA}
  \forall a: \nat @ C6 \in UClosed~\{ opPlusI, opBoxA~a \}
\end{theorem}

\begin{zproof}
prove;
use C6InPower;
with enabled (UClosedDef, C6UF) prove;
use InPair[\power (\power (\num \cross \num) \cross \power (\num \cross \num))][u := u, p := opPlusI, q := opBoxA~a];
split u = opBoxA~a;
cases;
prove;
next;
prove;
next;
\end{zproof}

\begin{theorem}{C6BClosedComp}
  C6 \in BClosed~\{ comp \}
\end{theorem}

\begin{zproof}
use C6InPower;
with enabled (BClosedDef, C6UF) prove;
\end{zproof}

\begin{theorem}{ExtASub6}
  \forall X: \nat \fun \nat; a: \nat | X \in C6 @ clos~\{ X \}~\{ opPlusI, opBoxA~a \}~\{ comp \} \subseteq C6
\end{theorem}

\begin{zproof}
prove;
use ClosLeast[I := \{ X \}, U := \{ opPlusI, opBoxA~a \}, B := \{ comp \}, Z := C6];
use C6InPower;
use C6UClosedA[a := a];
use C6BClosedComp;
prove;
\end{zproof}

\subsection{Zero for a decreasing $X$: $\bar 0 = X^a X^{\square(a)}$}

\begin{zed}
  Decr == \{ X: \nat \fun \nat | X~0 = 0 \land (\forall y: \nat | 0 < y @ X~y < y) \}
\end{zed}

\begin{zproof}[Decr\$domainCheck]
prove;
\end{zproof}

\begin{theorem}{DecrZero}
  \forall X: \nat \fun \nat | X \in Decr @ X~0 = 0
\end{theorem}

\begin{zproof}
invoke Decr;
prove;
\end{zproof}

\begin{theorem}{DecrLt}
  \forall X: \nat \fun \nat; y: \nat | X \in Decr \land 0 < y @ X~y < y \land X~0 = 0
\end{theorem}

\begin{zproof}
prove;
invoke Decr;
prove;
instantiate y\_\_0 == y;
prove;
\end{zproof}

\begin{zed}
  DecBoundSet == \{ n: \nat | \forall a: \nat; X: \nat \fun \nat | X \in Decr @ opBoxA~a~X~n \leq a \}
\end{zed}

\begin{zproof}[DecBoundSet\$domainCheck]
prove;
\end{zproof}

\begin{theorem}{rule DecBoundBase}
  0 \in DecBoundSet
\end{theorem}

\begin{zproof}
invoke DecBoundSet;
prove;
\end{zproof}

\begin{theorem}{rule DecBoundStep}
  \forall n: DecBoundSet @ 1 + n \in DecBoundSet
\end{theorem}

\begin{zproof}
invoke DecBoundSet;
prove;
instantiate a\_\_0 == a, X\_\_0 == X;
use DecrZero[X := X];
use DecrLt[X := X, y := opBoxA~a~X~n];
use UFNonneg[f := X, x := opBoxA~a~X~n];
use UFNonneg[f := opBoxA~a~X, x := n];
split opBoxA~a~X~n = 0;
cases;
prove;
next;
prove;
next;
\end{zproof}

\begin{theorem}{DecBoundSub}
  \nat \subseteq DecBoundSet
\end{theorem}

\begin{zproof}
apply natInduction to predicate \nat \subseteq DecBoundSet;
prove;
invoke DecBoundSet;
prove;
\end{zproof}

\begin{theorem}{DecBoundAll}
  \forall n: \nat @ \forall a: \nat; X: \nat \fun \nat | X \in Decr @ opBoxA~a~X~n \leq a
\end{theorem}

\begin{zproof}
prove;
use DecBoundSub;
invoke DecBoundSet;
with enabled (inPower) prove;
instantiate e == n;
prove;
instantiate a\_\_1 == a, X\_\_1 == X;
prove;
\end{zproof}

\begin{zed}
  DecIterSet == \{ k: \nat | \forall X: \nat \fun \nat; y: \nat | X \in Decr \land y \leq k @ iter[\nat]~k~X~y = 0 \}
\end{zed}

\begin{zproof}[DecIterSet\$domainCheck]
prove;
\end{zproof}

\begin{theorem}{rule DecIterBase}
  0 \in DecIterSet
\end{theorem}

\begin{zproof}
invoke DecIterSet;
prove;
\end{zproof}

\begin{theorem}{rule DecIterStep}
  \forall k: DecIterSet @ 1 + k \in DecIterSet
\end{theorem}

\begin{zproof}
invoke DecIterSet;
prove;
instantiate X\_\_0 == X, y\_\_0 == X~y;
use IterSucR[F := X, k := k, a := y];
use DecrZero[X := X];
use DecrLt[X := X, y := y];
use UFNonneg[f := X, x := y];
split y = 0;
cases;
prove;
next;
prove;
next;
\end{zproof}

\begin{theorem}{DecIterSub}
  \nat \subseteq DecIterSet
\end{theorem}

\begin{zproof}
apply natInduction to predicate \nat \subseteq DecIterSet;
prove;
invoke DecIterSet;
prove;
\end{zproof}

\begin{theorem}{DecIterAll}
  \forall k: \nat @ \forall X: \nat \fun \nat; y: \nat | X \in Decr \land y \leq k @ iter[\nat]~k~X~y = 0
\end{theorem}

\begin{zproof}
prove;
use DecIterSub;
invoke DecIterSet;
with enabled (inPower) prove;
instantiate e == k;
prove;
instantiate X\_\_1 == X, y\_\_1 == y;
prove;
\end{zproof}

\begin{theorem}{DecZeroEq}
  \forall X: \nat \fun \nat; a: \nat | X \in Decr @ comp(iter[\nat]~a~X, opBoxA~a~X) = cst~0
\end{theorem}

\begin{zproof}
apply UFExt to predicate comp(iter[\nat]~a~X, opBoxA~a~X) = cst~0;
prove;
use DecBoundAll[n := x, a := a, X := X];
use DecIterAll[k := a, X := X, y := opBoxA~a~X~x];
prove;
\end{zproof}

\begin{theorem}{ZeroInCB}
  \forall X: \nat \fun \nat; a: \nat | X \in Decr @ cst~0 \in clos~\{ X \}~\{ opPlusI, opBoxA~a \}~\{ comp \}
\end{theorem}

\begin{zproof}
prove;
use PowIn[I := \{ X \}, U := \{ opPlusI, opBoxA~a \}, B := \{ comp \}, F := X, G := opBoxA~a~X, k := a];
use DecZeroEq[X := X, a := a];
prove;
\end{zproof}

\subsection{The case $X = E$}

\begin{theorem}{ELt}
  \forall y: \nat | 0 < y @ funE~y < y
\end{theorem}

\begin{zproof}
prove;
use RtSpec[x := y];
use RtPos[y := y];
use timesMonotonic2[a := 1, b := 1, c := funRt~y, d := funRt~y];
with enabled (EVal) prove;
\end{zproof}

\begin{theorem}{EDecr}
  funE \in Decr
\end{theorem}

\begin{zproof}
invoke Decr;
prove;
use ELt[y := y];
prove;
\end{zproof}

\begin{theorem}{EMult}
  funE \in Mult
\end{theorem}

\begin{zproof}
invoke Mult;
prove;
instantiate r == i;
use ESqAdd[m := i, k := i];
prove;
\end{zproof}

\begin{theorem}{rule IterSVal}
  \forall k, x: \nat @ iter[\nat]~k~funS~x = x + k
\end{theorem}

\begin{zproof}
prove;
use IterSAll[k := k, x := x];
prove;
\end{zproof}

\begin{theorem}{RRec}
  \forall t: \nat @ funR~(1 + t) = 2 * funQ~(4 + funR~t) + funR~t + 1
\end{theorem}

\begin{zproof}
prove;
use RStep[n := t];
prove;
\end{zproof}

With $c = a^2 + a + 4$: $B = (S (D Q S^c)^+)^{\square(a)}$ satisfies $B(n) = R(n + a^2) - a^2 - a$,
$G = (S^{a+1} B)^{\square(a)}$ satisfies $G(n) = n^2 + 2an + a$, and
$\hat P = E\,S^{a^2 - a - 1} G^+$ is a predecessor.

\begin{zed}
  BEValSet == \{ n: \nat | \forall a: \nat @ opBoxA~a~(comp(funS, opPlusI~(comp(comp(funD, funQ), iter[\nat]~(a * a + a + 4)~funS))))~n = funR~(n + a * a) - a * a - a \}
\end{zed}

\begin{zproof}[BEValSet\$domainCheck]
prove;
\end{zproof}

\begin{theorem}{rule BEValBase}
  0 \in BEValSet
\end{theorem}

\begin{zproof}
invoke BEValSet;
prove;
\end{zproof}

\begin{theorem}{rule BEValStep}
  \forall n: BEValSet @ 1 + n \in BEValSet
\end{theorem}

\begin{zproof}
invoke BEValSet;
with disabled (QVal, RVal) prove;
instantiate a\_\_0 == a;
use RRec[t := n + a * a];
with disabled (QVal, RVal) prove;
\end{zproof}

\begin{theorem}{BEValSub}
  \nat \subseteq BEValSet
\end{theorem}

\begin{zproof}
apply natInduction to predicate \nat \subseteq BEValSet;
prove;
invoke BEValSet;
prove;
\end{zproof}

\begin{theorem}{BEValAll}
  \forall n: \nat @ \forall a: \nat @ opBoxA~a~(comp(funS, opPlusI~(comp(comp(funD, funQ), iter[\nat]~(a * a + a + 4)~funS))))~n = funR~(n + a * a) - a * a - a
\end{theorem}

\begin{zproof}
prove;
use BEValSub;
invoke BEValSet;
with enabled (inPower) prove;
instantiate e == n;
prove;
instantiate a\_\_1 == a;
prove;
\end{zproof}

\begin{zed}
  GEValSet == \{ n: \nat | \forall a: \nat @ opBoxA~a~(comp(iter[\nat]~(1 + a)~funS, opBoxA~a~(comp(funS, opPlusI~(comp(comp(funD, funQ), iter[\nat]~(a * a + a + 4)~funS))))))~n = n * n + 2 * a * n + a \}
\end{zed}

\begin{zproof}[GEValSet\$domainCheck]
prove;
\end{zproof}

\begin{theorem}{rule GEValBase}
  0 \in GEValSet
\end{theorem}

\begin{zproof}
invoke GEValSet;
prove;
\end{zproof}

\begin{theorem}{rule GEValStep}
  \forall n: GEValSet @ 1 + n \in GEValSet
\end{theorem}

\begin{zproof}
invoke GEValSet;
prove;
instantiate a\_\_0 == a;
use BEValAll[n := n * n + 2 * a * n + a, a := a];
use RtSqAdd[m := n + a, k := a];
prove;
\end{zproof}

\begin{theorem}{GEValSub}
  \nat \subseteq GEValSet
\end{theorem}

\begin{zproof}
apply natInduction to predicate \nat \subseteq GEValSet;
prove;
invoke GEValSet;
prove;
\end{zproof}

\begin{theorem}{GEValAll}
  \forall n: \nat @ \forall a: \nat @ opBoxA~a~(comp(iter[\nat]~(1 + a)~funS, opBoxA~a~(comp(funS, opPlusI~(comp(comp(funD, funQ), iter[\nat]~(a * a + a + 4)~funS))))))~n = n * n + 2 * a * n + a
\end{theorem}

\begin{zproof}
prove;
use GEValSub;
invoke GEValSet;
with enabled (inPower) prove;
instantiate e == n;
prove;
instantiate a\_\_1 == a;
prove;
\end{zproof}

\begin{theorem}{PEVal}
  \forall a, k: \nat | 2 \leq a @ comp(comp(funE, iter[\nat]~(a * a - a - 1)~funS), opPlusI~(opBoxA~a~(comp(iter[\nat]~(1 + a)~funS, opBoxA~a~(comp(funS, opPlusI~(comp(comp(funD, funQ), iter[\nat]~(a * a + a + 4)~funS))))))))~(1 + k) = k
\end{theorem}

\begin{zproof}
prove;
use timesMonotonic2[a := 2, b := a, c := a, d := a];
use GEValAll[n := 1 + k, a := a];
use ESqAdd[m := 1 + k + a, k := k];
prove;
\end{zproof}

\begin{theorem}{CBComp}
  \forall X: \nat \fun \nat; a: \nat; F, G: \power (\num \cross \num) | F \in clos~\{ X \}~\{ opPlusI, opBoxA~a \}~\{ comp \} \land G \in clos~\{ X \}~\{ opPlusI, opBoxA~a \}~\{ comp \} @ comp(F, G) \in clos~\{ X \}~\{ opPlusI, opBoxA~a \}~\{ comp \}
\end{theorem}

\begin{zproof}
prove;
use ClosElemUF[I := \{ X \}, U := \{ opPlusI, opBoxA~a \}, B := \{ comp \}, F := F];
use ClosElemUF[I := \{ X \}, U := \{ opPlusI, opBoxA~a \}, B := \{ comp \}, F := G];
prove;
\end{zproof}

\begin{theorem}{CBPlusI}
  \forall X: \nat \fun \nat; a: \nat; F: \power (\num \cross \num) | F \in clos~\{ X \}~\{ opPlusI, opBoxA~a \}~\{ comp \} @ opPlusI~F \in clos~\{ X \}~\{ opPlusI, opBoxA~a \}~\{ comp \}
\end{theorem}

\begin{zproof}
prove;
use ClosElemUF[I := \{ X \}, U := \{ opPlusI, opBoxA~a \}, B := \{ comp \}, F := F];
prove;
\end{zproof}

\begin{theorem}{CBBoxA}
  \forall X: \nat \fun \nat; a: \nat; F: \power (\num \cross \num) | F \in clos~\{ X \}~\{ opPlusI, opBoxA~a \}~\{ comp \} @ opBoxA~a~F \in clos~\{ X \}~\{ opPlusI, opBoxA~a \}~\{ comp \}
\end{theorem}

\begin{zproof}
prove;
use ClosElemUF[I := \{ X \}, U := \{ opPlusI, opBoxA~a \}, B := \{ comp \}, F := F];
prove;
\end{zproof}

\begin{theorem}{QExprInCE}
  \forall a: \nat | 1 \leq a @ comp(funO, funE) \in clos~\{ funE \}~\{ opPlusI, opBoxA~a \}~\{ comp \}
\end{theorem}

\begin{zproof}
prove;
use ZeroInCB[X := funE, a := a];
use EDecr;
use EMult;
use ExtAOne[X := funE, a := a];
prove;
\end{zproof}

\begin{theorem}{QInCE}
  \forall a: \nat | 1 \leq a @ funQ \in clos~\{ funE \}~\{ opPlusI, opBoxA~a \}~\{ comp \}
\end{theorem}

\begin{zproof}
prove;
use QExprInCE[a := a];
use QExpr;
prove;
\end{zproof}

\begin{theorem}{PEInCE}
  \forall a: \nat | 2 \leq a @ comp(comp(funE, iter[\nat]~(a * a - a - 1)~funS), opPlusI~(opBoxA~a~(comp(iter[\nat]~(1 + a)~funS, opBoxA~a~(comp(funS, opPlusI~(comp(comp(funD, funQ), iter[\nat]~(a * a + a + 4)~funS)))))))) \in clos~\{ funE \}~\{ opPlusI, opBoxA~a \}~\{ comp \}
\end{theorem}

\begin{zproof}
prove;
use timesMonotonic2[a := 2, b := a, c := a, d := a];
use ZeroInCB[X := funE, a := a];
use EDecr;
use EMult;
use ExtAI[X := funE, a := a];
use ExtAS[X := funE, a := a];
use QInCE[a := a];
prove;
use IterIn[I := \{ funE \}, U := \{ opPlusI, opBoxA~a \}, B := \{ comp \}, F := funS, k := a * a + a + 4];
use IterIn[I := \{ funE \}, U := \{ opPlusI, opBoxA~a \}, B := \{ comp \}, F := funS, k := 1 + a];
use IterIn[I := \{ funE \}, U := \{ opPlusI, opBoxA~a \}, B := \{ comp \}, F := funS, k := a * a - a - 1];
use CBComp[X := funE, a := a, F := funD, G := funQ];
prove;
use CBComp[X := funE, a := a, F := comp(funD, funQ), G := iter[\nat]~(a * a + a + 4)~funS];
prove;
use CBPlusI[X := funE, a := a, F := comp(comp(funD, funQ), iter[\nat]~(a * a + a + 4)~funS)];
prove;
\end{zproof}

\begin{theorem}{PEUF}
  \forall a: \nat | 2 \leq a @ comp(comp(funE, iter[\nat]~(a * a - a - 1)~funS), opPlusI~(opBoxA~a~(comp(iter[\nat]~(1 + a)~funS, opBoxA~a~(comp(funS, opPlusI~(comp(comp(funD, funQ), iter[\nat]~(a * a + a + 4)~funS)))))))) \in \nat \fun \nat
\end{theorem}

\begin{zproof}
prove;
use PEInCE[a := a];
use ClosInPow[I := \{ funE \}, U := \{ opPlusI, opBoxA~a \}, B := \{ comp \}];
apply inPower to predicate clos~\{ funE \}~\{ opPlusI, opBoxA~a \}~\{ comp \} \in \power (\nat \fun \nat);
instantiate e == comp(comp(funE, iter[\nat]~(a * a - a - 1)~funS), opPlusI~(opBoxA~a~(comp(iter[\nat]~(1 + a)~funS, opBoxA~a~(comp(funS, opPlusI~(comp(comp(funD, funQ), iter[\nat]~(a * a + a + 4)~funS))))))));
prove;
\end{zproof}

\begin{theorem}{PEPred}
  \forall a: \nat | 2 \leq a @ comp(comp(funE, iter[\nat]~(a * a - a - 1)~funS), opPlusI~(opBoxA~a~(comp(iter[\nat]~(1 + a)~funS, opBoxA~a~(comp(funS, opPlusI~(comp(comp(funD, funQ), iter[\nat]~(a * a + a + 4)~funS)))))))) \in Preds
\end{theorem}

\begin{zproof}
prove;
use PEInCE[a := a];
use ClosInPow[I := \{ funE \}, U := \{ opPlusI, opBoxA~a \}, B := \{ comp \}];
apply inPower to predicate clos~\{ funE \}~\{ opPlusI, opBoxA~a \}~\{ comp \} \in \power (\nat \fun \nat);
instantiate e == comp(comp(funE, iter[\nat]~(a * a - a - 1)~funS), opPlusI~(opBoxA~a~(comp(iter[\nat]~(1 + a)~funS, opBoxA~a~(comp(funS, opPlusI~(comp(comp(funD, funQ), iter[\nat]~(a * a + a + 4)~funS))))))));
prove;
invoke Preds;
prove;
use PEVal[a := a, k := x];
prove;
\end{zproof}

\begin{theorem}{ExtAE}
  \forall a: \nat | 2 \leq a @ clos~\{ funE \}~\{ opPlusI, opBoxA~a \}~\{ comp \} = C6
\end{theorem}

\begin{zproof}
prove;
use EDecr;
use EMult;
use ZeroInCB[X := funE, a := a];
use PEInCE[a := a];
use PEUF[a := a];
use PEPred[a := a];
use ExtASup[X := funE, P := comp(comp(funE, iter[\nat]~(a * a - a - 1)~funS), opPlusI~(opBoxA~a~(comp(iter[\nat]~(1 + a)~funS, opBoxA~a~(comp(funS, opPlusI~(comp(comp(funD, funQ), iter[\nat]~(a * a + a + 4)~funS)))))))), a := a];
use ExtASub6[X := funE, a := a];
use C6Def;
use ClosInPow[I := \{ funE \}, U := \{ opPlusI, opBoxA~a \}, B := \{ comp \}];
use C6InPower;
use SubAntisym[\nat \fun \nat][A := clos~\{ funE \}~\{ opPlusI, opBoxA~a \}~\{ comp \}, B := C6];
prove;
\end{zproof}


\subsection{The case $X = K$}

\begin{theorem}{KLt}
  \forall y: \nat | 0 < y @ funK~y < y
\end{theorem}

\begin{zproof}
prove;
use Decomp[x := y];
use UFNonneg[f := funV, x := y];
use UFNonneg[f := funK, x := y];
prove;
split funV~y = 0;
cases;
use AZero;
prove;
next;
use AMono[m := 1, n := funV~y];
use AZero;
prove;
next;
\end{zproof}

\begin{theorem}{KZero}
  funK~0 = 0
\end{theorem}

\begin{zproof}
use KAt[m := 0, t := 0, k := 0];
use AZero;
prove;
\end{zproof}

\begin{theorem}{KDecr}
  funK \in Decr
\end{theorem}

\begin{zproof}
invoke Decr;
use KZero;
prove;
use KLt[y := y];
prove;
\end{zproof}

$K(2i(i+1)) = i$, since $2i(i+1) = A(2i) + i$.

\begin{theorem}{KMult}
  funK \in Mult
\end{theorem}

\begin{zproof}
invoke Mult;
use KZero;
prove;
instantiate r == 2 * i;
use TwoA[x := 2 * i];
use KAt[m := 2 * i * (1 + i), t := 2 * i, k := i];
prove;
\end{zproof}

With $c = A(a)$: $B = (S^2 (K S^c)^+)^{\square(a)}$ satisfies $B(n) + A(a) + 1 = A(n + a + 1)$,
and $\hat P = K S^c B^+$ is a predecessor.

\begin{zed}
  BKValSet == \{ n: \nat | \forall a: \nat @ (opBoxA~a~(comp(iter[\nat]~2~funS, opPlusI~(comp(funK, iter[\nat]~(funA~a)~funS)))))~n + funA~a + 1 = funA~(n + a + 1) \}
\end{zed}

\begin{zproof}[BKValSet\$domainCheck]
prove;
\end{zproof}

\begin{theorem}{rule BKValBase}
  0 \in BKValSet
\end{theorem}

\begin{zproof}
invoke BKValSet;
prove;
use ASuc[x := a];
prove;
\end{zproof}

\begin{theorem}{rule BKValStep}
  \forall n: BKValSet @ 1 + n \in BKValSet
\end{theorem}

\begin{zproof}
invoke BKValSet;
prove;
instantiate a\_\_0 == a;
use ASuc[x := n + a];
use ASuc[x := 1 + n + a];
use KAt[m := (opBoxA~a~(comp(iter[\nat]~2~funS, opPlusI~(comp(funK, iter[\nat]~(funA~a)~funS)))))~n + funA~a, t := n + a, k := n + a];
prove;
\end{zproof}

\begin{theorem}{BKValSub}
  \nat \subseteq BKValSet
\end{theorem}

\begin{zproof}
apply natInduction to predicate \nat \subseteq BKValSet;
prove;
invoke BKValSet;
prove;
\end{zproof}

\begin{theorem}{BKValAll}
  \forall n, a: \nat @ (opBoxA~a~(comp(iter[\nat]~2~funS, opPlusI~(comp(funK, iter[\nat]~(funA~a)~funS)))))~n + funA~a + 1 = funA~(n + a + 1)
\end{theorem}

\begin{zproof}
prove;
use BKValSub;
invoke BKValSet;
with enabled (inPower) prove;
instantiate e == n;
prove;
instantiate a\_\_1 == a;
prove;
\end{zproof}

\begin{theorem}{PKVal}
  \forall a, k: \nat @ comp(comp(funK, iter[\nat]~(funA~a)~funS), opPlusI~(opBoxA~a~(comp(iter[\nat]~2~funS, opPlusI~(comp(funK, iter[\nat]~(funA~a)~funS))))))~(1 + k) = k
\end{theorem}

\begin{zproof}
prove;
use BKValAll[n := 1 + k, a := a];
use KAt[m := (opBoxA~a~(comp(iter[\nat]~2~funS, opPlusI~(comp(funK, iter[\nat]~(funA~a)~funS)))))~(1 + k) + (1 + k) + funA~a, t := k + a + 2, k := k];
prove;
\end{zproof}

\begin{theorem}{PKInCK}
  \forall a: \nat | 1 \leq a @ comp(comp(funK, iter[\nat]~(funA~a)~funS), opPlusI~(opBoxA~a~(comp(iter[\nat]~2~funS, opPlusI~(comp(funK, iter[\nat]~(funA~a)~funS)))))) \in clos~\{ funK \}~\{ opPlusI, opBoxA~a \}~\{ comp \}
\end{theorem}

\begin{zproof}
prove;
use ZeroInCB[X := funK, a := a];
use KDecr;
use KMult;
use ExtAI[X := funK, a := a];
use ExtAS[X := funK, a := a];
prove;
use IterIn[I := \{ funK \}, U := \{ opPlusI, opBoxA~a \}, B := \{ comp \}, F := funS, k := funA~a];
use IterIn[I := \{ funK \}, U := \{ opPlusI, opBoxA~a \}, B := \{ comp \}, F := funS, k := 2];
prove;
\end{zproof}

\begin{theorem}{PKUF}
  \forall a: \nat | 1 \leq a @ comp(comp(funK, iter[\nat]~(funA~a)~funS), opPlusI~(opBoxA~a~(comp(iter[\nat]~2~funS, opPlusI~(comp(funK, iter[\nat]~(funA~a)~funS)))))) \in \nat \fun \nat
\end{theorem}

\begin{zproof}
prove;
\end{zproof}

\begin{theorem}{PKPred}
  \forall a: \nat | 1 \leq a @ comp(comp(funK, iter[\nat]~(funA~a)~funS), opPlusI~(opBoxA~a~(comp(iter[\nat]~2~funS, opPlusI~(comp(funK, iter[\nat]~(funA~a)~funS)))))) \in Preds
\end{theorem}

\begin{zproof}
prove;
use PKUF[a := a];
invoke Preds;
prove;
use PKVal[a := a, k := x];
prove;
\end{zproof}

\begin{theorem}{ExtAK}
  \forall a: \nat | 1 \leq a @ clos~\{ funK \}~\{ opPlusI, opBoxA~a \}~\{ comp \} = C6
\end{theorem}

\begin{zproof}
prove;
use KDecr;
use KMult;
use ZeroInCB[X := funK, a := a];
use PKInCK[a := a];
use PKUF[a := a];
use PKPred[a := a];
use ExtASup[X := funK, P := comp(comp(funK, iter[\nat]~(funA~a)~funS), opPlusI~(opBoxA~a~(comp(iter[\nat]~2~funS, opPlusI~(comp(funK, iter[\nat]~(funA~a)~funS)))))), a := a];
use ExtASub6[X := funK, a := a];
use Theorem68K;
use C6KDef;
use ClosInPow[I := \{ funK \}, U := \{ opPlusI, opBoxA~a \}, B := \{ comp \}];
use C6InPower;
use SubAntisym[\nat \fun \nat][A := clos~\{ funK \}~\{ opPlusI, opBoxA~a \}~\{ comp \}, B := C6];
prove;
\end{zproof}


\subsection{The case $X = L$}

$\bar 0 = L L^+$, because $L(x + L(x)) = 0$.

\begin{theorem}{LZeroExpr}
  comp(funL, opPlusI~funL) = cst~0
\end{theorem}

\begin{zproof}
apply UFExt to predicate comp(funL, opPlusI~funL) = cst~0;
prove;
use Decomp[x := x];
prove;
use LAt[m := x, t := funV~x, k := funK~x];
prove;
use LAt[m := funA~(funV~x) + funV~x, t := funV~x, k := funV~x];
prove;
\end{zproof}

\begin{theorem}{LZeroExprInCL}
  \forall a: \nat @ comp(funL, opPlusI~funL) \in clos~\{ funL \}~\{ opPlusI, opBoxA~a \}~\{ comp \}
\end{theorem}

\begin{zproof}
prove;
\end{zproof}

\begin{theorem}{ZeroInCL}
  \forall a: \nat @ cst~0 \in clos~\{ funL \}~\{ opPlusI, opBoxA~a \}~\{ comp \}
\end{theorem}

\begin{zproof}
prove;
use LZeroExprInCL[a := a];
use LZeroExpr;
prove;
\end{zproof}

$L(2(i+1)(i+1)) = i$, since $2(i+1)(i+1) = A(2i+1) + i + 1$.

\begin{theorem}{LMult}
  funL \in Mult
\end{theorem}

\begin{zproof}
invoke Mult;
use LAt[m := 0, t := 0, k := 0];
use AZero;
prove;
instantiate r == 2 + 2 * i;
use TwoA[x := 1 + 2 * i];
use LAt[m := (2 + 2 * i) * (1 + i), t := 1 + 2 * i, k := 1 + i];
prove;
\end{zproof}

$K = L \bar 1^+ ((\bar 1^+)^+ L)^+$ (Theorem 6.8), so the case of $L$ reduces to that of $K$.

\begin{theorem}{KExprInCL}
  \forall a: \nat | 1 \leq a @ comp(comp(funL, opPlusI~(cst~1)), opPlusI~(comp(opPlusI~(opPlusI~(cst~1)), funL))) \in clos~\{ funL \}~\{ opPlusI, opBoxA~a \}~\{ comp \}
\end{theorem}

\begin{zproof}
prove;
use ZeroInCL[a := a];
use LMult;
use ExtAOne[X := funL, a := a];
prove;
\end{zproof}

\begin{theorem}{KInCL}
  \forall a: \nat | 1 \leq a @ funK \in clos~\{ funL \}~\{ opPlusI, opBoxA~a \}~\{ comp \}
\end{theorem}

\begin{zproof}
prove;
use KExprInCL[a := a];
use KExprL;
prove;
\end{zproof}

\begin{theorem}{CBUClosedA}
  \forall X: \nat \fun \nat; a: \nat @ clos~\{ X \}~\{ opPlusI, opBoxA~a \}~\{ comp \} \in UClosed~\{ opPlusI, opBoxA~a \}
\end{theorem}

\begin{zproof}
prove;
use ClosInPow[I := \{ X \}, U := \{ opPlusI, opBoxA~a \}, B := \{ comp \}];
with enabled (UClosedDef) prove;
\end{zproof}

\begin{theorem}{CKSubCL}
  \forall a: \nat | 1 \leq a @ clos~\{ funK \}~\{ opPlusI, opBoxA~a \}~\{ comp \} \subseteq clos~\{ funL \}~\{ opPlusI, opBoxA~a \}~\{ comp \}
\end{theorem}

\begin{zproof}
prove;
use KInCL[a := a];
use ClosLeast[I := \{ funK \}, U := \{ opPlusI, opBoxA~a \}, B := \{ comp \}, Z := clos~\{ funL \}~\{ opPlusI, opBoxA~a \}~\{ comp \}];
use ClosInPow[I := \{ funL \}, U := \{ opPlusI, opBoxA~a \}, B := \{ comp \}];
use CBUClosedA[X := funL, a := a];
use CBBClosed[X := funL, a := a];
prove;
\end{zproof}

\begin{theorem}{ExtAL}
  \forall a: \nat | 1 \leq a @ clos~\{ funL \}~\{ opPlusI, opBoxA~a \}~\{ comp \} = C6
\end{theorem}

\begin{zproof}
prove;
use ExtAK[a := a];
use CKSubCL[a := a];
use ExtASub6[X := funL, a := a];
use ClosInPow[I := \{ funL \}, U := \{ opPlusI, opBoxA~a \}, B := \{ comp \}];
use C6InPower;
use SubAntisym[\nat \fun \nat][A := clos~\{ funL \}~\{ opPlusI, opBoxA~a \}~\{ comp \}, B := C6];
prove;
\end{zproof}


\subsection{Every starting value $a \geq 1$ (Remark 6.9)}

For $a \geq 1$ the constructions for $K$ and $L$ above already work (\textsf{ExtAK}, \textsf{ExtAL}).
For $E$ and $a = 1$: $G = (S^2 B)^{\square(1)}$ (the $G$ above with $a = 1$) satisfies $G(n) = (n+1)^2$,
and $\hat P = E S^2 ((G^+)^+)^+$ is a predecessor, since $(x+2)^2 + 3(x+1) + 2 = (x+3)^2 + x$.

\begin{theorem}{PEOneVal}
  \forall k: \nat @ comp(comp(funE, iter[\nat]~2~funS), opPlusI~(opPlusI~(opPlusI~(opBoxA~1~(comp(iter[\nat]~2~funS, opBoxA~1~(comp(funS, opPlusI~(comp(comp(funD, funQ), iter[\nat]~6~funS))))))))))~(1 + k) = k
\end{theorem}

\begin{zproof}
prove;
use GEValAll[n := 1 + k, a := 1];
use ESqAdd[m := 3 + k, k := k];
prove;
\end{zproof}

\begin{theorem}{PEOneInCE}
  comp(comp(funE, iter[\nat]~2~funS), opPlusI~(opPlusI~(opPlusI~(opBoxA~1~(comp(iter[\nat]~2~funS, opBoxA~1~(comp(funS, opPlusI~(comp(comp(funD, funQ), iter[\nat]~6~funS)))))))))) \in clos~\{ funE \}~\{ opPlusI, opBoxA~1 \}~\{ comp \}
\end{theorem}

\begin{zproof}
use ZeroInCB[X := funE, a := 1];
use EDecr;
use EMult;
use ExtAI[X := funE, a := 1];
use ExtAS[X := funE, a := 1];
use QInCE[a := 1];
prove;
use IterIn[I := \{ funE \}, U := \{ opPlusI, opBoxA~1 \}, B := \{ comp \}, F := funS, k := 6];
use IterIn[I := \{ funE \}, U := \{ opPlusI, opBoxA~1 \}, B := \{ comp \}, F := funS, k := 2];
prove;
use CBComp[X := funE, a := 1, F := funD, G := funQ];
prove;
use CBComp[X := funE, a := 1, F := comp(funD, funQ), G := iter[\nat]~6~funS];
prove;
use CBPlusI[X := funE, a := 1, F := comp(comp(funD, funQ), iter[\nat]~6~funS)];
prove;
use CBComp[X := funE, a := 1, F := funS, G := opPlusI~(comp(comp(funD, funQ), iter[\nat]~6~funS))];
prove;
use CBBoxA[X := funE, a := 1, F := comp(funS, opPlusI~(comp(comp(funD, funQ), iter[\nat]~6~funS)))];
prove;
use CBComp[X := funE, a := 1, F := iter[\nat]~2~funS, G := opBoxA~1~(comp(funS, opPlusI~(comp(comp(funD, funQ), iter[\nat]~6~funS))))];
prove;
use CBBoxA[X := funE, a := 1, F := comp(iter[\nat]~2~funS, opBoxA~1~(comp(funS, opPlusI~(comp(comp(funD, funQ), iter[\nat]~6~funS)))))];
prove;
use CBPlusI[X := funE, a := 1, F := opBoxA~1~(comp(iter[\nat]~2~funS, opBoxA~1~(comp(funS, opPlusI~(comp(comp(funD, funQ), iter[\nat]~6~funS))))))];
prove;
use CBPlusI[X := funE, a := 1, F := opPlusI~(opBoxA~1~(comp(iter[\nat]~2~funS, opBoxA~1~(comp(funS, opPlusI~(comp(comp(funD, funQ), iter[\nat]~6~funS)))))))];
prove;
use CBPlusI[X := funE, a := 1, F := opPlusI~(opPlusI~(opBoxA~1~(comp(iter[\nat]~2~funS, opBoxA~1~(comp(funS, opPlusI~(comp(comp(funD, funQ), iter[\nat]~6~funS))))))))];
prove;
use CBComp[X := funE, a := 1, F := funE, G := iter[\nat]~2~funS];
prove;
use CBComp[X := funE, a := 1, F := comp(funE, iter[\nat]~2~funS), G := opPlusI~(opPlusI~(opPlusI~(opBoxA~1~(comp(iter[\nat]~2~funS, opBoxA~1~(comp(funS, opPlusI~(comp(comp(funD, funQ), iter[\nat]~6~funS)))))))))];
prove;
\end{zproof}

\begin{theorem}{PEOneUF}
  comp(comp(funE, iter[\nat]~2~funS), opPlusI~(opPlusI~(opPlusI~(opBoxA~1~(comp(iter[\nat]~2~funS, opBoxA~1~(comp(funS, opPlusI~(comp(comp(funD, funQ), iter[\nat]~6~funS)))))))))) \in \nat \fun \nat
\end{theorem}

\begin{zproof}
use PEOneInCE;
use ClosInPow[I := \{ funE \}, U := \{ opPlusI, opBoxA~1 \}, B := \{ comp \}];
apply inPower to predicate clos~\{ funE \}~\{ opPlusI, opBoxA~1 \}~\{ comp \} \in \power (\nat \fun \nat);
instantiate e == comp(comp(funE, iter[\nat]~2~funS), opPlusI~(opPlusI~(opPlusI~(opBoxA~1~(comp(iter[\nat]~2~funS, opBoxA~1~(comp(funS, opPlusI~(comp(comp(funD, funQ), iter[\nat]~6~funS))))))))));
prove;
\end{zproof}

\begin{theorem}{PEOnePred}
  comp(comp(funE, iter[\nat]~2~funS), opPlusI~(opPlusI~(opPlusI~(opBoxA~1~(comp(iter[\nat]~2~funS, opBoxA~1~(comp(funS, opPlusI~(comp(comp(funD, funQ), iter[\nat]~6~funS)))))))))) \in Preds
\end{theorem}

\begin{zproof}
use PEOneUF;
invoke Preds;
prove;
use PEOneVal[k := x];
prove;
\end{zproof}

\begin{theorem}{ExtAEOne}
  clos~\{ funE \}~\{ opPlusI, opBoxA~1 \}~\{ comp \} = C6
\end{theorem}

\begin{zproof}
use EDecr;
use EMult;
use ZeroInCB[X := funE, a := 1];
use PEOneInCE;
use PEOneUF;
use PEOnePred;
use ExtASup[X := funE, P := comp(comp(funE, iter[\nat]~2~funS), opPlusI~(opPlusI~(opPlusI~(opBoxA~1~(comp(iter[\nat]~2~funS, opBoxA~1~(comp(funS, opPlusI~(comp(comp(funD, funQ), iter[\nat]~6~funS)))))))))), a := 1];
use ExtASub6[X := funE, a := 1];
use C6Def;
use ClosInPow[I := \{ funE \}, U := \{ opPlusI, opBoxA~1 \}, B := \{ comp \}];
use C6InPower;
use SubAntisym[\nat \fun \nat][A := clos~\{ funE \}~\{ opPlusI, opBoxA~1 \}~\{ comp \}, B := C6];
prove;
\end{zproof}

\begin{theorem}{Remark69}
  clos~\{ funE \}~\{ opPlusI, opBoxA~1 \}~\{ comp \} = C6 \land clos~\{ funK \}~\{ opPlusI, opBoxA~1 \}~\{ comp \} = C6 \land clos~\{ funL \}~\{ opPlusI, opBoxA~1 \}~\{ comp \} = C6
\end{theorem}

\begin{zproof}
use ExtAEOne;
use ExtAK[a := 1];
use ExtAL[a := 1];
prove;
\end{zproof}

\begin{theorem}{ExtAAll}
  \forall a: \nat; X: \nat \fun \nat | 1 \leq a \land X \in \{ funE, funK, funL \} @ clos~\{ X \}~\{ opPlusI, opBoxA~a \}~\{ comp \} = C6
\end{theorem}

\begin{zproof}
prove;
use ExtAK[a := a];
use ExtAL[a := a];
split a = 1;
cases;
use ExtAEOne;
prove;
split X = funE;
cases;
prove;
next;
split X = funK;
cases;
prove;
next;
prove;
next;
next;
use ExtAE[a := a];
prove;
split X = funE;
cases;
prove;
next;
split X = funK;
cases;
prove;
next;
prove;
next;
next;
\end{zproof}


\begin{theorem}{rule UFInPow}
  \forall F: \nat \fun \nat @ F \in \power (\num \cross \num)
\end{theorem}

\begin{zproof}
prove;
\end{zproof}

\section{The shifted operator $F^{+b}$ (open question 2.2)}

For fixed $b$, $F^{+b}(x) = F(x) + x + b$. The basis $\langle E, F^{+b}, FG, F^\square \rangle$ generates
the same functions as Robinson's basis iff $b \geq 1$.

\begin{axdef}
  opPlusIb: \nat \fun ((\nat \fun \nat) \fun (\nat \fun \nat))
\where
  \Label{rule PlusIbVal} \forall b: \nat; F: \nat \fun \nat; x: \nat @ opPlusIb~b~F~x = F~x + x + b
\end{axdef}

\begin{zproof}[opPlusIb\$domainCheck]
prove;
\end{zproof}

\begin{theorem}{rule PlusIbInner}
  \forall b: \nat; F: \nat \fun \nat @ (\lambda x: \nat @ F~x + x + b) \in \nat \fun \nat
\end{theorem}

\begin{zproof}
apply FunIntro to predicate (\lambda x: \nat @ F~x + x + b) \in \nat \fun \nat;
with enabled (RelPow, inPower) prove by reduce;
\end{zproof}

\begin{theorem}{rule PlusIbMid}
  \forall b: \nat @ (\lambda F: \nat \fun \nat @ (\lambda x: \nat @ F~x + x + b)) \in (\nat \fun \nat) \fun (\nat \fun \nat)
\end{theorem}

\begin{zproof}
apply FunIntro to predicate (\lambda F: \nat \fun \nat @ (\lambda x: \nat @ F~x + x + b)) \in (\nat \fun \nat) \fun (\nat \fun \nat);
with enabled (RelPow, inPower) prove by reduce;
\end{zproof}

\begin{zed}
  PlusIbWit == (\lambda b: \nat @ (\lambda F: \nat \fun \nat @ (\lambda x: \nat @ F~x + x + b)))
\end{zed}

\begin{zproof}[PlusIbWit\$domainCheck]
prove;
\end{zproof}

\begin{theorem}{PlusIbWitFun}
  PlusIbWit \in \nat \fun ((\nat \fun \nat) \fun (\nat \fun \nat))
\end{theorem}

\begin{zproof}
invoke PlusIbWit;
apply FunIntro to predicate (\lambda b: \nat @ (\lambda F: \nat \fun \nat @ (\lambda x: \nat @ F~x + x + b))) \in \nat \fun ((\nat \fun \nat) \fun (\nat \fun \nat));
with enabled (RelPow, inPower) prove by reduce;
\end{zproof}

\begin{theorem}{PlusIbWitApp}
  \forall b: \nat @ PlusIbWit~b = (\lambda F: \nat \fun \nat @ (\lambda x: \nat @ F~x + x + b))
\end{theorem}

\begin{zproof}
use PlusIbWitFun;
use pairInFunction[\nat, (\nat \fun \nat) \fun (\nat \fun \nat)][f := PlusIbWit, x := b, y := (\lambda F: \nat \fun \nat @ (\lambda x: \nat @ F~x + x + b))];
invoke PlusIbWit;
prove by reduce;
\end{zproof}

\begin{theorem}{PlusIbWitApp2}
  \forall b: \nat; F: \nat \fun \nat @ PlusIbWit~b~F = (\lambda x: \nat @ F~x + x + b)
\end{theorem}

\begin{zproof}
prove;
use PlusIbWitApp[b := b];
use pairInFunction[\nat \fun \nat, \nat \fun \nat][f := (\lambda F: \nat \fun \nat @ (\lambda x: \nat @ F~x + x + b)), x := F, y := (\lambda x: \nat @ F~x + x + b)];
prove;
\end{zproof}

\begin{theorem}{rule PlusIbWitVal}
  \forall b: \nat; F: \nat \fun \nat; x: \nat @ PlusIbWit~b~F~x = F~x + x + b
\end{theorem}

\begin{zproof}
prove;
use PlusIbWitApp2[b := b, F := F];
use pairInFunction[\nat, \nat][f := (\lambda x: \nat @ F~x + x + b), x := x, y := F~x + x + b];
prove;
\end{zproof}

\begin{theorem}{PlusIbEx}
  \exists opPlusIb: \nat \fun ((\nat \fun \nat) \fun (\nat \fun \nat)) @
    \forall b: \nat; F: \nat \fun \nat; x: \nat @ opPlusIb~b~F~x = F~x + x + b
\end{theorem}

\begin{zproof}
instantiate opPlusIb == PlusIbWit;
use PlusIbWitFun;
prove;
\end{zproof}

\subsection{Shifts}

\begin{zed}
  Shift == \{ F: \nat \fun \nat; c: \nat | \forall y: \nat @ F~y = y + c \}
\end{zed}

\begin{zproof}[Shift\$domainCheck]
prove;
\end{zproof}

\begin{zed}
  ShiftPowSet == \{ k: \nat | \forall F: \nat \fun \nat; c, x: \nat | (F, c) \in Shift @ iter[\nat]~k~F~x = x + k * c \}
\end{zed}

\begin{zproof}[ShiftPowSet\$domainCheck]
prove;
\end{zproof}

\begin{theorem}{rule ShiftPowBase}
  0 \in ShiftPowSet
\end{theorem}

\begin{zproof}
invoke ShiftPowSet;
prove;
\end{zproof}

\begin{theorem}{rule ShiftPowStep}
  \forall k: ShiftPowSet @ 1 + k \in ShiftPowSet
\end{theorem}

\begin{zproof}
invoke ShiftPowSet;
prove;
instantiate F\_\_0 == F, c\_\_0 == c, x\_\_0 == x;
invoke Shift;
prove;
instantiate y\_\_0 == iter[\nat]~k~F~x;
prove;
\end{zproof}

\begin{theorem}{ShiftPowSub}
  \nat \subseteq ShiftPowSet
\end{theorem}

\begin{zproof}
apply natInduction to predicate \nat \subseteq ShiftPowSet;
prove;
invoke ShiftPowSet;
prove;
\end{zproof}

\begin{theorem}{ShiftPow}
  \forall k, c, x: \nat; F: \nat \fun \nat | (F, c) \in Shift @ iter[\nat]~k~F~x = x + k * c
\end{theorem}

\begin{zproof}
prove;
use ShiftPowSub;
invoke ShiftPowSet;
with enabled (inPower) prove;
instantiate e == k;
prove;
instantiate F\_\_1 == F, c\_\_1 == c, x\_\_1 == x;
prove;
\end{zproof}

\subsection{The construction, with $b = c + 1$}

$\bar 0 = E^\square$, $A = \bar 0^{+b}$ ($A(x) = x + b$), $C = E A^{b-1} \bar 0$ (the constant $b - 1$),
$B = C^{+b}$ ($B(x) = x + 2b - 1$), $T = B^{2b-2} A A$ ($T(x) = x + 4b^2 - 4b + 2$),
$H = A^3 (E T)^{+b}$, $G = H^\square$ ($G(n) = n^2 + 4bn$),
$\hat P = E A^{4b+1} ((G^{+b})^{+b})^{+b}$ (a predecessor), $\bar 1 = \hat P^{b-1} A \bar 0$, and
$F^+ = \hat P^b F^{+b}$.

\begin{theorem}{AShift}
  \forall c: \nat @ (opPlusIb~(1 + c)~(cst~0), 1 + c) \in Shift
\end{theorem}

\begin{zproof}
invoke Shift;
prove;
\end{zproof}

\begin{theorem}{rule CbVal}
  \forall c, x: \nat @ comp(funE, comp(iter[\nat]~c~(opPlusIb~(1 + c)~(cst~0)), cst~0))~x = c
\end{theorem}

\begin{zproof}
prove;
use AShift[c := c];
use ShiftPow[k := c, c := 1 + c, x := 0, F := opPlusIb~(1 + c)~(cst~0)];
use ESqAdd[m := c, k := c];
prove;
\end{zproof}

\begin{theorem}{BShift}
  \forall c: \nat @ (opPlusIb~(1 + c)~(comp(funE, comp(iter[\nat]~c~(opPlusIb~(1 + c)~(cst~0)), cst~0))), 1 + 2 * c) \in Shift
\end{theorem}

\begin{zproof}
invoke Shift;
prove;
\end{zproof}

\begin{theorem}{rule TbVal}
  \forall c, x: \nat @ comp(iter[\nat]~(2 * c)~(opPlusIb~(1 + c)~(comp(funE, comp(iter[\nat]~c~(opPlusIb~(1 + c)~(cst~0)), cst~0)))), comp(opPlusIb~(1 + c)~(cst~0), opPlusIb~(1 + c)~(cst~0)))~x = x + 4 * c * c + 4 * c + 2
\end{theorem}

\begin{zproof}
prove;
use BShift[c := c];
use ShiftPow[k := 2 * c, c := 1 + 2 * c, x := 2 + 2 * c + x, F := opPlusIb~(1 + c)~(comp(funE, comp(iter[\nat]~c~(opPlusIb~(1 + c)~(cst~0)), cst~0)))];
prove;
\end{zproof}

\begin{theorem}{rule HbVal}
  \forall c, y: \nat @ comp(iter[\nat]~3~(opPlusIb~(1 + c)~(cst~0)), opPlusIb~(1 + c)~(comp(funE, comp(iter[\nat]~(2 * c)~(opPlusIb~(1 + c)~(comp(funE, comp(iter[\nat]~c~(opPlusIb~(1 + c)~(cst~0)), cst~0)))), comp(opPlusIb~(1 + c)~(cst~0), opPlusIb~(1 + c)~(cst~0))))))~y = funE~(y + 4 * c * c + 4 * c + 2) + y + 4 * c + 4
\end{theorem}

\begin{zproof}
prove;
use AShift[c := c];
use ShiftPow[k := 3, c := 1 + c, x := funE~(y + 4 * c * c + 4 * c + 2) + y + 1 + c, F := opPlusIb~(1 + c)~(cst~0)];
prove;
\end{zproof}

\begin{zed}
  GbSet == \{ n: \nat | \forall c: \nat @ opBox~(comp(iter[\nat]~3~(opPlusIb~(1 + c)~(cst~0)), opPlusIb~(1 + c)~(comp(funE, comp(iter[\nat]~(2 * c)~(opPlusIb~(1 + c)~(comp(funE, comp(iter[\nat]~c~(opPlusIb~(1 + c)~(cst~0)), cst~0)))), comp(opPlusIb~(1 + c)~(cst~0), opPlusIb~(1 + c)~(cst~0)))))))~n = n * n + 4 * (1 + c) * n \}
\end{zed}

\begin{zproof}[GbSet\$domainCheck]
prove;
\end{zproof}

\begin{theorem}{rule GbBase}
  0 \in GbSet
\end{theorem}

\begin{zproof}
invoke GbSet;
prove;
\end{zproof}

\begin{theorem}{rule GbStep}
  \forall n: GbSet @ 1 + n \in GbSet
\end{theorem}

\begin{zproof}
invoke GbSet;
prove;
instantiate c\_\_0 == c;
use ESqAdd[m := n + 2 * c + 1, k := 2 * n + 1];
prove;
\end{zproof}

\begin{theorem}{GbSub}
  \nat \subseteq GbSet
\end{theorem}

\begin{zproof}
apply natInduction to predicate \nat \subseteq GbSet;
prove;
invoke GbSet;
prove;
\end{zproof}

\begin{theorem}{rule GbVal}
  \forall n, c: \nat @ opBox~(comp(iter[\nat]~3~(opPlusIb~(1 + c)~(cst~0)), opPlusIb~(1 + c)~(comp(funE, comp(iter[\nat]~(2 * c)~(opPlusIb~(1 + c)~(comp(funE, comp(iter[\nat]~c~(opPlusIb~(1 + c)~(cst~0)), cst~0)))), comp(opPlusIb~(1 + c)~(cst~0), opPlusIb~(1 + c)~(cst~0)))))))~n = n * n + 4 * (1 + c) * n
\end{theorem}

\begin{zproof}
prove;
use GbSub;
invoke GbSet;
with enabled (inPower) prove;
instantiate e == n;
prove;
instantiate c\_\_1 == c;
prove;
\end{zproof}

\begin{theorem}{PbVal}
  \forall c, k: \nat @ comp(funE, comp(iter[\nat]~(4 * c + 5)~(opPlusIb~(1 + c)~(cst~0)), opPlusIb~(1 + c)~(opPlusIb~(1 + c)~(opPlusIb~(1 + c)~(opBox~(comp(iter[\nat]~3~(opPlusIb~(1 + c)~(cst~0)), opPlusIb~(1 + c)~(comp(funE, comp(iter[\nat]~(2 * c)~(opPlusIb~(1 + c)~(comp(funE, comp(iter[\nat]~c~(opPlusIb~(1 + c)~(cst~0)), cst~0)))), comp(opPlusIb~(1 + c)~(cst~0), opPlusIb~(1 + c)~(cst~0))))))))))))~(1 + k) = k
\end{theorem}

\begin{zproof}
with disabled (BoxSuc, HbVal) prove;
use AShift[c := c];
use ShiftPow[k := 4 * c + 5, c := 1 + c, x := 6 + 3 * c + 3 * k + (1 + k) * (1 + k) + 4 * (1 + c) * (1 + k), F := opPlusIb~(1 + c)~(cst~0)];
use ESqAdd[m := k + 2 * c + 4, k := k];
with disabled (BoxSuc, HbVal) prove;
\end{zproof}


\begin{theorem}{PbPred}
  \forall c: \nat @ comp(funE, comp(iter[\nat]~(4 * c + 5)~(opPlusIb~(1 + c)~(cst~0)), opPlusIb~(1 + c)~(opPlusIb~(1 + c)~(opPlusIb~(1 + c)~(opBox~(comp(iter[\nat]~3~(opPlusIb~(1 + c)~(cst~0)), opPlusIb~(1 + c)~(comp(funE, comp(iter[\nat]~(2 * c)~(opPlusIb~(1 + c)~(comp(funE, comp(iter[\nat]~c~(opPlusIb~(1 + c)~(cst~0)), cst~0)))), comp(opPlusIb~(1 + c)~(cst~0), opPlusIb~(1 + c)~(cst~0)))))))))))) \in Preds
\end{theorem}

\begin{zproof}
prove;
invoke Preds;
prove;
use PbVal[c := c, k := x];
prove;
\end{zproof}

\begin{theorem}{OneEqVb}
  \forall c: \nat @ comp(iter[\nat]~c~(comp(funE, comp(iter[\nat]~(4 * c + 5)~(opPlusIb~(1 + c)~(cst~0)), opPlusIb~(1 + c)~(opPlusIb~(1 + c)~(opPlusIb~(1 + c)~(opBox~(comp(iter[\nat]~3~(opPlusIb~(1 + c)~(cst~0)), opPlusIb~(1 + c)~(comp(funE, comp(iter[\nat]~(2 * c)~(opPlusIb~(1 + c)~(comp(funE, comp(iter[\nat]~c~(opPlusIb~(1 + c)~(cst~0)), cst~0)))), comp(opPlusIb~(1 + c)~(cst~0), opPlusIb~(1 + c)~(cst~0))))))))))))), comp(opPlusIb~(1 + c)~(cst~0), cst~0)) = cst~1
\end{theorem}

\begin{zproof}
apply UFExt to predicate comp(iter[\nat]~c~(comp(funE, comp(iter[\nat]~(4 * c + 5)~(opPlusIb~(1 + c)~(cst~0)), opPlusIb~(1 + c)~(opPlusIb~(1 + c)~(opPlusIb~(1 + c)~(opBox~(comp(iter[\nat]~3~(opPlusIb~(1 + c)~(cst~0)), opPlusIb~(1 + c)~(comp(funE, comp(iter[\nat]~(2 * c)~(opPlusIb~(1 + c)~(comp(funE, comp(iter[\nat]~c~(opPlusIb~(1 + c)~(cst~0)), cst~0)))), comp(opPlusIb~(1 + c)~(cst~0), opPlusIb~(1 + c)~(cst~0))))))))))))), comp(opPlusIb~(1 + c)~(cst~0), cst~0)) = cst~1;
prove;
use PbPred[c := c];
use IterPAll[k := c, P := comp(funE, comp(iter[\nat]~(4 * c + 5)~(opPlusIb~(1 + c)~(cst~0)), opPlusIb~(1 + c)~(opPlusIb~(1 + c)~(opPlusIb~(1 + c)~(opBox~(comp(iter[\nat]~3~(opPlusIb~(1 + c)~(cst~0)), opPlusIb~(1 + c)~(comp(funE, comp(iter[\nat]~(2 * c)~(opPlusIb~(1 + c)~(comp(funE, comp(iter[\nat]~c~(opPlusIb~(1 + c)~(cst~0)), cst~0)))), comp(opPlusIb~(1 + c)~(cst~0), opPlusIb~(1 + c)~(cst~0)))))))))))), x := 1];
prove;
\end{zproof}

\begin{theorem}{PlusIEqVb}
  \forall c: \nat; F: \nat \fun \nat @ comp(iter[\nat]~(1 + c)~(comp(funE, comp(iter[\nat]~(4 * c + 5)~(opPlusIb~(1 + c)~(cst~0)), opPlusIb~(1 + c)~(opPlusIb~(1 + c)~(opPlusIb~(1 + c)~(opBox~(comp(iter[\nat]~3~(opPlusIb~(1 + c)~(cst~0)), opPlusIb~(1 + c)~(comp(funE, comp(iter[\nat]~(2 * c)~(opPlusIb~(1 + c)~(comp(funE, comp(iter[\nat]~c~(opPlusIb~(1 + c)~(cst~0)), cst~0)))), comp(opPlusIb~(1 + c)~(cst~0), opPlusIb~(1 + c)~(cst~0))))))))))))), opPlusIb~(1 + c)~F) = opPlusI~F
\end{theorem}

\begin{zproof}
apply UFExt to predicate comp(iter[\nat]~(1 + c)~(comp(funE, comp(iter[\nat]~(4 * c + 5)~(opPlusIb~(1 + c)~(cst~0)), opPlusIb~(1 + c)~(opPlusIb~(1 + c)~(opPlusIb~(1 + c)~(opBox~(comp(iter[\nat]~3~(opPlusIb~(1 + c)~(cst~0)), opPlusIb~(1 + c)~(comp(funE, comp(iter[\nat]~(2 * c)~(opPlusIb~(1 + c)~(comp(funE, comp(iter[\nat]~c~(opPlusIb~(1 + c)~(cst~0)), cst~0)))), comp(opPlusIb~(1 + c)~(cst~0), opPlusIb~(1 + c)~(cst~0))))))))))))), opPlusIb~(1 + c)~F) = opPlusI~F;
prove;
use PbPred[c := c];
use IterPAll[k := 1 + c, P := comp(funE, comp(iter[\nat]~(4 * c + 5)~(opPlusIb~(1 + c)~(cst~0)), opPlusIb~(1 + c)~(opPlusIb~(1 + c)~(opPlusIb~(1 + c)~(opBox~(comp(iter[\nat]~3~(opPlusIb~(1 + c)~(cst~0)), opPlusIb~(1 + c)~(comp(funE, comp(iter[\nat]~(2 * c)~(opPlusIb~(1 + c)~(comp(funE, comp(iter[\nat]~c~(opPlusIb~(1 + c)~(cst~0)), cst~0)))), comp(opPlusIb~(1 + c)~(cst~0), opPlusIb~(1 + c)~(cst~0)))))))))))), x := x + F~x];
prove;
\end{zproof}

\subsection{Closure steps in $\langle E, F^{+b}, FG, F^\square \rangle$}

\begin{theorem}{CVComp}
  \forall b: \nat; F, G: \power (\num \cross \num) | F \in clos~\{ funE \}~\{ opPlusIb~b, opBox \}~\{ comp \} \land G \in clos~\{ funE \}~\{ opPlusIb~b, opBox \}~\{ comp \} @ comp(F, G) \in clos~\{ funE \}~\{ opPlusIb~b, opBox \}~\{ comp \}
\end{theorem}

\begin{zproof}
prove;
use ClosElemUF[I := \{ funE \}, U := \{ opPlusIb~b, opBox \}, B := \{ comp \}, F := F];
use ClosElemUF[I := \{ funE \}, U := \{ opPlusIb~b, opBox \}, B := \{ comp \}, F := G];
prove;
\end{zproof}

\begin{theorem}{CVPlusIb}
  \forall b: \nat; F: \power (\num \cross \num) | F \in clos~\{ funE \}~\{ opPlusIb~b, opBox \}~\{ comp \} @ opPlusIb~b~F \in clos~\{ funE \}~\{ opPlusIb~b, opBox \}~\{ comp \}
\end{theorem}

\begin{zproof}
prove;
use ClosElemUF[I := \{ funE \}, U := \{ opPlusIb~b, opBox \}, B := \{ comp \}, F := F];
prove;
\end{zproof}

\begin{theorem}{CVBox}
  \forall b: \nat; F: \power (\num \cross \num) | F \in clos~\{ funE \}~\{ opPlusIb~b, opBox \}~\{ comp \} @ opBox~F \in clos~\{ funE \}~\{ opPlusIb~b, opBox \}~\{ comp \}
\end{theorem}

\begin{zproof}
prove;
use ClosElemUF[I := \{ funE \}, U := \{ opPlusIb~b, opBox \}, B := \{ comp \}, F := F];
prove;
\end{zproof}

\begin{theorem}{CVIter}
  \forall b, k: \nat; F, G: \power (\num \cross \num) | F \in clos~\{ funE \}~\{ opPlusIb~b, opBox \}~\{ comp \} \land G \in clos~\{ funE \}~\{ opPlusIb~b, opBox \}~\{ comp \} @ comp(iter[\nat]~k~F, G) \in clos~\{ funE \}~\{ opPlusIb~b, opBox \}~\{ comp \}
\end{theorem}

\begin{zproof}
prove;
use ClosElemUF[I := \{ funE \}, U := \{ opPlusIb~b, opBox \}, B := \{ comp \}, F := F];
use ClosElemUF[I := \{ funE \}, U := \{ opPlusIb~b, opBox \}, B := \{ comp \}, F := G];
use PowIn[I := \{ funE \}, U := \{ opPlusIb~b, opBox \}, B := \{ comp \}, F := F, G := G, k := k];
prove;
\end{zproof}

\begin{theorem}{ZeroExprInCV}
  \forall b: \nat @ opBox~funE \in clos~\{ funE \}~\{ opPlusIb~b, opBox \}~\{ comp \}
\end{theorem}

\begin{zproof}
prove;
\end{zproof}

\begin{theorem}{ZeroInCV}
  \forall b: \nat @ cst~0 \in clos~\{ funE \}~\{ opPlusIb~b, opBox \}~\{ comp \}
\end{theorem}

\begin{zproof}
prove;
use ZeroExprInCV[b := b];
use BoxEExpr;
prove;
equality substitute cst~0;
prove;
\end{zproof}

\begin{theorem}{PbInCV}
  \forall c: \nat @ comp(funE, comp(iter[\nat]~(4 * c + 5)~(opPlusIb~(1 + c)~(cst~0)), opPlusIb~(1 + c)~(opPlusIb~(1 + c)~(opPlusIb~(1 + c)~(opBox~(comp(iter[\nat]~3~(opPlusIb~(1 + c)~(cst~0)), opPlusIb~(1 + c)~(comp(funE, comp(iter[\nat]~(2 * c)~(opPlusIb~(1 + c)~(comp(funE, comp(iter[\nat]~c~(opPlusIb~(1 + c)~(cst~0)), cst~0)))), comp(opPlusIb~(1 + c)~(cst~0), opPlusIb~(1 + c)~(cst~0)))))))))))) \in clos~\{ funE \}~\{ opPlusIb~(1 + c), opBox \}~\{ comp \}
\end{theorem}

\begin{zproof}
prove;
use ZeroInCV[b := 1 + c];
prove;
use CVPlusIb[b := 1 + c, F := cst~0];
prove;
use CVIter[b := 1 + c, k := c, F := opPlusIb~(1 + c)~(cst~0), G := cst~0];
prove;
use CVComp[b := 1 + c, F := funE, G := comp(iter[\nat]~c~(opPlusIb~(1 + c)~(cst~0)), cst~0)];
prove;
use CVPlusIb[b := 1 + c, F := comp(funE, comp(iter[\nat]~c~(opPlusIb~(1 + c)~(cst~0)), cst~0))];
prove;
use CVComp[b := 1 + c, F := opPlusIb~(1 + c)~(cst~0), G := opPlusIb~(1 + c)~(cst~0)];
prove;
use CVIter[b := 1 + c, k := 2 * c, F := opPlusIb~(1 + c)~(comp(funE, comp(iter[\nat]~c~(opPlusIb~(1 + c)~(cst~0)), cst~0))), G := comp(opPlusIb~(1 + c)~(cst~0), opPlusIb~(1 + c)~(cst~0))];
prove;
use CVComp[b := 1 + c, F := funE, G := comp(iter[\nat]~(2 * c)~(opPlusIb~(1 + c)~(comp(funE, comp(iter[\nat]~c~(opPlusIb~(1 + c)~(cst~0)), cst~0)))), comp(opPlusIb~(1 + c)~(cst~0), opPlusIb~(1 + c)~(cst~0)))];
prove;
use CVPlusIb[b := 1 + c, F := comp(funE, comp(iter[\nat]~(2 * c)~(opPlusIb~(1 + c)~(comp(funE, comp(iter[\nat]~c~(opPlusIb~(1 + c)~(cst~0)), cst~0)))), comp(opPlusIb~(1 + c)~(cst~0), opPlusIb~(1 + c)~(cst~0))))];
prove;
use CVIter[b := 1 + c, k := 3, F := opPlusIb~(1 + c)~(cst~0), G := opPlusIb~(1 + c)~(comp(funE, comp(iter[\nat]~(2 * c)~(opPlusIb~(1 + c)~(comp(funE, comp(iter[\nat]~c~(opPlusIb~(1 + c)~(cst~0)), cst~0)))), comp(opPlusIb~(1 + c)~(cst~0), opPlusIb~(1 + c)~(cst~0)))))];
prove;
use CVBox[b := 1 + c, F := comp(iter[\nat]~3~(opPlusIb~(1 + c)~(cst~0)), opPlusIb~(1 + c)~(comp(funE, comp(iter[\nat]~(2 * c)~(opPlusIb~(1 + c)~(comp(funE, comp(iter[\nat]~c~(opPlusIb~(1 + c)~(cst~0)), cst~0)))), comp(opPlusIb~(1 + c)~(cst~0), opPlusIb~(1 + c)~(cst~0))))))];
prove;
use CVPlusIb[b := 1 + c, F := opBox~(comp(iter[\nat]~3~(opPlusIb~(1 + c)~(cst~0)), opPlusIb~(1 + c)~(comp(funE, comp(iter[\nat]~(2 * c)~(opPlusIb~(1 + c)~(comp(funE, comp(iter[\nat]~c~(opPlusIb~(1 + c)~(cst~0)), cst~0)))), comp(opPlusIb~(1 + c)~(cst~0), opPlusIb~(1 + c)~(cst~0)))))))];
prove;
use CVPlusIb[b := 1 + c, F := opPlusIb~(1 + c)~(opBox~(comp(iter[\nat]~3~(opPlusIb~(1 + c)~(cst~0)), opPlusIb~(1 + c)~(comp(funE, comp(iter[\nat]~(2 * c)~(opPlusIb~(1 + c)~(comp(funE, comp(iter[\nat]~c~(opPlusIb~(1 + c)~(cst~0)), cst~0)))), comp(opPlusIb~(1 + c)~(cst~0), opPlusIb~(1 + c)~(cst~0))))))))];
prove;
use CVPlusIb[b := 1 + c, F := opPlusIb~(1 + c)~(opPlusIb~(1 + c)~(opBox~(comp(iter[\nat]~3~(opPlusIb~(1 + c)~(cst~0)), opPlusIb~(1 + c)~(comp(funE, comp(iter[\nat]~(2 * c)~(opPlusIb~(1 + c)~(comp(funE, comp(iter[\nat]~c~(opPlusIb~(1 + c)~(cst~0)), cst~0)))), comp(opPlusIb~(1 + c)~(cst~0), opPlusIb~(1 + c)~(cst~0)))))))))];
prove;
use CVIter[b := 1 + c, k := 4 * c + 5, F := opPlusIb~(1 + c)~(cst~0), G := opPlusIb~(1 + c)~(opPlusIb~(1 + c)~(opPlusIb~(1 + c)~(opBox~(comp(iter[\nat]~3~(opPlusIb~(1 + c)~(cst~0)), opPlusIb~(1 + c)~(comp(funE, comp(iter[\nat]~(2 * c)~(opPlusIb~(1 + c)~(comp(funE, comp(iter[\nat]~c~(opPlusIb~(1 + c)~(cst~0)), cst~0)))), comp(opPlusIb~(1 + c)~(cst~0), opPlusIb~(1 + c)~(cst~0))))))))))];
prove;
\end{zproof}

\begin{theorem}{OneInCV}
  \forall c: \nat @ cst~1 \in clos~\{ funE \}~\{ opPlusIb~(1 + c), opBox \}~\{ comp \}
\end{theorem}

\begin{zproof}
prove;
use ZeroInCV[b := 1 + c];
use PbInCV[c := c];
prove;
use CVPlusIb[b := 1 + c, F := cst~0];
prove;
use CVComp[b := 1 + c, F := opPlusIb~(1 + c)~(cst~0), G := cst~0];
prove;
use CVIter[b := 1 + c, k := c, F := comp(funE, comp(iter[\nat]~(4 * c + 5)~(opPlusIb~(1 + c)~(cst~0)), opPlusIb~(1 + c)~(opPlusIb~(1 + c)~(opPlusIb~(1 + c)~(opBox~(comp(iter[\nat]~3~(opPlusIb~(1 + c)~(cst~0)), opPlusIb~(1 + c)~(comp(funE, comp(iter[\nat]~(2 * c)~(opPlusIb~(1 + c)~(comp(funE, comp(iter[\nat]~c~(opPlusIb~(1 + c)~(cst~0)), cst~0)))), comp(opPlusIb~(1 + c)~(cst~0), opPlusIb~(1 + c)~(cst~0)))))))))))), G := comp(opPlusIb~(1 + c)~(cst~0), cst~0)];
prove;
use OneEqVb[c := c];
prove;
\end{zproof}

\begin{theorem}{PlusIExprInCV}
  \forall c: \nat; F: \nat \fun \nat | F \in clos~\{ funE \}~\{ opPlusIb~(1 + c), opBox \}~\{ comp \} @ comp(iter[\nat]~(1 + c)~(comp(funE, comp(iter[\nat]~(4 * c + 5)~(opPlusIb~(1 + c)~(cst~0)), opPlusIb~(1 + c)~(opPlusIb~(1 + c)~(opPlusIb~(1 + c)~(opBox~(comp(iter[\nat]~3~(opPlusIb~(1 + c)~(cst~0)), opPlusIb~(1 + c)~(comp(funE, comp(iter[\nat]~(2 * c)~(opPlusIb~(1 + c)~(comp(funE, comp(iter[\nat]~c~(opPlusIb~(1 + c)~(cst~0)), cst~0)))), comp(opPlusIb~(1 + c)~(cst~0), opPlusIb~(1 + c)~(cst~0))))))))))))), opPlusIb~(1 + c)~F) \in clos~\{ funE \}~\{ opPlusIb~(1 + c), opBox \}~\{ comp \}
\end{theorem}

\begin{zproof}
prove;
use PbInCV[c := c];
use CVPlusIb[b := 1 + c, F := F];
prove;
use CVIter[b := 1 + c, k := 1 + c, F := comp(funE, comp(iter[\nat]~(4 * c + 5)~(opPlusIb~(1 + c)~(cst~0)), opPlusIb~(1 + c)~(opPlusIb~(1 + c)~(opPlusIb~(1 + c)~(opBox~(comp(iter[\nat]~3~(opPlusIb~(1 + c)~(cst~0)), opPlusIb~(1 + c)~(comp(funE, comp(iter[\nat]~(2 * c)~(opPlusIb~(1 + c)~(comp(funE, comp(iter[\nat]~c~(opPlusIb~(1 + c)~(cst~0)), cst~0)))), comp(opPlusIb~(1 + c)~(cst~0), opPlusIb~(1 + c)~(cst~0)))))))))))), G := opPlusIb~(1 + c)~F];
prove;
\end{zproof}

\begin{theorem}{PlusIInCV}
  \forall c: \nat; F: \nat \fun \nat | F \in clos~\{ funE \}~\{ opPlusIb~(1 + c), opBox \}~\{ comp \} @ opPlusI~F \in clos~\{ funE \}~\{ opPlusIb~(1 + c), opBox \}~\{ comp \}
\end{theorem}

\begin{zproof}
prove;
use PlusIExprInCV[c := c, F := F];
use PlusIEqVb[c := c, F := F];
prove;
\end{zproof}

\subsection{$\langle E, F^{+b}, FG, F^\square \rangle$ for $b \geq 1$}

\begin{theorem}{CVUClosed}
  \forall c: \nat @ clos~\{ funE \}~\{ opPlusIb~(1 + c), opBox \}~\{ comp \} \in UClosed~\{ opPlusI, opBox \}
\end{theorem}

\begin{zproof}
prove;
use ClosInPow[I := \{ funE \}, U := \{ opPlusIb~(1 + c), opBox \}, B := \{ comp \}];
with enabled (UClosedDef) prove;
use ClosElemUF[I := \{ funE \}, U := \{ opPlusIb~(1 + c), opBox \}, B := \{ comp \}, F := F];
split u = opPlusI;
cases;
use PlusIInCV[c := c, F := F];
prove;
next;
prove;
next;
\end{zproof}

\begin{theorem}{CVBClosed}
  \forall b: \nat @ clos~\{ funE \}~\{ opPlusIb~b, opBox \}~\{ comp \} \in BClosed~\{ comp \}
\end{theorem}

\begin{zproof}
prove;
use ClosInPow[I := \{ funE \}, U := \{ opPlusIb~b, opBox \}, B := \{ comp \}];
with enabled (BClosedDef) prove;
\end{zproof}

\begin{theorem}{C6SubCV}
  \forall c: \nat @ C6 \subseteq clos~\{ funE \}~\{ opPlusIb~(1 + c), opBox \}~\{ comp \}
\end{theorem}

\begin{zproof}
prove;
use C6Def;
use ClosLeast[I := \{ cst~1, funE \}, U := \{ opPlusI, opBox \}, B := \{ comp \}, Z := clos~\{ funE \}~\{ opPlusIb~(1 + c), opBox \}~\{ comp \}];
use ClosInPow[I := \{ funE \}, U := \{ opPlusIb~(1 + c), opBox \}, B := \{ comp \}];
use CVUClosed[c := c];
use CVBClosed[b := 1 + c];
use OneInCV[c := c];
prove;
\end{zproof}

$F^{+b} = (S^b F)^+$, so $C_6$ is closed under $F^{+b}$.

\begin{theorem}{PlusIbExpr}
  \forall b: \nat; F: \nat \fun \nat @ opPlusIb~b~F = opPlusI~(comp(iter[\nat]~b~funS, F))
\end{theorem}

\begin{zproof}
prove;
apply UFExt to predicate opPlusIb~b~F = opPlusI~(comp(iter[\nat]~b~funS, F));
prove;
\end{zproof}

\begin{theorem}{C6PlusIbExpr}
  \forall b: \nat; F: \nat \fun \nat | F \in C6 @ opPlusI~(comp(iter[\nat]~b~funS, F)) \in C6
\end{theorem}

\begin{zproof}
prove;
use C6Def;
use PowIn[I := \{ cst~1, funE \}, U := \{ opPlusI, opBox \}, B := \{ comp \}, F := funS, G := F, k := b];
prove;
\end{zproof}

\begin{theorem}{C6PlusIb}
  \forall b: \nat; F: \nat \fun \nat | F \in C6 @ opPlusIb~b~F \in C6
\end{theorem}

\begin{zproof}
prove;
use C6PlusIbExpr[b := b, F := F];
use PlusIbExpr[b := b, F := F];
prove;
\end{zproof}

\begin{theorem}{C6UClosedB}
  \forall b: \nat @ C6 \in UClosed~\{ opPlusIb~b, opBox \}
\end{theorem}

\begin{zproof}
prove;
use C6InPower;
with enabled (UClosedDef, C6UF) prove;
use InPair[\power (\power (\num \cross \num) \cross \power (\num \cross \num))][u := u, p := opPlusIb~b, q := opBox];
split u = opBox;
cases;
prove;
next;
use C6PlusIb[b := b, F := F];
prove;
next;
\end{zproof}

\begin{theorem}{CVSubC6}
  \forall b: \nat @ clos~\{ funE \}~\{ opPlusIb~b, opBox \}~\{ comp \} \subseteq C6
\end{theorem}

\begin{zproof}
prove;
use ClosLeast[I := \{ funE \}, U := \{ opPlusIb~b, opBox \}, B := \{ comp \}, Z := C6];
use C6InPower;
use C6UClosedB[b := b];
use C6BClosedComp;
prove;
\end{zproof}

\begin{theorem}{ExtVbSucc}
  \forall c: \nat @ clos~\{ funE \}~\{ opPlusIb~(1 + c), opBox \}~\{ comp \} = C6
\end{theorem}

\begin{zproof}
prove;
use C6SubCV[c := c];
use CVSubC6[b := 1 + c];
use ClosInPow[I := \{ funE \}, U := \{ opPlusIb~(1 + c), opBox \}, B := \{ comp \}];
use C6InPower;
use SubAntisym[\nat \fun \nat][A := clos~\{ funE \}~\{ opPlusIb~(1 + c), opBox \}~\{ comp \}, B := C6];
prove;
\end{zproof}

\begin{theorem}{ExtVb}
  \forall b: \nat | 1 \leq b @ clos~\{ funE \}~\{ opPlusIb~b, opBox \}~\{ comp \} = C6
\end{theorem}

\begin{zproof}
prove;
use ExtVbSucc[c := b - 1];
prove;
\end{zproof}

\subsection{$b = 0$: every generated function fixes $0$, so $\bar 1$ is not generated}

\begin{zed}
  FixZero == \{ F: \nat \fun \nat | F~0 = 0 \}
\end{zed}

\begin{zproof}[FixZero\$domainCheck]
prove;
\end{zproof}

\begin{theorem}{FixZeroInPower}
  FixZero \in \power (\nat \fun \nat)
\end{theorem}

\begin{zproof}
invoke FixZero;
with enabled (inPower) prove;
\end{zproof}

\begin{theorem}{FixZeroUClosed}
  FixZero \in UClosed~\{ opPlusIb~0, opBox \}
\end{theorem}

\begin{zproof}
use FixZeroInPower;
with enabled (UClosedDef) prove;
invoke FixZero;
prove;
use InPair[\power (\power (\num \cross \num) \cross \power (\num \cross \num))][u := u, p := opPlusIb~0, q := opBox];
split u = opBox;
cases;
prove;
next;
prove;
next;
\end{zproof}

\begin{theorem}{FixZeroBClosed}
  FixZero \in BClosed~\{ comp \}
\end{theorem}

\begin{zproof}
use FixZeroInPower;
with enabled (BClosedDef) prove;
invoke FixZero;
prove;
\end{zproof}

\begin{theorem}{CVZeroFix}
  clos~\{ funE \}~\{ opPlusIb~0, opBox \}~\{ comp \} \subseteq FixZero
\end{theorem}

\begin{zproof}
use ClosLeast[I := \{ funE \}, U := \{ opPlusIb~0, opBox \}, B := \{ comp \}, Z := FixZero];
use FixZeroInPower;
use FixZeroUClosed;
use FixZeroBClosed;
invoke FixZero;
prove;
\end{zproof}

\begin{theorem}{OneNotCVZero}
  \lnot cst~1 \in clos~\{ funE \}~\{ opPlusIb~0, opBox \}~\{ comp \}
\end{theorem}

\begin{zproof}
use CVZeroFix;
apply subsetDef to predicate clos~\{ funE \}~\{ opPlusIb~0, opBox \}~\{ comp \} \subseteq FixZero;
apply inPower to predicate clos~\{ funE \}~\{ opPlusIb~0, opBox \}~\{ comp \} \in \power FixZero;
instantiate e == cst~1;
invoke FixZero;
prove;
\end{zproof}

\begin{theorem}{OneInEqC6}
  \forall Y: \power (\nat \fun \nat) | Y = C6 @ cst~1 \in Y
\end{theorem}

\begin{zproof}
prove;
\end{zproof}

\begin{theorem}{ExtVbZero}
  \lnot clos~\{ funE \}~\{ opPlusIb~0, opBox \}~\{ comp \} = C6
\end{theorem}

\begin{zproof}
use OneNotCVZero;
use ClosInPow[I := \{ funE \}, U := \{ opPlusIb~0, opBox \}, B := \{ comp \}];
use OneInEqC6[Y := clos~\{ funE \}~\{ opPlusIb~0, opBox \}~\{ comp \}];
prove;
\end{zproof}

\begin{theorem}{ExtVbIff}
  \forall b: \nat @ clos~\{ funE \}~\{ opPlusIb~b, opBox \}~\{ comp \} = C6 \iff 1 \leq b
\end{theorem}

\begin{zproof}
prove;
split b = 0;
cases;
use ExtVbZero;
prove;
next;
use ExtVb[b := b];
prove;
next;
\end{zproof}

\section{General recursive functions (README \S 2.3)}

J.~Robinson (Proc.\ AMS 1 (1950), Sec.~4, Thm.~3) proved that the unary general recursive
functions are the closure of $\langle S, E, F + G, FG, F^{-1}\rangle$, where
$F^{-1}(x) = \mu y\,[F(y) = x]$ for $F$ onto.  Here $F + G$ is replaced by $F^+$.

\subsection{Least preimages and inversion}

Every value that is taken is taken at a least argument.

\begin{zed}
  LeastPreSet == \{ k: \nat | \forall F: \nat \fun \nat; x: \nat @ (\forall j: \nat | j \leq k @ \lnot F~j = x) \lor (\exists m: \nat @ F~m = x \land (\forall n: \nat | F~n = x @ m \leq n)) \}
\end{zed}

\begin{zproof}[LeastPreSet\$domainCheck]
prove;
\end{zproof}

\begin{theorem}{rule LeastPreBase}
  0 \in LeastPreSet
\end{theorem}

\begin{zproof}
invoke LeastPreSet;
prove;
instantiate m == 0;
prove;
\end{zproof}

\begin{theorem}{rule LeastPreStep}
  \forall k: LeastPreSet @ 1 + k \in LeastPreSet
\end{theorem}

\begin{zproof}
invoke LeastPreSet;
prove;
split j \leq k;
cases;
instantiate F\_\_0 == F, j\_\_0 == j;
prove;
instantiate m\_\_1 == m;
prove;
instantiate n\_\_1 == n;
prove;
next;
split \exists n: \nat | n \leq k @ F~n = F~j;
cases;
prove;
instantiate F\_\_0 == F, j\_\_0 == n;
prove;
instantiate m\_\_1 == m;
prove;
instantiate n\_\_2 == n\_\_0;
prove;
next;
prove;
instantiate m\_\_0 == j;
prove;
instantiate n\_\_0 == n;
prove;
next;
\end{zproof}

\begin{theorem}{LeastPreSub}
  \nat \subseteq LeastPreSet
\end{theorem}

\begin{zproof}
apply natInduction to predicate \nat \subseteq LeastPreSet;
prove;
invoke LeastPreSet;
prove;
\end{zproof}

\begin{theorem}{LeastPre}
  \forall F: \nat \fun \nat; y: \nat @ \exists m: \nat @ F~m = F~y \land (\forall n: \nat | F~n = F~y @ m \leq n)
\end{theorem}

\begin{zproof}
prove;
use LeastPreSub;
invoke LeastPreSet;
with enabled (inPower) prove;
instantiate e == y;
prove;
instantiate F\_\_1 == F, j\_\_0 == y;
prove;
instantiate m\_\_2 == m;
prove;
instantiate n\_\_2 == n;
prove;
\end{zproof}


$y$ is the least preimage of $x$ under $F$ (or $0$ if $x$ is not a value).

\begin{zed}
  IsLeast == \{ F: \nat \fun \nat; x, y: \nat | (\forall v: \nat | F~v = x @ y \leq v) \land (F~y = x \lor ((\forall w: \nat @ \lnot F~w = x) \land y = 0)) \}
\end{zed}

\begin{zproof}[IsLeast\$domainCheck]
prove;
\end{zproof}

\begin{theorem}{LinvUnique}
  \forall F: \nat \fun \nat; x, y, z: \nat | (F, x, y) \in IsLeast \land (F, x, z) \in IsLeast @ y = z
\end{theorem}

\begin{zproof}
invoke IsLeast;
prove;
instantiate v == z, v\_\_0 == y, w == z, w\_\_0 == y;
prove;
split F~y = x;
cases;
with normalization prove;
next;
prove;
next;
\end{zproof}

\begin{theorem}{LinvExists}
  \forall F: \nat \fun \nat; x: \nat @ \exists y: \nat @ (F, x, y) \in IsLeast
\end{theorem}

\begin{zproof}
invoke IsLeast;
prove;
split \exists w: \nat @ F~w = x;
cases;
prove;
use LeastPre[F := F, y := w];
prove;
instantiate y == m;
prove;
instantiate n == v;
prove;
next;
prove;
next;
\end{zproof}

\begin{theorem}{LinvInner}
  \forall F: \nat \fun \nat @ \{ x, y: \nat | (F, x, y) \in IsLeast @ (x, y) \} \in \nat \fun \nat
\end{theorem}

\begin{zproof}
apply FunIntro to predicate \{ x, y: \nat | (F, x, y) \in IsLeast @ (x, y) \} \in \nat \fun \nat;
with enabled (RelPow) prove;
use LinvExists[F := F, x := x];
use LinvUnique[F := F, x := x\_\_0, y := y, z := z];
prove;
instantiate y\_\_1 == y\_\_0;
prove;
\end{zproof}

\begin{zed}
  LinvWit == (\lambda F: \nat \fun \nat @ \{ x, y: \nat | (F, x, y) \in IsLeast @ (x, y) \})
\end{zed}

\begin{theorem}{LinvWitFun}
  LinvWit \in (\nat \fun \nat) \fun (\nat \fun \nat)
\end{theorem}

\begin{zproof}
invoke LinvWit;
apply FunIntro to predicate (\lambda F: \nat \fun \nat @ \{ x, y: \nat | (F, x, y) \in IsLeast @ (x, y) \}) \in (\nat \fun \nat) \fun (\nat \fun \nat);
with enabled (RelPow, inPower) prove;
use LinvInner[F := F];
use LinvInner[F := x];
prove;
\end{zproof}

\begin{theorem}{LinvWitApp}
  \forall F: \nat \fun \nat @ LinvWit~F = \{ x, y: \nat | (F, x, y) \in IsLeast @ (x, y) \}
\end{theorem}

\begin{zproof}
use LinvWitFun;
use pairInFunction[\nat \fun \nat, \nat \fun \nat][f := LinvWit, x := F, y := \{ x, y: \nat | (F, x, y) \in IsLeast @ (x, y) \}];
invoke LinvWit;
prove;
\end{zproof}

\begin{theorem}{LinvWitVal}
  \forall F: \nat \fun \nat; x: \nat @ (F, x, LinvWit~F~x) \in IsLeast
\end{theorem}

\begin{zproof}
prove;
use LinvWitApp[F := F];
use LinvInner[F := F];
use LinvExists[F := F, x := x];
prove;
use pairInFunction[\nat, \nat][f := \{ x, y: \nat | (F, x, y) \in IsLeast @ (x, y) \}, x := x, y := y];
prove;
\end{zproof}

The least preimage of $x$ (an arbitrary value if $x$ is not a value of $F$).

\begin{axdef}
  linv: (\nat \fun \nat) \fun (\nat \fun \nat)
\where
  \Label{LinvSpec} \forall F: \nat \fun \nat; x: \nat @ (F, x, linv~F~x) \in IsLeast
\end{axdef}

\begin{zproof}[linv\$domainCheck]
prove;
\end{zproof}

\begin{theorem}{LinvEx}
  \exists linv: (\nat \fun \nat) \fun (\nat \fun \nat) @ \forall F: \nat \fun \nat; x: \nat @ (F, x, linv~F~x) \in IsLeast
\end{theorem}

\begin{zproof}
instantiate linv == LinvWit;
use LinvWitFun;
prove;
use LinvWitVal[F := F, x := x];
prove;
\end{zproof}

\begin{theorem}{LinvProp}
  \forall F: \nat \fun \nat; x, y: \nat | F~y = x @ F~(linv~F~x) = x \land linv~F~x \leq y
\end{theorem}

\begin{zproof}
prove;
use LinvSpec[F := F, x := x];
invoke IsLeast;
prove;
instantiate v == y, w == y;
prove;
\end{zproof}

$F$ is onto.  Inversion: $F^{-1}(x) = \mu y\,[F(y) = x]$ if $F$ is onto, and $F^{-1} = F$
otherwise, so that a closure under \textit{opInv} is a closure under Robinson's inversion.

\begin{zed}
  Surj == \{ F: \nat \fun \nat | \forall x: \nat @ \exists y: \nat @ F~y = x \}
\end{zed}

\begin{zproof}[Surj\$domainCheck]
prove;
\end{zproof}

\begin{axdef}
  opInv: (\nat \fun \nat) \fun (\nat \fun \nat)
\where
  \Label{disabled rule InvVal} \forall F: \nat \fun \nat; x: \nat @ opInv~F~x = \IF F \in Surj \THEN linv~F~x \ELSE F~x
\end{axdef}

\begin{zproof}[opInv\$domainCheck]
prove;
\end{zproof}

\begin{theorem}{rule InvInner}
  \forall F: \nat \fun \nat @ (\lambda x: \nat @ \IF F \in Surj \THEN linv~F~x \ELSE F~x) \in \nat \fun \nat
\end{theorem}

\begin{zproof}
apply FunIntro to predicate (\lambda x: \nat @ \IF F \in Surj \THEN linv~F~x \ELSE F~x) \in \nat \fun \nat;
with enabled (RelPow, inPower) prove;
split F \in Surj;
cases;
prove;
next;
prove;
next;
\end{zproof}

\begin{zed}
  InvWit == (\lambda F: \nat \fun \nat @ (\lambda x: \nat @ \IF F \in Surj \THEN linv~F~x \ELSE F~x))
\end{zed}

\begin{zproof}[InvWit\$domainCheck]
prove;
\end{zproof}

\begin{theorem}{InvWitFun}
  InvWit \in (\nat \fun \nat) \fun (\nat \fun \nat)
\end{theorem}

\begin{zproof}
invoke InvWit;
apply FunIntro to predicate (\lambda F: \nat \fun \nat @ (\lambda x: \nat @ \IF F \in Surj \THEN linv~F~x \ELSE F~x)) \in (\nat \fun \nat) \fun (\nat \fun \nat);
with enabled (RelPow, inPower) prove;
\end{zproof}

\begin{theorem}{InvWitApp}
  \forall F: \nat \fun \nat @ InvWit~F = (\lambda x: \nat @ \IF F \in Surj \THEN linv~F~x \ELSE F~x)
\end{theorem}

\begin{zproof}
use InvWitFun;
use pairInFunction[\nat \fun \nat, \nat \fun \nat][f := InvWit, x := F, y := (\lambda x: \nat @ \IF F \in Surj \THEN linv~F~x \ELSE F~x)];
invoke InvWit;
prove by reduce;
\end{zproof}

\begin{theorem}{rule InvWitVal}
  \forall F: \nat \fun \nat; x: \nat @ InvWit~F~x = \IF F \in Surj \THEN linv~F~x \ELSE F~x
\end{theorem}

\begin{zproof}
prove;
use InvWitApp[F := F];
use pairInFunction[\nat, \nat][f := (\lambda x: \nat @ \IF F \in Surj \THEN linv~F~x \ELSE F~x), x := x, y := \IF F \in Surj \THEN linv~F~x \ELSE F~x];
prove;
\end{zproof}

\begin{theorem}{InvEx}
  \exists opInv: (\nat \fun \nat) \fun (\nat \fun \nat) @ \forall F: \nat \fun \nat; x: \nat @ opInv~F~x = \IF F \in Surj \THEN linv~F~x \ELSE F~x
\end{theorem}

\begin{zproof}
instantiate opInv == InvWit;
use InvWitFun;
prove;
\end{zproof}

\begin{theorem}{InvNot}
  \forall F: \nat \fun \nat | \lnot F \in Surj @ opInv~F = F
\end{theorem}

\begin{zproof}
apply UFExt to predicate opInv~F = F;
prove;
use InvVal[F := F, x := x];
prove;
\end{zproof}

\begin{theorem}{rule InvRight}
  \forall F: \nat \fun \nat; x: \nat | F \in Surj @ F~(opInv~F~x) = x
\end{theorem}

\begin{zproof}
use InvVal[F := F, x := x];
invoke Surj;
prove;
instantiate x\_\_0 == x;
prove;
use LinvProp[F := F, x := x, y := y];
prove;
\end{zproof}

\begin{theorem}{InvLeast}
  \forall F: \nat \fun \nat; x, y: \nat | F \in Surj \land F~y = x @ opInv~F~x \leq y
\end{theorem}

\begin{zproof}
use InvVal[F := F, x := x];
use LinvProp[F := F, x := x, y := y];
prove;
\end{zproof}

\begin{theorem}{InvFix0}
  \forall F: \nat \fun \nat | F~0 = 0 @ opInv~F~0 = 0
\end{theorem}

\begin{zproof}
use InvVal[F := F, x := 0];
use LinvProp[F := F, x := 0, y := 0];
use UFNonneg[f := linv~F, x := 0];
prove;
\end{zproof}

\subsection{Lemmas 6.2--6.6 in a rich set}

The expressions of Lemmas 6.2--6.6 only use $I, E, S, D, M_3, \mathrm{Sq}, \mathrm{Hf}, \mathrm{Rt}$,
$F^+$ and composition, so they work in every set of functions that contains these and is closed
under $F^+$ and composition.

\begin{zed}
  Rich == \{ Y: \power (\nat \fun \nat) | funI \in Y \land funE \in Y \land funS \in Y \land funD \in Y \land funM3 \in Y \land funSq \in Y \land funHf \in Y \land funRt \in Y \land (\forall F: Y @ opPlusI~F \in Y) \land (\forall F, G: Y @ comp(F, G) \in Y) \}
\end{zed}

\begin{zproof}[Rich\$domainCheck]
prove;
\end{zproof}

\begin{theorem}{rule RichI}
  \forall Y: Rich @ funI \in Y
\end{theorem}

\begin{zproof}
invoke Rich;
prove;
\end{zproof}

\begin{theorem}{rule RichE}
  \forall Y: Rich @ funE \in Y
\end{theorem}

\begin{zproof}
invoke Rich;
prove;
\end{zproof}

\begin{theorem}{rule RichS}
  \forall Y: Rich @ funS \in Y
\end{theorem}

\begin{zproof}
invoke Rich;
prove;
\end{zproof}

\begin{theorem}{rule RichD}
  \forall Y: Rich @ funD \in Y
\end{theorem}

\begin{zproof}
invoke Rich;
prove;
\end{zproof}

\begin{theorem}{rule RichM3}
  \forall Y: Rich @ funM3 \in Y
\end{theorem}

\begin{zproof}
invoke Rich;
prove;
\end{zproof}

\begin{theorem}{rule RichSq}
  \forall Y: Rich @ funSq \in Y
\end{theorem}

\begin{zproof}
invoke Rich;
prove;
\end{zproof}

\begin{theorem}{rule RichHf}
  \forall Y: Rich @ funHf \in Y
\end{theorem}

\begin{zproof}
invoke Rich;
prove;
\end{zproof}

\begin{theorem}{rule RichRt}
  \forall Y: Rich @ funRt \in Y
\end{theorem}

\begin{zproof}
invoke Rich;
prove;
\end{zproof}

\begin{theorem}{rule RichPlusI}
  \forall Y: Rich; F: \nat \fun \nat | F \in Y @ opPlusI~F \in Y
\end{theorem}

\begin{zproof}
invoke Rich;
prove;
instantiate F\_\_0 == F;
prove;
\end{zproof}

\begin{theorem}{rule RichComp}
  \forall Y: Rich; F, G: \nat \fun \nat | F \in Y \land G \in Y @ comp(F, G) \in Y
\end{theorem}

\begin{zproof}
invoke Rich;
prove;
instantiate F\_\_1 == F, G\_\_0 == G;
prove;
\end{zproof}

\begin{theorem}{rule RMinusIIn}
  \forall Y: Rich; F: \nat \fun \nat | F \in Y @ exprMinusI~F \in Y
\end{theorem}

\begin{zproof}
with enabled (exprMinusIApp) prove;
\end{zproof}

\begin{theorem}{rule ROplusIn}
  \forall Y: Rich; F, G: \nat \fun \nat | F \in Y \land G \in Y @ exprOplus(F, G) \in Y
\end{theorem}

\begin{zproof}
with enabled (exprOplusApp) prove;
\end{zproof}

\begin{theorem}{rule ROminusIn}
  \forall Y: Rich; F, G: \nat \fun \nat | F \in Y \land G \in Y @ exprOminus(F, G) \in Y
\end{theorem}

\begin{zproof}
with enabled (exprOminusApp) prove;
\end{zproof}

\begin{theorem}{rule RPhatIn}
  \forall Y: Rich @ exprPhat \in Y
\end{theorem}

\begin{zproof}
invoke exprPhat;
prove;
\end{zproof}

\begin{theorem}{rule ROtimesIn}
  \forall Y: Rich; F, G: \nat \fun \nat | F \in Y \land G \in Y @ exprOtimes(F, G) \in Y
\end{theorem}

\begin{zproof}
with enabled (exprOtimesApp) prove;
\end{zproof}

\begin{theorem}{rule RPlusExprIn}
  \forall Y: Rich; F, G: \nat \fun \nat | F \in Y \land G \in Y @ exprPlus(F, G) \in Y
\end{theorem}

\begin{zproof}
with enabled (exprPlusApp) prove;
\end{zproof}

\begin{theorem}{rule RMinusExprIn}
  \forall Y: Rich; F, G: \nat \fun \nat | F \in Y \land G \in Y @ exprMinus(F, G) \in Y
\end{theorem}

\begin{zproof}
with enabled (exprMinusApp) prove;
\end{zproof}

Lemma 6.6 in a rich set.

\begin{theorem}{rule RichPlus}
  \forall Y: Rich; F, G: \nat \fun \nat | F \in Y \land G \in Y @ opPlus(F, G) \in Y
\end{theorem}

\begin{zproof}
use PlusExprEq[F := F, G := G];
prove;
equality substitute opPlus(F, G);
prove;
\end{zproof}

\begin{theorem}{rule RichDiffE}
  \forall Y: Rich; F, G: \nat \fun \nat | F \in Y \land G \in Y @ opDiffE(F, G) \in Y
\end{theorem}

\begin{zproof}
use MinusExprEq[F := F, G := G];
prove;
equality substitute opDiffE(F, G);
prove;
\end{zproof}

$A$, $V$, $K$ and $L$ are in every rich set.

\begin{theorem}{RichAExpr}
  \forall Y: Rich @ comp(funHf, opPlusI~funSq) \in Y
\end{theorem}

\begin{zproof}
prove;
\end{zproof}

\begin{theorem}{rule RichA}
  \forall Y: Rich @ funA \in Y
\end{theorem}

\begin{zproof}
prove;
use RichAExpr[Y := Y];
use AExpr;
prove;
\end{zproof}

\begin{theorem}{RichVExpr}
  \forall Y: Rich @ comp(comp(comp(comp(comp(comp(funHf, exprPhat), funRt), funS), funD), funD), funD) \in Y
\end{theorem}

\begin{zproof}
prove;
\end{zproof}

\begin{theorem}{rule RichV}
  \forall Y: Rich @ funV \in Y
\end{theorem}

\begin{zproof}
prove;
use RichVExpr[Y := Y];
use VExpr;
prove;
\end{zproof}

\begin{theorem}{RichKExpr}
  \forall Y: Rich @ opDiffE(funI, comp(funA, funV)) \in Y
\end{theorem}

\begin{zproof}
prove;
\end{zproof}

\begin{theorem}{rule RichK}
  \forall Y: Rich @ funK \in Y
\end{theorem}

\begin{zproof}
prove;
use RichKExpr[Y := Y];
use KExpr;
prove;
\end{zproof}

\begin{theorem}{RichLExpr}
  \forall Y: Rich @ opDiffE(comp(funA, comp(funS, funV)), funS) \in Y
\end{theorem}

\begin{zproof}
prove;
\end{zproof}

\begin{theorem}{rule RichL}
  \forall Y: Rich @ funL \in Y
\end{theorem}

\begin{zproof}
prove;
use RichLExpr[Y := Y];
use LExpr;
prove;
\end{zproof}

\subsection{The bases with inversion}

J.~Robinson's basis $\langle S, E, F + G, FG, F^{-1}\rangle$ and $\langle \bar 1, X, F^+, FG, F^{-1}\rangle$
for $X \in \{E, K, L\}$.

\begin{zed}
  ClosJR == clos~\{ funS, funE \}~\{ opInv \}~\{ opPlus, comp \}
\end{zed}

\begin{zproof}[ClosJR\$domainCheck]
prove;
\end{zproof}

\begin[disabled]{theorem}{rule ClosJRDef}
  ClosJR = clos~\{ funS, funE \}~\{ opInv \}~\{ opPlus, comp \}
\end{theorem}

\begin{zproof}
invoke ClosJR;
prove;
\end{zproof}

\begin{theorem}{ClosJRInPower}
  ClosJR \in \power (\nat \fun \nat)
\end{theorem}

\begin{zproof}
use ClosInPow[I := \{ funS, funE \}, U := \{ opInv \}, B := \{ opPlus, comp \}];
with enabled (ClosJRDef) prove;
\end{zproof}

\begin[disabled]{theorem}{rule ClosJRUF}
  \forall F: \power (\num \cross \num) | F \in ClosJR @ F \in \nat \fun \nat
\end{theorem}

\begin{zproof}
prove;
use ClosElemUF[I := \{ funS, funE \}, U := \{ opInv \}, B := \{ opPlus, comp \}, F := F];
with enabled (ClosJRDef) prove;
\end{zproof}

\begin{theorem}{rule ClosJRInit0}
  funS \in ClosJR
\end{theorem}

\begin{zproof}
with enabled (ClosJRDef) prove;
\end{zproof}

\begin{theorem}{rule ClosJRInit1}
  funE \in ClosJR
\end{theorem}

\begin{zproof}
with enabled (ClosJRDef) prove;
\end{zproof}

\begin{theorem}{rule ClosJRUnop0}
  \forall F: \nat \fun \nat | F \in ClosJR @ opInv~F \in ClosJR
\end{theorem}

\begin{zproof}
with enabled (ClosJRDef) prove;
\end{zproof}

\begin{theorem}{rule ClosJRBinop0}
  \forall F, G: \nat \fun \nat | F \in ClosJR \land G \in ClosJR @ opPlus(F, G) \in ClosJR
\end{theorem}

\begin{zproof}
with enabled (ClosJRDef) prove;
\end{zproof}

\begin{theorem}{rule ClosJRBinop1}
  \forall F, G: \nat \fun \nat | F \in ClosJR \land G \in ClosJR @ comp(F, G) \in ClosJR
\end{theorem}

\begin{zproof}
with enabled (ClosJRDef) prove;
\end{zproof}

\begin{zed}
  CGE == clos~\{ cst~1, funE \}~\{ opPlusI, opInv \}~\{ comp \}
\end{zed}

\begin{zproof}[CGE\$domainCheck]
prove;
\end{zproof}

\begin[disabled]{theorem}{rule CGEDef}
  CGE = clos~\{ cst~1, funE \}~\{ opPlusI, opInv \}~\{ comp \}
\end{theorem}

\begin{zproof}
invoke CGE;
prove;
\end{zproof}

\begin{theorem}{CGEInPower}
  CGE \in \power (\nat \fun \nat)
\end{theorem}

\begin{zproof}
use ClosInPow[I := \{ cst~1, funE \}, U := \{ opPlusI, opInv \}, B := \{ comp \}];
with enabled (CGEDef) prove;
\end{zproof}

\begin[disabled]{theorem}{rule CGEUF}
  \forall F: \power (\num \cross \num) | F \in CGE @ F \in \nat \fun \nat
\end{theorem}

\begin{zproof}
prove;
use ClosElemUF[I := \{ cst~1, funE \}, U := \{ opPlusI, opInv \}, B := \{ comp \}, F := F];
with enabled (CGEDef) prove;
\end{zproof}

\begin{theorem}{rule CGEInit0}
  cst~1 \in CGE
\end{theorem}

\begin{zproof}
with enabled (CGEDef) prove;
\end{zproof}

\begin{theorem}{rule CGEInit1}
  funE \in CGE
\end{theorem}

\begin{zproof}
with enabled (CGEDef) prove;
\end{zproof}

\begin{theorem}{rule CGEUnop0}
  \forall F: \nat \fun \nat | F \in CGE @ opPlusI~F \in CGE
\end{theorem}

\begin{zproof}
with enabled (CGEDef) prove;
\end{zproof}

\begin{theorem}{rule CGEUnop1}
  \forall F: \nat \fun \nat | F \in CGE @ opInv~F \in CGE
\end{theorem}

\begin{zproof}
with enabled (CGEDef) prove;
\end{zproof}

\begin{theorem}{rule CGEBinop0}
  \forall F, G: \nat \fun \nat | F \in CGE \land G \in CGE @ comp(F, G) \in CGE
\end{theorem}

\begin{zproof}
with enabled (CGEDef) prove;
\end{zproof}


\subsection{The case $X = E$: arithmetic}

\begin{theorem}{EBound}
  \forall y: \nat @ funE~y \leq 2 * funRt~y
\end{theorem}

\begin{zproof}
prove;
use RtSpec[x := y];
with enabled (EVal) prove;
\end{zproof}

\begin{theorem}{EDecomp}
  \forall y: \nat @ funE~y + funRt~y * funRt~y = y
\end{theorem}

\begin{zproof}
prove;
use RtSpec[x := y];
with enabled (EVal) prove;
\end{zproof}

\begin{theorem}{TwoLeSq}
  \forall a: \nat @ 2 * a \leq a * a + 1
\end{theorem}

\begin{zproof}
prove;
split a = 0;
cases;
prove;
next;
use MulNonneg[a := a - 1, b := a - 1];
prove;
next;
\end{zproof}

\begin{theorem}{TwoELe}
  \forall y: \nat @ 2 * funE~y \leq y + 1
\end{theorem}

\begin{zproof}
prove;
use EBound[y := y];
use TwoLeSq[a := funRt~y];
use RtSpec[x := y];
use UFNonneg[f := funRt, x := y];
with enabled (EVal) prove;
\end{zproof}

\begin{theorem}{ESurj}
  funE \in Surj
\end{theorem}

\begin{zproof}
invoke Surj;
prove;
instantiate y == x * x + x;
use ESqAdd[m := x, k := x];
prove;
\end{zproof}

$E^{-1}(x) = c^2 + x$ where $c = \lceil x/2 \rceil$.

\begin{theorem}{EinvVal}
  \forall x, c: \nat | x \leq 2 * c \land 2 * c \leq x + 1 @ opInv~funE~x = c * c + x
\end{theorem}

\begin{zproof}
prove;
use ESurj;
use ESqAdd[m := c, k := x];
use InvLeast[F := funE, x := x, y := c * c + x];
use InvRight[F := funE, x := x];
use EBound[y := opInv~funE~x];
use EDecomp[y := opInv~funE~x];
use UFNonneg[f := funRt, x := opInv~funE~x];
use timesMonotonic2[a := c, b := c, c := funRt~(opInv~funE~x), d := funRt~(opInv~funE~x)];
prove;
\end{zproof}

\begin{theorem}{rule EinvD}
  \forall z: \nat @ opInv~funE~(2 * z) = z * z + 2 * z
\end{theorem}

\begin{zproof}
prove;
use EinvVal[x := 2 * z, c := z];
prove;
\end{zproof}

\begin{theorem}{rule EinvOdd}
  \forall m: \nat @ opInv~funE~(1 + 2 * m) = (m + 1) * (m + 1) + 2 * m + 1
\end{theorem}

\begin{zproof}
prove;
use EinvVal[x := 1 + 2 * m, c := m + 1];
prove;
\end{zproof}

\subsection{The case $X = E$: $S$, $\bar 0$, $I$, $D$, $M_3$, $\mathrm{Sq}$}

\begin{theorem}{IExprZero}
  opPlusI~(cst~0) = funI
\end{theorem}

\begin{zproof}
apply UFExt to predicate opPlusI~(cst~0) = funI;
prove;
\end{zproof}

\begin{theorem}{GESExpr}
  opPlusI~(cst~1) \in CGE
\end{theorem}

\begin{zproof}
prove;
\end{zproof}

\begin{theorem}{rule GES}
  funS \in CGE
\end{theorem}

\begin{zproof}
use GESExpr;
use SExpr;
prove;
\end{zproof}

\begin{theorem}{GEZeroExpr}
  comp(funE, cst~1) \in CGE
\end{theorem}

\begin{zproof}
prove;
\end{zproof}

\begin{theorem}{rule GEZero}
  cst~0 \in CGE
\end{theorem}

\begin{zproof}
use GEZeroExpr;
use ZeroExpr;
prove;
\end{zproof}

\begin{theorem}{GEIExpr}
  opPlusI~(cst~0) \in CGE
\end{theorem}

\begin{zproof}
prove;
\end{zproof}

\begin{theorem}{rule GEI}
  funI \in CGE
\end{theorem}

\begin{zproof}
use GEIExpr;
use IExprZero;
prove;
\end{zproof}

\begin{theorem}{GEDExpr}
  opPlusI~(opPlusI~(cst~0)) \in CGE
\end{theorem}

\begin{zproof}
prove;
\end{zproof}

\begin{theorem}{rule GED}
  funD \in CGE
\end{theorem}

\begin{zproof}
use GEDExpr;
use DExpr;
prove;
\end{zproof}

\begin{theorem}{GEM3Expr}
  opPlusI~funD \in CGE
\end{theorem}

\begin{zproof}
prove;
\end{zproof}

\begin{theorem}{rule GEM3}
  funM3 \in CGE
\end{theorem}

\begin{zproof}
use GEM3Expr;
use M3Expr;
prove;
\end{zproof}


$F^- = E\,(S S S M_3 F E)^+ S\,(E^{-1} D F^+)^+$, without $\mathrm{Sq}$.

\begin{zed}
  exprMinusIE == (\lambda F: \nat \fun \nat @ comp(comp(comp(funE, opPlusI~(comp(comp(comp(comp(comp(funS, funS), funS), funM3), F), funE))), funS), opPlusI~(comp(comp(opInv~funE, funD), opPlusI~F))))
\end{zed}

\begin{zproof}[exprMinusIE\$domainCheck]
prove;
\end{zproof}

\begin{theorem}{rule exprMinusIEFun}
  exprMinusIE \in (\nat \fun \nat) \fun (\nat \fun \nat)
\end{theorem}

\begin{zproof}
invoke exprMinusIE;
apply FunIntro to predicate (\lambda F: \nat \fun \nat @ comp(comp(comp(funE, opPlusI~(comp(comp(comp(comp(comp(funS, funS), funS), funM3), F), funE))), funS), opPlusI~(comp(comp(opInv~funE, funD), opPlusI~F)))) \in (\nat \fun \nat) \fun (\nat \fun \nat);
with enabled (RelPow, inPower) prove;
\end{zproof}

\begin[disabled]{theorem}{rule exprMinusIEApp}
  \forall F: \nat \fun \nat @ exprMinusIE~F = comp(comp(comp(funE, opPlusI~(comp(comp(comp(comp(comp(funS, funS), funS), funM3), F), funE))), funS), opPlusI~(comp(comp(opInv~funE, funD), opPlusI~F)))
\end{theorem}

\begin{zproof}
use exprMinusIEFun;
use pairInFunction[\nat \fun \nat, \nat \fun \nat][f := exprMinusIE, x := F, y := comp(comp(comp(funE, opPlusI~(comp(comp(comp(comp(comp(funS, funS), funS), funM3), F), funE))), funS), opPlusI~(comp(comp(opInv~funE, funD), opPlusI~F)))];
invoke exprMinusIE;
prove;
\end{zproof}


\begin{theorem}{rule MinusIEIn}
  \forall F: \nat \fun \nat | F \in CGE @ exprMinusIE~F \in CGE
\end{theorem}

\begin{zproof}
with enabled (exprMinusIEApp) prove;
\end{zproof}

\begin{theorem}{rule MinusIEVal}
  \forall F: \nat \fun \nat; x: \nat | x \leq F~x @ exprMinusIE~F~x = F~x - x
\end{theorem}

\begin{zproof}
with enabled (exprMinusIEApp) prove;
use EinvD[z := F~x + x];
use ESqAdd[m := F~x + x + 1, k := x];
use ESqAdd[m := F~x + x + 2, k := F~x - x];
prove;
\end{zproof}

\begin{theorem}{SqExprGE}
  exprMinusIE~(exprMinusIE~(comp(opInv~funE, funD))) = funSq
\end{theorem}

\begin{zproof}
apply UFExt to predicate exprMinusIE~(exprMinusIE~(comp(opInv~funE, funD))) = funSq;
prove;
\end{zproof}

\begin{theorem}{GESqExpr}
  exprMinusIE~(exprMinusIE~(comp(opInv~funE, funD))) \in CGE
\end{theorem}

\begin{zproof}
prove;
\end{zproof}

\begin{theorem}{rule GESq}
  funSq \in CGE
\end{theorem}

\begin{zproof}
use GESqExpr;
use SqExprGE;
prove;
\end{zproof}


$I - F = E\,(S\,(F^+)^+ E)^+ (\mathrm{Sq}\,F^+)^+$ where $F(x) \leq x$.

\begin{zed}
  exprIminusE == (\lambda F: \nat \fun \nat @ comp(comp(funE, opPlusI~(comp(comp(funS, opPlusI~(opPlusI~F)), funE))), opPlusI~(comp(funSq, opPlusI~F))))
\end{zed}

\begin{zproof}[exprIminusE\$domainCheck]
prove;
\end{zproof}

\begin{theorem}{rule exprIminusEFun}
  exprIminusE \in (\nat \fun \nat) \fun (\nat \fun \nat)
\end{theorem}

\begin{zproof}
invoke exprIminusE;
apply FunIntro to predicate (\lambda F: \nat \fun \nat @ comp(comp(funE, opPlusI~(comp(comp(funS, opPlusI~(opPlusI~F)), funE))), opPlusI~(comp(funSq, opPlusI~F)))) \in (\nat \fun \nat) \fun (\nat \fun \nat);
with enabled (RelPow, inPower) prove;
\end{zproof}

\begin[disabled]{theorem}{rule exprIminusEApp}
  \forall F: \nat \fun \nat @ exprIminusE~F = comp(comp(funE, opPlusI~(comp(comp(funS, opPlusI~(opPlusI~F)), funE))), opPlusI~(comp(funSq, opPlusI~F)))
\end{theorem}

\begin{zproof}
use exprIminusEFun;
use pairInFunction[\nat \fun \nat, \nat \fun \nat][f := exprIminusE, x := F, y := comp(comp(funE, opPlusI~(comp(comp(funS, opPlusI~(opPlusI~F)), funE))), opPlusI~(comp(funSq, opPlusI~F)))];
invoke exprIminusE;
prove;
\end{zproof}


\begin{theorem}{rule IminusEIn}
  \forall F: \nat \fun \nat | F \in CGE @ exprIminusE~F \in CGE
\end{theorem}

\begin{zproof}
with enabled (exprIminusEApp) prove;
\end{zproof}

\begin{theorem}{rule IminusEVal}
  \forall F: \nat \fun \nat; x: \nat | F~x \leq x @ exprIminusE~F~x = x - F~x
\end{theorem}

\begin{zproof}
with enabled (exprIminusEApp) prove;
use UFNonneg[f := F, x := x];
use ESqAdd[m := F~x + x, k := x];
use RobDiffE[a := x, b := F~x];
prove;
\end{zproof}

$W = E\,(S S D E E)^+ E^{-1}$: $W(2q) = 2E(2q) + 1$ and $W(2q + 1) = 2E(2q + 1)$.

\begin{zed}
  exprW == comp(comp(funE, opPlusI~(comp(comp(comp(comp(funS, funS), funD), funE), funE))), opInv~funE)
\end{zed}

\begin{zproof}[exprW\$domainCheck]
prove;
\end{zproof}

\begin{theorem}{WEven}
  \forall q: \nat @ exprW~(2 * q) = 2 * funE~(2 * q) + 1
\end{theorem}

\begin{zproof}
invoke exprW;
prove;
use ESqAdd[m := q, k := 2 * q];
use TwoELe[y := 2 * q];
use ESqAdd[m := q + 1, k := 2 * funE~(2 * q) + 1];
prove;
\end{zproof}

\begin{theorem}{WOdd}
  \forall q: \nat @ exprW~(1 + 2 * q) = 2 * funE~(1 + 2 * q)
\end{theorem}

\begin{zproof}
invoke exprW;
prove;
use ESqAdd[m := q + 1, k := 1 + 2 * q];
use TwoELe[y := 1 + 2 * q];
use ESqAdd[m := q + 2, k := 2 * funE~(1 + 2 * q)];
prove;
\end{zproof}

\begin{theorem}{EAt}
  \forall y, m, k: \nat | y = m * m + k \land k \leq 2 * m @ funE~y = k
\end{theorem}

\begin{zproof}
prove;
use ESqAdd[m := m, k := k];
prove;
\end{zproof}

\begin{theorem}{WFun}
  exprW \in \nat \fun \nat
\end{theorem}

\begin{zproof}
invoke exprW;
prove;
\end{zproof}

\begin{theorem}{WSurj}
  exprW \in Surj
\end{theorem}

\begin{zproof}
invoke Surj;
use WFun;
prove;
use DivModOf[x := x, d := 2];
use TwoA[x := x \div 2];
use TwoA[x := 1 + x \div 2];
prove;
split x \mod 2 = 0;
cases;
instantiate y == 1 + 2 * (funA~(x \div 2) + x \div 2);
use WOdd[q := funA~(x \div 2) + x \div 2];
use EAt[y := 1 + 2 * (funA~(x \div 2) + x \div 2), m := 1 + x \div 2, k := x \div 2];
prove;
next;
instantiate y\_\_0 == 2 * (funA~(1 + x \div 2) + x \div 2 + 1);
use WEven[q := funA~(1 + x \div 2) + x \div 2 + 1];
use EAt[y := 2 * (funA~(1 + x \div 2) + x \div 2 + 1), m := 2 + x \div 2, k := x \div 2];
prove;
next;
\end{zproof}

\subsection{The case $X = E$: $\mathrm{Hf} = E\,W^{-1}$ and $\mathrm{Rt}$}

\begin{theorem}{WPar}
  \forall y: \nat @ (y \mod 2 = 0 \implies exprW~y = 2 * funE~y + 1) \land (y \mod 2 = 1 \implies exprW~y = 2 * funE~y)
\end{theorem}

\begin{zproof}
prove;
use DivModOf[x := y, d := 2];
use WEven[q := y \div 2];
use WOdd[q := y \div 2];
prove;
\end{zproof}

\begin{theorem}{HfExprGE}
  comp(funE, opInv~exprW) = funHf
\end{theorem}

\begin{zproof}
apply UFExt to predicate comp(funE, opInv~exprW) = funHf;
prove;
use WSurj;
use WFun;
use InvRight[F := exprW, x := x];
use UFNonneg[f := opInv~exprW, x := x];
use WPar[y := opInv~exprW~x];
use DivModOf[x := opInv~exprW~x, d := 2];
use UFNonneg[f := funE, x := opInv~exprW~x];
prove;
split opInv~exprW~x \mod 2 = 0;
cases;
use DivMod[x := x, d := 2, q := funE~(opInv~exprW~x), r := 1];
prove;
next;
use DivMod[x := x, d := 2, q := funE~(opInv~exprW~x), r := 0];
prove;
next;
\end{zproof}

\begin{theorem}{GEHfExpr}
  comp(funE, opInv~exprW) \in CGE
\end{theorem}

\begin{zproof}
use WFun;
invoke exprW;
prove;
\end{zproof}

\begin{theorem}{rule GEHf}
  funHf \in CGE
\end{theorem}

\begin{zproof}
use GEHfExpr;
use HfExprGE;
prove;
\end{zproof}


$\mathrm{Rt}(4x + 1) \in \{2\,\mathrm{Rt}(x), 2\,\mathrm{Rt}(x) + 1\}$, so
$\mathrm{Rt} = \mathrm{Hf}\,\mathrm{Hf}\,S\,S\,E\,(I - \bar 1)(I - E)\,S\,M_3^+$.

\begin{theorem}{HfRt4}
  \forall x: \nat @ funHf~(funRt~(1 + 4 * x)) = funRt~x
\end{theorem}

\begin{zproof}
prove;
use RtSpec[x := x];
use UFNonneg[f := funRt, x := x];
split 1 + 4 * x < (2 * funRt~x + 1) * (2 * funRt~x + 1);
cases;
use RtI[r := 2 * funRt~x, x := 1 + 4 * x];
use DivMod[x := 2 * funRt~x, d := 2, q := funRt~x, r := 0];
prove;
next;
use RtI[r := 2 * funRt~x + 1, x := 1 + 4 * x];
use DivMod[x := 2 * funRt~x + 1, d := 2, q := funRt~x, r := 1];
prove;
next;
\end{zproof}

\begin{theorem}{rule IminusEE}
  \forall z: \nat @ exprIminusE~funE~z = funRt~z * funRt~z
\end{theorem}

\begin{zproof}
prove;
use ELe[x := z];
use IminusEVal[F := funE, x := z];
use EDecomp[y := z];
prove;
\end{zproof}

\begin{theorem}{IminusEOne}
  \forall n: \nat | 1 \leq n @ exprIminusE~(cst~1)~(n * n) = n * n - 1
\end{theorem}

\begin{zproof}
prove;
\end{zproof}

\begin{theorem}{ESqPred}
  \forall n: \nat | 1 \leq n @ funE~(n * n - 1) = 2 * n - 2
\end{theorem}

\begin{zproof}
prove;
use EAt[y := n * n - 1, m := n - 1, k := 2 * n - 2];
prove;
\end{zproof}

\begin{theorem}{HfTwo}
  \forall n: \nat @ funHf~(2 * n) = n
\end{theorem}

\begin{zproof}
prove;
\end{zproof}

\begin{zed}
  exprRtGE == comp(comp(comp(comp(comp(comp(comp(comp(funHf, funHf), funS), funS), funE), exprIminusE~(cst~1)), exprIminusE~funE), funS), opPlusI~funM3)
\end{zed}

\begin{zproof}[exprRtGE\$domainCheck]
prove;
\end{zproof}

\begin{theorem}{RtExprGE}
  exprRtGE = funRt
\end{theorem}

\begin{zproof}
invoke exprRtGE;
apply UFExt to predicate comp(comp(comp(comp(comp(comp(comp(comp(funHf, funHf), funS), funS), funE), exprIminusE~(cst~1)), exprIminusE~funE), funS), opPlusI~funM3) = funRt;
prove;
use RtPos[y := 1 + 4 * x];
use IminusEOne[n := funRt~(1 + 4 * x)];
use ESqPred[n := funRt~(1 + 4 * x)];
use HfTwo[n := funRt~(1 + 4 * x)];
use HfRt4[x := x];
prove;
\end{zproof}

\begin{theorem}{GERtExpr}
  exprRtGE \in CGE
\end{theorem}

\begin{zproof}
invoke exprRtGE;
prove;
\end{zproof}

\begin{theorem}{rule GERt}
  funRt \in CGE
\end{theorem}

\begin{zproof}
use GERtExpr;
use RtExprGE;
prove;
\end{zproof}

\begin{theorem}{GERich}
  CGE \in Rich
\end{theorem}

\begin{zproof}
use CGEInPower;
invoke Rich;
with enabled (CGEUF) prove;
\end{zproof}

\subsection{The case $X = E$: the closure is Robinson's}

\begin{theorem}{JRUClosedGE}
  CGE \in UClosed~\{ opInv \}
\end{theorem}

\begin{zproof}
use CGEInPower;
with enabled (UClosedDef, CGEUF) prove;
\end{zproof}

\begin{theorem}{JRBClosedGE}
  CGE \in BClosed~\{ opPlus, comp \}
\end{theorem}

\begin{zproof}
use CGEInPower;
use GERich;
with enabled (BClosedDef, CGEUF) prove;
split b = opPlus;
cases;
prove;
next;
prove;
next;
\end{zproof}

\begin{theorem}{JRSubGE}
  ClosJR \subseteq CGE
\end{theorem}

\begin{zproof}
use ClosLeast[I := \{ funS, funE \}, U := \{ opInv \}, B := \{ opPlus, comp \}, Z := CGE];
use CGEInPower;
use JRUClosedGE;
use JRBClosedGE;
with enabled (ClosJRDef) prove;
\end{zproof}

In Robinson's closure: $I = E\,E^{-1}$, $\bar 0 = E\,S\,E^{-1}(I + I)$, $\bar 1 = S\,\bar 0$ and
$F^+ = F + I$.

\begin{theorem}{EEinv}
  comp(funE, opInv~funE) = funI
\end{theorem}

\begin{zproof}
apply UFExt to predicate comp(funE, opInv~funE) = funI;
prove;
use ESurj;
use InvRight[F := funE, x := x];
prove;
\end{zproof}

\begin{theorem}{JRIExpr}
  comp(funE, opInv~funE) \in ClosJR
\end{theorem}

\begin{zproof}
prove;
\end{zproof}

\begin{theorem}{rule JRI}
  funI \in ClosJR
\end{theorem}

\begin{zproof}
use JRIExpr;
use EEinv;
prove;
\end{zproof}


\begin{theorem}{ZeroExprJR}
  comp(comp(comp(funE, funS), opInv~funE), opPlus(funI, funI)) = cst~0
\end{theorem}

\begin{zproof}
apply UFExt to predicate comp(comp(comp(funE, funS), opInv~funE), opPlus(funI, funI)) = cst~0;
prove;
use EinvD[z := x];
use EAt[y := 1 + x * x + 2 * x, m := x + 1, k := 0];
prove;
\end{zproof}

\begin{theorem}{JRZeroExpr}
  comp(comp(comp(funE, funS), opInv~funE), opPlus(funI, funI)) \in ClosJR
\end{theorem}

\begin{zproof}
prove;
\end{zproof}

\begin{theorem}{rule JRZero}
  cst~0 \in ClosJR
\end{theorem}

\begin{zproof}
use JRZeroExpr;
use ZeroExprJR;
prove;
\end{zproof}

\begin{theorem}{JROneExpr}
  comp(funS, cst~0) \in ClosJR
\end{theorem}

\begin{zproof}
prove;
\end{zproof}

\begin{theorem}{rule JROne}
  cst~1 \in ClosJR
\end{theorem}

\begin{zproof}
use JROneExpr;
use OneExpr;
prove;
\end{zproof}


\begin{theorem}{rule JRPlusI}
  \forall F: \nat \fun \nat | F \in ClosJR @ opPlusI~F \in ClosJR
\end{theorem}

\begin{zproof}
use PlusIEq[F := F];
prove;
\end{zproof}

\begin{theorem}{GEUClosedJR}
  ClosJR \in UClosed~\{ opPlusI, opInv \}
\end{theorem}

\begin{zproof}
use ClosJRInPower;
with enabled (UClosedDef, ClosJRUF) prove;
split u = opPlusI;
cases;
prove;
next;
prove;
next;
\end{zproof}

\begin{theorem}{GEBClosedJR}
  ClosJR \in BClosed~\{ comp \}
\end{theorem}

\begin{zproof}
use ClosJRInPower;
with enabled (BClosedDef, ClosJRUF) prove;
\end{zproof}

\begin{theorem}{GESubJR}
  CGE \subseteq ClosJR
\end{theorem}

\begin{zproof}
use ClosLeast[I := \{ cst~1, funE \}, U := \{ opPlusI, opInv \}, B := \{ comp \}, Z := ClosJR];
use ClosJRInPower;
use GEUClosedJR;
use GEBClosedJR;
with enabled (CGEDef) prove;
\end{zproof}

\begin{theorem}{GREqJRE}
  CGE = ClosJR
\end{theorem}

\begin{zproof}
use GESubJR;
use JRSubGE;
use CGEInPower;
use ClosJRInPower;
use SubAntisym[\nat \fun \nat][A := CGE, B := ClosJR];
prove;
\end{zproof}

\subsection{The case $X = K$}

\begin{theorem}{Halve}
  \forall a, b: \num | 2 * a = 2 * b @ a = b
\end{theorem}

\begin{zproof}
prove;
\end{zproof}


$Z = K^{-1}$, $A = K S S S ((Z Z^+ K)^+ Z K)^+ Z$, $F^- = K (S S S M_3 F K)^+ (A F^+)^+$,
$\mathrm{Sq} = (((D Z)^-)^-)^-$, $P = K (K S)^{-1}$, $U = I - A K P (I - K)$ ($U = V + K$),
\begin{theorem}{Div2Bound}
  \forall x: \nat @ 2 * (x \div 2) \leq x \land x \leq 2 * (x \div 2) + 1
\end{theorem}

\begin{zproof}
use DivModOf[x := x, d := 2];
prove;
\end{zproof}

$\mathrm{Hf} = K U^{-1}$, $B = I - (I - \mathrm{Sq} (I - \mathrm{Hf}) U)(I - K)$
($B = K + \lceil V/2 \rceil^2$) and $E = K B^{-1}$.

\begin{zed}
  CGK == clos~\{ cst~1, funK \}~\{ opPlusI, opInv \}~\{ comp \}
\end{zed}

\begin{zproof}[CGK\$domainCheck]
prove;
\end{zproof}

\begin[disabled]{theorem}{rule CGKDef}
  CGK = clos~\{ cst~1, funK \}~\{ opPlusI, opInv \}~\{ comp \}
\end{theorem}

\begin{zproof}
invoke CGK;
prove;
\end{zproof}

\begin{theorem}{CGKInPower}
  CGK \in \power (\nat \fun \nat)
\end{theorem}

\begin{zproof}
use ClosInPow[I := \{ cst~1, funK \}, U := \{ opPlusI, opInv \}, B := \{ comp \}];
with enabled (CGKDef) prove;
\end{zproof}

\begin[disabled]{theorem}{rule CGKUF}
  \forall F: \power (\num \cross \num) | F \in CGK @ F \in \nat \fun \nat
\end{theorem}

\begin{zproof}
prove;
use ClosElemUF[I := \{ cst~1, funK \}, U := \{ opPlusI, opInv \}, B := \{ comp \}, F := F];
with enabled (CGKDef) prove;
\end{zproof}

\begin{theorem}{rule CGKInit0}
  cst~1 \in CGK
\end{theorem}

\begin{zproof}
with enabled (CGKDef) prove;
\end{zproof}

\begin{theorem}{rule CGKInit1}
  funK \in CGK
\end{theorem}

\begin{zproof}
with enabled (CGKDef) prove;
\end{zproof}

\begin{theorem}{rule CGKUnop0}
  \forall F: \nat \fun \nat | F \in CGK @ opPlusI~F \in CGK
\end{theorem}

\begin{zproof}
with enabled (CGKDef) prove;
\end{zproof}

\begin{theorem}{rule CGKUnop1}
  \forall F: \nat \fun \nat | F \in CGK @ opInv~F \in CGK
\end{theorem}

\begin{zproof}
with enabled (CGKDef) prove;
\end{zproof}

\begin{theorem}{rule CGKBinop0}
  \forall F, G: \nat \fun \nat | F \in CGK \land G \in CGK @ comp(F, G) \in CGK
\end{theorem}

\begin{zproof}
with enabled (CGKDef) prove;
\end{zproof}


\begin{theorem}{ZeroExprK}
  comp(funK, cst~1) = cst~0
\end{theorem}

\begin{zproof}
apply UFExt to predicate comp(funK, cst~1) = cst~0;
prove;
use KAt[m := 1, t := 1, k := 0];
use AZero;
prove;
\end{zproof}

\begin{theorem}{GKSExpr}
  opPlusI~(cst~1) \in CGK
\end{theorem}

\begin{zproof}
prove;
\end{zproof}

\begin{theorem}{rule GKS}
  funS \in CGK
\end{theorem}

\begin{zproof}
use GKSExpr;
use SExpr;
prove;
\end{zproof}

\begin{theorem}{GKZeroExpr}
  comp(funK, cst~1) \in CGK
\end{theorem}

\begin{zproof}
prove;
\end{zproof}

\begin{theorem}{rule GKZero}
  cst~0 \in CGK
\end{theorem}

\begin{zproof}
use GKZeroExpr;
use ZeroExprK;
prove;
\end{zproof}

\begin{theorem}{GKIExpr}
  opPlusI~(cst~0) \in CGK
\end{theorem}

\begin{zproof}
prove;
\end{zproof}

\begin{theorem}{rule GKI}
  funI \in CGK
\end{theorem}

\begin{zproof}
use GKIExpr;
use IExprZero;
prove;
\end{zproof}

\begin{theorem}{GKDExpr}
  opPlusI~(opPlusI~(cst~0)) \in CGK
\end{theorem}

\begin{zproof}
prove;
\end{zproof}

\begin{theorem}{rule GKD}
  funD \in CGK
\end{theorem}

\begin{zproof}
use GKDExpr;
use DExpr;
prove;
\end{zproof}

\begin{theorem}{GKM3Expr}
  opPlusI~funD \in CGK
\end{theorem}

\begin{zproof}
prove;
\end{zproof}

\begin{theorem}{rule GKM3}
  funM3 \in CGK
\end{theorem}

\begin{zproof}
use GKM3Expr;
use M3Expr;
prove;
\end{zproof}


\begin{theorem}{KSurj}
  funK \in Surj
\end{theorem}

\begin{zproof}
invoke Surj;
prove;
instantiate y == funA~x + x;
use KAt[m := funA~x + x, t := x, k := x];
use UFNonneg[f := funA, x := x];
prove;
\end{zproof}

\begin{theorem}{KinvZ}
  opInv~funK = funZ
\end{theorem}

\begin{zproof}
apply UFExt to predicate opInv~funK = funZ;
prove;
use KSurj;
use KAt[m := funA~x + x, t := x, k := x];
use UFNonneg[f := funA, x := x];
use InvLeast[F := funK, x := x, y := funA~x + x];
use InvRight[F := funK, x := x];
use UFNonneg[f := opInv~funK, x := x];
use Decomp[x := opInv~funK~x];
use UFNonneg[f := funV, x := opInv~funK~x];
use AMono[m := x, n := funV~(opInv~funK~x)];
prove;
\end{zproof}

\begin{theorem}{GKZExpr}
  opInv~funK \in CGK
\end{theorem}

\begin{zproof}
prove;
\end{zproof}

\begin{theorem}{rule GKZ}
  funZ \in CGK
\end{theorem}

\begin{zproof}
use GKZExpr;
use KinvZ;
prove;
\end{zproof}


\begin{theorem}{KZ}
  \forall x: \nat @ funK~(funA~x + x) = x
\end{theorem}

\begin{zproof}
prove;
use KAt[m := funA~x + x, t := x, k := x];
use UFNonneg[f := funA, x := x];
prove;
\end{zproof}

\begin{zed}
  exprAK == comp(comp(comp(comp(comp(funK, funS), funS), funS), opPlusI~(comp(comp(opPlusI~(comp(comp(funZ, opPlusI~funZ), funK)), funZ), funK))), funZ)
\end{zed}

\begin{zproof}[exprAK\$domainCheck]
prove;
\end{zproof}

\begin{theorem}{AExprK}
  exprAK = funA
\end{theorem}

\begin{zproof}
invoke exprAK;
apply UFExt to predicate comp(comp(comp(comp(comp(funK, funS), funS), funS), opPlusI~(comp(comp(opPlusI~(comp(comp(funZ, opPlusI~funZ), funK)), funZ), funK))), funZ) = funA;
prove;
use KZ[x := x];
use UFNonneg[f := funA, x := x];
use RobDiffK[a := funA~x + x, b := x];
prove;
\end{zproof}

\begin{theorem}{GKAExpr}
  exprAK \in CGK
\end{theorem}

\begin{zproof}
invoke exprAK;
prove;
\end{zproof}

\begin{theorem}{rule GKA}
  funA \in CGK
\end{theorem}

\begin{zproof}
use GKAExpr;
use AExprK;
prove;
\end{zproof}

\begin{zed}
  exprMinusIK == (\lambda F: \nat \fun \nat @ comp(comp(funK, opPlusI~(comp(comp(comp(comp(comp(funS, funS), funS), funM3), F), funK))), opPlusI~(comp(funA, opPlusI~F))))
\end{zed}

\begin{zproof}[exprMinusIK\$domainCheck]
prove;
\end{zproof}

\begin{theorem}{rule exprMinusIKFun}
  exprMinusIK \in (\nat \fun \nat) \fun (\nat \fun \nat)
\end{theorem}

\begin{zproof}
invoke exprMinusIK;
apply FunIntro to predicate (\lambda F: \nat \fun \nat @ comp(comp(funK, opPlusI~(comp(comp(comp(comp(comp(funS, funS), funS), funM3), F), funK))), opPlusI~(comp(funA, opPlusI~F)))) \in (\nat \fun \nat) \fun (\nat \fun \nat);
with enabled (RelPow, inPower) prove;
\end{zproof}

\begin[disabled]{theorem}{rule exprMinusIKApp}
  \forall F: \nat \fun \nat @ exprMinusIK~F = comp(comp(funK, opPlusI~(comp(comp(comp(comp(comp(funS, funS), funS), funM3), F), funK))), opPlusI~(comp(funA, opPlusI~F)))
\end{theorem}

\begin{zproof}
use exprMinusIKFun;
use pairInFunction[\nat \fun \nat, \nat \fun \nat][f := exprMinusIK, x := F, y := comp(comp(funK, opPlusI~(comp(comp(comp(comp(comp(funS, funS), funS), funM3), F), funK))), opPlusI~(comp(funA, opPlusI~F)))];
invoke exprMinusIK;
prove;
\end{zproof}


\begin{theorem}{rule MinusIKIn}
  \forall F: \nat \fun \nat | F \in CGK @ exprMinusIK~F \in CGK
\end{theorem}

\begin{zproof}
with enabled (exprMinusIKApp) prove;
\end{zproof}

\begin{theorem}{rule MinusIKVal}
  \forall F: \nat \fun \nat; x: \nat | x \leq F~x @ exprMinusIK~F~x = F~x - x
\end{theorem}

\begin{zproof}
with enabled (exprMinusIKApp) prove;
use UFNonneg[f := F, x := x];
use UFNonneg[f := funA, x := F~x + x];
use KAt[m := funA~(F~x + x) + x, t := F~x + x, k := x];
use ASuc[x := F~x + x];
use ASuc[x := 1 + F~x + x];
use KAt[m := 3 + (funA~(F~x + x) + x + 3 * F~x), t := 2 + F~x + x, k := F~x - x];
prove;
\end{zproof}

\begin{theorem}{SqExprGK}
  exprMinusIK~(exprMinusIK~(exprMinusIK~(comp(funD, funZ)))) = funSq
\end{theorem}

\begin{zproof}
apply UFExt to predicate exprMinusIK~(exprMinusIK~(exprMinusIK~(comp(funD, funZ)))) = funSq;
prove;
use TwoA[x := x];
prove;
\end{zproof}

\begin{theorem}{GKSqExpr}
  exprMinusIK~(exprMinusIK~(exprMinusIK~(comp(funD, funZ)))) \in CGK
\end{theorem}

\begin{zproof}
prove;
\end{zproof}

\begin{theorem}{rule GKSq}
  funSq \in CGK
\end{theorem}

\begin{zproof}
use GKSqExpr;
use SqExprGK;
prove;
\end{zproof}


The predecessor $P = K (K S)^{-1}$.

\begin{theorem}{KSSurj}
  comp(funK, funS) \in Surj
\end{theorem}

\begin{zproof}
invoke Surj;
prove;
instantiate y == funA~x + 2 * x;
use ASuc[x := x];
use UFNonneg[f := funA, x := x];
use KAt[m := 1 + funA~x + 2 * x, t := 1 + x, k := x];
prove;
\end{zproof}

\begin{zed}
  exprPK == comp(funK, opInv~(comp(funK, funS)))
\end{zed}

\begin{zproof}[exprPK\$domainCheck]
prove;
\end{zproof}

\begin{theorem}{rule PKFun}
  exprPK \in \nat \fun \nat
\end{theorem}

\begin{zproof}
invoke exprPK;
prove;
\end{zproof}

\begin{theorem}{KSInvRight}
  \forall x: \nat @ funK~(1 + opInv~(comp(funK, funS))~x) = x
\end{theorem}

\begin{zproof}
prove;
use KSSurj;
use InvRight[F := comp(funK, funS), x := x];
with disabled (InvRight) prove;
\end{zproof}

\begin{theorem}{PKZero}
  exprPK~0 = 0
\end{theorem}

\begin{zproof}
invoke exprPK;
prove;
use KSSurj;
use KAt[m := 1, t := 1, k := 0];
use AZero;
use InvLeast[F := comp(funK, funS), x := 0, y := 0];
use UFNonneg[f := opInv~(comp(funK, funS)), x := 0];
use KZero;
prove;
\end{zproof}

\begin{theorem}{PKSuc}
  \forall x: \nat @ exprPK~(1 + x) = x
\end{theorem}

\begin{zproof}
invoke exprPK;
prove;
use KSInvRight[x := 1 + x];
use UFNonneg[f := opInv~(comp(funK, funS)), x := 1 + x];
use Decomp[x := 1 + opInv~(comp(funK, funS))~(1 + x)];
use UFNonneg[f := funA, x := funV~(1 + opInv~(comp(funK, funS))~(1 + x))];
use KAt[m := opInv~(comp(funK, funS))~(1 + x), t := funV~(1 + opInv~(comp(funK, funS))~(1 + x)), k := x];
prove;
\end{zproof}

\begin{theorem}{GKPExpr}
  comp(funK, opInv~(comp(funK, funS))) \in CGK
\end{theorem}

\begin{zproof}
prove;
\end{zproof}

\begin{theorem}{rule GKP}
  exprPK \in CGK
\end{theorem}

\begin{zproof}
use GKPExpr;
invoke exprPK;
prove;
\end{zproof}


\begin{theorem}{rule GKIminusF}
  \forall F: \nat \fun \nat | F \in CGK @ exprIminusF~F \in CGK
\end{theorem}

\begin{zproof}
with enabled (exprIminusFApp) prove;
\end{zproof}

\begin{theorem}{rule IminusKVal}
  \forall y: \nat @ exprIminusF~funK~y = funA~(funV~y)
\end{theorem}

\begin{zproof}
prove;
use Decomp[x := y];
use UFNonneg[f := funA, x := funV~y];
use IminusFVal[F := funK, x := y];
prove;
\end{zproof}

\begin{zed}
  exprU == exprIminusF~(comp(comp(comp(funA, funK), exprPK), exprIminusF~funK))
\end{zed}

\begin{zproof}[exprU\$domainCheck]
use PKFun;
prove;
\end{zproof}

\begin{theorem}{UIn}
  exprU \in CGK
\end{theorem}

\begin{zproof}
invoke exprU;
use GKP;
prove;
\end{zproof}

\begin{theorem}{UPairPos}
  \forall t, k: \nat | k \leq 1 + t @ exprU~(funA~(1 + t) + k) = 1 + t + k
\end{theorem}

\begin{zproof}
invoke exprU;
prove;
use Diag[k := k, t := 1 + t];
use ASuc[x := t];
use UFNonneg[f := funA, x := t];
use PKSuc[x := funA~t + t];
use KAt[m := funA~t + t, t := t, k := t];
use IminusFVal[F := comp(comp(comp(funA, funK), exprPK), exprIminusF~funK), x := funA~(1 + t) + k];
prove;
\end{zproof}

\begin{theorem}{UZero}
  exprU~0 = 0
\end{theorem}

\begin{zproof}
use IminusFVal[F := comp(comp(comp(funA, funK), exprPK), exprIminusF~funK), x := 0];
use IminusKVal[y := 0];
use Diag[k := 0, t := 0];
use AZero;
use PKZero;
use KZero;
invoke exprU;
prove;
\end{zproof}

\begin{theorem}{UPair}
  \forall k, t: \nat | k \leq t @ exprU~(funA~t + k) = t + k
\end{theorem}

\begin{zproof}
prove;
split t = 0;
cases;
use UZero;
use AZero;
prove;
next;
use UPairPos[t := t - 1, k := k];
prove;
next;
\end{zproof}

\begin{theorem}{UVal}
  \forall y: \nat @ exprU~y = funV~y + funK~y
\end{theorem}

\begin{zproof}
prove;
use Decomp[x := y];
use UPair[k := funK~y, t := funV~y];
prove;
\end{zproof}

\begin{theorem}{rule UFun}
  exprU \in \nat \fun \nat
\end{theorem}

\begin{zproof}
invoke exprU;
prove;
\end{zproof}

\begin{theorem}{USurj}
  exprU \in Surj
\end{theorem}

\begin{zproof}
invoke Surj;
use UFun;
prove;
use DivModOf[x := x, d := 2];
split x \mod 2 = 0;
cases;
instantiate y == funA~(x \div 2) + x \div 2;
use UPair[k := x \div 2, t := x \div 2];
prove;
next;
instantiate y\_\_0 == funA~(1 + x \div 2) + x \div 2;
use UPair[k := x \div 2, t := 1 + x \div 2];
prove;
next;
\end{zproof}

$\mathrm{Hf} = K U^{-1}$: the least $y = A(t) + k$ with $t + k = x$ has $k = \lfloor x/2 \rfloor$.

\begin{theorem}{HfMin}
  \forall x, y, t, k: \nat | y = funA~t + k \land k \leq t \land t + k = x \land y \leq funA~(x - x \div 2) + x \div 2 @ k = x \div 2
\end{theorem}

\begin{zproof}
prove;
use Div2Bound[x := x];
use ASuc[x := x - x \div 2];
use UFNonneg[f := funA, x := x - x \div 2];
use AMono[m := 1 + (x - x \div 2), n := t];
prove;
\end{zproof}

\begin{theorem}{HfExprK}
  comp(funK, opInv~exprU) = funHf
\end{theorem}

\begin{zproof}
apply UFExt to predicate comp(funK, opInv~exprU) = funHf;
prove;
use USurj;
use InvRight[F := exprU, x := x];
use UFNonneg[f := opInv~exprU, x := x];
use UVal[y := opInv~exprU~x];
use Decomp[x := opInv~exprU~x];
use Div2Bound[x := x];
use UPair[k := x \div 2, t := x - x \div 2];
use InvLeast[F := exprU, x := x, y := funA~(x - x \div 2) + x \div 2];
use HfMin[x := x, y := opInv~exprU~x, t := funV~(opInv~exprU~x), k := funK~(opInv~exprU~x)];
prove;
\end{zproof}

\begin{theorem}{GKHfExpr}
  comp(funK, opInv~exprU) \in CGK
\end{theorem}

\begin{zproof}
use UIn;
use UFun;
prove;
\end{zproof}

\begin{theorem}{rule GKHf}
  funHf \in CGK
\end{theorem}

\begin{zproof}
use GKHfExpr;
use HfExprK;
prove;
\end{zproof}

$B = K + \lceil V/2 \rceil^2$ and $E = K B^{-1}$.

\begin{theorem}{AHalf}
  \forall h, r: \nat | r = 0 \lor r = 1 @ funA~(2 * h + r) = (h + r) * (h + r) + h * h + h
\end{theorem}

\begin{zproof}
prove;
use TwoA[x := 2 * h + r];
use Halve[a := funA~(2 * h + r), b := (h + r) * (h + r) + h * h + h];
split r = 0;
cases;
prove;
next;
prove;
next;
\end{zproof}

\begin{theorem}{rule IminusHfEven}
  \forall h: \nat @ exprIminusF~funHf~(2 * h) = h
\end{theorem}

\begin{zproof}
prove;
use DivMod[x := 2 * h, d := 2, q := h, r := 0];
use IminusFVal[F := funHf, x := 2 * h];
prove;
\end{zproof}

\begin{theorem}{rule IminusHfOdd}
  \forall h: \nat @ exprIminusF~funHf~(1 + 2 * h) = 1 + h
\end{theorem}

\begin{zproof}
prove;
use DivMod[x := 1 + 2 * h, d := 2, q := h, r := 1];
use IminusFVal[F := funHf, x := 1 + 2 * h];
prove;
\end{zproof}

\begin{zed}
  exprB == exprIminusF~(comp(exprIminusF~(comp(comp(funSq, exprIminusF~funHf), exprU)), exprIminusF~funK))
\end{zed}

\begin{zproof}[exprB\$domainCheck]
prove;
\end{zproof}

\begin{theorem}{BIn}
  exprB \in CGK
\end{theorem}

\begin{zproof}
invoke exprB;
use UIn;
use UFun;
prove;
\end{zproof}

\begin{theorem}{rule BFun}
  exprB \in \nat \fun \nat
\end{theorem}

\begin{zproof}
invoke exprB;
use UFun;
prove;
\end{zproof}

\begin{theorem}{BVal}
  \forall h, r, k: \nat | (r = 0 \lor r = 1) \land k \leq 2 * h + r @ exprB~(funA~(2 * h + r) + k) = k + (h + r) * (h + r)
\end{theorem}

\begin{zproof}
invoke exprB;
prove;
use Diag[k := k, t := 2 * h + r];
use AHalf[h := h, r := r];
use MulNonneg[a := h, b := h];
use UPair[k := 0, t := 2 * h + r];
use IminusFVal[F := comp(comp(funSq, exprIminusF~funHf), exprU), x := funA~(2 * h + r)];
use IminusFVal[F := comp(exprIminusF~(comp(comp(funSq, exprIminusF~funHf), exprU)), exprIminusF~funK), x := funA~(2 * h + r) + k];
prove;
split r = 0;
cases;
prove;
next;
prove;
next;
\end{zproof}

\begin{theorem}{BSurj}
  exprB \in Surj
\end{theorem}

\begin{zproof}
invoke Surj;
use BFun;
prove;
instantiate y == funA~(2 * funRt~x) + funE~x;
use EBound[y := x];
use UFNonneg[f := funA, x := 2 * funRt~x];
use Diag[k := funE~x, t := 2 * funRt~x];
use BVal[h := funRt~x, r := 0, k := funE~x];
use EDecomp[y := x];
prove;
\end{zproof}

\begin{theorem}{Mod2Cases}
  \forall x: \nat @ x = 2 * (x \div 2) + x \mod 2 \land (x \mod 2 = 0 \lor x \mod 2 = 1)
\end{theorem}

\begin{zproof}
prove;
\end{zproof}

\begin{theorem}{Div2Nonneg}
  \forall x: \nat @ 0 \leq x \div 2
\end{theorem}

\begin{zproof}
prove;
\end{zproof}

\begin{theorem}{EKey}
  \forall x, y, t, k, q, r: \nat | y = funA~t + k \land k \leq t \land t = 2 * q + r \land (r = 0 \lor r = 1) \land exprB~y = x @ funE~x = k
\end{theorem}

\begin{zproof}
prove;
use BVal[h := q, r := r, k := k];
use EAt[y := x, m := q + r, k := k];
prove;
\end{zproof}

\begin{theorem}{EExprKB}
  comp(funK, opInv~exprB) = funE
\end{theorem}

\begin{zproof}
apply UFExt to predicate comp(funK, opInv~exprB) = funE;
prove;
use BSurj;
use InvRight[F := exprB, x := x];
use UFNonneg[f := opInv~exprB, x := x];
use Decomp[x := opInv~exprB~x];
use Mod2Cases[x := funV~(opInv~exprB~x)];
use Div2Nonneg[x := funV~(opInv~exprB~x)];
use EKey[x := x, y := opInv~exprB~x, t := funV~(opInv~exprB~x), k := funK~(opInv~exprB~x), q := funV~(opInv~exprB~x) \div 2, r := funV~(opInv~exprB~x) \mod 2];
prove;
\end{zproof}

\begin{theorem}{GKEExpr}
  comp(funK, opInv~exprB) \in CGK
\end{theorem}

\begin{zproof}
use BIn;
use BFun;
prove;
\end{zproof}

\begin{theorem}{rule GKE}
  funE \in CGK
\end{theorem}

\begin{zproof}
use GKEExpr;
use EExprKB;
prove;
\end{zproof}

\subsection{The cases $X = K$ and $X = L$: the closure is Robinson's}

\begin{theorem}{CGEUClosedG}
  CGE \in UClosed~\{ opPlusI, opInv \}
\end{theorem}

\begin{zproof}
use CGEInPower;
with enabled (UClosedDef, CGEUF) prove;
split u = opPlusI;
cases;
prove;
next;
prove;
next;
\end{zproof}

\begin{theorem}{CGEBClosedG}
  CGE \in BClosed~\{ comp \}
\end{theorem}

\begin{zproof}
use CGEInPower;
with enabled (BClosedDef, CGEUF) prove;
\end{zproof}

\begin{theorem}{CGKUClosedG}
  CGK \in UClosed~\{ opPlusI, opInv \}
\end{theorem}

\begin{zproof}
use CGKInPower;
with enabled (UClosedDef, CGKUF) prove;
split u = opPlusI;
cases;
prove;
next;
prove;
next;
\end{zproof}

\begin{theorem}{CGKBClosedG}
  CGK \in BClosed~\{ comp \}
\end{theorem}

\begin{zproof}
use CGKInPower;
with enabled (BClosedDef, CGKUF) prove;
\end{zproof}


\begin{theorem}{rule GEK}
  funK \in CGE
\end{theorem}

\begin{zproof}
use GERich;
use RichK[Y := CGE];
prove;
\end{zproof}

\begin{theorem}{rule GEL}
  funL \in CGE
\end{theorem}

\begin{zproof}
use GERich;
use RichL[Y := CGE];
prove;
\end{zproof}

\begin{theorem}{GKSubGE}
  CGK \subseteq CGE
\end{theorem}

\begin{zproof}
use ClosLeast[I := \{ cst~1, funK \}, U := \{ opPlusI, opInv \}, B := \{ comp \}, Z := CGE];
use CGEInPower;
use CGEUClosedG;
use CGEBClosedG;
with enabled (CGKDef) prove;
\end{zproof}

\begin{theorem}{GESubGK}
  CGE \subseteq CGK
\end{theorem}

\begin{zproof}
use ClosLeast[I := \{ cst~1, funE \}, U := \{ opPlusI, opInv \}, B := \{ comp \}, Z := CGK];
use CGKInPower;
use CGKUClosedG;
use CGKBClosedG;
with enabled (CGEDef) prove;
\end{zproof}

\begin{zed}
  CGL == clos~\{ cst~1, funL \}~\{ opPlusI, opInv \}~\{ comp \}
\end{zed}

\begin{zproof}[CGL\$domainCheck]
prove;
\end{zproof}

\begin[disabled]{theorem}{rule CGLDef}
  CGL = clos~\{ cst~1, funL \}~\{ opPlusI, opInv \}~\{ comp \}
\end{theorem}

\begin{zproof}
invoke CGL;
prove;
\end{zproof}

\begin{theorem}{CGLInPower}
  CGL \in \power (\nat \fun \nat)
\end{theorem}

\begin{zproof}
use ClosInPow[I := \{ cst~1, funL \}, U := \{ opPlusI, opInv \}, B := \{ comp \}];
with enabled (CGLDef) prove;
\end{zproof}

\begin[disabled]{theorem}{rule CGLUF}
  \forall F: \power (\num \cross \num) | F \in CGL @ F \in \nat \fun \nat
\end{theorem}

\begin{zproof}
prove;
use ClosElemUF[I := \{ cst~1, funL \}, U := \{ opPlusI, opInv \}, B := \{ comp \}, F := F];
with enabled (CGLDef) prove;
\end{zproof}

\begin{theorem}{rule CGLInit0}
  cst~1 \in CGL
\end{theorem}

\begin{zproof}
with enabled (CGLDef) prove;
\end{zproof}

\begin{theorem}{rule CGLInit1}
  funL \in CGL
\end{theorem}

\begin{zproof}
with enabled (CGLDef) prove;
\end{zproof}

\begin{theorem}{rule CGLUnop0}
  \forall F: \nat \fun \nat | F \in CGL @ opPlusI~F \in CGL
\end{theorem}

\begin{zproof}
with enabled (CGLDef) prove;
\end{zproof}

\begin{theorem}{rule CGLUnop1}
  \forall F: \nat \fun \nat | F \in CGL @ opInv~F \in CGL
\end{theorem}

\begin{zproof}
with enabled (CGLDef) prove;
\end{zproof}

\begin{theorem}{rule CGLBinop0}
  \forall F, G: \nat \fun \nat | F \in CGL \land G \in CGL @ comp(F, G) \in CGL
\end{theorem}

\begin{zproof}
with enabled (CGLDef) prove;
\end{zproof}

\begin{theorem}{CGLUClosedG}
  CGL \in UClosed~\{ opPlusI, opInv \}
\end{theorem}

\begin{zproof}
use CGLInPower;
with enabled (UClosedDef, CGLUF) prove;
split u = opPlusI;
cases;
prove;
next;
prove;
next;
\end{zproof}

\begin{theorem}{CGLBClosedG}
  CGL \in BClosed~\{ comp \}
\end{theorem}

\begin{zproof}
use CGLInPower;
with enabled (BClosedDef, CGLUF) prove;
\end{zproof}

\begin{theorem}{GLKExpr}
  comp(comp(funL, opPlusI~(cst~1)), opPlusI~(comp(opPlusI~(opPlusI~(cst~1)), funL))) \in CGL
\end{theorem}

\begin{zproof}
prove;
\end{zproof}

\begin{theorem}{rule GLK}
  funK \in CGL
\end{theorem}

\begin{zproof}
use GLKExpr;
use KExprL;
prove;
\end{zproof}

\begin{theorem}{GKSubGL}
  CGK \subseteq CGL
\end{theorem}

\begin{zproof}
use ClosLeast[I := \{ cst~1, funK \}, U := \{ opPlusI, opInv \}, B := \{ comp \}, Z := CGL];
use CGLInPower;
use CGLUClosedG;
use CGLBClosedG;
with enabled (CGKDef) prove;
\end{zproof}

\begin{theorem}{GLSubGE}
  CGL \subseteq CGE
\end{theorem}

\begin{zproof}
use ClosLeast[I := \{ cst~1, funL \}, U := \{ opPlusI, opInv \}, B := \{ comp \}, Z := CGE];
use CGEInPower;
use CGEUClosedG;
use CGEBClosedG;
with enabled (CGLDef) prove;
\end{zproof}


\begin{theorem}{rule GLE}
  funE \in CGL
\end{theorem}

\begin{zproof}
use GKSubGL;
use GKE;
apply subsetDef to predicate CGK \subseteq CGL;
apply inPower to predicate CGK \in \power CGL;
instantiate e == funE;
prove;
\end{zproof}

\begin{theorem}{GESubGL}
  CGE \subseteq CGL
\end{theorem}

\begin{zproof}
use ClosLeast[I := \{ cst~1, funE \}, U := \{ opPlusI, opInv \}, B := \{ comp \}, Z := CGL];
use CGLInPower;
use CGLUClosedG;
use CGLBClosedG;
with enabled (CGEDef) prove;
\end{zproof}


\begin{theorem}{GREqJR}
  CGE = ClosJR \land CGK = ClosJR \land CGL = ClosJR
\end{theorem}

\begin{zproof}
use GREqJRE;
use GKSubGE;
use GESubGK;
use GKSubGL;
use GLSubGE;
use CGEInPower;
use CGKInPower;
use CGLInPower;
use SubAntisym[\nat \fun \nat][A := CGK, B := CGE];
use SubAntisym[\nat \fun \nat][A := CGL, B := CGE];
use GESubGL;
prove;
\end{zproof}

\subsection{The constant cannot be dropped}

Without $\bar 1$, every generated function fixes $0$: $\langle X, F^+, FG, F^{-1}\rangle$ does not
generate Robinson's closure for $X \in \{E, K, L\}$.

\begin{theorem}{rule L0}
  funL~0 = 0
\end{theorem}

\begin{zproof}
use LAt[m := 0, t := 0, k := 0];
use AZero;
prove;
\end{zproof}

\begin{theorem}{FixZeroUClosedInv}
  FixZero \in UClosed~\{ opPlusI, opInv \}
\end{theorem}

\begin{zproof}
use FixZeroInPower;
with enabled (UClosedDef) prove;
invoke FixZero;
prove;
split u = opPlusI;
cases;
prove;
next;
use InvFix0[F := F];
prove;
next;
\end{zproof}

\begin{theorem}{GR0Fix}
  \forall X: \nat \fun \nat | X~0 = 0 @ clos~\{ X \}~\{ opPlusI, opInv \}~\{ comp \} \subseteq FixZero
\end{theorem}

\begin{zproof}
prove;
use ClosLeast[I := \{ X \}, U := \{ opPlusI, opInv \}, B := \{ comp \}, Z := FixZero];
use FixZeroInPower;
use FixZeroUClosedInv;
use FixZeroBClosed;
invoke FixZero;
prove;
\end{zproof}

\begin{theorem}{GR0NotOne}
  \forall X: \nat \fun \nat | X~0 = 0 @ \lnot cst~1 \in clos~\{ X \}~\{ opPlusI, opInv \}~\{ comp \}
\end{theorem}

\begin{zproof}
use GR0Fix[X := X];
apply subsetDef to predicate clos~\{ X \}~\{ opPlusI, opInv \}~\{ comp \} \subseteq FixZero;
apply inPower to predicate clos~\{ X \}~\{ opPlusI, opInv \}~\{ comp \} \in \power FixZero;
instantiate e == cst~1;
invoke FixZero;
prove;
\end{zproof}

\begin{theorem}{OneInJREq}
  \forall Y: \power (\nat \fun \nat) | Y = ClosJR @ cst~1 \in Y
\end{theorem}

\begin{zproof}
prove;
\end{zproof}

\begin{theorem}{GR0Neq}
  \forall X: \nat \fun \nat | X~0 = 0 @ \lnot clos~\{ X \}~\{ opPlusI, opInv \}~\{ comp \} = ClosJR
\end{theorem}

\begin{zproof}
use GR0NotOne[X := X];
use JROne;
prove;
split clos~\{ X \}~\{ opPlusI, opInv \}~\{ comp \} = ClosJR;
cases;
equality substitute clos~\{ X \}~(\{ opPlusI \} \cup \{ opInv \})~\{ comp \};
prove;
next;
prove;
next;
\end{zproof}

\begin{theorem}{ExtGR0}
  \lnot clos~\{ funE \}~\{ opPlusI, opInv \}~\{ comp \} = ClosJR \land \lnot clos~\{ funK \}~\{ opPlusI, opInv \}~\{ comp \} = ClosJR \land \lnot clos~\{ funL \}~\{ opPlusI, opInv \}~\{ comp \} = ClosJR
\end{theorem}

\begin{zproof}
use GR0Neq[X := funE];
use GR0Neq[X := funK];
use GR0Neq[X := funL];
use KZero;
prove;
\end{zproof}

\subsection{Remark: $\bar 1$ and $S$ are interchangeable}

With $S$ in place of $\bar 1$: $\bar 0 = (E S E^{-1})(E S E^{-1})$, $\bar 0 = K S K^{-1}$ and
$\bar 0 = L (L S)^{-1} S$, and $\bar 1 = S \bar 0$.

\begin{theorem}{ESEEven}
  \forall x, q: \nat | x = 2 * q @ comp(comp(funE, funS), opInv~funE)~x = 0
\end{theorem}

\begin{zproof}
prove;
use EAt[y := 1 + q * q + 2 * q, m := q + 1, k := 0];
prove;
\end{zproof}

\begin{theorem}{ESEOdd}
  \forall x, q: \nat | x = 1 + 2 * q @ comp(comp(funE, funS), opInv~funE)~x = 2 + 2 * q
\end{theorem}

\begin{zproof}
prove;
use EAt[y := 2 + (q + 1) * (q + 1) + 2 * q, m := q + 1, k := 2 + 2 * q];
prove;
\end{zproof}

\begin{theorem}{ZeroExprSE}
  comp(comp(comp(funE, funS), opInv~funE), comp(comp(funE, funS), opInv~funE)) = cst~0
\end{theorem}

\begin{zproof}
apply UFExt to predicate comp(comp(comp(funE, funS), opInv~funE), comp(comp(funE, funS), opInv~funE)) = cst~0;
prove;
use Mod2Cases[x := x];
use Div2Nonneg[x := x];
use ESEEven[x := 0, q := 0];
split x \mod 2 = 0;
cases;
use ESEEven[x := x, q := x \div 2];
prove;
next;
use ESEOdd[x := x, q := x \div 2];
use ESEEven[x := 2 + 2 * (x \div 2), q := 1 + x \div 2];
prove;
next;
\end{zproof}

\begin{theorem}{ZeroExprSK}
  comp(comp(funK, funS), opInv~funK) = cst~0
\end{theorem}

\begin{zproof}
apply UFExt to predicate comp(comp(funK, funS), opInv~funK) = cst~0;
prove;
use KinvZ;
use ASuc[x := x];
use UFNonneg[f := funA, x := x];
use KAt[m := 1 + funA~x + x, t := 1 + x, k := 0];
prove;
\end{zproof}

\begin{theorem}{LSSurj}
  comp(funL, funS) \in Surj
\end{theorem}

\begin{zproof}
invoke Surj;
prove;
instantiate y == funA~(1 + x);
use UFNonneg[f := funA, x := 1 + x];
use LAt[m := 1 + funA~(1 + x), t := 1 + x, k := 1];
prove;
\end{zproof}

\begin{theorem}{LSLow}
  \forall x, y, t, k: \nat | 1 + y = funA~t + k \land k \leq t \land t - k = 1 + x @ funA~x + x \leq y
\end{theorem}

\begin{zproof}
prove;
use AMono[m := 1 + x, n := t];
use ASuc[x := x];
prove;
\end{zproof}

\begin{theorem}{LSInvSuc}
  \forall x: \nat @ opInv~(comp(funL, funS))~(1 + x) = funA~x + x
\end{theorem}

\begin{zproof}
prove;
use LSSurj;
use InvRight[F := comp(funL, funS), x := 1 + x];
use UFNonneg[f := opInv~(comp(funL, funS)), x := 1 + x];
use Decomp[x := 1 + opInv~(comp(funL, funS))~(1 + x)];
use LAt[m := 1 + opInv~(comp(funL, funS))~(1 + x), t := funV~(1 + opInv~(comp(funL, funS))~(1 + x)), k := funK~(1 + opInv~(comp(funL, funS))~(1 + x))];
use LSLow[x := x, y := opInv~(comp(funL, funS))~(1 + x), t := funV~(1 + opInv~(comp(funL, funS))~(1 + x)), k := funK~(1 + opInv~(comp(funL, funS))~(1 + x))];
use ASuc[x := x];
use UFNonneg[f := funA, x := x];
use LAt[m := 1 + funA~x + x, t := 1 + x, k := 0];
use InvLeast[F := comp(funL, funS), x := 1 + x, y := funA~x + x];
with disabled (InvRight) prove;
\end{zproof}

\begin{theorem}{ZeroExprSL}
  comp(comp(funL, opInv~(comp(funL, funS))), funS) = cst~0
\end{theorem}

\begin{zproof}
apply UFExt to predicate comp(comp(funL, opInv~(comp(funL, funS))), funS) = cst~0;
prove;
use LSInvSuc[x := x];
use UFNonneg[f := funA, x := x];
use LAt[m := funA~x + x, t := x, k := x];
prove;
\end{zproof}

\begin{zed}
  CGSE == clos~\{ funS, funE \}~\{ opPlusI, opInv \}~\{ comp \}
\end{zed}

\begin{zproof}[CGSE\$domainCheck]
prove;
\end{zproof}

\begin[disabled]{theorem}{rule CGSEDef}
  CGSE = clos~\{ funS, funE \}~\{ opPlusI, opInv \}~\{ comp \}
\end{theorem}

\begin{zproof}
invoke CGSE;
prove;
\end{zproof}

\begin{theorem}{CGSEInPower}
  CGSE \in \power (\nat \fun \nat)
\end{theorem}

\begin{zproof}
use ClosInPow[I := \{ funS, funE \}, U := \{ opPlusI, opInv \}, B := \{ comp \}];
with enabled (CGSEDef) prove;
\end{zproof}

\begin[disabled]{theorem}{rule CGSEUF}
  \forall F: \power (\num \cross \num) | F \in CGSE @ F \in \nat \fun \nat
\end{theorem}

\begin{zproof}
prove;
use ClosElemUF[I := \{ funS, funE \}, U := \{ opPlusI, opInv \}, B := \{ comp \}, F := F];
with enabled (CGSEDef) prove;
\end{zproof}

\begin{theorem}{rule CGSEInit0}
  funS \in CGSE
\end{theorem}

\begin{zproof}
with enabled (CGSEDef) prove;
\end{zproof}

\begin{theorem}{rule CGSEInit1}
  funE \in CGSE
\end{theorem}

\begin{zproof}
with enabled (CGSEDef) prove;
\end{zproof}

\begin{theorem}{rule CGSEUnop0}
  \forall F: \nat \fun \nat | F \in CGSE @ opPlusI~F \in CGSE
\end{theorem}

\begin{zproof}
with enabled (CGSEDef) prove;
\end{zproof}

\begin{theorem}{rule CGSEUnop1}
  \forall F: \nat \fun \nat | F \in CGSE @ opInv~F \in CGSE
\end{theorem}

\begin{zproof}
with enabled (CGSEDef) prove;
\end{zproof}

\begin{theorem}{rule CGSEBinop0}
  \forall F, G: \nat \fun \nat | F \in CGSE \land G \in CGSE @ comp(F, G) \in CGSE
\end{theorem}

\begin{zproof}
with enabled (CGSEDef) prove;
\end{zproof}

\begin{theorem}{CGSEUClosedG}
  CGSE \in UClosed~\{ opPlusI, opInv \}
\end{theorem}

\begin{zproof}
use CGSEInPower;
with enabled (UClosedDef, CGSEUF) prove;
split u = opPlusI;
cases;
prove;
next;
prove;
next;
\end{zproof}

\begin{theorem}{CGSEBClosedG}
  CGSE \in BClosed~\{ comp \}
\end{theorem}

\begin{zproof}
use CGSEInPower;
with enabled (BClosedDef, CGSEUF) prove;
\end{zproof}

\begin{theorem}{GSEZeroExpr}
  comp(comp(comp(funE, funS), opInv~funE), comp(comp(funE, funS), opInv~funE)) \in CGSE
\end{theorem}

\begin{zproof}
prove;
\end{zproof}

\begin{theorem}{rule GSEZero}
  cst~0 \in CGSE
\end{theorem}

\begin{zproof}
use GSEZeroExpr;
use ZeroExprSE;
prove;
\end{zproof}

\begin{theorem}{GSEOneExpr}
  comp(funS, cst~0) \in CGSE
\end{theorem}

\begin{zproof}
prove;
\end{zproof}

\begin{theorem}{rule GSEOne}
  cst~1 \in CGSE
\end{theorem}

\begin{zproof}
use GSEOneExpr;
use OneExpr;
prove;
\end{zproof}

\begin{theorem}{GRSESub1}
  CGSE \subseteq CGE
\end{theorem}

\begin{zproof}
use ClosLeast[I := \{ funS, funE \}, U := \{ opPlusI, opInv \}, B := \{ comp \}, Z := CGE];
use CGEInPower;
use CGEUClosedG;
use CGEBClosedG;
with enabled (CGSEDef) prove;
\end{zproof}

\begin{theorem}{GRSESub2}
  CGE \subseteq CGSE
\end{theorem}

\begin{zproof}
use ClosLeast[I := \{ cst~1, funE \}, U := \{ opPlusI, opInv \}, B := \{ comp \}, Z := CGSE];
use CGSEInPower;
use CGSEUClosedG;
use CGSEBClosedG;
with enabled (CGEDef) prove;
\end{zproof}

\begin{theorem}{GRSE}
  CGSE = CGE
\end{theorem}

\begin{zproof}
use GRSESub1;
use GRSESub2;
use CGSEInPower;
use CGEInPower;
use SubAntisym[\nat \fun \nat][A := CGSE, B := CGE];
prove;
\end{zproof}

\begin{zed}
  CGSK == clos~\{ funS, funK \}~\{ opPlusI, opInv \}~\{ comp \}
\end{zed}

\begin{zproof}[CGSK\$domainCheck]
prove;
\end{zproof}

\begin[disabled]{theorem}{rule CGSKDef}
  CGSK = clos~\{ funS, funK \}~\{ opPlusI, opInv \}~\{ comp \}
\end{theorem}

\begin{zproof}
invoke CGSK;
prove;
\end{zproof}

\begin{theorem}{CGSKInPower}
  CGSK \in \power (\nat \fun \nat)
\end{theorem}

\begin{zproof}
use ClosInPow[I := \{ funS, funK \}, U := \{ opPlusI, opInv \}, B := \{ comp \}];
with enabled (CGSKDef) prove;
\end{zproof}

\begin[disabled]{theorem}{rule CGSKUF}
  \forall F: \power (\num \cross \num) | F \in CGSK @ F \in \nat \fun \nat
\end{theorem}

\begin{zproof}
prove;
use ClosElemUF[I := \{ funS, funK \}, U := \{ opPlusI, opInv \}, B := \{ comp \}, F := F];
with enabled (CGSKDef) prove;
\end{zproof}

\begin{theorem}{rule CGSKInit0}
  funS \in CGSK
\end{theorem}

\begin{zproof}
with enabled (CGSKDef) prove;
\end{zproof}

\begin{theorem}{rule CGSKInit1}
  funK \in CGSK
\end{theorem}

\begin{zproof}
with enabled (CGSKDef) prove;
\end{zproof}

\begin{theorem}{rule CGSKUnop0}
  \forall F: \nat \fun \nat | F \in CGSK @ opPlusI~F \in CGSK
\end{theorem}

\begin{zproof}
with enabled (CGSKDef) prove;
\end{zproof}

\begin{theorem}{rule CGSKUnop1}
  \forall F: \nat \fun \nat | F \in CGSK @ opInv~F \in CGSK
\end{theorem}

\begin{zproof}
with enabled (CGSKDef) prove;
\end{zproof}

\begin{theorem}{rule CGSKBinop0}
  \forall F, G: \nat \fun \nat | F \in CGSK \land G \in CGSK @ comp(F, G) \in CGSK
\end{theorem}

\begin{zproof}
with enabled (CGSKDef) prove;
\end{zproof}

\begin{theorem}{CGSKUClosedG}
  CGSK \in UClosed~\{ opPlusI, opInv \}
\end{theorem}

\begin{zproof}
use CGSKInPower;
with enabled (UClosedDef, CGSKUF) prove;
split u = opPlusI;
cases;
prove;
next;
prove;
next;
\end{zproof}

\begin{theorem}{CGSKBClosedG}
  CGSK \in BClosed~\{ comp \}
\end{theorem}

\begin{zproof}
use CGSKInPower;
with enabled (BClosedDef, CGSKUF) prove;
\end{zproof}

\begin{theorem}{GSKZeroExpr}
  comp(comp(funK, funS), opInv~funK) \in CGSK
\end{theorem}

\begin{zproof}
prove;
\end{zproof}

\begin{theorem}{rule GSKZero}
  cst~0 \in CGSK
\end{theorem}

\begin{zproof}
use GSKZeroExpr;
use ZeroExprSK;
prove;
\end{zproof}

\begin{theorem}{GSKOneExpr}
  comp(funS, cst~0) \in CGSK
\end{theorem}

\begin{zproof}
prove;
\end{zproof}

\begin{theorem}{rule GSKOne}
  cst~1 \in CGSK
\end{theorem}

\begin{zproof}
use GSKOneExpr;
use OneExpr;
prove;
\end{zproof}

\begin{theorem}{GRSKSub1}
  CGSK \subseteq CGK
\end{theorem}

\begin{zproof}
use ClosLeast[I := \{ funS, funK \}, U := \{ opPlusI, opInv \}, B := \{ comp \}, Z := CGK];
use CGKInPower;
use CGKUClosedG;
use CGKBClosedG;
with enabled (CGSKDef) prove;
\end{zproof}

\begin{theorem}{GRSKSub2}
  CGK \subseteq CGSK
\end{theorem}

\begin{zproof}
use ClosLeast[I := \{ cst~1, funK \}, U := \{ opPlusI, opInv \}, B := \{ comp \}, Z := CGSK];
use CGSKInPower;
use CGSKUClosedG;
use CGSKBClosedG;
with enabled (CGKDef) prove;
\end{zproof}

\begin{theorem}{GRSK}
  CGSK = CGK
\end{theorem}

\begin{zproof}
use GRSKSub1;
use GRSKSub2;
use CGSKInPower;
use CGKInPower;
use SubAntisym[\nat \fun \nat][A := CGSK, B := CGK];
prove;
\end{zproof}

\begin{zed}
  CGSL == clos~\{ funS, funL \}~\{ opPlusI, opInv \}~\{ comp \}
\end{zed}

\begin{zproof}[CGSL\$domainCheck]
prove;
\end{zproof}

\begin[disabled]{theorem}{rule CGSLDef}
  CGSL = clos~\{ funS, funL \}~\{ opPlusI, opInv \}~\{ comp \}
\end{theorem}

\begin{zproof}
invoke CGSL;
prove;
\end{zproof}

\begin{theorem}{CGSLInPower}
  CGSL \in \power (\nat \fun \nat)
\end{theorem}

\begin{zproof}
use ClosInPow[I := \{ funS, funL \}, U := \{ opPlusI, opInv \}, B := \{ comp \}];
with enabled (CGSLDef) prove;
\end{zproof}

\begin[disabled]{theorem}{rule CGSLUF}
  \forall F: \power (\num \cross \num) | F \in CGSL @ F \in \nat \fun \nat
\end{theorem}

\begin{zproof}
prove;
use ClosElemUF[I := \{ funS, funL \}, U := \{ opPlusI, opInv \}, B := \{ comp \}, F := F];
with enabled (CGSLDef) prove;
\end{zproof}

\begin{theorem}{rule CGSLInit0}
  funS \in CGSL
\end{theorem}

\begin{zproof}
with enabled (CGSLDef) prove;
\end{zproof}

\begin{theorem}{rule CGSLInit1}
  funL \in CGSL
\end{theorem}

\begin{zproof}
with enabled (CGSLDef) prove;
\end{zproof}

\begin{theorem}{rule CGSLUnop0}
  \forall F: \nat \fun \nat | F \in CGSL @ opPlusI~F \in CGSL
\end{theorem}

\begin{zproof}
with enabled (CGSLDef) prove;
\end{zproof}

\begin{theorem}{rule CGSLUnop1}
  \forall F: \nat \fun \nat | F \in CGSL @ opInv~F \in CGSL
\end{theorem}

\begin{zproof}
with enabled (CGSLDef) prove;
\end{zproof}

\begin{theorem}{rule CGSLBinop0}
  \forall F, G: \nat \fun \nat | F \in CGSL \land G \in CGSL @ comp(F, G) \in CGSL
\end{theorem}

\begin{zproof}
with enabled (CGSLDef) prove;
\end{zproof}

\begin{theorem}{CGSLUClosedG}
  CGSL \in UClosed~\{ opPlusI, opInv \}
\end{theorem}

\begin{zproof}
use CGSLInPower;
with enabled (UClosedDef, CGSLUF) prove;
split u = opPlusI;
cases;
prove;
next;
prove;
next;
\end{zproof}

\begin{theorem}{CGSLBClosedG}
  CGSL \in BClosed~\{ comp \}
\end{theorem}

\begin{zproof}
use CGSLInPower;
with enabled (BClosedDef, CGSLUF) prove;
\end{zproof}

\begin{theorem}{GSLZeroExpr}
  comp(comp(funL, opInv~(comp(funL, funS))), funS) \in CGSL
\end{theorem}

\begin{zproof}
prove;
\end{zproof}

\begin{theorem}{rule GSLZero}
  cst~0 \in CGSL
\end{theorem}

\begin{zproof}
use GSLZeroExpr;
use ZeroExprSL;
prove;
\end{zproof}

\begin{theorem}{GSLOneExpr}
  comp(funS, cst~0) \in CGSL
\end{theorem}

\begin{zproof}
prove;
\end{zproof}

\begin{theorem}{rule GSLOne}
  cst~1 \in CGSL
\end{theorem}

\begin{zproof}
use GSLOneExpr;
use OneExpr;
prove;
\end{zproof}

\begin{theorem}{GLSExpr}
  opPlusI~(cst~1) \in CGL
\end{theorem}

\begin{zproof}
prove;
\end{zproof}

\begin{theorem}{rule GLS}
  funS \in CGL
\end{theorem}

\begin{zproof}
use GLSExpr;
use SExpr;
prove;
\end{zproof}

\begin{theorem}{GRSLSub1}
  CGSL \subseteq CGL
\end{theorem}

\begin{zproof}
use ClosLeast[I := \{ funS, funL \}, U := \{ opPlusI, opInv \}, B := \{ comp \}, Z := CGL];
use CGLInPower;
use CGLUClosedG;
use CGLBClosedG;
with enabled (CGSLDef) prove;
\end{zproof}

\begin{theorem}{GRSLSub2}
  CGL \subseteq CGSL
\end{theorem}

\begin{zproof}
use ClosLeast[I := \{ cst~1, funL \}, U := \{ opPlusI, opInv \}, B := \{ comp \}, Z := CGSL];
use CGSLInPower;
use CGSLUClosedG;
use CGSLBClosedG;
with enabled (CGLDef) prove;
\end{zproof}

\begin{theorem}{GRSL}
  CGSL = CGL
\end{theorem}

\begin{zproof}
use GRSLSub1;
use GRSLSub2;
use CGSLInPower;
use CGLInPower;
use SubAntisym[\nat \fun \nat][A := CGSL, B := CGL];
prove;
\end{zproof}


\begin{theorem}{GRSEq}
  CGSE = ClosJR \land CGSK = ClosJR \land CGSL = ClosJR
\end{theorem}

\begin{zproof}
use GRSE;
use GRSK;
use GRSL;
use GREqJR;
prove;
\end{zproof}

\subsection{The shifted operator with inversion}

$\langle X, F^{+b}, FG, F^{-1}\rangle$ generates Robinson's closure iff $b \geq 1$
($X \in \{E, K, L\}$).  A successor $S$ is generated first; a predecessor $P$ then gives
$F^+ = P^b F^{+b}$ and $\bar 1$.

\begin{theorem}{rule CbX}
  \forall X: \nat \fun \nat; b: \nat @ X \in clos~\{ X \}~\{ opPlusIb~b, opInv \}~\{ comp \}
\end{theorem}

\begin{zproof}
prove;
\end{zproof}

\begin{theorem}{CbComp}
  \forall X: \nat \fun \nat; b: \nat; F, G: \power (\num \cross \num) | F \in clos~\{ X \}~\{ opPlusIb~b, opInv \}~\{ comp \} \land G \in clos~\{ X \}~\{ opPlusIb~b, opInv \}~\{ comp \} @ comp(F, G) \in clos~\{ X \}~\{ opPlusIb~b, opInv \}~\{ comp \}
\end{theorem}

\begin{zproof}
prove;
use ClosElemUF[I := \{ X \}, U := \{ opPlusIb~b, opInv \}, B := \{ comp \}, F := F];
use ClosElemUF[I := \{ X \}, U := \{ opPlusIb~b, opInv \}, B := \{ comp \}, F := G];
prove;
\end{zproof}

\begin{theorem}{CbUn}
  \forall X: \nat \fun \nat; b: \nat; F: \power (\num \cross \num) | F \in clos~\{ X \}~\{ opPlusIb~b, opInv \}~\{ comp \} @ opPlusIb~b~F \in clos~\{ X \}~\{ opPlusIb~b, opInv \}~\{ comp \} \land opInv~F \in clos~\{ X \}~\{ opPlusIb~b, opInv \}~\{ comp \}
\end{theorem}

\begin{zproof}
prove;
use ClosElemUF[I := \{ X \}, U := \{ opPlusIb~b, opInv \}, B := \{ comp \}, F := F];
prove;
\end{zproof}

\begin{theorem}{CbIter}
  \forall X: \nat \fun \nat; b, k: \nat; F, G: \nat \fun \nat | F \in clos~\{ X \}~\{ opPlusIb~b, opInv \}~\{ comp \} \land G \in clos~\{ X \}~\{ opPlusIb~b, opInv \}~\{ comp \} @ comp(iter[\nat]~k~F, G) \in clos~\{ X \}~\{ opPlusIb~b, opInv \}~\{ comp \}
\end{theorem}

\begin{zproof}
prove;
use PowIn[I := \{ X \}, U := \{ opPlusIb~b, opInv \}, B := \{ comp \}, F := F, G := G, k := k];
prove;
\end{zproof}

$F^+ = P^b F^{+b}$ for a predecessor $P$.

\begin{theorem}{PlusIFromPred}
  \forall b: \nat; P, F: \nat \fun \nat | P \in Preds @ comp(iter[\nat]~b~P, opPlusIb~b~F) = opPlusI~F
\end{theorem}

\begin{zproof}
apply UFExt to predicate comp(iter[\nat]~b~P, opPlusIb~b~F) = opPlusI~F;
prove;
use IterPAll[k := b, P := P, x := F~x + x];
prove;
\end{zproof}

\begin{theorem}{CbPlusI}
  \forall X: \nat \fun \nat; b: \nat; P, F: \nat \fun \nat | P \in Preds \land P \in clos~\{ X \}~\{ opPlusIb~b, opInv \}~\{ comp \} \land F \in clos~\{ X \}~\{ opPlusIb~b, opInv \}~\{ comp \} @ opPlusI~F \in clos~\{ X \}~\{ opPlusIb~b, opInv \}~\{ comp \}
\end{theorem}

\begin{zproof}
prove;
use CbUn[X := X, b := b, F := F];
use CbIter[X := X, b := b, k := b, F := P, G := opPlusIb~b~F];
use PlusIFromPred[b := b, P := P, F := F];
prove;
\end{zproof}

With $\bar 1$ and a predecessor, the closure contains $\langle \bar 1, X, F^+, FG, F^{-1}\rangle$.

\begin{theorem}{CbUClosed}
  \forall X: \nat \fun \nat; b: \nat; P: \nat \fun \nat | P \in Preds \land P \in clos~\{ X \}~\{ opPlusIb~b, opInv \}~\{ comp \} @ clos~\{ X \}~\{ opPlusIb~b, opInv \}~\{ comp \} \in UClosed~\{ opPlusI, opInv \}
\end{theorem}

\begin{zproof}
prove;
use ClosInPow[I := \{ X \}, U := \{ opPlusIb~b, opInv \}, B := \{ comp \}];
with enabled (UClosedDef) prove;
use ClosElemUF[I := \{ X \}, U := \{ opPlusIb~b, opInv \}, B := \{ comp \}, F := F];
prove;
split u = opPlusI;
cases;
use CbPlusI[X := X, b := b, P := P, F := F];
prove;
next;
use CbUn[X := X, b := b, F := F];
prove;
next;
\end{zproof}

\begin{theorem}{CbBClosed}
  \forall X: \nat \fun \nat; b: \nat @ clos~\{ X \}~\{ opPlusIb~b, opInv \}~\{ comp \} \in BClosed~\{ comp \}
\end{theorem}

\begin{zproof}
prove;
use ClosInPow[I := \{ X \}, U := \{ opPlusIb~b, opInv \}, B := \{ comp \}];
with enabled (BClosedDef) prove;
\end{zproof}

\begin{theorem}{CbSub}
  \forall X: \nat \fun \nat; b: \nat; P: \nat \fun \nat | cst~1 \in clos~\{ X \}~\{ opPlusIb~b, opInv \}~\{ comp \} \land P \in Preds \land P \in clos~\{ X \}~\{ opPlusIb~b, opInv \}~\{ comp \} @ clos~\{ cst~1, X \}~\{ opPlusI, opInv \}~\{ comp \} \subseteq clos~\{ X \}~\{ opPlusIb~b, opInv \}~\{ comp \}
\end{theorem}

\begin{zproof}
prove;
use ClosLeast[I := \{ cst~1, X \}, U := \{ opPlusI, opInv \}, B := \{ comp \}, Z := clos~\{ X \}~\{ opPlusIb~b, opInv \}~\{ comp \}];
use ClosInPow[I := \{ X \}, U := \{ opPlusIb~b, opInv \}, B := \{ comp \}];
use CbUClosed[X := X, b := b, P := P];
use CbBClosed[X := X, b := b];
prove;
\end{zproof}

Soundness: $F^{+b} = (S^b F)^+$.

\begin{theorem}{GEPlusIb}
  \forall b: \nat; F: \nat \fun \nat | F \in CGE @ opPlusIb~b~F \in CGE
\end{theorem}

\begin{zproof}
prove;
use PlusIbExpr[b := b, F := F];
use GES;
use PowIn[I := \{ cst~1, funE \}, U := \{ opPlusI, opInv \}, B := \{ comp \}, F := funS, G := F, k := b];
with enabled (CGEDef) prove;
\end{zproof}

\begin{theorem}{GEUClosedb}
  \forall b: \nat @ CGE \in UClosed~\{ opPlusIb~b, opInv \}
\end{theorem}

\begin{zproof}
prove;
use CGEInPower;
with enabled (UClosedDef, CGEUF) prove;
use InPair[\power (\power (\num \cross \num) \cross \power (\num \cross \num))][u := u, p := opPlusIb~b, q := opInv];
split u = opInv;
cases;
prove;
next;
use GEPlusIb[b := b, F := F];
prove;
next;
\end{zproof}

\begin{theorem}{GRbSubGE}
  \forall X: \nat \fun \nat; b: \nat | X \in CGE @ clos~\{ X \}~\{ opPlusIb~b, opInv \}~\{ comp \} \subseteq CGE
\end{theorem}

\begin{zproof}
prove;
use ClosLeast[I := \{ X \}, U := \{ opPlusIb~b, opInv \}, B := \{ comp \}, Z := CGE];
use CGEInPower;
use CGEBClosedG;
use GEUClosedb[b := b];
prove;
\end{zproof}

\begin{theorem}{ShiftVal}
  \forall W: \nat \fun \nat; c, y: \nat | (W, c) \in Shift @ W~y = y + c
\end{theorem}

\begin{zproof}
invoke Shift;
prove;
instantiate y\_\_0 == y;
prove;
\end{zproof}

Subtraction by inversion with $E$: if $V(y) = y + c$ then $E\,(E V)^{-1}$ subtracts $c$.

\begin{theorem}{ESubSurj}
  \forall W: \nat \fun \nat; c: \nat | (W, c) \in Shift @ comp(funE, W) \in Surj
\end{theorem}

\begin{zproof}
invoke Surj;
prove;
instantiate y == (x + c) * (x + c) + x - c;
use LeSq[n := x + c];
use ShiftVal[W := W, c := c, y := (x + c) * (x + c) + x - c];
use EAt[y := (x + c) * (x + c) + x, m := x + c, k := x];
prove;
\end{zproof}

\begin{theorem}{ESubCore}
  \forall y, c, x: \nat | funE~(y + c) = x \land c \leq x @ funE~y = x - c
\end{theorem}

\begin{zproof}
prove;
use EDecomp[y := y + c];
use EBound[y := y + c];
use EAt[y := y, m := funRt~(y + c), k := x - c];
prove;
\end{zproof}

\begin{theorem}{ESubVal}
  \forall W: \nat \fun \nat; c, x: \nat | (W, c) \in Shift \land c \leq x @ funE~(opInv~(comp(funE, W))~x) = x - c
\end{theorem}

\begin{zproof}
prove;
use ESubSurj[W := W, c := c];
use InvRight[F := comp(funE, W), x := x];
use UFNonneg[f := opInv~(comp(funE, W)), x := x];
use ShiftVal[W := W, c := c, y := opInv~(comp(funE, W))~x];
use ESubCore[y := opInv~(comp(funE, W))~x, c := c, x := x];
with disabled (InvRight) prove;
\end{zproof}


\subsection{$\langle E, F^{+b}, FG, F^{-1}\rangle$}

$E(E^{-1}(2m) + x + c + 1) = x + c$ and $E(E^{-1}(2m + 1) + x + c + 2) = x + c$.

\begin{theorem}{EShiftE}
  \forall m, x, c: \nat | x + c \leq 2 * m + 2 @ funE~(opInv~funE~(2 * m) + x + c + 1) = x + c
\end{theorem}

\begin{zproof}
prove;
use EAt[y := m * m + 2 * m + x + c + 1, m := m + 1, k := x + c];
prove;
\end{zproof}

\begin{theorem}{EShiftO}
  \forall m, x, c: \nat | x + c \leq 2 * m + 4 @ funE~(opInv~funE~(1 + 2 * m) + x + c + 2) = x + c
\end{theorem}

\begin{zproof}
prove;
use EAt[y := (m + 1) * (m + 1) + 2 * m + 1 + x + c + 2, m := m + 2, k := x + c];
prove;
\end{zproof}

For $b \geq 2$, $W_1 = E\,(E^{-1} T)^{+b}$ and $W_2 = E\,(E^{-1} W_1 T)^{+b}$ ($T = I^{+b}$) are the
shifts $x + b - 1$ and $x + b - 2$, in an order given by the parity of $b$.

\begin{theorem}{W1EvenVal}
  \forall p, x: \nat @ comp(funE, opPlusIb~(2 + 2 * p)~(comp(opInv~funE, opPlusIb~(2 + 2 * p)~funI)))~x = x + 2 * p + 1
\end{theorem}

\begin{zproof}
prove;
use EShiftE[m := x + p + 1, x := x, c := 2 * p + 1];
prove;
\end{zproof}

\begin{theorem}{W2EvenVal}
  \forall p, x: \nat @ comp(funE, opPlusIb~(2 + 2 * p)~(comp(comp(opInv~funE, comp(funE, opPlusIb~(2 + 2 * p)~(comp(opInv~funE, opPlusIb~(2 + 2 * p)~funI)))), opPlusIb~(2 + 2 * p)~funI)))~x = x + 2 * p
\end{theorem}

\begin{zproof}
prove;
use W1EvenVal[p := p, x := 2 + 2 * p + 2 * x];
use EShiftO[m := x + 2 * p + 1, x := x, c := 2 * p];
prove;
\end{zproof}

\begin{theorem}{W1OddVal}
  \forall p, x: \nat @ comp(funE, opPlusIb~(3 + 2 * p)~(comp(opInv~funE, opPlusIb~(3 + 2 * p)~funI)))~x = x + 2 * p + 1
\end{theorem}

\begin{zproof}
prove;
use EShiftO[m := x + p + 1, x := x, c := 2 * p + 1];
prove;
\end{zproof}

\begin{theorem}{W2OddVal}
  \forall p, x: \nat @ comp(funE, opPlusIb~(3 + 2 * p)~(comp(comp(opInv~funE, comp(funE, opPlusIb~(3 + 2 * p)~(comp(opInv~funE, opPlusIb~(3 + 2 * p)~funI)))), opPlusIb~(3 + 2 * p)~funI)))~x = x + 2 * p + 2
\end{theorem}

\begin{zproof}
prove;
use W1OddVal[p := p, x := 3 + 2 * p + 2 * x];
use EShiftE[m := x + 2 * p + 2, x := x, c := 2 * p + 2];
prove;
\end{zproof}

\begin{theorem}{CbEIn}
  \forall b: \nat @ funI \in clos~\{ funE \}~\{ opPlusIb~b, opInv \}~\{ comp \} \land comp(funE, opPlusIb~b~(comp(opInv~funE, opPlusIb~b~funI))) \in clos~\{ funE \}~\{ opPlusIb~b, opInv \}~\{ comp \} \land comp(funE, opPlusIb~b~(comp(comp(opInv~funE, comp(funE, opPlusIb~b~(comp(opInv~funE, opPlusIb~b~funI)))), opPlusIb~b~funI))) \in clos~\{ funE \}~\{ opPlusIb~b, opInv \}~\{ comp \}
\end{theorem}

\begin{zproof}
prove;
use CbUn[X := funE, b := b, F := funE];
use CbComp[X := funE, b := b, F := funE, G := opInv~funE];
use EEinv;
prove;
\end{zproof}

$S = E\,(E\,W_2)^{-1} W_1$ (or with $W_1$ and $W_2$ exchanged).

\begin{theorem}{SFromShifts}
  \forall G, H: \nat \fun \nat; d: \nat | (G, d + 1) \in Shift \land (H, d) \in Shift @ comp(comp(funE, opInv~(comp(funE, H))), G) = funS
\end{theorem}

\begin{zproof}
apply UFExt to predicate comp(comp(funE, opInv~(comp(funE, H))), G) = funS;
prove;
use ShiftVal[W := G, c := d + 1, y := x];
use ESubVal[W := H, c := d, x := 1 + x + d];
prove;
\end{zproof}

\begin{theorem}{SInCbFromShifts}
  \forall X: \nat \fun \nat; b, d: \nat; G, H: \nat \fun \nat | funE \in clos~\{ X \}~\{ opPlusIb~b, opInv \}~\{ comp \} \land G \in clos~\{ X \}~\{ opPlusIb~b, opInv \}~\{ comp \} \land H \in clos~\{ X \}~\{ opPlusIb~b, opInv \}~\{ comp \} \land (G, d + 1) \in Shift \land (H, d) \in Shift @ funS \in clos~\{ X \}~\{ opPlusIb~b, opInv \}~\{ comp \}
\end{theorem}

\begin{zproof}
use SFromShifts[G := G, H := H, d := d];
prove;
equality substitute funS;
use CbComp[X := X, b := b, F := funE, G := H];
use CbUn[X := X, b := b, F := comp(funE, H)];
use CbComp[X := X, b := b, F := funE, G := opInv~(comp(funE, H))];
use CbComp[X := X, b := b, F := comp(funE, opInv~(comp(funE, H))), G := G];
prove;
\end{zproof}

\begin{theorem}{W12Shifts}
  \forall p: \nat @ (comp(funE, opPlusIb~(2 + 2 * p)~(comp(opInv~funE, opPlusIb~(2 + 2 * p)~funI))), 2 * p + 1) \in Shift \land (comp(funE, opPlusIb~(2 + 2 * p)~(comp(comp(opInv~funE, comp(funE, opPlusIb~(2 + 2 * p)~(comp(opInv~funE, opPlusIb~(2 + 2 * p)~funI)))), opPlusIb~(2 + 2 * p)~funI))), 2 * p) \in Shift \land (comp(funE, opPlusIb~(3 + 2 * p)~(comp(comp(opInv~funE, comp(funE, opPlusIb~(3 + 2 * p)~(comp(opInv~funE, opPlusIb~(3 + 2 * p)~funI)))), opPlusIb~(3 + 2 * p)~funI))), 2 * p + 2) \in Shift \land (comp(funE, opPlusIb~(3 + 2 * p)~(comp(opInv~funE, opPlusIb~(3 + 2 * p)~funI))), 2 * p + 1) \in Shift
\end{theorem}

\begin{zproof}
invoke Shift;
prove;
use W1EvenVal[p := p, x := y];
use W2EvenVal[p := p, x := y\_\_0];
use W2OddVal[p := p, x := y\_\_1];
use W1OddVal[p := p, x := y\_\_2];
prove;
\end{zproof}

\begin{theorem}{CbESEven}
  \forall p: \nat @ funS \in clos~\{ funE \}~\{ opPlusIb~(2 + 2 * p), opInv \}~\{ comp \}
\end{theorem}

\begin{zproof}
prove;
use CbEIn[b := 2 + 2 * p];
use W12Shifts[p := p];
use SInCbFromShifts[X := funE, b := 2 + 2 * p, d := 2 * p, G := comp(funE, opPlusIb~(2 + 2 * p)~(comp(opInv~funE, opPlusIb~(2 + 2 * p)~funI))), H := comp(funE, opPlusIb~(2 + 2 * p)~(comp(comp(opInv~funE, comp(funE, opPlusIb~(2 + 2 * p)~(comp(opInv~funE, opPlusIb~(2 + 2 * p)~funI)))), opPlusIb~(2 + 2 * p)~funI)))];
prove;
\end{zproof}

\begin{theorem}{CbESOdd}
  \forall p: \nat @ funS \in clos~\{ funE \}~\{ opPlusIb~(3 + 2 * p), opInv \}~\{ comp \}
\end{theorem}

\begin{zproof}
prove;
use W12Shifts[p := p];
use SInCbFromShifts[X := funE, b := 3 + 2 * p, d := 2 * p + 1, G := comp(funE, opPlusIb~(3 + 2 * p)~(comp(comp(opInv~funE, comp(funE, opPlusIb~(3 + 2 * p)~(comp(opInv~funE, opPlusIb~(3 + 2 * p)~funI)))), opPlusIb~(3 + 2 * p)~funI))), H := comp(funE, opPlusIb~(3 + 2 * p)~(comp(opInv~funE, opPlusIb~(3 + 2 * p)~funI)))];
prove;
use CbEIn[b := 3 + 2 * p];
with disabled (ClosBinop, ClosUnop, ClosInit) prove;
\end{zproof}


For $b = 1$: $\hat P = E\,(E^{-1} T)^{+1}$ ($\hat P(y + 1) = y$, $\hat P(0) = 2$) and
$S = E\,\hat P^{-1} (E^{-1} \hat P T)^{+1}$.

\begin{theorem}{PE1Suc}
  \forall y: \nat @ comp(funE, opPlusIb~1~(comp(opInv~funE, opPlusIb~1~funI)))~(1 + y) = y
\end{theorem}

\begin{zproof}
prove;
use EShiftO[m := 1 + y, x := y, c := 0];
prove;
\end{zproof}

\begin{theorem}{PE1Zero}
  comp(funE, opPlusIb~1~(comp(opInv~funE, opPlusIb~1~funI)))~0 = 2
\end{theorem}

\begin{zproof}
prove;
use EinvOdd[m := 0];
use EAt[y := 3, m := 1, k := 2];
prove;
\end{zproof}

\begin{theorem}{PE1Core}
  \forall v, y: \nat | comp(funE, opPlusIb~1~(comp(opInv~funE, opPlusIb~1~funI)))~y = v \land \lnot v = 2 @ y = 1 + v
\end{theorem}

\begin{zproof}
prove;
use PE1Zero;
use PE1Suc[y := y - 1];
split y = 0;
cases;
prove;
next;
prove;
next;
\end{zproof}

\begin{theorem}{PE1Surj}
  comp(funE, opPlusIb~1~(comp(opInv~funE, opPlusIb~1~funI))) \in Surj
\end{theorem}

\begin{zproof}
invoke Surj;
with disabled (EinvOdd, EinvD) prove;
use PE1Suc[y := x];
instantiate y == 1 + x;
with disabled (EinvOdd, EinvD) prove;
\end{zproof}

\begin{theorem}{PE1Inv}
  \forall v: \nat | \lnot v = 2 @ opInv~(comp(funE, opPlusIb~1~(comp(opInv~funE, opPlusIb~1~funI))))~v = 1 + v
\end{theorem}

\begin{zproof}
prove;
use PE1Surj;
use InvRight[F := comp(funE, opPlusIb~1~(comp(opInv~funE, opPlusIb~1~funI))), x := v];
use UFNonneg[f := opInv~(comp(funE, opPlusIb~1~(comp(opInv~funE, opPlusIb~1~funI)))), x := v];
use PE1Core[v := v, y := opInv~(comp(funE, opPlusIb~1~(comp(opInv~funE, opPlusIb~1~funI))))~v];
with disabled (InvRight) prove;
\end{zproof}

\begin{theorem}{SE1Val}
  comp(comp(funE, opInv~(comp(funE, opPlusIb~1~(comp(opInv~funE, opPlusIb~1~funI))))), opPlusIb~1~(comp(comp(opInv~funE, comp(funE, opPlusIb~1~(comp(opInv~funE, opPlusIb~1~funI)))), opPlusIb~1~funI))) = funS
\end{theorem}

\begin{zproof}
apply UFExt to predicate comp(comp(funE, opInv~(comp(funE, opPlusIb~1~(comp(opInv~funE, opPlusIb~1~funI))))), opPlusIb~1~(comp(comp(opInv~funE, comp(funE, opPlusIb~1~(comp(opInv~funE, opPlusIb~1~funI)))), opPlusIb~1~funI))) = funS;
prove;
use PE1Suc[y := 2 * x];
use MulNonneg[a := x, b := x];
use PE1Inv[v := 1 + x * x + 3 * x];
use EAt[y := 2 + x * x + 3 * x, m := x + 1, k := x + 1];
prove;
split x = 0;
cases;
prove;
next;
prove;
next;
\end{zproof}

\begin{theorem}{SE1In}
  comp(comp(funE, opInv~(comp(funE, opPlusIb~1~(comp(opInv~funE, opPlusIb~1~funI))))), opPlusIb~1~(comp(comp(opInv~funE, comp(funE, opPlusIb~1~(comp(opInv~funE, opPlusIb~1~funI)))), opPlusIb~1~funI))) \in clos~\{ funE \}~\{ opPlusIb~1, opInv \}~\{ comp \}
\end{theorem}

\begin{zproof}
use CbEIn[b := 1];
prove;
\end{zproof}

\begin{theorem}{CbESOne}
  funS \in clos~\{ funE \}~\{ opPlusIb~1, opInv \}~\{ comp \}
\end{theorem}

\begin{zproof}
use SE1In;
use SE1Val;
prove;
\end{zproof}

From $S$: the predecessor $E\,(E S)^{-1}$, then $F^+$, $\bar 0 = E S E^{-1} I^+$ and
$\bar 1 = S \bar 0$.

\begin{theorem}{SShift}
  (funS, 1) \in Shift
\end{theorem}

\begin{zproof}
invoke Shift;
prove;
\end{zproof}

\begin{theorem}{PEPreds}
  comp(funE, opInv~(comp(funE, funS))) \in Preds
\end{theorem}

\begin{zproof}
invoke Preds;
prove;
use SShift;
use ESubVal[W := funS, c := 1, x := 1 + x];
prove;
\end{zproof}

\begin{theorem}{ZeroExprE2}
  comp(comp(comp(funE, funS), opInv~funE), opPlusI~funI) = cst~0
\end{theorem}

\begin{zproof}
apply UFExt to predicate comp(comp(comp(funE, funS), opInv~funE), opPlusI~funI) = cst~0;
prove;
use EAt[y := 1 + x * x + 2 * x, m := x + 1, k := 0];
prove;
\end{zproof}

\begin{theorem}{CbEPI}
  \forall b: \nat | funS \in clos~\{ funE \}~\{ opPlusIb~b, opInv \}~\{ comp \} @ comp(funE, opInv~(comp(funE, funS))) \in clos~\{ funE \}~\{ opPlusIb~b, opInv \}~\{ comp \} \land opPlusI~funI \in clos~\{ funE \}~\{ opPlusIb~b, opInv \}~\{ comp \}
\end{theorem}

\begin{zproof}
prove;
use CbEIn[b := b];
use PEPreds;
use CbPlusI[X := funE, b := b, P := comp(funE, opInv~(comp(funE, funS))), F := funI];
prove;
\end{zproof}

\begin{theorem}{CbEZeroIn}
  \forall b: \nat | funS \in clos~\{ funE \}~\{ opPlusIb~b, opInv \}~\{ comp \} @ comp(comp(comp(funE, funS), opInv~funE), opPlusI~funI) \in clos~\{ funE \}~\{ opPlusIb~b, opInv \}~\{ comp \}
\end{theorem}

\begin{zproof}
prove;
use CbEPI[b := b];
prove;
\end{zproof}

\begin{theorem}{CbEOne}
  \forall b: \nat | funS \in clos~\{ funE \}~\{ opPlusIb~b, opInv \}~\{ comp \} @ cst~0 \in clos~\{ funE \}~\{ opPlusIb~b, opInv \}~\{ comp \} \land cst~1 \in clos~\{ funE \}~\{ opPlusIb~b, opInv \}~\{ comp \}
\end{theorem}

\begin{zproof}
prove;
use CbEZeroIn[b := b];
use ZeroExprE2;
prove;
use OneExpr;
prove;
\end{zproof}

\begin{theorem}{CbEFromS}
  \forall b: \nat | funS \in clos~\{ funE \}~\{ opPlusIb~b, opInv \}~\{ comp \} @ clos~\{ funE \}~\{ opPlusIb~b, opInv \}~\{ comp \} = ClosJR
\end{theorem}

\begin{zproof}
prove;
use CbEOne[b := b];
use CbEPI[b := b];
use PEPreds;
use CbSub[X := funE, b := b, P := comp(funE, opInv~(comp(funE, funS)))];
use GRbSubGE[X := funE, b := b];
use GREqJRE;
use ClosInPow[I := \{ funE \}, U := \{ opPlusIb~b, opInv \}, B := \{ comp \}];
use CGEInPower;
use SubAntisym[\nat \fun \nat][A := clos~\{ funE \}~\{ opPlusIb~b, opInv \}~\{ comp \}, B := CGE];
with enabled (CGEDef) prove;
\end{zproof}


\subsection{$\langle K, F^{+b}, FG, F^{-1}\rangle$}

\begin{theorem}{rule CbKZ}
  \forall b: \nat @ funZ \in clos~\{ funK \}~\{ opPlusIb~b, opInv \}~\{ comp \}
\end{theorem}

\begin{zproof}
prove;
use CbUn[X := funK, b := b, F := funK];
use KinvZ;
prove;
\end{zproof}

\begin{theorem}{KZI}
  comp(funK, funZ) = funI
\end{theorem}

\begin{zproof}
apply UFExt to predicate comp(funK, funZ) = funI;
prove;
use KZ[x := x];
prove;
\end{zproof}

\begin{theorem}{CbKIExpr}
  \forall b: \nat @ comp(funK, funZ) \in clos~\{ funK \}~\{ opPlusIb~b, opInv \}~\{ comp \}
\end{theorem}

\begin{zproof}
prove;
\end{zproof}

\begin{theorem}{rule CbKI}
  \forall b: \nat @ funI \in clos~\{ funK \}~\{ opPlusIb~b, opInv \}~\{ comp \}
\end{theorem}

\begin{zproof}
prove;
use CbKIExpr[b := b];
use KZI;
prove;
\end{zproof}

\begin{theorem}{ALeCore}
  \forall a, n, q: \num | n \leq q \land 2 * a = n + q @ n \leq a
\end{theorem}

\begin{zproof}
prove;
\end{zproof}

\begin{theorem}{ALe}
  \forall n: \nat @ n \leq funA~n
\end{theorem}

\begin{zproof}
prove;
use TwoA[x := n];
use LeSq[n := n];
use ALeCore[a := funA~n, n := n, q := n * n];
prove;
\end{zproof}

\begin{theorem}{ShiftComp}
  \forall F, G: \nat \fun \nat; a, c: \nat | (F, a) \in Shift \land (G, c) \in Shift @ (comp(F, G), a + c) \in Shift
\end{theorem}

\begin{zproof}
prove;
invoke Shift;
prove;
instantiate y\_\_0 == G~y, y\_\_1 == y;
prove;
\end{zproof}

Subtraction by inversion with $K$.

\begin{theorem}{KSubSurj}
  \forall W: \nat \fun \nat; c: \nat | (W, c) \in Shift @ comp(funK, W) \in Surj
\end{theorem}

\begin{zproof}
invoke Surj;
prove;
instantiate y == funA~(x + c) + x - c;
use ALe[n := x + c];
use ShiftVal[W := W, c := c, y := funA~(x + c) + x - c];
use KAt[m := funA~(x + c) + x, t := x + c, k := x];
prove;
\end{zproof}

\begin{theorem}{KSubCore}
  \forall y, c, x: \nat | funK~(y + c) = x \land c \leq x @ funK~y = x - c
\end{theorem}

\begin{zproof}
prove;
use Decomp[x := y + c];
use KAt[m := y, t := funV~(y + c), k := x - c];
prove;
\end{zproof}

\begin{theorem}{KSubVal}
  \forall W: \nat \fun \nat; c, x: \nat | (W, c) \in Shift \land c \leq x @ funK~(opInv~(comp(funK, W))~x) = x - c
\end{theorem}

\begin{zproof}
prove;
use KSubSurj[W := W, c := c];
use InvRight[F := comp(funK, W), x := x];
use UFNonneg[f := opInv~(comp(funK, W)), x := x];
use ShiftVal[W := W, c := c, y := opInv~(comp(funK, W))~x];
use KSubCore[y := opInv~(comp(funK, W))~x, c := c, x := x];
with disabled (InvRight) prove;
\end{zproof}

$U = K\,(Z T)^{+b}$ is the shift $b - 1$.

\begin{theorem}{UKShift}
  \forall c: \nat @ (comp(funK, opPlusIb~(1 + c)~(comp(funZ, opPlusIb~(1 + c)~funI))), c) \in Shift
\end{theorem}

\begin{zproof}
invoke Shift;
prove;
use ASuc[x := 1 + c + 2 * y];
use UFNonneg[f := funA, x := 1 + c + 2 * y];
use KAt[m := funA~(2 + c + 2 * y) + y + c, t := 2 + c + 2 * y, k := y + c];
prove;
\end{zproof}

\begin{theorem}{rule CbKUIn}
  \forall b: \nat @ comp(funK, opPlusIb~b~(comp(funZ, opPlusIb~b~funI))) \in clos~\{ funK \}~\{ opPlusIb~b, opInv \}~\{ comp \}
\end{theorem}

\begin{zproof}
prove;
\end{zproof}

For $b = d + 2$: the constant $C = K\,(Z P_1)^{+b} U^3 = 2d$ with $P_1 = K\,(K U)^{-1}$, and
$S = K\,(K\,C^{+b})^{-1} U^3$.

\begin{theorem}{CKVal}
  \forall d, x: \nat @ comp(comp(funK, opPlusIb~(2 + d)~(comp(funZ, comp(funK, opInv~(comp(funK, comp(funK, opPlusIb~(2 + d)~(comp(funZ, opPlusIb~(2 + d)~funI))))))))), comp(comp(comp(funK, opPlusIb~(2 + d)~(comp(funZ, opPlusIb~(2 + d)~funI))), comp(funK, opPlusIb~(2 + d)~(comp(funZ, opPlusIb~(2 + d)~funI)))), comp(funK, opPlusIb~(2 + d)~(comp(funZ, opPlusIb~(2 + d)~funI)))))~x = 2 * d
\end{theorem}

\begin{zproof}
prove;
use UKShift[c := 1 + d];
use ShiftVal[W := comp(funK, opPlusIb~(2 + d)~(comp(funZ, opPlusIb~(2 + d)~funI))), c := 1 + d, y := x];
use ShiftVal[W := comp(funK, opPlusIb~(2 + d)~(comp(funZ, opPlusIb~(2 + d)~funI))), c := 1 + d, y := 1 + d + x];
use ShiftVal[W := comp(funK, opPlusIb~(2 + d)~(comp(funZ, opPlusIb~(2 + d)~funI))), c := 1 + d, y := 2 + 2 * d + x];
use KSubVal[W := comp(funK, opPlusIb~(2 + d)~(comp(funZ, opPlusIb~(2 + d)~funI))), c := 1 + d, x := 3 + 3 * d + x];
use ASuc[x := 2 + 2 * d + x];
use ASuc[x := 3 + 2 * d + x];
use UFNonneg[f := funA, x := 2 + 2 * d + x];
use KAt[m := funA~(4 + 2 * d + x) + 2 * d, t := 4 + 2 * d + x, k := 2 * d];
prove;
\end{zproof}

\begin{theorem}{KShifts}
  \forall d: \nat @ (comp(comp(comp(funK, opPlusIb~(2 + d)~(comp(funZ, opPlusIb~(2 + d)~funI))), comp(funK, opPlusIb~(2 + d)~(comp(funZ, opPlusIb~(2 + d)~funI)))), comp(funK, opPlusIb~(2 + d)~(comp(funZ, opPlusIb~(2 + d)~funI)))), 3 * d + 3) \in Shift \land (opPlusIb~(2 + d)~(comp(comp(funK, opPlusIb~(2 + d)~(comp(funZ, comp(funK, opInv~(comp(funK, comp(funK, opPlusIb~(2 + d)~(comp(funZ, opPlusIb~(2 + d)~funI))))))))), comp(comp(comp(funK, opPlusIb~(2 + d)~(comp(funZ, opPlusIb~(2 + d)~funI))), comp(funK, opPlusIb~(2 + d)~(comp(funZ, opPlusIb~(2 + d)~funI)))), comp(funK, opPlusIb~(2 + d)~(comp(funZ, opPlusIb~(2 + d)~funI)))))), 3 * d + 2) \in Shift
\end{theorem}

\begin{zproof}
prove;
use UKShift[c := 1 + d];
use ShiftComp[F := comp(funK, opPlusIb~(2 + d)~(comp(funZ, opPlusIb~(2 + d)~funI))), G := comp(funK, opPlusIb~(2 + d)~(comp(funZ, opPlusIb~(2 + d)~funI))), a := 1 + d, c := 1 + d];
use ShiftComp[F := comp(comp(funK, opPlusIb~(2 + d)~(comp(funZ, opPlusIb~(2 + d)~funI))), comp(funK, opPlusIb~(2 + d)~(comp(funZ, opPlusIb~(2 + d)~funI)))), G := comp(funK, opPlusIb~(2 + d)~(comp(funZ, opPlusIb~(2 + d)~funI))), a := 2 + 2 * d, c := 1 + d];
prove;
invoke Shift;
prove;
use CKVal[d := d, x := y];
prove;
\end{zproof}

\begin{theorem}{SFromShiftsK}
  \forall G, H: \nat \fun \nat; d: \nat | (G, d + 1) \in Shift \land (H, d) \in Shift @ comp(comp(funK, opInv~(comp(funK, H))), G) = funS
\end{theorem}

\begin{zproof}
apply UFExt to predicate comp(comp(funK, opInv~(comp(funK, H))), G) = funS;
prove;
use ShiftVal[W := G, c := d + 1, y := x];
use KSubVal[W := H, c := d, x := 1 + x + d];
prove;
\end{zproof}

\begin{theorem}{rule CbKU3In}
  \forall b: \nat @ comp(comp(comp(funK, opPlusIb~b~(comp(funZ, opPlusIb~b~funI))), comp(funK, opPlusIb~b~(comp(funZ, opPlusIb~b~funI)))), comp(funK, opPlusIb~b~(comp(funZ, opPlusIb~b~funI)))) \in clos~\{ funK \}~\{ opPlusIb~b, opInv \}~\{ comp \}
\end{theorem}

\begin{zproof}
prove;
\end{zproof}

\begin{theorem}{rule CbKP1In}
  \forall b: \nat @ comp(funK, opInv~(comp(funK, comp(funK, opPlusIb~b~(comp(funZ, opPlusIb~b~funI)))))) \in clos~\{ funK \}~\{ opPlusIb~b, opInv \}~\{ comp \}
\end{theorem}

\begin{zproof}
prove;
\end{zproof}

\begin{theorem}{rule CbKCIn}
  \forall b: \nat @ comp(comp(funK, opPlusIb~b~(comp(funZ, comp(funK, opInv~(comp(funK, comp(funK, opPlusIb~b~(comp(funZ, opPlusIb~b~funI))))))))), comp(comp(comp(funK, opPlusIb~b~(comp(funZ, opPlusIb~b~funI))), comp(funK, opPlusIb~b~(comp(funZ, opPlusIb~b~funI)))), comp(funK, opPlusIb~b~(comp(funZ, opPlusIb~b~funI))))) \in clos~\{ funK \}~\{ opPlusIb~b, opInv \}~\{ comp \}
\end{theorem}

\begin{zproof}
prove;
\end{zproof}

\begin{theorem}{CbKSTwoIn}
  \forall d: \nat @ comp(comp(funK, opInv~(comp(funK, opPlusIb~(2 + d)~(comp(comp(funK, opPlusIb~(2 + d)~(comp(funZ, comp(funK, opInv~(comp(funK, comp(funK, opPlusIb~(2 + d)~(comp(funZ, opPlusIb~(2 + d)~funI))))))))), comp(comp(comp(funK, opPlusIb~(2 + d)~(comp(funZ, opPlusIb~(2 + d)~funI))), comp(funK, opPlusIb~(2 + d)~(comp(funZ, opPlusIb~(2 + d)~funI)))), comp(funK, opPlusIb~(2 + d)~(comp(funZ, opPlusIb~(2 + d)~funI))))))))), comp(comp(comp(funK, opPlusIb~(2 + d)~(comp(funZ, opPlusIb~(2 + d)~funI))), comp(funK, opPlusIb~(2 + d)~(comp(funZ, opPlusIb~(2 + d)~funI)))), comp(funK, opPlusIb~(2 + d)~(comp(funZ, opPlusIb~(2 + d)~funI))))) \in clos~\{ funK \}~\{ opPlusIb~(2 + d), opInv \}~\{ comp \}
\end{theorem}

\begin{zproof}
prove;
\end{zproof}

\begin{theorem}{CbKSTwo}
  \forall d: \nat @ funS \in clos~\{ funK \}~\{ opPlusIb~(2 + d), opInv \}~\{ comp \}
\end{theorem}

\begin{zproof}
prove;
use CbKSTwoIn[d := d];
use KShifts[d := d];
use SFromShiftsK[G := comp(comp(comp(funK, opPlusIb~(2 + d)~(comp(funZ, opPlusIb~(2 + d)~funI))), comp(funK, opPlusIb~(2 + d)~(comp(funZ, opPlusIb~(2 + d)~funI)))), comp(funK, opPlusIb~(2 + d)~(comp(funZ, opPlusIb~(2 + d)~funI)))), H := opPlusIb~(2 + d)~(comp(comp(funK, opPlusIb~(2 + d)~(comp(funZ, comp(funK, opInv~(comp(funK, comp(funK, opPlusIb~(2 + d)~(comp(funZ, opPlusIb~(2 + d)~funI))))))))), comp(comp(comp(funK, opPlusIb~(2 + d)~(comp(funZ, opPlusIb~(2 + d)~funI))), comp(funK, opPlusIb~(2 + d)~(comp(funZ, opPlusIb~(2 + d)~funI)))), comp(funK, opPlusIb~(2 + d)~(comp(funZ, opPlusIb~(2 + d)~funI)))))), d := 3 * d + 2];
prove;
\end{zproof}

For $b = 1$: $\hat P = K\,(Z^{+1})^{+1}$ ($\hat P(y + 1) = y$, $\hat P(0) = 1$) and
$S = K\,\hat P^{-1} (Z T)^{+1}$.

\begin{theorem}{PK1Suc}
  \forall y: \nat @ comp(funK, opPlusIb~1~(opPlusIb~1~funZ))~(1 + y) = y
\end{theorem}

\begin{zproof}
prove;
use ASuc[x := 1 + y];
use ASuc[x := 2 + y];
use UFNonneg[f := funA, x := 1 + y];
use KAt[m := funA~(3 + y) + y, t := 3 + y, k := y];
prove;
\end{zproof}

\begin{theorem}{PK1Zero}
  comp(funK, opPlusIb~1~(opPlusIb~1~funZ))~0 = 1
\end{theorem}

\begin{zproof}
prove;
use AZero;
use KAt[m := 2, t := 1, k := 1];
prove;
\end{zproof}

\begin{theorem}{PK1Core}
  \forall v, y: \nat | comp(funK, opPlusIb~1~(opPlusIb~1~funZ))~y = v \land \lnot v = 1 @ y = 1 + v
\end{theorem}

\begin{zproof}
prove;
use PK1Zero;
use PK1Suc[y := y - 1];
split y = 0;
cases;
prove;
next;
prove;
next;
\end{zproof}

\begin{theorem}{PK1Surj}
  comp(funK, opPlusIb~1~(opPlusIb~1~funZ)) \in Surj
\end{theorem}

\begin{zproof}
invoke Surj;
prove;
use PK1Suc[y := x];
instantiate y == 1 + x;
prove;
\end{zproof}

\begin{theorem}{PK1Inv}
  \forall v: \nat | \lnot v = 1 @ opInv~(comp(funK, opPlusIb~1~(opPlusIb~1~funZ)))~v = 1 + v
\end{theorem}

\begin{zproof}
prove;
use PK1Surj;
use InvRight[F := comp(funK, opPlusIb~1~(opPlusIb~1~funZ)), x := v];
use UFNonneg[f := opInv~(comp(funK, opPlusIb~1~(opPlusIb~1~funZ))), x := v];
use PK1Core[v := v, y := opInv~(comp(funK, opPlusIb~1~(opPlusIb~1~funZ)))~v];
with disabled (InvRight) prove;
\end{zproof}

\begin{theorem}{SK1Val}
  comp(comp(funK, opInv~(comp(funK, opPlusIb~1~(opPlusIb~1~funZ)))), opPlusIb~1~(comp(funZ, opPlusIb~1~funI))) = funS
\end{theorem}

\begin{zproof}
apply UFExt to predicate comp(comp(funK, opInv~(comp(funK, opPlusIb~1~(opPlusIb~1~funZ)))), opPlusIb~1~(comp(funZ, opPlusIb~1~funI))) = funS;
prove;
use ASuc[x := 1 + 2 * x];
use UFNonneg[f := funA, x := 1 + 2 * x];
use ALe[n := 2 + 2 * x];
use PK1Inv[v := funA~(2 + 2 * x) + x];
use KAt[m := 1 + funA~(2 + 2 * x) + x, t := 2 + 2 * x, k := 1 + x];
prove;
\end{zproof}

\begin{theorem}{SK1In}
  comp(comp(funK, opInv~(comp(funK, opPlusIb~1~(opPlusIb~1~funZ)))), opPlusIb~1~(comp(funZ, opPlusIb~1~funI))) \in clos~\{ funK \}~\{ opPlusIb~1, opInv \}~\{ comp \}
\end{theorem}

\begin{zproof}
prove;
use CbUn[X := funK, b := 1, F := funZ];
use CbUn[X := funK, b := 1, F := opPlusIb~1~funZ];
use CbComp[X := funK, b := 1, F := funK, G := opPlusIb~1~(opPlusIb~1~funZ)];
prove;
\end{zproof}

\begin{theorem}{CbKSOne}
  funS \in clos~\{ funK \}~\{ opPlusIb~1, opInv \}~\{ comp \}
\end{theorem}

\begin{zproof}
use SK1In;
use SK1Val;
prove;
\end{zproof}

From $S$: $P = K (K S)^{-1}$, $F^+ = P^b F^{+b}$, $A$, $\bar 0 = K A$ and $\bar 1 = S \bar 0$.

\begin{theorem}{PKPreds}
  exprPK \in Preds
\end{theorem}

\begin{zproof}
invoke Preds;
prove;
use PKSuc[x := x];
prove;
\end{zproof}

\begin{theorem}{rule CbKPlusI}
  \forall b: \nat; F: \nat \fun \nat | funS \in clos~\{ funK \}~\{ opPlusIb~b, opInv \}~\{ comp \} \land F \in clos~\{ funK \}~\{ opPlusIb~b, opInv \}~\{ comp \} @ opPlusI~F \in clos~\{ funK \}~\{ opPlusIb~b, opInv \}~\{ comp \}
\end{theorem}

\begin{zproof}
prove;
use PKPreds;
use CbComp[X := funK, b := b, F := funK, G := funS];
use CbUn[X := funK, b := b, F := comp(funK, funS)];
use CbComp[X := funK, b := b, F := funK, G := opInv~(comp(funK, funS))];
invoke exprPK;
use CbPlusI[X := funK, b := b, P := comp(funK, opInv~(comp(funK, funS))), F := F];
prove;
\end{zproof}

\begin{theorem}{ZeroExprKA}
  comp(funK, funA) = cst~0
\end{theorem}

\begin{zproof}
apply UFExt to predicate comp(funK, funA) = cst~0;
prove;
use UFNonneg[f := funA, x := x];
use KAt[m := funA~x, t := x, k := 0];
prove;
\end{zproof}

\begin{theorem}{CbKAIn}
  \forall b: \nat | funS \in clos~\{ funK \}~\{ opPlusIb~b, opInv \}~\{ comp \} @ exprAK \in clos~\{ funK \}~\{ opPlusIb~b, opInv \}~\{ comp \}
\end{theorem}

\begin{zproof}
invoke exprAK;
prove;
\end{zproof}

\begin{theorem}{CbKOne}
  \forall b: \nat | funS \in clos~\{ funK \}~\{ opPlusIb~b, opInv \}~\{ comp \} @ cst~0 \in clos~\{ funK \}~\{ opPlusIb~b, opInv \}~\{ comp \} \land cst~1 \in clos~\{ funK \}~\{ opPlusIb~b, opInv \}~\{ comp \}
\end{theorem}

\begin{zproof}
prove;
use CbKAIn[b := b];
use AExprK;
prove;
use ZeroExprKA;
prove;
use OneExpr;
prove;
\end{zproof}

\begin{theorem}{CbKPIn}
  \forall b: \nat | funS \in clos~\{ funK \}~\{ opPlusIb~b, opInv \}~\{ comp \} @ exprPK \in clos~\{ funK \}~\{ opPlusIb~b, opInv \}~\{ comp \}
\end{theorem}

\begin{zproof}
invoke exprPK;
prove;
\end{zproof}

\begin{theorem}{CbKFromS}
  \forall b: \nat | funS \in clos~\{ funK \}~\{ opPlusIb~b, opInv \}~\{ comp \} @ clos~\{ funK \}~\{ opPlusIb~b, opInv \}~\{ comp \} = ClosJR
\end{theorem}

\begin{zproof}
prove;
use CbKOne[b := b];
use CbKPIn[b := b];
use PKPreds;
use CbSub[X := funK, b := b, P := exprPK];
use GRbSubGE[X := funK, b := b];
use GEK;
use GREqJR;
use GKSubGE;
use GESubGK;
use ClosInPow[I := \{ funK \}, U := \{ opPlusIb~b, opInv \}, B := \{ comp \}];
use CGEInPower;
use CGKInPower;
use SubAntisym[\nat \fun \nat][A := clos~\{ funK \}~\{ opPlusIb~b, opInv \}~\{ comp \}, B := CGE];
with enabled (CGKDef) prove;
\end{zproof}


\subsection{$\langle L, F^{+b}, FG, F^{-1}\rangle$}

$L^{-1} = A$; $U = L\,(\dots(A^{+b})\dots)^{+b}$ ($2b - 2$ times) is the shift $b - 1$,
$L A^{+b}$ subtracts $b - 2$, and $S = L A^{+b} U$.

\begin{theorem}{LSurj}
  funL \in Surj
\end{theorem}

\begin{zproof}
invoke Surj;
prove;
instantiate y == funA~x;
use UFNonneg[f := funA, x := x];
use LAt[m := funA~x, t := x, k := 0];
prove;
\end{zproof}

\begin{theorem}{LLowCore}
  \forall x, y, t, k: \nat | y = funA~t + k \land k \leq t \land t - k = x @ funA~x \leq y
\end{theorem}

\begin{zproof}
prove;
use AMono[m := x, n := t];
prove;
\end{zproof}

\begin{theorem}{LinvA}
  opInv~funL = funA
\end{theorem}

\begin{zproof}
apply UFExt to predicate opInv~funL = funA;
prove;
use LSurj;
use UFNonneg[f := funA, x := x];
use LAt[m := funA~x, t := x, k := 0];
use InvLeast[F := funL, x := x, y := funA~x];
use InvRight[F := funL, x := x];
use UFNonneg[f := opInv~funL, x := x];
use Decomp[x := opInv~funL~x];
use LAt[m := opInv~funL~x, t := funV~(opInv~funL~x), k := funK~(opInv~funL~x)];
use LLowCore[x := x, y := opInv~funL~x, t := funV~(opInv~funL~x), k := funK~(opInv~funL~x)];
with disabled (InvRight) prove;
split funA~x \leq opInv~funL~x;
cases;
prove;
next;
with normalization prove;
next;
\end{zproof}

\begin{theorem}{rule CbLA}
  \forall b: \nat @ funA \in clos~\{ funL \}~\{ opPlusIb~b, opInv \}~\{ comp \}
\end{theorem}

\begin{zproof}
prove;
use CbUn[X := funL, b := b, F := funL];
use LinvA;
prove;
\end{zproof}

$A + n(I + \bar b)$, i.e.\ $A^{+b}$ iterated $n$ times, is in the closure.

\begin{zed}
  IterLSet == \{ n: \nat | \forall b: \nat @ opPlus(funA, comp(funM~n, opPlus(funI, cst~b))) \in clos~\{ funL \}~\{ opPlusIb~b, opInv \}~\{ comp \} \}
\end{zed}

\begin{zproof}[IterLSet\$domainCheck]
prove;
\end{zproof}

\begin{theorem}{IterLEqZero}
  \forall b: \nat @ opPlus(funA, comp(funM~0, opPlus(funI, cst~b))) = funA
\end{theorem}

\begin{zproof}
prove;
apply UFExt to predicate opPlus(funA, comp(funM~0, opPlus(funI, cst~b))) = funA;
prove;
\end{zproof}

\begin{theorem}{IterLEqSuc}
  \forall n, b: \nat @ opPlus(funA, comp(funM~(1 + n), opPlus(funI, cst~b))) = opPlusIb~b~(opPlus(funA, comp(funM~n, opPlus(funI, cst~b))))
\end{theorem}

\begin{zproof}
prove;
apply UFExt to predicate opPlus(funA, comp(funM~(1 + n), opPlus(funI, cst~b))) = opPlusIb~b~(opPlus(funA, comp(funM~n, opPlus(funI, cst~b))));
prove;
\end{zproof}

\begin{theorem}{rule IterLBase}
  0 \in IterLSet
\end{theorem}

\begin{zproof}
invoke IterLSet;
prove;
use IterLEqZero[b := b];
prove;
\end{zproof}

\begin{theorem}{rule IterLStep}
  \forall n: IterLSet @ 1 + n \in IterLSet
\end{theorem}

\begin{zproof}
invoke IterLSet;
prove;
instantiate b\_\_0 == b;
use IterLEqSuc[n := n, b := b];
use CbUn[X := funL, b := b, F := opPlus(funA, comp(funM~n, opPlus(funI, cst~b)))];
prove;
\end{zproof}

\begin{theorem}{IterLSub}
  \nat \subseteq IterLSet
\end{theorem}

\begin{zproof}
apply natInduction to predicate \nat \subseteq IterLSet;
prove;
invoke IterLSet;
prove;
\end{zproof}

\begin{theorem}{rule IterLAll}
  \forall n, b: \nat @ opPlus(funA, comp(funM~n, opPlus(funI, cst~b))) \in clos~\{ funL \}~\{ opPlusIb~b, opInv \}~\{ comp \}
\end{theorem}

\begin{zproof}
prove;
use IterLSub;
invoke IterLSet;
with enabled (inPower) prove;
instantiate e == n;
prove;
instantiate b\_\_1 == b;
prove;
\end{zproof}


\begin{theorem}{AAddTwo}
  \forall x, c: \nat @ funA~(x + 2 * c) = funA~x + 2 * c * x + 2 * c * c + c
\end{theorem}

\begin{zproof}
prove;
use TwoA[x := x + 2 * c];
use TwoA[x := x];
use Halve[a := funA~(x + 2 * c), b := funA~x + 2 * c * x + 2 * c * c + c];
prove;
\end{zproof}

\begin{theorem}{ULShift}
  \forall c: \nat @ (comp(funL, opPlus(funA, comp(funM~(2 * c), opPlus(funI, cst~(1 + c))))), c) \in Shift
\end{theorem}

\begin{zproof}
invoke Shift;
prove;
use AAddTwo[x := y, c := c];
use UFNonneg[f := funA, x := y + 2 * c];
use LAt[m := funA~(y + 2 * c) + c, t := y + 2 * c, k := c];
prove;
\end{zproof}

\begin{theorem}{SigmaS}
  \forall c: \nat; U: \nat \fun \nat | (U, c) \in Shift @ comp(comp(funL, opPlusIb~(1 + c)~funA), U) = funS
\end{theorem}

\begin{zproof}
apply UFExt to predicate comp(comp(funL, opPlusIb~(1 + c)~funA), U) = funS;
prove;
use ShiftVal[W := U, c := c, y := x];
use ASuc[x := x + c];
use UFNonneg[f := funA, x := x + c];
use LAt[m := funA~(1 + x + c) + c, t := 1 + x + c, k := c];
prove;
\end{zproof}

\begin{theorem}{CbLSIn}
  \forall c: \nat @ comp(comp(funL, opPlusIb~(1 + c)~funA), comp(funL, opPlus(funA, comp(funM~(2 * c), opPlus(funI, cst~(1 + c)))))) \in clos~\{ funL \}~\{ opPlusIb~(1 + c), opInv \}~\{ comp \}
\end{theorem}

\begin{zproof}
prove;
\end{zproof}

\begin{theorem}{CbLS}
  \forall c: \nat @ funS \in clos~\{ funL \}~\{ opPlusIb~(1 + c), opInv \}~\{ comp \}
\end{theorem}

\begin{zproof}
prove;
use CbLSIn[c := c];
use ULShift[c := c];
use SigmaS[c := c, U := comp(funL, opPlus(funA, comp(funM~(2 * c), opPlus(funI, cst~(1 + c)))))];
prove;
\end{zproof}

From $S$: the predecessor $L S A$, $F^+$, $\bar 0 = L A^+$ and $\bar 1 = S \bar 0$.

\begin{theorem}{PLPreds}
  comp(comp(funL, funS), funA) \in Preds
\end{theorem}

\begin{zproof}
invoke Preds;
prove;
use UFNonneg[f := funA, x := 1 + x];
use LAt[m := 1 + funA~(1 + x), t := 1 + x, k := 1];
prove;
\end{zproof}

\begin{theorem}{ZeroExprLA}
  comp(funL, opPlusI~funA) = cst~0
\end{theorem}

\begin{zproof}
apply UFExt to predicate comp(funL, opPlusI~funA) = cst~0;
prove;
use UFNonneg[f := funA, x := x];
use LAt[m := funA~x + x, t := x, k := x];
prove;
\end{zproof}

\begin{theorem}{CbLPI}
  \forall b: \nat | funS \in clos~\{ funL \}~\{ opPlusIb~b, opInv \}~\{ comp \} @ comp(comp(funL, funS), funA) \in clos~\{ funL \}~\{ opPlusIb~b, opInv \}~\{ comp \} \land opPlusI~funA \in clos~\{ funL \}~\{ opPlusIb~b, opInv \}~\{ comp \}
\end{theorem}

\begin{zproof}
prove;
use PLPreds;
use CbPlusI[X := funL, b := b, P := comp(comp(funL, funS), funA), F := funA];
prove;
\end{zproof}

\begin{theorem}{CbLOne}
  \forall b: \nat | funS \in clos~\{ funL \}~\{ opPlusIb~b, opInv \}~\{ comp \} @ cst~0 \in clos~\{ funL \}~\{ opPlusIb~b, opInv \}~\{ comp \} \land cst~1 \in clos~\{ funL \}~\{ opPlusIb~b, opInv \}~\{ comp \}
\end{theorem}

\begin{zproof}
prove;
use CbLPI[b := b];
use ZeroExprLA;
prove;
use OneExpr;
prove;
\end{zproof}

\begin{theorem}{CbLFromS}
  \forall b: \nat | funS \in clos~\{ funL \}~\{ opPlusIb~b, opInv \}~\{ comp \} @ clos~\{ funL \}~\{ opPlusIb~b, opInv \}~\{ comp \} = ClosJR
\end{theorem}

\begin{zproof}
prove;
use CbLOne[b := b];
use CbLPI[b := b];
use PLPreds;
use CbSub[X := funL, b := b, P := comp(comp(funL, funS), funA)];
use GRbSubGE[X := funL, b := b];
use GEL;
use GREqJR;
use GLSubGE;
use ClosInPow[I := \{ funL \}, U := \{ opPlusIb~b, opInv \}, B := \{ comp \}];
use CGEInPower;
use CGLInPower;
use SubAntisym[\nat \fun \nat][A := clos~\{ funL \}~\{ opPlusIb~b, opInv \}~\{ comp \}, B := CGE];
with enabled (CGLDef) prove;
\end{zproof}

\subsection{The shifted operator with inversion: the result}

\begin{theorem}{CbES}
  \forall b: \nat | 1 \leq b @ funS \in clos~\{ funE \}~\{ opPlusIb~b, opInv \}~\{ comp \}
\end{theorem}

\begin{zproof}
prove;
use Mod2Cases[x := b];
use Div2Nonneg[x := b];
split b = 1;
cases;
use CbESOne;
prove;
next;
split b \mod 2 = 0;
cases;
use CbESEven[p := b \div 2 - 1];
prove;
next;
use CbESOdd[p := b \div 2 - 1];
prove;
next;
next;
\end{zproof}

\begin{theorem}{CbKS}
  \forall b: \nat | 1 \leq b @ funS \in clos~\{ funK \}~\{ opPlusIb~b, opInv \}~\{ comp \}
\end{theorem}

\begin{zproof}
prove;
split b = 1;
cases;
use CbKSOne;
prove;
next;
use CbKSTwo[d := b - 2];
prove;
next;
\end{zproof}

\begin{theorem}{CbLSAll}
  \forall b: \nat | 1 \leq b @ funS \in clos~\{ funL \}~\{ opPlusIb~b, opInv \}~\{ comp \}
\end{theorem}

\begin{zproof}
prove;
use CbLS[c := b - 1];
prove;
\end{zproof}

\begin{theorem}{ExtGRb}
  \forall b: \nat | 1 \leq b @ clos~\{ funE \}~\{ opPlusIb~b, opInv \}~\{ comp \} = ClosJR \land clos~\{ funK \}~\{ opPlusIb~b, opInv \}~\{ comp \} = ClosJR \land clos~\{ funL \}~\{ opPlusIb~b, opInv \}~\{ comp \} = ClosJR
\end{theorem}

\begin{zproof}
prove;
use CbES[b := b];
use CbKS[b := b];
use CbLSAll[b := b];
use CbEFromS[b := b];
use CbKFromS[b := b];
use CbLFromS[b := b];
prove;
\end{zproof}

$b = 0$: every generated function fixes $0$.

\begin{theorem}{FixZeroUClosedIb}
  FixZero \in UClosed~\{ opPlusIb~0, opInv \}
\end{theorem}

\begin{zproof}
use FixZeroInPower;
with enabled (UClosedDef) prove;
invoke FixZero;
prove;
use InPair[\power (\power (\num \cross \num) \cross \power (\num \cross \num))][u := u, p := opPlusIb~0, q := opInv];
split u = opInv;
cases;
use InvFix0[F := F];
prove;
next;
prove;
next;
\end{zproof}

\begin{theorem}{GRbZeroFix}
  \forall X: \nat \fun \nat | X~0 = 0 @ clos~\{ X \}~\{ opPlusIb~0, opInv \}~\{ comp \} \subseteq FixZero
\end{theorem}

\begin{zproof}
prove;
use ClosLeast[I := \{ X \}, U := \{ opPlusIb~0, opInv \}, B := \{ comp \}, Z := FixZero];
use FixZeroInPower;
use FixZeroUClosedIb;
use FixZeroBClosed;
invoke FixZero;
prove;
\end{zproof}

\begin{theorem}{GRbZeroNotOne}
  \forall X: \nat \fun \nat | X~0 = 0 @ \lnot cst~1 \in clos~\{ X \}~\{ opPlusIb~0, opInv \}~\{ comp \}
\end{theorem}

\begin{zproof}
use GRbZeroFix[X := X];
apply subsetDef to predicate clos~\{ X \}~\{ opPlusIb~0, opInv \}~\{ comp \} \subseteq FixZero;
apply inPower to predicate clos~\{ X \}~\{ opPlusIb~0, opInv \}~\{ comp \} \in \power FixZero;
instantiate e == cst~1;
invoke FixZero;
prove;
\end{zproof}

\begin{theorem}{GRbZeroNeq}
  \forall X: \nat \fun \nat | X~0 = 0 @ \lnot clos~\{ X \}~\{ opPlusIb~0, opInv \}~\{ comp \} = ClosJR
\end{theorem}

\begin{zproof}
use GRbZeroNotOne[X := X];
use JROne;
prove;
split clos~\{ X \}~\{ opPlusIb~0, opInv \}~\{ comp \} = ClosJR;
cases;
equality substitute clos~\{ X \}~(\{ opPlusIb~0 \} \cup \{ opInv \})~\{ comp \};
prove;
next;
prove;
next;
\end{zproof}

\begin{theorem}{ExtGRbIff}
  \forall b: \nat @ (clos~\{ funE \}~\{ opPlusIb~b, opInv \}~\{ comp \} = ClosJR \iff 1 \leq b) \land (clos~\{ funK \}~\{ opPlusIb~b, opInv \}~\{ comp \} = ClosJR \iff 1 \leq b) \land (clos~\{ funL \}~\{ opPlusIb~b, opInv \}~\{ comp \} = ClosJR \iff 1 \leq b)
\end{theorem}

\begin{zproof}
prove;
split b = 0;
cases;
use GRbZeroNeq[X := funE];
use GRbZeroNeq[X := funK];
use GRbZeroNeq[X := funL];
use KZero;
prove;
next;
use ExtGRb[b := b];
prove;
next;
\end{zproof}


\subsection{Remark: any constant $\bar n$ with $n \geq 1$}

For $X \in \{E, K, L\}$ and a set $U$ of unary operators with $F^+ \in U$,
$\langle \bar n, X, U, FG\rangle$ and $\langle \bar 1, X, U, FG\rangle$ have the same closure if
$n \geq 1$.  Since $\bar n^+(x) = x + n$, every constant $c + kn$ is generated from $c$, and $X$
steps down from $c$ to $c - 1$:
$E(c + (nc^2 + 2c)n) = c - 1$, $K(c + (2nc^2 + 3c)n) = c - 1$,
$L(c + (2nc^2 + 2c)n) = nc - c$ and $L((nc - c) + 2nc^2 n) = c - 1$.

\begin{theorem}{CnComp}
  \forall X: \nat \fun \nat; U: \power ((\nat \fun \nat) \fun (\nat \fun \nat)); n: \nat; F, G: \nat \fun \nat | F \in clos~\{ cst~n, X \}~U~\{ comp \} \land G \in clos~\{ cst~n, X \}~U~\{ comp \} @ comp(F, G) \in clos~\{ cst~n, X \}~U~\{ comp \}
\end{theorem}

\begin{zproof}
prove;
\end{zproof}

\begin{theorem}{CstComp}
  \forall F: \nat \fun \nat; m: \nat @ comp(F, cst~m) = cst~(F~m)
\end{theorem}

\begin{zproof}
prove;
apply UFExt to predicate comp(F, cst~m) = cst~(F~m);
prove;
\end{zproof}

\begin{theorem}{CstShiftEq}
  \forall n, c, k: \nat @ comp(iter[\nat]~k~(opPlusI~(cst~n)), cst~c) = cst~(c + k * n)
\end{theorem}

\begin{zproof}
prove;
apply UFExt to predicate comp(iter[\nat]~k~(opPlusI~(cst~n)), cst~c) = cst~(c + k * n);
prove;
use ShiftPow[k := k, c := n, x := c, F := opPlusI~(cst~n)];
invoke Shift;
prove;
\end{zproof}

\begin{theorem}{CstShift}
  \forall X: \nat \fun \nat; U: \power ((\nat \fun \nat) \fun (\nat \fun \nat)); n, c, k: \nat | opPlusI \in U \land cst~c \in clos~\{ cst~n, X \}~U~\{ comp \} @ cst~(c + k * n) \in clos~\{ cst~n, X \}~U~\{ comp \} \land cst~(X~(c + k * n)) \in clos~\{ cst~n, X \}~U~\{ comp \}
\end{theorem}

\begin{zproof}
prove;
use PowIn[I := \{ cst~n, X \}, U := U, B := \{ comp \}, F := opPlusI~(cst~n), G := cst~c, k := k];
use CstShiftEq[n := n, c := c, k := k];
prove;
use CstComp[F := X, m := c + k * n];
prove;
\end{zproof}

The descent, with $n = n_1 + 1$ and $c = c_1 + 1$.

\begin{theorem}{EDesc}
  \forall n1, c1: \nat @ funE~((1 + c1) + ((1 + n1) * (1 + c1) * (1 + c1) + 2 * (1 + c1)) * (1 + n1)) = c1
\end{theorem}

\begin{zproof}
prove;
use MulNonneg[a := n1, b := c1];
use EAt[y := (1 + c1) + ((1 + n1) * (1 + c1) * (1 + c1) + 2 * (1 + c1)) * (1 + n1), m := (1 + n1) * (1 + c1) + 1, k := c1];
prove;
\end{zproof}

\begin{theorem}{KDesc}
  \forall n1, c1: \nat @ funK~((1 + c1) + (2 * (1 + n1) * (1 + c1) * (1 + c1) + 3 * (1 + c1)) * (1 + n1)) = c1
\end{theorem}

\begin{zproof}
prove;
use MulNonneg[a := n1, b := c1];
use TwoA[x := 2 * (1 + n1) * (1 + c1) + 1];
use Halve[a := funA~(2 * (1 + n1) * (1 + c1) + 1), b := (2 * (1 + n1) * (1 + c1) + 1) * ((1 + n1) * (1 + c1) + 1)];
use KAt[m := (1 + c1) + (2 * (1 + n1) * (1 + c1) * (1 + c1) + 3 * (1 + c1)) * (1 + n1), t := 2 * (1 + n1) * (1 + c1) + 1, k := c1];
prove;
\end{zproof}

\begin{theorem}{LDescOne}
  \forall n1, c: \nat @ funL~(c + (2 * (1 + n1) * c * c + 2 * c) * (1 + n1)) = n1 * c
\end{theorem}

\begin{zproof}
prove;
use MulNonneg[a := n1, b := c];
use TwoA[x := 2 * (1 + n1) * c];
use Halve[a := funA~(2 * (1 + n1) * c), b := 2 * (1 + n1) * (1 + n1) * c * c + (1 + n1) * c];
use LAt[m := c + (2 * (1 + n1) * c * c + 2 * c) * (1 + n1), t := 2 * (1 + n1) * c, k := (1 + n1) * c + c];
prove;
\end{zproof}

\begin{theorem}{LDescTwo}
  \forall n1, c1: \nat @ funL~(n1 * (1 + c1) + (2 * (1 + n1) * (1 + c1) * (1 + c1)) * (1 + n1)) = c1
\end{theorem}

\begin{zproof}
prove;
use MulNonneg[a := n1, b := c1];
use TwoA[x := (2 * n1 + 1) * (1 + c1) + c1];
use Halve[a := funA~((2 * n1 + 1) * (1 + c1) + c1), b := n1 * (1 + c1) + (2 * (1 + n1) * (1 + c1) * (1 + c1)) * (1 + n1) - (2 * n1 + 1) * (1 + c1)];
use LAt[m := n1 * (1 + c1) + (2 * (1 + n1) * (1 + c1) * (1 + c1)) * (1 + n1), t := (2 * n1 + 1) * (1 + c1) + c1, k := (2 * n1 + 1) * (1 + c1)];
prove;
\end{zproof}


\begin{theorem}{StepE}
  \forall U: \power ((\nat \fun \nat) \fun (\nat \fun \nat)); n, c: \nat | opPlusI \in U \land 1 \leq n \land 1 \leq c \land cst~c \in clos~\{ cst~n, funE \}~U~\{ comp \} @ cst~(c - 1) \in clos~\{ cst~n, funE \}~U~\{ comp \}
\end{theorem}

\begin{zproof}
prove;
use CstShift[X := funE, U := U, n := n, c := c, k := n * c * c + 2 * c];
use EDesc[n1 := n - 1, c1 := c - 1];
prove;
\end{zproof}

\begin{zed}
  DescESet == \{ j: \nat | \forall U: \power ((\nat \fun \nat) \fun (\nat \fun \nat)); n: \nat | opPlusI \in U \land j + 1 \leq n @ cst~(n - j) \in clos~\{ cst~n, funE \}~U~\{ comp \} \}
\end{zed}

\begin{zproof}[DescESet\$domainCheck]
prove;
\end{zproof}

\begin{theorem}{rule DescEBase}
  0 \in DescESet
\end{theorem}

\begin{zproof}
invoke DescESet;
prove;
\end{zproof}

\begin{theorem}{rule DescEStep}
  \forall j: DescESet @ 1 + j \in DescESet
\end{theorem}

\begin{zproof}
invoke DescESet;
prove;
instantiate U\_\_0 == U, n\_\_0 == n;
use StepE[U := U, n := n, c := n - j];
prove;
\end{zproof}

\begin{theorem}{DescESub}
  \nat \subseteq DescESet
\end{theorem}

\begin{zproof}
apply natInduction to predicate \nat \subseteq DescESet;
prove;
invoke DescESet;
prove;
\end{zproof}

\begin{theorem}{OneInE}
  \forall U: \power ((\nat \fun \nat) \fun (\nat \fun \nat)); m: \nat | opPlusI \in U @ cst~1 \in clos~\{ cst~(1 + m), funE \}~U~\{ comp \}
\end{theorem}

\begin{zproof}
prove;
use DescESub;
invoke DescESet;
with enabled (inPower) prove;
instantiate e\_\_0 == m, U\_\_0 == U, n == 1 + m;
with enabled (inPower) prove;
\end{zproof}


\begin{theorem}{StepK}
  \forall U: \power ((\nat \fun \nat) \fun (\nat \fun \nat)); n, c: \nat | opPlusI \in U \land 1 \leq n \land 1 \leq c \land cst~c \in clos~\{ cst~n, funK \}~U~\{ comp \} @ cst~(c - 1) \in clos~\{ cst~n, funK \}~U~\{ comp \}
\end{theorem}

\begin{zproof}
prove;
use CstShift[X := funK, U := U, n := n, c := c, k := 2 * n * c * c + 3 * c];
use KDesc[n1 := n - 1, c1 := c - 1];
prove;
\end{zproof}

\begin{zed}
  DescKSet == \{ j: \nat | \forall U: \power ((\nat \fun \nat) \fun (\nat \fun \nat)); n: \nat | opPlusI \in U \land j + 1 \leq n @ cst~(n - j) \in clos~\{ cst~n, funK \}~U~\{ comp \} \}
\end{zed}

\begin{zproof}[DescKSet\$domainCheck]
prove;
\end{zproof}

\begin{theorem}{rule DescKBase}
  0 \in DescKSet
\end{theorem}

\begin{zproof}
invoke DescKSet;
prove;
\end{zproof}

\begin{theorem}{rule DescKStep}
  \forall j: DescKSet @ 1 + j \in DescKSet
\end{theorem}

\begin{zproof}
invoke DescKSet;
prove;
instantiate U\_\_0 == U, n\_\_0 == n;
use StepK[U := U, n := n, c := n - j];
prove;
\end{zproof}

\begin{theorem}{DescKSub}
  \nat \subseteq DescKSet
\end{theorem}

\begin{zproof}
apply natInduction to predicate \nat \subseteq DescKSet;
prove;
invoke DescKSet;
prove;
\end{zproof}

\begin{theorem}{OneInK}
  \forall U: \power ((\nat \fun \nat) \fun (\nat \fun \nat)); m: \nat | opPlusI \in U @ cst~1 \in clos~\{ cst~(1 + m), funK \}~U~\{ comp \}
\end{theorem}

\begin{zproof}
prove;
use DescKSub;
invoke DescKSet;
with enabled (inPower) prove;
instantiate e\_\_0 == m, U\_\_0 == U, n == 1 + m;
with enabled (inPower) prove;
\end{zproof}


\begin{theorem}{StepL}
  \forall U: \power ((\nat \fun \nat) \fun (\nat \fun \nat)); n, c: \nat | opPlusI \in U \land 1 \leq n \land 1 \leq c \land cst~c \in clos~\{ cst~n, funL \}~U~\{ comp \} @ cst~(c - 1) \in clos~\{ cst~n, funL \}~U~\{ comp \}
\end{theorem}

\begin{zproof}
prove;
use CstShift[X := funL, U := U, n := n, c := c, k := 2 * n * c * c + 2 * c];
use LDescOne[n1 := n - 1, c := c];
prove;
use CstShift[X := funL, U := U, n := n, c := (n - 1) * c, k := 2 * n * c * c];
use LDescTwo[n1 := n - 1, c1 := c - 1];
prove;
\end{zproof}

\begin{zed}
  DescLSet == \{ j: \nat | \forall U: \power ((\nat \fun \nat) \fun (\nat \fun \nat)); n: \nat | opPlusI \in U \land j + 1 \leq n @ cst~(n - j) \in clos~\{ cst~n, funL \}~U~\{ comp \} \}
\end{zed}

\begin{zproof}[DescLSet\$domainCheck]
prove;
\end{zproof}

\begin{theorem}{rule DescLBase}
  0 \in DescLSet
\end{theorem}

\begin{zproof}
invoke DescLSet;
prove;
\end{zproof}

\begin{theorem}{rule DescLStep}
  \forall j: DescLSet @ 1 + j \in DescLSet
\end{theorem}

\begin{zproof}
invoke DescLSet;
prove;
instantiate U\_\_0 == U, n\_\_0 == n;
use StepL[U := U, n := n, c := n - j];
prove;
\end{zproof}

\begin{theorem}{DescLSub}
  \nat \subseteq DescLSet
\end{theorem}

\begin{zproof}
apply natInduction to predicate \nat \subseteq DescLSet;
prove;
invoke DescLSet;
prove;
\end{zproof}

\begin{theorem}{OneInL}
  \forall U: \power ((\nat \fun \nat) \fun (\nat \fun \nat)); m: \nat | opPlusI \in U @ cst~1 \in clos~\{ cst~(1 + m), funL \}~U~\{ comp \}
\end{theorem}

\begin{zproof}
prove;
use DescLSub;
invoke DescLSet;
with enabled (inPower) prove;
instantiate e\_\_0 == m, U\_\_0 == U, n == 1 + m;
with enabled (inPower) prove;
\end{zproof}


With $\bar 1$ in the closure, the bases with $\bar n$ and with $\bar 1$ are equivalent.

\begin{theorem}{CnUClosed}
  \forall X: \nat \fun \nat; U: \power ((\nat \fun \nat) \fun (\nat \fun \nat)); n: \nat @ clos~\{ cst~n, X \}~U~\{ comp \} \in UClosed~U
\end{theorem}

\begin{zproof}
prove;
use ClosInPow[I := \{ cst~n, X \}, U := U, B := \{ comp \}];
with enabled (UClosedDef) prove;
\end{zproof}

\begin{theorem}{CnBClosed}
  \forall X: \nat \fun \nat; U: \power ((\nat \fun \nat) \fun (\nat \fun \nat)); n: \nat @ clos~\{ cst~n, X \}~U~\{ comp \} \in BClosed~\{ comp \}
\end{theorem}

\begin{zproof}
prove;
use ClosInPow[I := \{ cst~n, X \}, U := U, B := \{ comp \}];
with enabled (BClosedDef) prove;
\end{zproof}

\begin{theorem}{CEqCOne}
  \forall X: \nat \fun \nat; U: \power ((\nat \fun \nat) \fun (\nat \fun \nat)); n: \nat | opPlusI \in U \land 1 \leq n \land cst~1 \in clos~\{ cst~n, X \}~U~\{ comp \} @ clos~\{ cst~n, X \}~U~\{ comp \} = clos~\{ cst~1, X \}~U~\{ comp \}
\end{theorem}

\begin{zproof}
prove;
use CstShift[X := X, U := U, n := 1, c := 1, k := n - 1];
use ClosLeast[I := \{ cst~1, X \}, U := U, B := \{ comp \}, Z := clos~\{ cst~n, X \}~U~\{ comp \}];
use ClosLeast[I := \{ cst~n, X \}, U := U, B := \{ comp \}, Z := clos~\{ cst~1, X \}~U~\{ comp \}];
use ClosInPow[I := \{ cst~n, X \}, U := U, B := \{ comp \}];
use ClosInPow[I := \{ cst~1, X \}, U := U, B := \{ comp \}];
use CnUClosed[X := X, U := U, n := n];
use CnUClosed[X := X, U := U, n := 1];
use CnBClosed[X := X, U := U, n := n];
use CnBClosed[X := X, U := U, n := 1];
use SubAntisym[\nat \fun \nat][A := clos~\{ cst~n, X \}~U~\{ comp \}, B := clos~\{ cst~1, X \}~U~\{ comp \}];
prove;
\end{zproof}

\begin{theorem}{CEq}
  \forall U: \power ((\nat \fun \nat) \fun (\nat \fun \nat)); n: \nat | opPlusI \in U \land 1 \leq n @ clos~\{ cst~n, funE \}~U~\{ comp \} = clos~\{ cst~1, funE \}~U~\{ comp \} \land clos~\{ cst~n, funK \}~U~\{ comp \} = clos~\{ cst~1, funK \}~U~\{ comp \} \land clos~\{ cst~n, funL \}~U~\{ comp \} = clos~\{ cst~1, funL \}~U~\{ comp \}
\end{theorem}

\begin{zproof}
prove;
use OneInE[U := U, m := n - 1];
use OneInK[U := U, m := n - 1];
use OneInL[U := U, m := n - 1];
use CEqCOne[X := funE, U := U, n := n];
use CEqCOne[X := funK, U := U, n := n];
use CEqCOne[X := funL, U := U, n := n];
prove;
\end{zproof}

With $\bar 0$, every generated function fixes $0$ (for $U = \{F^+, F^\square\}$ or $\{F^+, F^{-1}\}$).

\begin{theorem}{FixZeroUClosedBox}
  FixZero \in UClosed~\{ opPlusI, opBox \}
\end{theorem}

\begin{zproof}
use FixZeroInPower;
with enabled (UClosedDef) prove;
invoke FixZero;
prove;
split u = opBox;
cases;
prove;
next;
prove;
next;
\end{zproof}

\begin{theorem}{CZeroFix}
  \forall X: \nat \fun \nat; U: \power ((\nat \fun \nat) \fun (\nat \fun \nat)) | X~0 = 0 \land FixZero \in UClosed~U @ \lnot cst~1 \in clos~\{ cst~0, X \}~U~\{ comp \}
\end{theorem}

\begin{zproof}
prove;
use ClosLeast[I := \{ cst~0, X \}, U := U, B := \{ comp \}, Z := FixZero];
use FixZeroInPower;
use FixZeroBClosed;
invoke FixZero;
prove;
apply inPower to predicate clos~(\{ X \} \cup \{ cst~0 \})~U~\{ comp \} \in \power \{ F: \nat \fun \nat | F~0 = 0 \};
instantiate e == cst~1;
prove;
\end{zproof}


\begin{theorem}{CnBoxNeq}
  \forall X: \nat \fun \nat | X~0 = 0 @ \lnot clos~\{ cst~0, X \}~\{ opPlusI, opBox \}~\{ comp \} = C6
\end{theorem}

\begin{zproof}
use CZeroFix[X := X, U := \{ opPlusI, opBox \}];
use FixZeroUClosedBox;
prove;
split clos~(\{ X \} \cup \{ cst~0 \})~(\{ opPlusI \} \cup \{ opBox \})~\{ comp \} = C6;
cases;
equality substitute clos~(\{ X \} \cup \{ cst~0 \})~(\{ opPlusI \} \cup \{ opBox \})~\{ comp \};
prove;
next;
prove;
next;
\end{zproof}

\begin{theorem}{CnBox}
  \forall n: \nat @ (clos~\{ cst~n, funE \}~\{ opPlusI, opBox \}~\{ comp \} = C6 \iff 1 \leq n) \land (clos~\{ cst~n, funK \}~\{ opPlusI, opBox \}~\{ comp \} = C6 \iff 1 \leq n) \land (clos~\{ cst~n, funL \}~\{ opPlusI, opBox \}~\{ comp \} = C6 \iff 1 \leq n)
\end{theorem}

\begin{zproof}
prove;
split n = 0;
cases;
use CnBoxNeq[X := funE];
use CnBoxNeq[X := funK];
use CnBoxNeq[X := funL];
use KZero;
prove;
next;
use CEq[U := \{ opPlusI, opBox \}, n := n];
use Theorem68;
with enabled (C6Def, C6KDef, C6LDef) prove;
next;
\end{zproof}


\begin{theorem}{CnInvNeq}
  \forall X: \nat \fun \nat | X~0 = 0 @ \lnot clos~\{ cst~0, X \}~\{ opPlusI, opInv \}~\{ comp \} = ClosJR
\end{theorem}

\begin{zproof}
use CZeroFix[X := X, U := \{ opPlusI, opInv \}];
use FixZeroUClosedInv;
prove;
split clos~(\{ X \} \cup \{ cst~0 \})~(\{ opPlusI \} \cup \{ opInv \})~\{ comp \} = ClosJR;
cases;
equality substitute clos~(\{ X \} \cup \{ cst~0 \})~(\{ opPlusI \} \cup \{ opInv \})~\{ comp \};
prove;
next;
prove;
next;
\end{zproof}

\begin{theorem}{CnInv}
  \forall n: \nat @ (clos~\{ cst~n, funE \}~\{ opPlusI, opInv \}~\{ comp \} = ClosJR \iff 1 \leq n) \land (clos~\{ cst~n, funK \}~\{ opPlusI, opInv \}~\{ comp \} = ClosJR \iff 1 \leq n) \land (clos~\{ cst~n, funL \}~\{ opPlusI, opInv \}~\{ comp \} = ClosJR \iff 1 \leq n)
\end{theorem}

\begin{zproof}
prove;
split n = 0;
cases;
use CnInvNeq[X := funE];
use CnInvNeq[X := funK];
use CnInvNeq[X := funL];
use KZero;
prove;
next;
use CEq[U := \{ opPlusI, opInv \}, n := n];
use GREqJR;
with enabled (CGEDef, CGKDef, CGLDef) prove;
next;
\end{zproof}


\end{document}
